\input{preamble}
\begin{document}
\input{portada}

\section*{Introducción}

El presente documento se configura como un marco normativo y evaluativo de carácter fundamental, específicamente diseñado para sistematizar el proceso de dictamen de proyectos de tesis de pregrado en la Escuela Profesional de Ingeniería Civil. Su objetivo central es la estandarización de la revisión académica mediante la aplicación de criterios rigurosos que reduzcan la subjetividad y aseguren el cumplimiento de la normativa legal vigente.  

Al establecer directrices explícitas y coherentes, este manual garantiza que los protocolos de investigación respeten la integridad científica y los principios éticos, a la vez que alcanzan los elevados estándares de calidad técnica exigidos en el contexto universitario contemporáneo. De este modo, se constituye en un instrumento de referencia esencial para tesistas, asesores y miembros de jurados evaluadores, proporcionando un marco de actuación uniforme y verificable a lo largo de todo el proceso de formulación, revisión y aprobación de proyectos de tesis.

La estructura de este manual se fundamenta en un marco de evaluación multidimensional que trasciende la mera revisión metodológica, incorporando de manera sistemática la pertinencia institucional y el impacto social del proyecto. Mediante rúbricas exhaustivas, se verifica el grado de alineación de la investigación con la misión y visión de la universidad, así como con el perfil de egreso del ingeniero civil. Asimismo, se establece la exigencia de una articulación sólida con las demandas globales y locales, valorando la contribución empírica de la tesis al logro de los Objetivos de Desarrollo Sostenible (ODS) y a la implementación del Plan Estratégico de Desarrollo Nacional. En consecuencia, se fomenta el desarrollo de una ciencia aplicada orientada a responder de manera efectiva y contextualizada a los desafíos del entorno regional y nacional.

En su núcleo científico, el documento plantea la necesidad de un ajuste epistemológico riguroso y de una coherencia metodológica estricta. Se establece que la problematización del fenómeno debe trascender la mera aplicación profesional para configurarse como una verdadera brecha en el estado del conocimiento. El manual expone con detalle la exigencia de formular hipótesis empíricamente contrastables, de operacionalizar las variables con precisión metrológica y de diseñar esquemas experimentales o analíticos susceptibles de reproducción independiente. Este umbral de exigencia técnica asegura que las investigaciones en geotecnia, estructuras o pavimentos fundamenten sus conclusiones en principios físicos y mecánicos verificables, mediante la aplicación de un control sistemático de las fuentes de incertidumbre y la garantía de la validez interna y externa de los resultados empíricos obtenidos y proyectados.

De manera complementaria, la viabilidad operativa y la calidad formal del documento son sometidas a un análisis exhaustivo y sistemático. La guía establece criterios estrictos para la evaluación del presupuesto y del cronograma, garantizando que la asignación de recursos, la planificación de ensayos de laboratorio y los tiempos de ejecución sean técnicamente factibles y se ajusten a condiciones reales de operación. Paralelamente, se exige un estándar elevado en la redacción académica y en la presentación del documento, que incluye el uso preciso y unívoco de la terminología propia de la ingeniería, una organización jerárquica interna consistente y el cumplimiento estricto de las normas tipográficas y de citación vigentes. Esta formalización resulta esencial para asegurar la trazabilidad de los datos y la transparencia en la comunicación científica.

Como colofón, la implementación de este manual de rúbricas constituye un avance cualitativo mayor hacia la consolidación de estándares de excelencia académica en la formación de ingenieros civiles. Al proporcionar un instrumento de evaluación con criterios objetivos, reproducibles y escalables, la facultad posibilita la detección temprana de debilidades estructurales en los protocolos de investigación, lo que a su vez facilita la adopción de medidas correctivas en fases iniciales del proceso. 

Esta herramienta normativa no solo racionaliza y sistematiza los procedimientos administrativos de revisión y aprobación, sino que también promueve una cultura de investigación caracterizada por la integridad, la transparencia y la rigurosidad metodológica, orientada a la producción de resultados de alto impacto. En consecuencia y a través de el, la Universidad Andina del Cusco se fortalece y operacionaliza su compromiso institucional con la generación de conocimiento innovador, socialmente pertinente y metodológicamente sólido, contribuyendo de manera sustantiva al desarrollo tecnológico y sostenible del país.

\clearpage
\section{Cumplimiento de Normativa}

El propósito de esta sección es verificar el cumplimiento integral de las obligaciones normativas derivadas del corpus legal peruano vigente al año 2026 en materia de formulación de tesis e investigación. La base normativa está al amparo de:

\vspace{10pt}
\begingroup
\normalsize
\begin{enumerate}[leftmargin=*, nosep, label={\arabic*.}]
  \item Constitución Política del Perú (1993), Art.~18 --- Autonomía universitaria.
  \item Ley N.\textsuperscript{o}~30220, Ley Universitaria (2014, con modif.\ al 2026) --- Arts.~5.7, 8.3, 45, 47, 48, 53, 87.5, 94, 95, 100.13.
  \item Ley N.\textsuperscript{o}~31803 (2023) --- Modifica Arts.~87.5 y 100.13 de la Ley 30220 (gratuidad de asesoría).
  \item Ley N.\textsuperscript{o}~31250 (2021) --- Ley del Sistema Nacional de Ciencia, Tecnología e Innovación (SINACTI).
  \item Ley N.\textsuperscript{o}~30035 (2013) --- Repositorio Nacional Digital de CTI de Acceso Abierto.
  \item Decreto Legislativo N.\textsuperscript{o}~822 --- Ley sobre el Derecho de Autor.
  \item Resolución de Presidencia N.\textsuperscript{o}~035-2024-CONCYTEC-P --- Código Nacional de Integridad Científica.
  \item Directiva N.\textsuperscript{o}~001-2020-CONCYTEC-P --- Directiva ALICIA (Repositorio de Acceso Abierto).
  \item Directiva de Gestión de Datos de Investigación (CONCYTEC/SINACTI).
  \item RCD N.\textsuperscript{o}~084-2022-SUNEDU/CD --- Texto Integrado del Reglamento RENATI.
  \item RCD N.\textsuperscript{o}~0003-2026-SUNEDU-CD - Reglamento de Registro Nacional de Grados y Títulos (modif.\ 2026), incl.\ Art.~17-A.
  \item RD N.\textsuperscript{o}~00001-2026-SUNEDU-DS-DIRGRATU --- Supresión de modalidad de estudios en diplomas.
  \item Reglamento de Infracciones y Sanciones de SUNEDU (2019).
  \item RM N.\textsuperscript{o}~244-2025-MINEDU --- Norma Técnica para EESP.
\end{enumerate}
\endgroup

\clearpage
\subsection{Fundamento Legal y Legitimidad Institucional del Proyecto}

\textit{A continuación se presenta una rúbrica que sintetiza el marco normativo y el respaldo institucional del proyecto. La tabla organiza los criterios y sus indicadores verificables, y permite valorar el grado de cumplimiento mediante una escala de 0 a 3, facilitando una lectura comparativa y transparente de la legitimidad institucional.} Los criterios e indicadores verificables de la legitimidad institucional del proyecto se indican en la Tabla \ref{tab:Rúbrica Legitimidad Institucional}

\begingroup
\scriptsize
\begin{longtable}{|C{0.6cm}|L{3cm}|L{7.8cm}|C{0.4cm}|C{0.4cm}|C{0.4cm}|C{0.4cm}|}
\caption{Rúbrica de Evaluación del Fundamento Legal y Legitimidad Institucional del Proyecto}
\label{tab:Rúbrica Legitimidad Institucional} \\
\hline
\rowcolor{headerblue}
\thd{N.\textsuperscript{o}} & \thd{Criterio} & \thd{Indicador verificable} & \thd{3} & \thd{2} & \thd{1} & \thd{0} \\
\hline
\endfirsthead
\hline
\rowcolor{headerblue}
\thd{N.\textsuperscript{o}} & \thd{Criterio} & \thd{Indicador verificable} & \thd{3} & \thd{2} & \thd{1} & \thd{0} \\
\hline
\endhead
1.1 & Correspondencia entre el tipo de trabajo y el nivel académico postulado & El proyecto se autoidentifica correctamente como \enquote{tesis} (título profesional, maestría o doctorado) o \enquote{trabajo de investigación} (bachiller), conforme al Art.~45 de la Ley 30220. Para doctorado, se explicita la orientación de \enquote{máxima rigurosidad académica y carácter original}. & & & & \\ \hline
\rowcolor{lightgray}
1.2 & Designación formal de asesor & El proyecto acredita la designación oficial del asesor mediante resolución o constancia institucional. & & & & \\ \hline
1.3 & Adecuación al Reglamento de Grados y Títulos institucional vigente & La estructura del proyecto sigue la plantilla y los requisitos formales establecidos en el Reglamento interno de la universidad. & & & & \\ \hline
1.4 & Identificación de la línea de investigación institucional & El proyecto se inscribe explícitamente en una línea o sublínea de investigación aprobada por la universidad, conforme al principio de pertinencia del Art.~48 de la Ley 30220 (investigación que responde a necesidades de la sociedad y la realidad nacional). & & & & \\ \hline
\multicolumn{3}{|r|}{\textbf{Subtotal Sección~I (máx.~12 puntos):}} & \multicolumn{4}{c|}{/12} \\ \hline
\end{longtable}
\endgroup

\clearpage
\subsection{Integridad Científica y Ética de la Investigación}

\textit{Debe ajustarse a la siguiente base normativa: Res. de Presidencia N.\textsuperscript{o}~035-2024-CONCYTEC-P (Código Nacional de Integridad Científica); Ley N.\textsuperscript{o}~31250 (SINACTI); D.\,Leg.~822 (Ley sobre el Derecho de Autor); Reglamento RENATI, RCD N.\textsuperscript{o}~084-2022-SUNEDU/CD, Art.~8.} Los criterios pueden verse en la tabla \ref{tab:integridad}.

\begingroup
\scriptsize
\begin{longtable}{|C{0.6cm}|L{3cm}|L{7.5cm}|C{0.4cm}|C{0.4cm}|C{0.4cm}|C{0.4cm}|}
\caption{Rúbrica de la Integridad Científica y Ética del Proyecto}
\label{tab:integridad} \\
\hline
\rowcolor{headerblue}
\thd{N.\textsuperscript{o}} & \thd{Criterio} & \thd{Indicador verificable} & \thd{3} & \thd{2} & \thd{1} & \thd{0} \\
\hline
\endfirsthead
\hline
\rowcolor{headerblue}
\thd{N.\textsuperscript{o}} & \thd{Criterio} & \thd{Indicador verificable} & \thd{3} & \thd{2} & \thd{1} & \thd{0} \\
\hline
\endhead
2.1 & Declaración jurada de originalidad y autenticidad & El proyecto incorpora una declaración jurada de originalidad firmada solidariamente por el autor y el asesor, conforme al Art.~8 del Reglamento RENATI. Se indica la disposición a someterse a verificación de similitud. & & & & \\ \hline
\rowcolor{lightgray}
2.2 & Ausencia de indicios de plagio, fabricación o falsificación en el proyecto & El marco teórico y los antecedentes del proyecto no presentan indicios de apropiación ilícita de textos, ideas o datos ajenos. Se citan correctamente las fuentes. No se detectan tipologías de Mala Conducta Científica definidas en la Res.~035-2024-CONCYTEC-P (plagio, fabricación, falsificación, autoría inmerecida). & & & & \\ \hline
2.3 & Sistema de citación y referenciación normalizado & El proyecto emplea de manera uniforme y rigurosa un sistema de citación reconocido (APA, Vancouver, ISO~690, u otro exigido por el reglamento institucional). Las referencias bibliográficas son verificables y rastreables. & & & & \\ \hline
\rowcolor{lightgray}
2.4 & Compromiso explícito con el Código Nacional de Integridad Científica & El proyecto declara adhesión a los principios del Código Nacional de Integridad Científica (Res.~035-2024-CONCYTEC-P) y reconoce la jurisdicción del Comité de Integridad Científica (CIC) de la universidad. & & & & \\ \hline
2.5 & Corresponsabilidad del asesor documentada & Existe constancia de que el asesor ha revisado y avalado el contenido del proyecto. Se acredita el cumplimiento de la supervisión directa y corresponsabilidad exigida al docente asesor por el Código de Integridad y los Arts.~87.5 y 95.2 de la Ley 30220. & & & & \\ \hline
\rowcolor{lightgray}
2.6 &
Conocimiento del régimen sancionador aplicable & El proyecto o la documentación adjunta evidencia que el autor y el asesor conocen las consecuencias legales del fraude académico: nulidad de grado (Art.~17-A, RCD 0003-2026), sanciones institucionales (RIS-SUNEDU), y sanciones al docente (Arts.~94--95, Ley 30220). &
& & & \\ \hline
\multicolumn{3}{|r|}{\textbf{Subtotal Sección~II (máx.~18 puntos):}} & \multicolumn{4}{c|}{/18} \\ \hline
\end{longtable}
\endgroup

\clearpage
\subsection{Plan de Gestión de Datos y Ciencia Abierta}

\textit{Cumple con la base normativa: Ley N.\textsuperscript{o}~31250 (SINACTI), Título Preliminar; Directiva de Gestión de Datos de Investigación (CONCYTEC); Ley N.\textsuperscript{o}~30035 (Repositorio Nacional Digital / ALICIA); Directiva N.\textsuperscript{o}~001-2020-CONCYTEC-P.} Los criterios para evaluar esta sección se muestran en la Tabla \ref{tab:Gestión de datos}.

\begingroup
\scriptsize
\begin{longtable}{|C{0.6cm}|L{2.8cm}|L{8.1cm}|C{0.4cm}|C{0.4cm}|C{0.4cm}|C{0.4cm}|}
\caption{Rúbrica para Evaluar el Plan de Gestión de Datos y Ciencia Abierta}
\label{tab:Gestión de datos} \\
\hline
\rowcolor{headerblue}
\thd{N.\textsuperscript{o}} & \thd{Criterio} & \thd{Indicador verificable} & \thd{3} & \thd{2} & \thd{1} & \thd{0} \\
\hline
\endfirsthead
\hline
\rowcolor{headerblue}
\thd{N.\textsuperscript{o}} & \thd{Criterio} & \thd{Indicador verificable} & \thd{3} & \thd{2} & \thd{1} & \thd{0} \\
\hline
\endhead
3.1 & Formulación del Plan de Gestión de Datos (PGD) & El proyecto incluye un PGD o sección equivalente que detalla: (a)~tipo de datos a recopilar; (b)~formatos de almacenamiento (CSV, PDF, etc.); (c)~mecanismos de seguridad de la información; (d)~responsable de custodia. Obligatorio en posgrado y proyectos con financiamiento público. & & & & \\ \hline
\rowcolor{lightgray}
3.2 & Trazabilidad de la evidencia empírica planificada & El proyecto describe los mecanismos concretos para garantizar la trazabilidad absoluta de los datos primarios: bitácoras, actas de validación, registros fotográficos, hojas de cálculo numeradas u otros instrumentos que impidan la fabricación posterior de datos. & & & & \\ \hline
3.3 & Compromiso de depósito en repositorio institucional & El proyecto declara el compromiso de depositar la tesis finalizada y los datos de soporte en el repositorio institucional de la universidad, conforme a la Ley 30035 y la Directiva ALICIA, para su cosecha por el repositorio nacional. & & & & \\ \hline
\rowcolor{lightgray}
3.4 & Licenciamiento y acceso abierto & El proyecto identifica el tipo de licencia prevista (p.\,ej., Creative Commons BY-NC) o, en su defecto, justifica legalmente la restricción de acceso (patentes, propiedad industrial, datos sensibles) conforme al Reglamento RENATI. & & & & \\ \hline
3.5 & Preservación de metadatos & El proyecto prevé la generación de metadatos descriptivos que permitan la interoperabilidad del registro con el sistema ALICIA y el RENATI (autor, título, descriptores, tipo de investigación, institución). & & & & \\ \hline
\multicolumn{3}{|r|}{\textbf{Subtotal Sección~III (máx.~15 puntos):}} & \multicolumn{4}{c|}{/15} \\ \hline
\end{longtable}
\endgroup

\clearpage
La Tabla~\ref{tab:consolidado-normativa} resume los puntajes obtenidos en cada sección de la evaluación de cumplimiento normativo. La Tabla~\ref{tab:escala-calificacion} establece los niveles de calificación del cumplimiento normativo según el puntaje total obtenido.

\begin{table}[h]
\caption{Cuadro consolidado de resultados de la evaluación de cumplimiento normativo}\label{tab:consolidado-normativa}
\vspace{6pt}
\centering
\scriptsize
\begin{tabularx}{\textwidth}{|C{0.8cm}|L{7.5cm}|C{1.8cm}|C{1.8cm}|X|}
\hline
\rowcolor{headerblue}
\thd{Dim.} & \thd{Denominación} & \thd{Pts. Máx.} & \thd{Pts. Obt.} & \thd{Obs.} \\ \hline
I & Fundamento Legal y Legitimidad Institucional del Proyecto & 12 &  & \\ \hline
\rowcolor{lightgray}
II & Integridad Científica y Ética de la Investigación & 18 &  & \\ \hline
III & Plan de Gestión de Datos y Ciencia Abierta & 15 &  & \\ \hline
\rowcolor{white}
& \textbf{TOTAL GENERAL} & \textbf{45} &  & \\ \hline
\end{tabularx}
\end{table}

\begin{table}[htbp]
\centering
\scriptsize
\caption{Escala de calificación global del cumplimiento normativo}\label{tab:escala-calificacion}
\begin{tabularx}{\textwidth}{|C{2cm}|C{3.5cm}|X|}
  \hline \rowcolor{headerblue} \thd{Rango de puntaje} & \thd{Nivel de cumplimiento} & \thd{Descriptor} \\ \hline
  \textbf{37--45} & \textbf{Cumplimiento \newline excelente} & El proyecto satisface de manera integral y rigurosa las exigencias normativas. La documentación es completa, coherente y plenamente verificable. \\ \hline
  \rowcolor{lightgray} \textbf{28--36} & \textbf{Cumplimiento \newline satisfactorio} & El proyecto cumple la mayoría de las exigencias normativas con evidencia suficiente. Se detectan omisiones menores que no comprometen la viabilidad legal del proyecto. \\ \hline
  \textbf{19--27} & \textbf{Cumplimiento \newline parcial} & El proyecto presenta deficiencias normativas apreciables. Requiere correcciones sustantivas antes de su aprobación para garantizar la conformidad legal e institucional. \\ \hline
  \rowcolor{lightgray} \textbf{0--18} & \textbf{Cumplimiento \newline insuficiente} & El proyecto incumple de manera generalizada las exigencias normativas. No reúne las condiciones mínimas para su aprobación y requiere reformulación integral. \\ \hline
\end{tabularx}
\end{table}

\clearpage
\section{Respuesta a las Necesidades del Planeta}

{\footnotesize Los Objetivos de Desarrollo Sostenible (ODS), adoptados por la ONU en 2015, son un llamado global para erradicar la pobreza, proteger el planeta y garantizar paz y prosperidad para todas las personas antes de 2030. Los 17 ODS son interdependientes: las acciones en un ámbito afectan a los demás y el desarrollo debe equilibrar la sostenibilidad social, económica y ambiental. Esta rúbrica permite evaluar la contribución potencial de un proyecto de tesis según los cinco ejes de la Agenda 2030.: \textbf{Personas}, \textbf{Planeta}, \textbf{Prosperidad}, \textbf{Paz} y \textbf{Alianzas}. Los 17 Objetivos de Desarrollo Sostenible se muestran en la tabla \ref{tab:ods}. El agrupamiento por Ejes de la Agenda 2030 se aprecia en la tabla \ref{tab:ejes-agenda2030}. Por último, la escala de valoración se aprecia en la tabla \ref{tab:escala-valoracion-ods}. La referencia normativa está con base en: Naciones Unidas, \textit{Transformar nuestro mundo: la Agenda 2030 para el Desarrollo Sostenible}, Resolución A/RES/70/1, 25 de septiembre de 2015. PNUD \url{https://www.undp.org/sustainable-development-goals}.\par}

\begin{table}[htbp]
\centering
\scriptsize
\caption{Objetivos de Desarrollo Sostenible (ODS) de la Agenda 2030.}
\label{tab:ods}
\begin{tabularx}{\textwidth}{|C{1.3cm}|X|C{1.3cm}|X|}
\hline
\rowcolor{headerblue}
\thd{ODS} & \thd{Denominación} & \thd{ODS} & \thd{Denominación} \\ \hline
1 & Fin de la pobreza & 10 & Reducción de las desigualdades \\ \hline
\rowcolor{lightgray}
2 & Hambre cero & 11 & Ciudades y comunidades sostenibles \\ \hline
3 & Salud y bienestar & 12 & Producción y consumo responsables \\ \hline
\rowcolor{lightgray}
4 & Educación de calidad & 13 & Acción por el clima \\ \hline
5 & Igualdad de género & 14 & Vida submarina \\ \hline
\rowcolor{lightgray}
6 & Agua limpia y saneamiento & 15 & Vida de ecosistemas terrestres \\ \hline
7 & Energía asequible y no contaminante & 16 & Paz, justicia e instituciones sólidas \\ \hline
\rowcolor{lightgray}
8 & Trabajo decente y crecimiento económico & 17 & Alianzas para lograr los objetivos \\ \hline
9 & Industria, innovación e infraestructura & & \\ \hline
\end{tabularx}
\end{table}

\begin{table}[htbp]
\centering
\scriptsize
\caption{Agrupamiento de los ODS por ejes de la Agenda 2030.}
\label{tab:ejes-agenda2030}
\begin{tabularx}{\textwidth}{|L{2.5cm}|L{3.8cm}|X|}
\hline
\rowcolor{headerblue}
\thd{Eje} & \thd{ODS asociados} & \thd{Propósito} \\ \hline
\textbf{Personas} & ODS 1, 2, 3, 4, 5 & Poner fin a la pobreza y el hambre, garantizar vida sana, educación de calidad e igualdad de género. \\ \hline
\rowcolor{lightgray}
\textbf{Prosperidad} & ODS 7, 8, 9, 10, 11 & Garantizar vidas prósperas y plenas en armonía con la naturaleza; infraestructura resiliente, innovación, energía, empleo y ciudades sostenibles. \\ \hline
\textbf{Planeta} & ODS 6, 12, 13, 14, 15 & Proteger los recursos naturales y el clima: agua, producción responsable, acción climática, ecosistemas marinos y terrestres. \\ \hline
\rowcolor{lightgray}
\textbf{Paz y Alianzas} & ODS 16, 17 & Fomentar sociedades pacíficas e inclusivas y revitalizar alianzas globales para el desarrollo sostenible. \\ \hline
\end{tabularx}
\end{table}

\begin{table}[htbp]
\centering
\scriptsize
\caption{Escala de valoración para la contribución del proyecto a los ODS o ejes evaluados.}
\label{tab:escala-valoracion-ods}
\begin{tabularx}{\textwidth}{|C{2.8cm}|C{1cm}|X|}
\hline
\rowcolor{headerblue}
\thd{Nivel} & \thd{Pts.} & \thd{Descriptor operativo} \\ \hline
\textcolor{csuf}{\textbf{Excelente}} & \textbf{3} & El proyecto evidencia contribución explícita, fundamentada y verificable al ODS o eje evaluado. La articulación es orgánica dentro de la estructura argumentativa y metodológica del proyecto, con indicadores o metas específicas del ODS identificados. \\ \hline
\rowcolor{lightgray}
\textcolor{cace}{\textbf{Aceptable}} & \textbf{2} & La contribución es identificable pero parcial o implícita. El proyecto aborda el ODS sin fundamentación suficiente o sin vincular metas específicas. Susceptible de fortalecimiento. \\ \hline
\textcolor{cdef}{\textbf{Deficiente}} & \textbf{1} & El vínculo con el ODS es tangencial o meramente declarativo. No hay sustento argumentativo ni coherencia con el diseño del proyecto. \\ \hline
\rowcolor{lightgray}
\textcolor{cins}{\textbf{Ausente}} & \textbf{0} & El proyecto no evidencia relación alguna con el ODS o eje evaluado. \\ \hline
\end{tabularx}
\end{table}

\clearpage
\subsection{Identificación y Fundamentación de ODS Pertinentes}

\textit{Se evalúa si el proyecto identifica de manera explícita, justificada y coherente los ODS a los que contribuye, estableciendo la trazabilidad entre el problema de investigación y las metas globales de la Agenda 2030.}. Los criterios de este ítem se revelan en la tabla \ref{tab:dimension1-ods}.

\begingroup
\scriptsize
\begin{longtable}{|C{0.55cm}|L{3.8cm}|L{5.8cm}|C{0.65cm}|C{0.65cm}|C{0.65cm}|C{0.65cm}|}
\caption{Rúbrica de evaluación: identificación y fundamentación de ODS pertinentes (Dimensión~I).}\label{tab:dimension1-ods}\\
\hline
\rowcolor{headerblue}
\thd{N.\textsuperscript{o}} & \thd{Criterio} & \thd{Indicador verificable} & \thd{3} & \thd{2} & \thd{1} & \thd{0} \\
\hline\endfirsthead
\hline
\rowcolor{headerblue}
\thd{N.\textsuperscript{o}} & \thd{Criterio} & \thd{Indicador verificable} & \thd{3} & \thd{2} & \thd{1} & \thd{0} \\
\hline\endhead
1.1 & Identificación explícita de ODS vinculados al proyecto & El proyecto declara de manera explícita al menos un ODS al que contribuye. La identificación se basa en el análisis del problema, los objetivos y los resultados esperados, no en una mención genérica o nominal. Se cita el número y la denominación oficial del ODS. &
& & & \\ \hline
\rowcolor{lightgray}
1.2 & Vinculación con metas específicas (\textit{targets}) del ODS & El proyecto identifica las metas específicas (\textit{targets}) del ODS seleccionado (p.\,ej., meta~6.1, meta~9.1, meta~11.5) y argumenta cómo la investigación contribuye a su consecución. El vínculo no es genérico sino operacionalizado con indicadores o resultados concretos. &
& & & \\ \hline
1.3 & Justificación de la pertinencia del ODS en el contexto local/regional & El proyecto argumenta por qué el ODS seleccionado es pertinente para el contexto local, regional o nacional en el que se desarrolla la investigación. Se presentan datos, estadísticas o diagnósticos que evidencian la brecha o necesidad que el proyecto pretende atender en relación con el ODS. &
& & & \\ \hline
\rowcolor{lightgray}
1.4 & Coherencia entre el ODS seleccionado y el diseño de investigación & Existe correspondencia lógica entre el ODS declarado y el planteamiento del problema, los objetivos, las variables y la metodología del proyecto. El ODS no es un añadido retórico sino un eje articulador verificable en la estructura del proyecto. &
& & & \\ \hline
\multicolumn{3}{|r|}{\cellcolor{white}\textbf{Subtotal Dimensión~I (máx.~12 puntos):}} & \multicolumn{4}{c|}{\cellcolor{white}\textbf{/12}} \\ \hline
\end{longtable}
\endgroup

\clearpage
\subsection{Eje del Planeta: Sostenibilidad Ambiental}

\textit{Se evalúa la contribución del proyecto a la protección de los recursos naturales y el clima: agua limpia y saneamiento (ODS~6), producción y consumo responsables (ODS~12), acción por el clima (ODS~13), vida submarina (ODS~14) y vida de ecosistemas terrestres (ODS~15).}. Los criterios de evaluación de este ítem se enuncian en la tabla \ref{tab:dimension2-planeta}.

\begingroup
\scriptsize
\begin{longtable}{|C{0.55cm}|L{3.8cm}|L{5.8cm}|C{0.65cm}|C{0.65cm}|C{0.65cm}|C{0.65cm}|}
\caption{Rúbrica de evaluación: eje Planeta -- sostenibilidad ambiental (Dimensión~II).}\label{tab:dimension2-planeta}\\
\hline
\rowcolor{headerblue}
\thd{N.\textsuperscript{o}} & \thd{Criterio} & \thd{Indicador verificable} & \thd{3} & \thd{2} & \thd{1} & \thd{0} \\
\hline\endfirsthead
\hline
\rowcolor{headerblue}
\thd{N.\textsuperscript{o}} & \thd{Criterio} & \thd{Indicador verificable} & \thd{3} & \thd{2} & \thd{1} & \thd{0} \\
\hline\endhead
2.1 & Agua limpia y saneamiento \odsref{ODS\,6} & El proyecto contribuye al acceso a agua potable segura, a la gestión sostenible de recursos hídricos, al tratamiento de aguas residuales, a la protección de ecosistemas acuáticos o a la infraestructura de saneamiento. Se identifican metas específicas del ODS~6 atendidas por la investigación. &
& & & \\ \hline
\rowcolor{lightgray}
2.2 & Producción y consumo responsables \odsref{ODS\,12} & El proyecto promueve la gestión eficiente de recursos naturales, la reducción de residuos, el reciclaje, la reutilización de materiales, o la adopción de prácticas constructivas o productivas con menor huella ambiental. Se evidencia un enfoque de ciclo de vida o economía circular. &
& & & \\ \hline
2.3 & Acción por el clima \odsref{ODS\,13} & El proyecto aborda la mitigación o adaptación al cambio climático: reducción de emisiones de gases de efecto invernadero, resiliencia de infraestructura ante eventos climáticos extremos, gestión de riesgos de desastres, o modelamiento de escenarios climáticos aplicados al ámbito de estudio. &
& & & \\ \hline
\rowcolor{lightgray}
2.4 & Vida de ecosistemas terrestres y/o submarinos \odsref{ODS\,14--15} & El proyecto considera la protección, restauración o uso sostenible de ecosistemas terrestres o acuáticos: conservación de suelos, gestión de cuencas, prevención de la degradación del terreno, protección de biodiversidad, o evaluación de impactos de infraestructura sobre ecosistemas. &
& & & \\ \hline
\multicolumn{3}{|r|}{\cellcolor{white}\textbf{Subtotal Dimensión~II (máx.~12 puntos):}} & \multicolumn{4}{c|}{\cellcolor{white}\textbf{/12}} \\ \hline
\end{longtable}
\endgroup

\clearpage
\subsection{Eje Prosperidad: Infraestructura, Innovación y Ciudades}

\textit{Se evalúa la contribución del proyecto a la construcción de infraestructura resiliente, la promoción de la industrialización inclusiva, la innovación tecnológica, la energía sostenible, el crecimiento económico, la reducción de desigualdades y las ciudades sostenibles.}. Los criterios para su evaluación se representan en la tabla \ref{tab:dimension3-prosperidad}.

\begingroup
\scriptsize
\begin{longtable}{|C{0.55cm}|L{3.8cm}|L{5.8cm}|C{0.65cm}|C{0.65cm}|C{0.65cm}|C{0.65cm}|}
\caption{Rúbrica de evaluación: eje Prosperidad -- infraestructura, innovación y ciudades (Dimensión~III).}\label{tab:dimension3-prosperidad}\\
\hline
\rowcolor{headerblue}
\thd{N.\textsuperscript{o}} & \thd{Criterio} & \thd{Indicador verificable} & \thd{3} & \thd{2} & \thd{1} & \thd{0} \\
\hline\endfirsthead
\hline
\rowcolor{headerblue}
\thd{N.\textsuperscript{o}} & \thd{Criterio} & \thd{Indicador verificable} & \thd{3} & \thd{2} & \thd{1} & \thd{0} \\
\hline\endhead
3.1 & Industria, innovación e infraestructura \odsref{ODS\,9} & El proyecto contribuye a la construcción de infraestructura resiliente, la promoción de la industrialización inclusiva y sostenible, o el fomento de la innovación tecnológica. Se identifica el aporte a la modernización de infraestructura, la mejora de procesos productivos o la investigación aplicada con impacto industrial o constructivo. &
& & & \\ \hline
\rowcolor{lightgray}
3.2 & Ciudades y comunidades sostenibles \odsref{ODS\,11} & El proyecto aborda desafíos urbanos: vivienda segura y asequible, sistemas de transporte sostenibles, planificación urbana inclusiva, protección del patrimonio cultural, reducción del impacto ambiental de las ciudades, gestión de residuos sólidos, o resiliencia ante desastres en zonas urbanas. &
& & & \\ \hline
3.3 & Energía asequible y no contaminante \odsref{ODS\,7} & El proyecto contribuye al acceso a energía asequible, fiable, sostenible y moderna: eficiencia energética en edificaciones o infraestructura, energías renovables, reducción de la dependencia de combustibles fósiles, o evaluación energética de sistemas constructivos. &
& & & \\ \hline
\rowcolor{lightgray}
3.4 & Trabajo decente, crecimiento económico y reducción de desigualdades \odsref{ODS\,8, 10} & El proyecto tiene potencial para generar empleo, mejorar la productividad laboral, promover el crecimiento económico sostenible, o reducir desigualdades en el acceso a servicios, infraestructura o vivienda. Se identifica el beneficiario socioeconómico directo o indirecto de los resultados esperados. &
& & & \\ \hline
\multicolumn{3}{|r|}{\cellcolor{white}\textbf{Subtotal Dimensión~III (máx.~12 puntos):}} & \multicolumn{4}{c|}{\cellcolor{white}\textbf{/12}} \\ \hline
\end{longtable}
\endgroup

\clearpage
\subsection{Eje Personas: Bienestar Social y Desarrollo Humano}

\textit{Se evalúa la contribución del proyecto al bienestar humano: erradicación de la pobreza (ODS~1), seguridad alimentaria (ODS~2), salud y bienestar (ODS~3), educación de calidad (ODS~4) e igualdad de género (ODS~5).}. El resumen de criterios destinados a su evaluación se presenta en la tabla \ref{tab:dimension4-personas}.

\begingroup
\scriptsize
\begin{longtable}{|C{0.55cm}|L{3.8cm}|L{5.8cm}|C{0.65cm}|C{0.65cm}|C{0.65cm}|C{0.65cm}|}
\caption{Rúbrica de evaluación: eje Personas -- bienestar social y desarrollo humano (Dimensión~IV).}\label{tab:dimension4-personas}\\
\hline
\rowcolor{headerblue}
\thd{N.\textsuperscript{o}} & \thd{Criterio} & \thd{Indicador verificable} & \thd{3} & \thd{2} & \thd{1} & \thd{0} \\
\hline\endfirsthead
\hline
\rowcolor{headerblue}
\thd{N.\textsuperscript{o}} & \thd{Criterio} & \thd{Indicador verificable} & \thd{3} & \thd{2} & \thd{1} & \thd{0} \\
\hline\endhead
4.1 & Fin de la pobreza y seguridad alimentaria \odsref{ODS\,1--2} & El proyecto contribuye, directa o indirectamente, a mejorar las condiciones de vida de poblaciones vulnerables: acceso a servicios básicos, infraestructura para comunidades rurales o periurbanas, seguridad alimentaria mediante infraestructura agrícola, de riego o almacenamiento, o fortalecimiento de la resiliencia ante desastres de los grupos más expuestos. &
& & & \\ \hline
\rowcolor{lightgray}
4.2 & Salud y bienestar \odsref{ODS\,3} & El proyecto atiende necesidades vinculadas a la salud pública: infraestructura sanitaria, calidad del agua potable, control de contaminación atmosférica o de suelos, reducción de riesgos para la salud derivados de deficiencias en edificaciones, vías de transporte, o gestión de residuos peligrosos. &
& & & \\ \hline
4.3 & Educación de calidad \odsref{ODS\,4} & El proyecto contribuye a la mejora de la infraestructura educativa, genera conocimiento transferible al ámbito académico, o fortalece la formación técnica y científica mediante la investigación aplicada. Se considera también la contribución a la educación para el desarrollo sostenible. &
& & & \\ \hline
\rowcolor{lightgray}
4.4 & Igualdad de género, inclusión y equidad \odsref{ODS\,5} & El proyecto incorpora la perspectiva de género, accesibilidad universal o inclusión social en su diseño: consideración de necesidades diferenciadas, participación equitativa, diseño universal de infraestructura, o análisis de impacto diferencial de las soluciones propuestas sobre grupos históricamente excluidos. &
& & & \\ \hline
\multicolumn{3}{|r|}{\cellcolor{white}\textbf{Subtotal Dimensión~IV (máx.~12 puntos):}} & \multicolumn{4}{c|}{\cellcolor{white}\textbf{/12}} \\ \hline
\end{longtable}
\endgroup

\clearpage
\subsection{Coherencia Integral, Transversalidad y Profundidad de la Contribución}

\textit{Se evalúa la articulación transversal, la profundidad argumentativa y el potencial de impacto real del proyecto respecto de los ODS, incluyendo la interconexión entre objetivos y la vinculación con alianzas e institucionalidad (ODS~16--17).}. Los criterios de clarificación de este aspecto puntual se evidencian en la tabla \ref{tab:dimension5-coherencia}.

\begingroup
\scriptsize
\begin{longtable}{|C{0.55cm}|L{3.8cm}|L{5.8cm}|C{0.65cm}|C{0.65cm}|C{0.65cm}|C{0.65cm}|}
\caption{Rúbrica de evaluación: coherencia integral, transversalidad y profundidad de la contribución (Dimensión~V).}\label{tab:dimension5-coherencia}\\
\hline
\rowcolor{headerblue}
\thd{N.\textsuperscript{o}} & \thd{Criterio} & \thd{Indicador verificable} & \thd{3} & \thd{2} & \thd{1} & \thd{0} \\
\hline\endfirsthead
\hline
\rowcolor{headerblue}
\thd{N.\textsuperscript{o}} & \thd{Criterio} & \thd{Indicador verificable} & \thd{3} & \thd{2} & \thd{1} & \thd{0} \\
\hline\endhead
5.1 & Interconexión entre múltiples ODS & El proyecto reconoce y articula la interconexión entre dos o más ODS: demuestra que los resultados en un área afectan los resultados en otra (p.\,ej., infraestructura resiliente [ODS~9] que mejora la salud [ODS~3] y reduce la pobreza [ODS~1]). La integralidad de la Agenda 2030 se refleja en el diseño del proyecto. &
& & & \\ \hline
\rowcolor{lightgray}
5.2 & Paz, justicia, instituciones sólidas y alianzas \odsref{ODS\,16--17} & El proyecto fortalece instituciones, promueve el acceso a la justicia, contribuye a la transparencia, fomenta alianzas interinstitucionales, o prevé mecanismos de transferencia de conocimiento a organismos públicos, comunidades o sector productivo. Se consideran los convenios, autorizaciones o colaboraciones planificadas. &
& & & \\ \hline
5.3 & Transversalidad del alineamiento ODS en el proyecto & El alineamiento con los ODS no se limita a una sección aislada del proyecto (p.\,ej., solo la justificación), sino que permea de manera transversal el planteamiento del problema, los objetivos, el marco teórico, el diseño metodológico y los resultados esperados. &
& & & \\ \hline
\rowcolor{lightgray}
5.4 & Potencial de impacto verificable y escalabilidad & El proyecto demuestra que su contribución al ODS es medible o verificable: se prevén indicadores de impacto, métricas de desempeño, o resultados cuantificables vinculados a las metas del ODS. Se valora el potencial de escalabilidad o replicabilidad de la solución propuesta en otros contextos. &
& & & \\ \hline
\multicolumn{3}{|r|}{\cellcolor{white}\textbf{Subtotal Dimensión~V (máx.~12 puntos):}} & \multicolumn{4}{c|}{\cellcolor{white}\textbf{/12}} \\ \hline
\end{longtable}
\endgroup

\clearpage
Tanto el cuadro consolidado de resultados como la escala de calificación de alineamiento con los ODS se aprecian en las tablas \ref{tab:cuadro-consolidado} y \ref{tab:escala-calificacion-ods}.

\begin{table}[htbp]
\centering
\scriptsize
\caption{Cuadro consolidado de resultados por dimensión de evaluación ODS.}
\label{tab:cuadro-consolidado}
\begin{tabularx}{\textwidth}{|C{0.8cm}|L{7.5cm}|C{1.8cm}|C{1.8cm}|X|}
\hline
\rowcolor{headerblue}
\thd{Dim.} & \thd{Denominación} & \thd{Pts. Máx.} & \thd{Pts. Obt.} & \thd{Obs.} \\ \hline
I & Identificación y Fundamentación de ODS Pertinentes & 12 & & \\ \hline
\rowcolor{lightgray}
II & Eje Planeta: Sostenibilidad Ambiental (ODS 6, 12, 13, 14, 15) & 12 & & \\ \hline
III & Eje Prosperidad: Infraestructura, Innovación, Ciudades (ODS 7, 8, 9, 10, 11) & 12 & & \\ \hline
\rowcolor{lightgray}
IV & Eje Personas: Bienestar Social y Desarrollo Humano (ODS 1--5) & 12 & & \\ \hline
V & Coherencia Integral, Transversalidad y Profundidad & 12 & & \\ \hline
\rowcolor{white}
& \textbf{TOTAL GENERAL} & \textbf{60} & & \\ \hline
\end{tabularx}
\end{table}

\begin{table}[htbp]
\centering
\scriptsize
\caption{Escala de calificación del alineamiento del proyecto con los ODS.}
\label{tab:escala-calificacion-ods}
\begin{tabularx}{\textwidth}{|C{2.5cm}|C{3cm}|X|}
\hline
\rowcolor{headerblue}
\thd{Rango} & \thd{Calificación} & \thd{Significado} \\ \hline
\textbf{54--60}\newline(90--100\,\%) &
\textcolor{csuf}{\textbf{ALINEAMIENTO\newline PLENO}} &
El proyecto demuestra contribución integral, explícita y fundamentada a los ODS. La investigación constituye un aporte ejemplar a la Agenda 2030, con trazabilidad verificable a metas específicas y articulación transversal entre múltiples objetivos. \\ \hline
\rowcolor{lightgray}
\textbf{42--53}\newline(70--89\,\%) &
\textcolor{cace}{\textbf{ALINEAMIENTO\newline SUFICIENTE}} &
El proyecto evidencia contribución sustancial a los ODS, con áreas específicas susceptibles de fortalecimiento. Los vínculos principales están demostrados y la pertinencia global es verificable. \\ \hline
\textbf{30--41}\newline(50--69\,\%) &
\textcolor{cdef}{\textbf{ALINEAMIENTO\newline PARCIAL}} &
La contribución es incompleta o predominantemente declarativa. El proyecto requiere fortalecimiento de la justificación y articulación con las metas de los ODS. \\ \hline
\rowcolor{lightgray}
\textbf{0--29}\newline($<$50\,\%) &
\textcolor{cins}{\textbf{ALINEAMIENTO\newline INSUFICIENTE}} &
El proyecto no demuestra contribución verificable a los ODS. La desconexión con la Agenda 2030 compromete la pertinencia global de la investigación frente a las necesidades urgentes del desarrollo sostenible. \\ \hline
\end{tabularx}
\end{table}

\clearpage
\section{Respuesta a las Necesidades del País}

{\footnotesize El nivel de complejidad metodológica del proyecto de tesis debe alinearse con la ley universitaria vigente y responder a las necesidades del país. El Plan Estratégico de Desarrollo Nacional (PEDN) al 2050, aprobado por el CEPLAN, órgano rector del SINAPLAN, es el principal instrumento de planeamiento estratégico del Perú. El PEDN sitúa a las personas y su vida en comunidad en el centro de su enfoque y define una estrategia de mediano y largo plazo para un desarrollo armónico y sostenible. Se organiza en cuatro Objetivos Nacionales (ON) —Desarrollo de las personas, Territorio sostenible, Competitividad e innovación, y Democracia y paz— que expresan cambios en el bienestar de la población bajo un enfoque sistémico. Esta rúbrica permite al dictaminante evaluar la contribución potencial de un proyecto de tesis a estos Objetivos Nacionales, a sus objetivos específicos y a los retos estructurales del PEDN, articulados con la Visión del Perú al 2050 y las 35 Políticas de Estado del Acuerdo Nacional. La Referencia normativa considerada es: Centro Nacional de Planeamiento Estratégico (CEPLAN), \textit{Plan Estratégico de Desarrollo Nacional al 2050}, aprobado por D.\,S.\,N.\textsuperscript{o}~095-2022-PCM, actualizado en noviembre de 2024. Visión del Perú al 2050, aprobada por consenso en el Foro del Acuerdo Nacional, abril de 2019. Fuente: \url{https://peru2050.ceplan.gob.pe}. En esta sección se presentan, en primer lugar, los Objetivos Nacionales del PEDN al 2050 (Tabla~\ref{tab:pedn-on}) y, seguidamente, los ejes que estructuran la Visión del Perú al 2050 (Tabla~\ref{tab:vision-peru-2050-ejes}), como marco de referencia para el análisis. Finalmente, se incluye la escala de valoración que orienta la interpretación de los puntajes y la consistencia del dictamen (Tabla~\ref{tab:escala-valoracion-pedn}).\par}

\begin{table}[htbp]
\centering
\scriptsize
\caption{Objetivos Nacionales (ON) del Plan Estratégico de Desarrollo Nacional (PEDN) 2050.}
\label{tab:pedn-on}
\begin{tabularx}{\textwidth}{|C{0.8cm}|L{4.5cm}|X|}
\hline
\rowcolor{headerblue}
\thd{ON} & \thd{Denominación} & \thd{Propósito} \\ \hline
1 & Desarrollo de las personas & Alcanzar el pleno desarrollo de las capacidades de las personas sin dejar a nadie atrás. \\ \hline
\rowcolor{lightgray}
2 & Territorio sostenible & Gestionar el territorio de manera sostenible a fin de prevenir y reducir los riesgos y amenazas que afectan a las personas y sus medios de vida. \\ \hline
3 & Competitividad e innovación & Elevar los niveles de competitividad y productividad en el país. \\ \hline
\rowcolor{lightgray}
4 & Democracia y paz & Garantizar una sociedad justa, democrática, pacífica y un Estado efectivo al servicio de las personas. \\ \hline
\end{tabularx}
\end{table}

\begin{table}[htbp]
\centering
\scriptsize
\caption{Ejes de la Visión del Perú al 2050.}
\label{tab:vision-peru-2050-ejes}
\begin{tabularx}{\textwidth}{|C{0.8cm}|X|}
\hline
\rowcolor{headerblue}
\thd{Eje} & \thd{Descripción} \\ \hline
1 & Las personas alcanzan su potencial en igualdad de oportunidades y sin discriminación para gozar de una vida plena. \\ \hline
\rowcolor{lightgray}
2 & Gestión sostenible de la naturaleza y medidas frente al cambio climático. \\ \hline
3 & Desarrollo sostenible con empleo digno y en armonía con la naturaleza. \\ \hline
\rowcolor{lightgray}
4 & Sociedad democrática, pacífica, respetuosa de los derechos humanos y libre del temor y de la violencia. \\ \hline
5 & Estado moderno, eficiente, transparente y descentralizado que garantiza una sociedad justa e inclusiva, sin corrupción y sin dejar a nadie atrás. \\ \hline
\end{tabularx}
\end{table}

\begin{table}[htbp]
\centering
\scriptsize
\caption{Escala de valoración del alineamiento del proyecto con el PEDN al 2050.}
\label{tab:escala-valoracion-pedn}
\begin{tabularx}{\textwidth}{|C{2.8cm}|C{1cm}|X|}
\hline
\rowcolor{headerblue}
\thd{Nivel} & \thd{Pts.} & \thd{Descriptor operativo} \\ \hline
\textcolor{csuf}{\textbf{Excelente}} & \textbf{3} & El proyecto evidencia contribución explícita, fundamentada y verificable al Objetivo Nacional o eje evaluado. La articulación es orgánica dentro de la estructura argumentativa y metodológica del proyecto, con vinculación a objetivos específicos, retos estructurales o indicadores del PEDN. \\ \hline
\rowcolor{lightgray}
\textcolor{cace}{\textbf{Aceptable}} & \textbf{2} & La contribución es identificable pero parcial o implícita. El proyecto aborda el ON sin fundamentación suficiente o sin vincular objetivos específicos ni retos del PEDN. Susceptible de fortalecimiento. \\ \hline
\textcolor{cdef}{\textbf{Deficiente}} & \textbf{1} & El vínculo con el ON es tangencial o meramente declarativo. No hay sustento argumentativo ni coherencia con el diseño del proyecto. \\ \hline
\rowcolor{lightgray}
\textcolor{cins}{\textbf{Ausente}} & \textbf{0} & El proyecto no evidencia relación alguna con el ON o eje evaluado del PEDN al 2050. \\ \hline
\end{tabularx}
\end{table}

\clearpage
\subsection{Identificación y Fundamentación del Alineamiento con el PEDN al 2050}

\textit{Se evalúa si el proyecto identifica de manera explícita, justificada y coherente los Objetivos Nacionales, los objetivos específicos y los retos estructurales del PEDN al 2050 a los que contribuye, estableciendo la trazabilidad entre el problema de investigación y las metas de desarrollo nacional.} Los criterios para la evaluación se mencionan en la tabla \ref{tab:dimension1-pedn}. 

\begingroup
\scriptsize
\begin{longtable}{|C{0.55cm}|L{3.8cm}|L{5.8cm}|C{0.65cm}|C{0.65cm}|C{0.65cm}|C{0.65cm}|}
\caption{Rúbrica de evaluación: identificación y fundamentación del alineamiento con el PEDN al 2050 (Dimensión~I).}\label{tab:dimension1-pedn}\\
\hline
\rowcolor{headerblue}
\thd{N.\textsuperscript{o}} & \thd{Criterio} & \thd{Indicador verificable} & \thd{3} & \thd{2} & \thd{1} & \thd{0} \\
\hline\endfirsthead
\hline
\rowcolor{headerblue}
\thd{N.\textsuperscript{o}} & \thd{Criterio} & \thd{Indicador verificable} & \thd{3} & \thd{2} & \thd{1} & \thd{0} \\
\hline\endhead
1.1 & Identificación explícita de ON y objetivos específicos vinculados & El proyecto declara de manera explícita al menos un Objetivo Nacional y su(s) objetivo(s) específico(s) a los que contribuye. La identificación se basa en el análisis del problema, los objetivos y los resultados esperados, no en una mención genérica. Se cita la numeración y denominación oficial del PEDN. &
& & & \\ \hline
\rowcolor{lightgray}
1.2 & Vinculación con retos estructurales del PEDN & El proyecto identifica al menos uno de los 16~retos estructurales del PEDN (p.\,ej., Reto~4: insuficiente acceso a agua y saneamiento; Reto~11: deficiente infraestructura) y argumenta cómo la investigación contribuye a abordar dicho reto. El vínculo no es genérico sino operacionalizado con variables o resultados concretos. &
& & & \\ \hline
1.3 & Justificación de la pertinencia en el contexto local/regional & El proyecto argumenta por qué el ON seleccionado es pertinente para el contexto local, regional o nacional en el que se desarrolla la investigación. Se presentan datos, estadísticas o diagnósticos que evidencian la brecha o necesidad que el proyecto pretende atender en relación con las metas del PEDN al 2050. &
& & & \\ \hline
\rowcolor{lightgray}
1.4 & Coherencia entre el ON y el diseño de investigación & Existe correspondencia lógica entre el ON declarado y el planteamiento del problema, los objetivos, las variables y la metodología del proyecto. El ON no es un añadido retórico sino un eje articulador verificable en la estructura del proyecto. &
& & & \\ \hline
\multicolumn{3}{|r|}{\cellcolor{white}\textbf{Subtotal Dimensión~I (máx.~12 puntos):}} & \multicolumn{4}{c|}{\cellcolor{white}\textbf{/12}} \\ \hline
\end{longtable}
\endgroup

\clearpage
\subsection{ON~1: Desarrollo de las Personas (OE 1.1--1.6)}
\textit{Se evalúa la contribución del proyecto al pleno desarrollo de las capacidades de las personas: educación de calidad (OE~1.1), salud universal (OE~1.2), vivienda digna y saneamiento (OE~1.3), transporte seguro (OE~1.4), igualdad de oportunidades (OE~1.5) y poblaciones de frontera (OE~1.6). Retos estructurales asociados: 1--4.}. Para dicha evaluación, se presentan los criterios de la tabla \ref{tab:dimension2-on1-pedn}.

\begingroup
\scriptsize
\begin{longtable}{|C{0.55cm}|L{3.8cm}|L{5.8cm}|C{0.65cm}|C{0.65cm}|C{0.65cm}|C{0.65cm}|}
\caption{Rúbrica de evaluación: ON~1 -- desarrollo de las personas (Dimensión~II).}\label{tab:dimension2-on1-pedn}\\
\hline
\rowcolor{headerblue}
\thd{N.\textsuperscript{o}} & \thd{Criterio} & \thd{Indicador verificable} & \thd{3} & \thd{2} & \thd{1} & \thd{0} \\
\hline\endfirsthead
\hline
\rowcolor{headerblue}
\thd{N.\textsuperscript{o}} & \thd{Criterio} & \thd{Indicador verificable} & \thd{3} & \thd{2} & \thd{1} & \thd{0} \\
\hline\endhead
2.1 & Educación de calidad e inclusiva \onref{OE\,1.1} & El proyecto contribuye a la formación educativa de calidad, al desarrollo de competencias, al uso de tecnologías educativas o digitales, o al fortalecimiento de la investigación aplicada con impacto en la sostenibilidad y competitividad del país. Se identifica la vinculación con el Reto~2 (bajo rendimiento escolar). &
& & & \\ \hline
\rowcolor{lightgray}
2.2 & Salud, vivienda digna y servicios básicos \onref{OE\,1.2--1.3} & El proyecto contribuye a reducir la mortalidad o morbilidad, al acceso universal a la salud, a la salud digital, a la provisión de vivienda digna, accesible y resiliente, o al acceso a agua potable y saneamiento básico, especialmente en zonas rurales y periurbanas. Se identifica la vinculación con los Retos~3 y~4. &
& & & \\ \hline
2.3 & Transporte seguro y sostenible \onref{OE\,1.4} & El proyecto aborda la reducción del tiempo, costo o inseguridad en el traslado de personas en entornos urbanos o rurales, mediante sistemas de transporte seguros, accesibles, conectados, de calidad y con sostenibilidad ambiental y social. &
& & & \\ \hline
\rowcolor{lightgray}
2.4 & Igualdad de oportunidades e inclusión social \onref{OE\,1.5--1.6} & El proyecto garantiza la igualdad de oportunidades, la inclusión social de mujeres y grupos vulnerables, el respeto y valoración de la diversidad cultural, étnica y de género, o la mejora de calidad de vida de poblaciones de frontera. Se identifica la vinculación con el Reto~1 (pobreza persistente). &
& & & \\ \hline
\multicolumn{3}{|r|}{\cellcolor{white}\textbf{Subtotal Dimensión~II (máx.~12 puntos):}} & \multicolumn{4}{c|}{\cellcolor{white}\textbf{/12}} \\ \hline
\end{longtable}
\endgroup

\clearpage
\subsection{ON~2: Territorio Sostenible (OE 2.1--2.8)}

\textit{Se evalúa la contribución del proyecto a la gestión sostenible del territorio: ordenamiento territorial (OE~2.1), gestión del riesgo de desastres (OE~2.2), desarrollo urbano y rural (OE~2.3), diversidad biológica (OE~2.4), recursos hídricos (OE~2.5), calidad ambiental (OE~2.6), cambio climático (OE~2.7) y defensa nacional (OE~2.8). Retos estructurales asociados: 5--7.} Los criterios para esta evaluación se enuncian en la tabla \ref{tab:dimension3-on2-pedn}.

\begingroup
\scriptsize
\begin{longtable}{|C{0.55cm}|L{3.8cm}|L{5.8cm}|C{0.65cm}|C{0.65cm}|C{0.65cm}|C{0.65cm}|}
\caption{Rúbrica de evaluación: ON~2 -- territorio sostenible (Dimensión~III).}\label{tab:dimension3-on2-pedn}\\
\hline
\rowcolor{headerblue}
\thd{N.\textsuperscript{o}} & \thd{Criterio} & \thd{Indicador verificable} & \thd{3} & \thd{2} & \thd{1} & \thd{0} \\
\hline\endfirsthead
\hline
\rowcolor{headerblue}
\thd{N.\textsuperscript{o}} & \thd{Criterio} & \thd{Indicador verificable} & \thd{3} & \thd{2} & \thd{1} & \thd{0} \\
\hline\endhead
3.1 & Ordenamiento territorial y gestión del riesgo de desastres \onref{OE\,2.1--2.2} & El proyecto contribuye a la gestión del territorio con visión estratégica, al manejo sostenible de recursos naturales, a la reducción de la vulnerabilidad ante desastres, a la comprensión del riesgo, o a la mejora del uso y ocupación del territorio. Se identifica la vinculación con los Retos~5 y~7. &
& & & \\ \hline
\rowcolor{lightgray}
3.2 & Ciudades sostenibles y diversidad biológica \onref{OE\,2.3--2.4} & El proyecto promueve ciudades y centros poblados sostenibles, equitativos y competitivos, o la sostenibilidad de los servicios ecosistémicos mediante la gestión integrada de recursos naturales y ecosistemas. Se consideran acciones de planificación urbana, protección de biodiversidad y conservación de suelos. &
& & & \\ \hline
3.3 & Recursos hídricos y calidad ambiental \onref{OE\,2.5--2.6} & El proyecto aborda la disponibilidad, calidad y sostenibilidad de los recursos hídricos, la gestión de residuos sólidos, la calidad del aire, el tratamiento de aguas residuales, o el control y gobernanza ambiental, en un contexto de estrés hídrico y deterioro ambiental. &
& & & \\ \hline
\rowcolor{lightgray}
3.4 & Resiliencia ante el cambio climático y protección de intereses nacionales \onref{OE\,2.7--2.8} & El proyecto aumenta la resiliencia y adaptación ante el cambio climático, contribuye al tránsito hacia una economía baja en carbono, al monitoreo de fenómenos geológicos o hidroclimáticos, o a la defensa de los intereses nacionales destinados a la paz y seguridad internacional. Se identifica la vinculación con el Reto~6. &
& & & \\ \hline
\multicolumn{3}{|r|}{\cellcolor{white}\textbf{Subtotal Dimensión~III (máx.~12 puntos):}} & \multicolumn{4}{c|}{\cellcolor{white}\textbf{/12}} \\ \hline
\end{longtable}
\endgroup

\clearpage
\subsection{ON~3: Competitividad e Innovación (OE 3.1--3.7)}

\textit{Se evalúa la contribución del proyecto a elevar la competitividad y productividad del país: estabilidad macroeconómica (OE~3.1), empleo decente (OE~3.2), sectores productivos (OE~3.3), MIPYME (OE~3.4), ciencia, tecnología e innovación (OE~3.5), infraestructura nacional (OE~3.6) y competencia de mercado (OE~3.7). Retos estructurales asociados: 8--12.} Los criterios para esta evaluación se enuncian en la tabla \ref{tab:dimension4-on3-pedn}.

\begingroup
\scriptsize
\begin{longtable}{|C{0.55cm}|L{3.8cm}|L{5.8cm}|C{0.65cm}|C{0.65cm}|C{0.65cm}|C{0.65cm}|}
\caption{Rúbrica de evaluación: ON~3 -- competitividad e innovación (Dimensión~IV).}\label{tab:dimension4-on3-pedn}\\
\hline
\rowcolor{headerblue}
\thd{N.\textsuperscript{o}} & \thd{Criterio} & \thd{Indicador verificable} & \thd{3} & \thd{2} & \thd{1} & \thd{0} \\
\hline\endfirsthead
\hline
\rowcolor{headerblue}
\thd{N.\textsuperscript{o}} & \thd{Criterio} & \thd{Indicador verificable} & \thd{3} & \thd{2} & \thd{1} & \thd{0} \\
\hline\endhead
4.1 & Ciencia, tecnología, innovación y transformación digital \onref{OE\,3.5} & El proyecto eleva la capacidad científica y de innovación tecnológica del país, promueve la investigación, creación, adaptación o transferencia tecnológica, impulsa la transformación digital, o favorece la articulación academia--Estado--sector productivo. Se identifica la vinculación con el Reto~12 (escasas políticas de innovación). &
& & & \\ \hline
\rowcolor{lightgray}
4.2 & Infraestructura, conectividad y sectores productivos \onref{OE\,3.3, 3.6} & El proyecto contribuye a elevar la competitividad y productividad de los sectores económicos, a la diversificación productiva, a la generación de valor agregado, o a elevar la conectividad del país a través de infraestructura moderna, sostenible, resiliente y de calidad. Se identifica la vinculación con los Retos~10 y~11. &
& & & \\ \hline
4.3 & Empleo decente, formalización y MIPYME \onref{OE\,3.2, 3.4} & El proyecto tiene potencial para incrementar los niveles de empleo decente, productivo y formal, mejorar las condiciones laborales, fomentar el emprendimiento, la creatividad y la innovación en micro, pequeñas y medianas empresas, o facilitar el acceso a servicios financieros y entornos digitales. Se identifica la vinculación con los Retos~8 y~9. &
& & & \\ \hline
\rowcolor{lightgray}
4.4 & Competencia de mercado y estabilidad macroeconómica \onref{OE\,3.1, 3.7} & El proyecto contribuye a la estabilidad macroeconómica, a garantizar un mercado competitivo con regulación eficiente, a la protección de los derechos de los consumidores, o promueve condiciones para el ingreso libre de nuevos competidores, en el marco de una economía verde y baja en carbono. &
& & & \\ \hline
\multicolumn{3}{|r|}{\cellcolor{white}\textbf{Subtotal Dimensión~IV (máx.~12 puntos):}} & \multicolumn{4}{c|}{\cellcolor{white}\textbf{/12}} \\ \hline
\end{longtable}
\endgroup

\clearpage
\subsection{ON~4: Democracia y Paz; Coherencia Integral y Transversalidad}

\textit{Se evalúa la contribución del proyecto al ON~4 ---sistema político (OE~4.1), identidad cultural (OE~4.2), justicia (OE~4.3), seguridad institucional (OE~4.4), gestión pública (OE~4.5) y descentralización (OE~4.6)--- y la articulación transversal del proyecto con los lineamientos y la Visión del Perú al 2050. Retos asociados: 13--16.} Los criterios para esta evaluación se enuncian en la tabla \ref{tab:dimension5-on4-pedn}.

\begingroup
\scriptsize
\begin{longtable}{|C{0.55cm}|L{3.8cm}|L{5.8cm}|C{0.65cm}|C{0.65cm}|C{0.65cm}|C{0.65cm}|}
\caption{Rúbrica de evaluación: ON~4 -- democracia y paz; coherencia integral y transversalidad (Dimensión~V).}\label{tab:dimension5-on4-pedn}\\
\hline
\rowcolor{headerblue}
\thd{N.\textsuperscript{o}} & \thd{Criterio} & \thd{Indicador verificable} & \thd{3} & \thd{2} & \thd{1} & \thd{0} \\
\hline\endfirsthead
\hline
\rowcolor{headerblue}
\thd{N.\textsuperscript{o}} & \thd{Criterio} & \thd{Indicador verificable} & \thd{3} & \thd{2} & \thd{1} & \thd{0} \\
\hline\endhead
5.1 & Democracia, justicia, institucionalidad y gestión pública \onref{OE\,4.1--4.5} & El proyecto fortalece la institucionalidad democrática, el acceso a la justicia, la identidad y diversidad cultural, la efectividad de la gestión pública, la lucha contra la corrupción, o la transformación digital del Estado. Se identifica la vinculación con los Retos~13--16 (desempeño institucional, discriminación, corrupción, representación política). &
& & & \\ \hline
\rowcolor{lightgray}
5.2 & Descentralización, enfoque territorial y articulación multinivel \onref{OE\,4.6} & El proyecto contribuye al proceso de descentralización y ordenamiento territorial, adopta un enfoque territorial que considere las dinámicas y potencialidades del territorio, o promueve la articulación entre los tres niveles de gobierno (nacional, regional, local) para la implementación del PEDN. &
& & & \\ \hline
5.3 & Transversalidad del alineamiento PEDN en el proyecto & El alineamiento con el PEDN al 2050 no se limita a una sección aislada del proyecto (p.\,ej., solo la justificación), sino que permea de manera transversal el planteamiento del problema, los objetivos, el marco teórico, el diseño metodológico y los resultados esperados. Se verifican los enfoques de derechos humanos, igualdad de género, interculturalidad y perspectiva de discapacidad. &
& & & \\ \hline
\rowcolor{lightgray}
5.4 & Interconexión entre ON y potencial de impacto verificable & El proyecto reconoce y articula la interconexión entre dos o más Objetivos Nacionales, demostrando el enfoque sistémico del PEDN. Además, la contribución al PEDN es medible o verificable: se prevén indicadores de impacto, métricas de desempeño o resultados cuantificables vinculados a las metas del plan. Se valora el potencial de escalabilidad o replicabilidad. &
& & & \\ \hline
\multicolumn{3}{|r|}{\cellcolor{white}\textbf{Subtotal Dimensión~V (máx.~12 puntos):}} & \multicolumn{4}{c|}{\cellcolor{white}\textbf{/12}} \\ \hline
\end{longtable}
\endgroup

\clearpage
Tanto el cuadro Consolidado de Resultados como la Escala de Calificación de Alineamiento con el PEDN al 2050 se aprecian en las Tablas~\ref{tab:cuadro-consolidado-pedn} y~\ref{tab:escala-calificacion-pedn}, respectivamente.

\begin{table}[htbp]
\centering
\scriptsize
\caption{Cuadro consolidado de resultados por dimensión de evaluación (PEDN al 2050).}
\label{tab:cuadro-consolidado-pedn}
\begin{tabularx}{\textwidth}{|C{0.8cm}|L{7.5cm}|C{1.8cm}|C{1.8cm}|X|}
\hline
\rowcolor{headerblue}
\thd{Dim.} & \thd{Denominación} & \thd{Pts. Máx.} & \thd{Pts. Obt.} & \thd{Obs.} \\ \hline
I & Identificación y Fundamentación del Alineamiento con el PEDN al 2050 & 12 & & \\ \hline
\rowcolor{lightgray}
II & ON~1: Desarrollo de las Personas (OE 1.1--1.6) & 12 & & \\ \hline
III & ON~2: Territorio Sostenible (OE 2.1--2.8) & 12 & & \\ \hline
\rowcolor{lightgray}
IV & ON~3: Competitividad e Innovación (OE 3.1--3.7) & 12 & & \\ \hline
V & ON~4: Democracia y Paz; Coherencia Integral y Transversalidad & 12 & & \\ \hline
\rowcolor{white}
& \textbf{TOTAL GENERAL} & \textbf{60} & & \\ \hline
\end{tabularx}
\end{table}

\begin{table}[htbp]
\centering
\scriptsize
\caption{Escala de calificación del alineamiento del proyecto con el PEDN al 2050.}
\label{tab:escala-calificacion-pedn}
\begin{tabularx}{\textwidth}{|C{2.5cm}|C{3cm}|X|}
\hline
\rowcolor{headerblue}
\thd{Rango} & \thd{Calificación} & \thd{Significado} \\ \hline
\textbf{54--60}\newline(90--100\,\%) &
\textcolor{csuf}{\textbf{ALINEAMIENTO\newline PLENO}} &
El proyecto demuestra contribución integral, explícita y fundamentada al PEDN al 2050. La investigación constituye un aporte ejemplar al desarrollo nacional, con trazabilidad verificable a objetivos específicos, retos estructurales e indicadores del plan, y articulación transversal entre múltiples Objetivos Nacionales. \\ \hline
\rowcolor{lightgray}
\textbf{42--53}\newline(70--89\,\%) &
\textcolor{cace}{\textbf{ALINEAMIENTO\newline SUFICIENTE}} &
El proyecto evidencia contribución sustancial al PEDN al 2050, con áreas específicas susceptibles de fortalecimiento. Los vínculos principales con los Objetivos Nacionales están demostrados y la pertinencia para el desarrollo nacional es verificable. \\ \hline
\textbf{30--41}\newline(50--69\,\%) &
\textcolor{cdef}{\textbf{ALINEAMIENTO\newline PARCIAL}} &
La contribución es incompleta o predominantemente declarativa. El proyecto requiere fortalecimiento de la justificación y articulación con los objetivos específicos y retos estructurales del PEDN al 2050. \\ \hline
\rowcolor{lightgray}
\textbf{0--29}\newline($<$50\,\%) &
\textcolor{cins}{\textbf{ALINEAMIENTO\newline INSUFICIENTE}} &
El proyecto no demuestra contribución verificable al PEDN al 2050. La desconexión con los Objetivos Nacionales compromete la pertinencia de la investigación frente a las necesidades estratégicas del desarrollo sostenible del país. \\ \hline
\end{tabularx}
\end{table}


\clearpage
\section{Ajuste Epistemológico del Proyecto de Tesis}

{\footnotesize El ajuste epistemológico de un proyecto de tesis exige que exista una alineación lógica y verificable entre los supuestos filosóficos del investigador, los objetivos de la investigación y los métodos seleccionados para alcanzarlos. La epistemología constituye un fundamento científico indispensable que orienta la adquisición, validación y aplicación del conocimiento en la investigación de la Ingeniería Civil. Un proyecto epistemológicamente coherente asegura que los compromisos ontológicos sobre la naturaleza de los fenómenos investigados, la orientación epistemológica sobre la validez del conocimiento y las implementaciones metodológicas se articulen de forma integrada y justificable. La presente rúbrica permite al dictaminante evaluar si el proyecto de tesis presenta coherencia y justificación filosófico-metodológica, flexibilidad e innovación, identificación de controversias y limitaciones, y comunicación científica mejorada en el contexto de las especialidades de la Ingeniería Civil: estabilización de suelos, pavimentos, cimentaciones, mecánica de materiales, investigación de nuevos materiales, entre otras. Las dimensiones de evaluación epistemológica y su escala de valoración se aprecian en las Tablas~\ref{tab:dim-epistemologica} y~\ref{tab:escala-epistemologica}. Se toma como referencia de fundamentación el Capítulo 2 del libro \enquote{Metodología de la Investigación para la Estabilización Geotécnica de Suelos} (Arbulú \& Solorzano, 2026), con énfasis en la sección \textit{Importancia del Ajuste Epistemológico en Investigación Geotécnica} (subsecciones: Coherencia y Justificación, Flexibilidad e Innovación, Identificación de Controversias y Limitaciones, y Comunicación Mejorada).\par} 

\begin{table}[htbp]
\centering
\scriptsize
\caption{Dimensiones de evaluación del ajuste epistemológico del proyecto de tesis.}
\label{tab:dim-epistemologica}
\begin{tabularx}{\textwidth}{|C{0.8cm}|L{4.5cm}|X|}
\hline
\rowcolor{headerblue}
\thd{Dim.} & \thd{Denominación} & \thd{Alcance} \\ \hline
I & Coherencia ontológica y pertinencia del campo & Verifica que la posición ontológica (objetivista, constructivista o mixta) sea consistente con la naturaleza del problema de Ingeniería Civil abordado y con las especialidades de la carrera. \\ \hline
\rowcolor{lightgray}
II & Orientación epistemológica y validez del conocimiento & Evalúa si la postura epistemológica adoptada (positivista, pragmática, interpretativista) es apropiada para el tipo de conocimiento que el proyecto pretende generar y si los criterios de validez son explícitos. \\ \hline
III & Coherencia metodológica y justificación del diseño & Examina la alineación entre los supuestos filosóficos declarados, el tipo de razonamiento empleado (deductivo, inductivo, abductivo, retroductivo) y los métodos de investigación seleccionados. \\ \hline
\rowcolor{lightgray}
IV & Axiología, valores e innovación en la investigación & Evalúa si el proyecto reconoce explícitamente los valores que orientan la investigación (sostenibilidad, seguridad, eficiencia) y si el marco filosófico permite flexibilidad e innovación metodológica. \\ \hline
V & Identificación de limitaciones y comunicación científica & Verifica si el proyecto identifica controversias epistemológicas, reconoce limitaciones del enfoque adoptado y comunica de manera explícita su posición paradigmática para facilitar la colaboración interdisciplinaria. \\ \hline
\end{tabularx}
\end{table}

\begin{table}[htbp]
\centering
\scriptsize
\caption{Escala de valoración del ajuste epistemológico del proyecto de tesis.}
\label{tab:escala-epistemologica}
\begin{tabularx}{\textwidth}{|C{2.8cm}|C{1cm}|X|}
\hline
\rowcolor{headerblue}
\thd{Nivel} & \thd{Pts.} & \thd{Descriptor operativo} \\ \hline
\textcolor{epgreen}{\textbf{Excelente}} & \textbf{3} & El proyecto evidencia ajuste epistemológico explícito, fundamentado y verificable en el criterio evaluado. La articulación entre los supuestos filosóficos, los objetivos y la metodología es orgánica, coherente y plenamente justificada dentro de la estructura del proyecto. \\ \hline
\rowcolor{lightgray}
\textcolor{epblue}{\textbf{Aceptable}} & \textbf{2} & El ajuste epistemológico es identificable pero parcial o implícito. El proyecto aborda el criterio sin fundamentación suficiente o sin articular de manera explícita la conexión entre sus componentes filosóficos y metodológicos. Susceptible de fortalecimiento. \\ \hline
\textcolor{eporange}{\textbf{Deficiente}} & \textbf{1} & El vínculo epistemológico es tangencial o meramente declarativo. No hay sustento argumentativo ni coherencia verificable entre la posición filosófica y el diseño de la investigación. \\ \hline
\rowcolor{lightgray}
\textcolor{epred}{\textbf{Ausente}} & \textbf{0} & El proyecto no evidencia consideración alguna del criterio evaluado. Existe desconexión total entre los fundamentos filosóficos y la implementación de la investigación. \\ \hline
\end{tabularx}
\end{table}

\clearpage
\subsection{Coherencia Ontológica y Pertinencia del Ámbito}

\textit{Se evalúa si el proyecto adopta una posición ontológica (objetivista, constructivista o mixta) consistente con la naturaleza del problema de Ingeniería Civil investigado y si la temática se inscribe dentro de las especialidades propias de la carrera: estabilización de suelos, pavimentos, cimentaciones, mecánica de materiales, investigación de nuevos materiales, entre otras.} Los criterios para esta evaluación se presentan en la Tabla~\ref{tab:coherencia-ontologica}.

\begingroup
\scriptsize
\begin{longtable}{|C{0.55cm}|L{3.8cm}|L{5.8cm}|C{0.65cm}|C{0.65cm}|C{0.65cm}|C{0.65cm}|}
\caption{Rúbrica de evaluación: coherencia ontológica y pertinencia del área.}\label{tab:coherencia-ontologica}\\
\hline
\rowcolor{headerblue}
\thd{N.\textsuperscript{o}} & \thd{Criterio} & \thd{Indicador verificable} & \thd{3} & \thd{2} & \thd{1} & \thd{0} \\
\hline\endfirsthead
\hline
\rowcolor{headerblue}
\thd{N.\textsuperscript{o}} & \thd{Criterio} & \thd{Indicador verificable} & \thd{3} & \thd{2} & \thd{1} & \thd{0} \\
\hline\endhead
1.1 & Pertinencia del problema dentro de la Ingeniería Civil & El problema de investigación se inscribe de manera verificable dentro de las áreas disciplinares de la Ingeniería Civil (estabilización de suelos, pavimentos, cimentaciones, mecánica de materiales, investigación de nuevos materiales, geotecnia computacional, mitigación de peligros, sostenibilidad de infraestructura, entre otras). El planteamiento no es genérico ni ajeno a la carrera. &
& & & \\ \hline
\rowcolor{lightgray}
1.2 & Posición ontológica coherente con el objeto de estudio & El proyecto asume, explícita o implícitamente, una posición ontológica (objetivista para propiedades cuantificables de geomateriales; constructivista para dimensiones socioambientales; o mixta para problemas interdisciplinares) que es congruente con la naturaleza del fenómeno investigado. &
& & & \\ \hline
1.3 & Supuestos sobre la realidad investigada & Los supuestos sobre la existencia independiente y mensurabilidad de las propiedades del suelo, roca, material o sistema estructural son coherentes con los procedimientos experimentales propuestos. Si el proyecto incluye dimensiones sociales (percepción de riesgo, sostenibilidad), se reconoce la construcción social de dichos fenómenos. &
& & & \\ \hline
\rowcolor{lightgray}
1.4 & Delimitación ontológica del alcance de la investigación & Los límites del estudio (qué se considera real, medible y cognoscible) están claramente delimitados y son consistentes con la posición ontológica adoptada. El proyecto no pretende cuantificar fenómenos inherentemente cualitativos sin justificación ni viceversa. &
& & & \\ \hline
\multicolumn{3}{|r|}{\cellcolor{white}\textbf{Subtotal Dimensión~I (máx.~12 puntos):}} & \multicolumn{4}{c|}{\cellcolor{white}\textbf{/12}} \\ \hline
\end{longtable}
\endgroup

\clearpage
\subsection{Orientación Epistemológica y Validez del Conocimiento}

\textit{Se evalúa si la postura epistemológica adoptada por el proyecto (positivista-realista, pragmática o interpretativista) es apropiada para el tipo de conocimiento que se pretende generar, y si los criterios de validez, confiabilidad y generalización son explícitos y consistentes con dicha postura.} Los criterios para esta evaluación se presentan en la Tabla~\ref{tab:validez-epistemologica}.

\begingroup
\scriptsize
\begin{longtable}{|C{0.55cm}|L{3.8cm}|L{5.8cm}|C{0.65cm}|C{0.65cm}|C{0.65cm}|C{0.65cm}|}
\caption{Rúbrica de evaluación: orientación epistemológica y validez del conocimiento.}\label{tab:validez-epistemologica}\\
\hline
\rowcolor{headerblue}
\thd{N.\textsuperscript{o}} & \thd{Criterio} & \thd{Indicador verificable} & \thd{3} & \thd{2} & \thd{1} & \thd{0} \\
\hline\endfirsthead
\hline
\rowcolor{headerblue}
\thd{N.\textsuperscript{o}} & \thd{Criterio} & \thd{Indicador verificable} & \thd{3} & \thd{2} & \thd{1} & \thd{0} \\
\hline\endhead
2.1 & Posición epistemológica apropiada al problema \epref{Positivismo / Pragmatismo / Interpretativismo} & El proyecto adopta una orientación epistemológica coherente con la naturaleza de su problema: positivista-realista para la caracterización empírica de materiales y comportamiento mecánico; pragmática para problemas aplicados que requieren flexibilidad metodológica; interpretativista para dimensiones sociotécnicas. La elección es justificable y no arbitraria. &
& & & \\ \hline
\rowcolor{lightgray}
2.2 & Criterios explícitos de validez del conocimiento & El proyecto establece, explícita o implícitamente, los criterios mediante los cuales se determinará la validez de los resultados: reproducibilidad experimental, significancia estadística, triangulación metodológica, saturación teórica, o coherencia interpretativa, según corresponda a la postura epistemológica adoptada. &
& & & \\ \hline
2.3 & Admisibilidad de la evidencia empírica & Los tipos de datos admitidos como evidencia (datos de laboratorio, datos de campo, mediciones geofísicas, simulaciones numéricas, percepciones de actores, juicio experto) son consistentes con la posición epistemológica. El proyecto no admite como evidencia datos cuya naturaleza contradice los supuestos epistemológicos declarados. &
& & & \\ \hline
\rowcolor{lightgray}
2.4 & Nivel de generalización coherente con el diseño & Las pretensiones de generalización de los resultados (universalidad, transferibilidad o particularidad) son proporcionales al alcance del diseño de investigación y a la postura epistemológica. El proyecto no extrapola conclusiones más allá de lo que la base empírica y el enfoque epistemológico permiten justificar. &
& & & \\ \hline
\multicolumn{3}{|r|}{\cellcolor{white}\textbf{Subtotal Dimensión~II (máx.~12 puntos):}} & \multicolumn{4}{c|}{\cellcolor{white}\textbf{/12}} \\ \hline
\end{longtable}
\endgroup

\clearpage
\subsection{Coherencia Metodológica y Justificación del Diseño}

\textit{Se evalúa la alineación entre los supuestos filosóficos, el tipo de preguntas de investigación (descriptivas, relacionales, causales, normativas), el enfoque de datos (cuantitativo, cualitativo, mixto), el razonamiento lógico empleado (deductivo, inductivo, abductivo, retroductivo) y los métodos e instrumentos seleccionados.} Los criterios para esta evaluación se presentan en la Tabla~\ref{tab:coherencia-metodologica}.

\begingroup
\scriptsize
\begin{longtable}{|C{0.55cm}|L{3.8cm}|L{5.8cm}|C{0.65cm}|C{0.65cm}|C{0.65cm}|C{0.65cm}|}
\caption{Rúbrica de evaluación: coherencia metodológica y justificación del diseño.}\label{tab:coherencia-metodologica}\\
\hline
\rowcolor{headerblue}
\thd{N.\textsuperscript{o}} & \thd{Criterio} & \thd{Indicador verificable} & \thd{3} & \thd{2} & \thd{1} & \thd{0} \\
\hline\endfirsthead
\hline
\rowcolor{headerblue}
\thd{N.\textsuperscript{o}} & \thd{Criterio} & \thd{Indicador verificable} & \thd{3} & \thd{2} & \thd{1} & \thd{0} \\
\hline\endhead
3.1 & Tipo de preguntas de investigación coherente con el marco filosófico & Las preguntas de investigación (descriptivas para caracterización de materiales, relacionales/causales para mecanismos de comportamiento, normativas para optimización de diseños) son apropiadas para la posición ontológica y epistemológica adoptada. No existe disonancia entre lo que se pregunta y lo que los supuestos filosóficos permiten conocer. &
& & & \\ \hline
\rowcolor{lightgray}
3.2 & Enfoque de datos alineado con la epistemología \epref{Cuantitativo / Cualitativo / Mixto} & El enfoque de recolección y análisis de datos (cuantitativo mediante ensayos estandarizados y modelización numérica; cualitativo mediante estudios de caso o consulta a actores; mixto para investigaciones interdisciplinares) es congruente con la postura epistemológica y con el tipo de evidencia que el proyecto requiere para validar sus conclusiones. &
& & & \\ \hline
3.3 & Razonamiento lógico apropiado y justificado \epref{Deductivo / Inductivo / Abductivo / Retroductivo} & El modo de razonamiento empleado es verificable: deductivo al contrastar teorías establecidas de mecánica de suelos con predicciones específicas; inductivo al desarrollar modelos a partir de datos experimentales; abductivo al inferir explicaciones plausibles de fenómenos complejos; retroductivo al identificar mecanismos causales subyacentes. El razonamiento seleccionado es consistente con el estado del conocimiento previo y la naturaleza del problema. &
& & & \\ \hline
\rowcolor{lightgray}
3.4 & Métodos e instrumentos consistentes con el marco jerárquico & Los métodos específicos seleccionados (ensayos SPT, CPT, triaxial, corte directo, prospecciones geofísicas, simulaciones MEF/MED, RSM, LCA/LCCA, consulta a actores, juicio experto, entre otros) se derivan lógicamente de la cadena ontología--epistemología--enfoque--razonamiento y no son seleccionados de manera arbitraria o por conveniencia sin justificación epistemológica. &
& & & \\ \hline
\multicolumn{3}{|r|}{\cellcolor{white}\textbf{Subtotal Dimensión~III (máx.~12 puntos):}} & \multicolumn{4}{c|}{\cellcolor{white}\textbf{/12}} \\ \hline
\end{longtable}
\endgroup

\clearpage
\subsection{Axiología, Valores e Innovación en la Investigación}
\textit{Se evalúa si el proyecto reconoce explícitamente los valores que orientan la investigación (sostenibilidad, seguridad, eficiencia, equidad intergeneracional), si la postura axiológica es coherente con el diseño metodológico y si el marco filosófico permite flexibilidad e innovación para abordar problemas emergentes de la Ingeniería Civil.} Los criterios para esta evaluación se presentan en la Tabla~\ref{tab:axiologia-valores}.

\begingroup
\scriptsize
\begin{longtable}{|C{0.55cm}|L{3.8cm}|L{5.8cm}|C{0.65cm}|C{0.65cm}|C{0.65cm}|C{0.65cm}|}
\caption{Rúbrica de evaluación: axiología, valores e innovación en la investigación.}\label{tab:axiologia-valores}\\
\hline
\rowcolor{headerblue}
\thd{N.\textsuperscript{o}} & \thd{Criterio} & \thd{Indicador verificable} & \thd{3} & \thd{2} & \thd{1} & \thd{0} \\
\hline\endfirsthead
\hline
\rowcolor{headerblue}
\thd{N.\textsuperscript{o}} & \thd{Criterio} & \thd{Indicador verificable} & \thd{3} & \thd{2} & \thd{1} & \thd{0} \\
\hline\endhead
4.1 & Reconocimiento explícito de los valores que orientan la investigación & El proyecto identifica los valores que motivan y orientan la investigación: sostenibilidad ambiental, seguridad estructural, protección de la vida humana, eficiencia económica, equidad intergeneracional, conservación de recursos o responsabilidad geoética. La posición axiológica no permanece oculta ni se asume como neutral sin justificación. &
& & & \\ \hline
\rowcolor{lightgray}
4.2 & Coherencia entre los valores declarados y las capacidades analíticas del diseño & Los compromisos valorativos se traducen en capacidades analíticas tangibles dentro del diseño de la investigación. Si el proyecto declara valores de sostenibilidad, el diseño incluye herramientas capaces de evaluar dimensiones ambientales y sociales (LCA, LCCA, consulta a actores), no solo métricas de desempeño técnico convencional. &
& & & \\ \hline
4.3 & Flexibilidad metodológica para abordar la especificidad del problema & El marco filosófico adoptado permite adaptabilidad metodológica en función de las características específicas del problema, las condiciones geológicas locales, las restricciones contextuales y los recursos disponibles, en lugar de conformarse a protocolos prescritos de manera rígida e injustificada. &
& & & \\ \hline
\rowcolor{lightgray}
4.4 & Integración de innovación o tecnologías emergentes con justificación epistemológica & Si el proyecto incorpora tecnologías emergentes (inteligencia artificial, materiales a nanoescala, métodos computacionales avanzados, materiales reciclados o alternativos), su integración se justifica epistemológicamente dentro del marco filosófico adoptado y no constituye un añadido superficial sin articulación con los supuestos de la investigación. &
& & & \\ \hline
\multicolumn{3}{|r|}{\cellcolor{white}\textbf{Subtotal Dimensión~IV (máx.~12 puntos):}} & \multicolumn{4}{c|}{\cellcolor{white}\textbf{/12}} \\ \hline
\end{longtable}
\endgroup

\clearpage
\subsection{Identificación de Limitaciones y Comunicación Científica}
\textit{Se evalúa si el proyecto identifica las controversias epistemológicas y las limitaciones inherentes al enfoque adoptado, si reconoce las tensiones entre dimensiones técnicas y sociocontextuales, y si comunica de manera explícita su posición paradigmática para facilitar la evaluación crítica y la colaboración interdisciplinaria.} Los criterios para esta evaluación se presentan en la Tabla~\ref{tab:limitaciones-comunicacion}.

\begingroup
\scriptsize
\begin{longtable}{|C{0.55cm}|L{3.8cm}|L{5.8cm}|C{0.65cm}|C{0.65cm}|C{0.65cm}|C{0.65cm}|}
\caption{Rúbrica de evaluación: identificación de limitaciones y comunicación científica.}\label{tab:limitaciones-comunicacion}\\
\hline
\rowcolor{headerblue}
\thd{N.\textsuperscript{o}} & \thd{Criterio} & \thd{Indicador verificable} & \thd{3} & \thd{2} & \thd{1} & \thd{0} \\
\hline\endfirsthead
\hline
\rowcolor{headerblue}
\thd{N.\textsuperscript{o}} & \thd{Criterio} & \thd{Indicador verificable} & \thd{3} & \thd{2} & \thd{1} & \thd{0} \\
\hline\endhead
5.1 & Identificación de controversias epistemológicas & El proyecto reconoce las tensiones potenciales entre los componentes de su paradigma de investigación: limitaciones de un enfoque estrictamente objetivista en proyectos con implicaciones sociales, insuficiencia de métodos puramente cuantitativos para capturar dimensiones cualitativas del problema, o restricciones de generalización derivadas del diseño. &
& & & \\ \hline
\rowcolor{lightgray}
5.2 & Reconocimiento explícito de las limitaciones del enfoque adoptado & El proyecto presenta una autocrítica epistemológica: identifica qué aspectos del problema no puede abordar con el enfoque elegido, qué dimensiones quedan fuera del alcance de la investigación por restricciones filosóficas o metodológicas, y qué supuestos podrían ser cuestionables. La estrategia resultante es más sofisticada y no pretende una completitud que el diseño no puede sustentar. &
& & & \\ \hline
5.3 & Comunicación explícita de la posición paradigmática & El proyecto explicita, al menos de forma implícita verificable, la posición paradigmática subyacente (ontología, epistemología, axiología, metodología), facilitando la evaluación crítica por parte de dictaminantes, asesores y pares. Esta claridad es particularmente relevante en proyectos interdisciplinares que involucran profesionales de tradiciones filosóficas divergentes. &
& & & \\ \hline
\rowcolor{lightgray}
5.4 & Coherencia integral a lo largo de todas las secciones del proyecto & El ajuste epistemológico no se limita a una sección aislada (p.\,ej., solo el marco teórico o la metodología), sino que permea de manera transversal el planteamiento del problema, los objetivos, el marco teórico, el diseño metodológico, los resultados esperados y las conclusiones anticipadas. La coherencia filosófico-metodológica es un eje articulador verificable en la estructura completa del proyecto. &
& & & \\ \hline
\multicolumn{3}{|r|}{\cellcolor{white}\textbf{Subtotal Dimensión~V (máx.~12 puntos):}} & \multicolumn{4}{c|}{\cellcolor{white}\textbf{/12}} \\ \hline
\end{longtable}
\endgroup

\newpage
El cuadro consolidado de resultados y la escala de calificación del ajuste epistemológico se presentan en las Tablas~\ref{tab:cuadro-consolidado-epistemologia} y~\ref{tab:escala-calificacion-epistemologia}.

\begin{table}[htbp]
\centering
\scriptsize
\caption{Cuadro consolidado de resultados por dimensión de evaluación epistemológica.}
\label{tab:cuadro-consolidado-epistemologia}
\begin{tabularx}{\textwidth}{|C{0.8cm}|L{7.5cm}|C{1.8cm}|C{1.8cm}|X|}
\hline
\rowcolor{headerblue}
\thd{Dim.} & \thd{Denominación} & \thd{Pts. Máx.} & \thd{Pts. Obt.} & \thd{Obs.} \\ \hline
I & Coherencia Ontológica y Pertinencia para el Campo de Conocimiento & 12 & & \\ \hline
\rowcolor{lightgray}
II & Orientación Epistemológica y Validez del Conocimiento & 12 & & \\ \hline
III & Coherencia Metodológica y Justificación del Diseño & 12 & & \\ \hline
\rowcolor{lightgray}
IV & Axiología, Valores e Innovación en la Investigación & 12 & & \\ \hline
V & Identificación de Limitaciones y Comunicación Científica & 12 & & \\ \hline
\rowcolor{white}
& \textbf{TOTAL GENERAL} & \textbf{60} & & \\ \hline
\end{tabularx}
\end{table}

\begin{table}[htbp]
\centering
\scriptsize
\caption{Escala de calificación del ajuste epistemológico del proyecto.}
\label{tab:escala-calificacion-epistemologia}
\begin{tabularx}{\textwidth}{|C{2.5cm}|C{3cm}|X|}
\hline
\rowcolor{headerblue}
\thd{Rango} & \thd{Calificación} & \thd{Significado} \\ \hline
\textbf{54--60}\newline(90--100\,\%) &
\textcolor{epgreen}{\textbf{AJUSTE\newline PLENO}} &
El proyecto demuestra coherencia epistemológica integral entre su posición ontológica, orientación epistemológica, postura axiológica, enfoque de investigación y diseño metodológico. Las relaciones lógicas entre los fundamentos filosóficos y las implementaciones metodológicas son explícitas, justificadas y verificables en todas las secciones del proyecto. \\ \hline
\rowcolor{lightgray}
\textbf{42--53}\newline(70--89\,\%) &
\textcolor{epblue}{\textbf{AJUSTE\newline SUFICIENTE}} &
El proyecto evidencia coherencia epistemológica sustancial, con áreas específicas susceptibles de fortalecimiento. Los vínculos principales entre filosofía y metodología están demostrados, aunque algunos componentes del marco jerárquico permanecen implícitos o requieren mayor articulación. \\ \hline
\textbf{30--41}\newline(50--69\,\%) &
\textcolor{eporange}{\textbf{AJUSTE\newline PARCIAL}} &
La coherencia epistemológica es incompleta o predominantemente declarativa. Se evidencian desconexiones significativas entre los supuestos filosóficos y las decisiones metodológicas. El proyecto requiere fortalecimiento sustancial de la justificación filosófico-metodológica. \\ \hline
\rowcolor{lightgray}
\textbf{0--29}\newline($<$50\,\%) &
\textcolor{epred}{\textbf{AJUSTE\newline INSUFICIENTE}} &
El proyecto no demuestra coherencia epistemológica verificable. Existe desconexión entre los fundamentos filosóficos y la implementación de la investigación, comprometiendo la credibilidad, la justificación metodológica y la solidez conceptual del estudio. Se requiere una reestructuración fundamental del diseño de investigación. \\ \hline
\end{tabularx}
\end{table}

\clearpage
\section{Evaluación de la Formulación del Problema de Investigación}

{\footnotesize
La problematización es el proceso intelectual mediante el cual el tesista identifica, delimita, contextualiza y estructura un problema de investigación, distinguiéndolo de uno de aplicación profesional y dándole sustento empírico y lógico previo a la hipótesis, la justificación y las preguntas de investigación. En Ingeniería Civil, exige rigor técnico en la descripción de fenómenos geotécnicos, estructurales, viales o de materiales, y sigue la secuencia lógica \textbf{Diagnóstico}~$\rightarrow$~\textbf{Pronóstico}~$\rightarrow$~\textbf{Control del Pronóstico}. El \textit{diagnóstico} describe la situación problemática actual con evidencia verificable; el \textit{pronóstico} proyecta las consecuencias de no intervenir; y el \textit{control del pronóstico} anticipa la estrategia de investigación para abordar el problema. Esta rúbrica permite evaluar la calidad, consistencia y completitud de la problematización, verificando que el planteamiento sea técnicamente preciso, no ambiguo, debidamente sustentado y lógicamente articulado como antecedente de la hipótesis, la justificación y la formulación del problema. La secuencia Diagnóstico--Pronóstico--Control del Pronóstico es una estructura lógico-argumental habitual en la tradición académica latinoamericana para formular el planteamiento del problema en proyectos de investigación de pregrado en Ingeniería Civil; su evaluación mediante esta rúbrica se centra en la consistencia interna, la precisión técnica y la suficiencia argumental como antecedentes de la hipótesis, la justificación y la formulación del problema de investigación.\par
}
Las dimensiones de evaluación de la problematización se resumen en la Tabla~\ref{tab:dim-problematizacion}, mientras que la escala de valoración y sus descriptores operativos se presentan en la Tabla~\ref{tab:escala-problematizacion}.

\begin{table}[htbp]
\centering
\scriptsize
\caption{Dimensiones de evaluación de la problematización y alcance de cada dimensión.}
\label{tab:dim-problematizacion}
\begin{tabularx}{\textwidth}{|C{0.8cm}|L{4.5cm}|X|}
\hline
\rowcolor{headerblue}
\thd{Dim.} & \thd{Denominación} & \thd{Alcance} \\ \hline
I & Identificación y delimitación del problema & Verifica la existencia de un problema de conocimiento claramente identificado, la delimitación espacial, temporal, temática y poblacional, la consistencia con el título del proyecto y la precisión de la unidad de análisis. \\ \hline
\rowcolor{lightgray}
II & Contextualización y evidencia del problema & Evalúa el sustento empírico, estadístico o bibliográfico del problema, la ausencia de afirmaciones sin respaldo, la coherencia causal con el abordaje metodológico y la precisión técnica del diagnóstico. \\ \hline
III & Diagnóstico técnico de la situación problemática & Examina la rigurosidad del diagnóstico: descripción verificable del estado actual del fenómeno, cuantificación de magnitudes y brechas, identificación de causas y factores, y coherencia con el dominio profesional de la Ingeniería Civil. \\ \hline
\rowcolor{lightgray}
IV & Pronóstico y consecuencias de la no intervención & Evalúa la proyección de consecuencias técnicas, estructurales, económicas o ambientales de no abordar el problema, la plausibilidad de los escenarios proyectados y la trazabilidad lógica desde el diagnóstico. \\ \hline
V & Control del pronóstico y articulación pre-hipótesis & Verifica la anticipación de la estrategia de investigación que permitirá abordar el problema, la coherencia interna de la cadena argumental diagnóstico--pronóstico--control, y la suficiencia de la problematización como antecedente de la hipótesis, la justificación y la pregunta de investigación. \\ \hline
\end{tabularx}
\end{table}

\begin{table}[htbp]
\centering
\scriptsize
\caption{Escala de valoración (puntaje) y descriptor operativo para la evaluación de la problematización.}
\label{tab:escala-problematizacion}
\begin{tabularx}{\textwidth}{|C{2.8cm}|C{1cm}|X|}
\hline
\rowcolor{headerblue}
\thd{Nivel} & \thd{Pts.} & \thd{Descriptor operativo} \\ \hline
\textcolor{csuf}{\textbf{Excelente}} & \textbf{3} & El proyecto satisface plenamente el criterio evaluado. La problematización es explícita, técnicamente precisa, debidamente sustentada con evidencia verificable y lógicamente articulada. No se identifican ambigüedades ni vacíos argumentativos. \\ \hline
\rowcolor{lightgray}
\textcolor{cace}{\textbf{Aceptable}} & \textbf{2} & El criterio es identificable pero se cumple de manera parcial o incompleta. La problematización aborda el aspecto evaluado sin la profundidad, la precisión o el sustento suficiente. Susceptible de fortalecimiento con correcciones específicas. \\ \hline
\textcolor{cdef}{\textbf{Deficiente}} & \textbf{1} & El cumplimiento del criterio es tangencial o meramente declarativo. La argumentación es superficial, carece de sustento verificable o presenta inconsistencias lógicas que debilitan la estructura de la problematización. \\ \hline
\rowcolor{lightgray}
\textcolor{cins}{\textbf{Ausente}} & \textbf{0} & El proyecto no evidencia consideración alguna del criterio evaluado. Existe omisión total del aspecto evaluado o una desconexión que invalida la problematización en ese componente. \\ \hline
\end{tabularx}
\end{table}

\clearpage
\subsection{Identificación y Delimitación del Problema}

\textit{Se evalúa si el proyecto identifica un problema de conocimiento genuino (diferenciado de un problema de aplicación profesional), si la delimitación es precisa y si existe consistencia entre el título, el problema formulado y la unidad de análisis. Esta dimensión verifica los requisitos mínimos de formulación antes de evaluar el sustento y la articulación argumental.} Véase la Tabla~\ref{tab:rubrica-ident-delim}.

{\scriptsize
\begin{longtable}{|C{0.55cm}|L{3.8cm}|L{5.8cm}|C{0.65cm}|C{0.65cm}|C{0.65cm}|C{0.65cm}|}
\caption{Rúbrica de evaluación de la Dimensión~I: identificación y delimitación del problema.}\label{tab:rubrica-ident-delim}\\
\hline
\rowcolor{headerblue}
\thd{N.\textsuperscript{o}} & \thd{Criterio} & \thd{Indicador verificable} & \thd{3} & \thd{2} & \thd{1} & \thd{0} \\
\hline\endfirsthead
\hline
\rowcolor{headerblue}
\thd{N.\textsuperscript{o}} & \thd{Criterio} & \thd{Indicador verificable} & \thd{3} & \thd{2} & \thd{1} & \thd{0} \\
\hline\endhead
1.1 & Existencia de un problema de conocimiento claramente identificado \epref{Problema de conocimiento vs.\ problema de aplicación} & El proyecto formula un problema de conocimiento genuino: una brecha, vacío, contradicción o insuficiencia en el estado del arte que requiere investigación para ser resuelta. El problema no se reduce a una necesidad de aplicación profesional (p.\,ej., \enquote{diseñar un pavimento} no es un problema de conocimiento; \enquote{determinar el efecto de un aditivo polimérico sobre la resistencia a la fatiga de mezclas asfálticas en condiciones de altura} sí lo es). Se identifica con claridad qué se desconoce y por qué es necesario investigarlo. &
& & & \\ \hline
\rowcolor{lightgray}
1.2 & Delimitación espacial, temporal, temática y poblacional del problema & El problema está delimitado en las cuatro dimensiones pertinentes: \textbf{espacial} (localización geográfica, tramo vial, zona geotécnica, sitio de estudio); \textbf{temporal} (periodo de análisis, horizonte de datos, vigencia del estudio); \textbf{temática} (subárea específica dentro de la Ingeniería Civil: estabilización, pavimentos, cimentaciones, mecánica de materiales, nuevos materiales, etc.); y \textbf{poblacional o de unidad muestral} (tipo de suelo, tipo de estructura, material específico, muestras, probetas). La delimitación es operativa y no meramente enunciativa. &
& & & \\ \hline
1.3 & Consistencia entre el título del proyecto y el problema formulado & Existe correspondencia biunívoca verificable entre el título del proyecto y el problema de investigación: las variables, el alcance, la delimitación y el enfoque expresados en el título se reflejan fielmente en la descripción del problema. No hay variables en el título ausentes en el problema ni viceversa; no hay discrepancias de alcance, de unidad de análisis o de enfoque entre ambos elementos. &
& & & \\ \hline
\rowcolor{lightgray}
1.4 & Precisión y ausencia de ambigüedades en la unidad de análisis & La unidad de análisis (material, suelo, estructura, mezcla, probeta, tramo, sistema) está definida sin ambigüedades: se especifica su naturaleza, clasificación técnica, procedencia y condición. No se emplean términos genéricos que admitan interpretaciones múltiples (p.\,ej., \enquote{suelo de la zona} sin clasificación SUCS/AASHTO; \enquote{material reciclado} sin especificación del tipo; \enquote{mezcla mejorada} sin definir el agente estabilizador ni las dosificaciones). &
& & & \\ \hline
\multicolumn{3}{|r|}{\cellcolor{white}\textbf{Subtotal Dimensión~I (máx.~12 puntos):}} & \multicolumn{4}{c|}{\cellcolor{white}\textbf{/12}} \\ \hline
\end{longtable}
}

\clearpage
\subsection{Contextualización y Evidencia del Problema}

\textit{Se evalúa si la descripción del problema se sustenta en evidencia empírica, estadística o bibliográfica verificable, si las afirmaciones están debidamente respaldadas y si existe coherencia entre las causas identificadas y el abordaje propuesto. Esta dimensión verifica la calidad del sustento argumental antes de evaluar la estructura diagnóstica.} Los criterios se aprecian en la tabla~\ref{tab:rubrica-dim2-contextualizacion-evidencia}.

\scriptsize
\begin{longtable}{|C{0.55cm}|L{2cm}|L{8cm}|C{0.3cm}|C{0.3cm}|C{0.3cm}|C{0.3cm}|}
\caption{Rúbrica de evaluación de la Dimensión II: Contextualización y evidencia del problema}\label{tab:rubrica-dim2-contextualizacion-evidencia}\\
\hline
\rowcolor{headerblue}
\thd{N.\textsuperscript{o}} & \thd{Criterio} & \thd{Indicador verificable} & \thd{3} & \thd{2} & \thd{1} & \thd{0} \\
\hline\endfirsthead
\hline
\rowcolor{headerblue}
\thd{N.\textsuperscript{o}} & \thd{Criterio} & \thd{Indicador verificable} & \thd{3} & \thd{2} & \thd{1} & \thd{0} \\
\hline\endhead
2.1 & Sustento empírico, estadístico o bibliográfico del problema en el contexto de estudio \epref{Evidencia contextual} & El problema se sustenta con datos verificables provenientes de al menos una de las siguientes fuentes: estadísticas oficiales (inventarios viales, registros de fallas, datos del SINPAD/CENEPRED/INGEMMET), resultados de ensayos previos (SPT, CBR, granulometría, límites de Atterberg), estudios publicados en revistas indexadas, informes técnicos institucionales, o bases de datos geotécnicas. El sustento es específico para el contexto espacial y temporal declarado y no se limita a generalidades globales descontextualizadas. &
& & & \\ \hline
\rowcolor{lightgray}
2.2 & Ausencia de afirmaciones sin respaldo presentadas como hechos o premisas de conocimiento común & La descripción del problema no contiene afirmaciones fácticas carentes de referencia bibliográfica, fuente estadística o evidencia empírica. No se presentan como \enquote{hechos conocidos} o \enquote{premisas de conocimiento común} proposiciones que requieren sustento técnico (p.\,ej., \enquote{los suelos arcillosos de la región son altamente expansivos} sin ensayos de expansión libre o índice de plasticidad documentado; \enquote{el tráfico vehicular ha aumentado considerablemente} sin datos de conteo vehicular o IMDA). Cada afirmación factual relevante para la construcción del problema es trazable a una fuente verificable. &
& & & \\ \hline
2.3 & Coherencia entre las causas identificadas y el abordaje metodológico propuesto & Las causas o factores identificados como generadores del problema son coherentes con la estrategia de investigación que el proyecto anticipa. Si el problema identifica, por ejemplo, la insuficiencia de un aditivo convencional como causa de baja durabilidad, el abordaje propuesto debe incluir la evaluación experimental de alternativas al aditivo convencional, no un estudio meramente descriptivo de las propiedades del material existente. No hay desconexión entre \enquote{lo que causa el problema} y \enquote{lo que el proyecto propone investigar}. &
& & & \\ \hline
\rowcolor{lightgray}
2.4 & Escalamiento lógico de la contextualización: global $\rightarrow$ nacional $\rightarrow$ regional $\rightarrow$ local & La contextualización del problema sigue una progresión lógica desde el contexto global o internacional, pasando por el contexto nacional y regional, hasta llegar al contexto local y específico del estudio. Cada nivel de escalamiento aporta información pertinente que converge hacia la justificación de la necesidad de investigación en el sitio o material específico. La progresión no es un recurso retórico de relleno sino un argumento de convergencia que demuestra la pertinencia local del problema. &
& & & \\ \hline
\multicolumn{3}{|r|}{\cellcolor{white}\textbf{Subtotal Dimensión~II (máx.~12 puntos):}} & \multicolumn{4}{c|}{\cellcolor{white}\textbf{/12}} \\ \hline
\end{longtable}
\normalsize

\clearpage
\subsection{Diagnóstico Técnico de la Situación Problemática}
\textit{Se evalúa la rigurosidad del diagnóstico como primer componente de la secuencia Diagnóstico--Pronóstico--Control del Pronóstico. El diagnóstico debe describir el estado actual del fenómeno o sistema investigado con precisión técnica, cuantificación verificable, identificación causal y coherencia con el dominio especializado de la Ingeniería Civil.} Los criterios se aprecian en la tabla~\ref{tab:rubrica-dim3-diagnostico-tecnico}.

\scriptsize
\begin{longtable}{|C{0.55cm}|L{2cm}|L{8cm}|C{0.3cm}|C{0.3cm}|C{0.3cm}|C{0.3cm}|}
\caption{Rúbrica de evaluación de la Dimensión~III: Diagnóstico técnico de la situación problemática}\label{tab:rubrica-dim3-diagnostico-tecnico}\\
\hline
\rowcolor{headerblue}
\thd{N.\textsuperscript{o}} & \thd{Criterio} & \thd{Indicador verificable} & \thd{3} & \thd{2} & \thd{1} & \thd{0} \\
\hline\endfirsthead
\hline
\rowcolor{headerblue}
\thd{N.\textsuperscript{o}} & \thd{Criterio} & \thd{Indicador verificable} & \thd{3} & \thd{2} & \thd{1} & \thd{0} \\
\hline\endhead
3.1 & Descripción verificable del estado actual del fenómeno o sistema \epref{Diagnóstico} & El diagnóstico presenta una descripción técnica verificable de la situación problemática actual: condición del suelo, estado del pavimento, comportamiento de la cimentación, propiedades del material, o desempeño de la infraestructura. La descripción se basa en observaciones documentadas, mediciones reportadas o estudios previos identificables, y no en apreciaciones subjetivas o anecdóticas. Se distingue claramente entre lo observado/medido y lo interpretado. &
& & & \\ \hline
\rowcolor{lightgray}
3.2 & Cuantificación de magnitudes, brechas o deficiencias & El diagnóstico cuantifica, en la medida de lo posible, la magnitud del problema con unidades técnicas apropiadas: valores de CBR insuficientes respecto del mínimo normativo, índices de plasticidad que exceden umbrales de trabajabilidad, espesores de capa inferiores a los requeridos por diseño, valores de asentamiento superiores a los admisibles, porcentajes de deterioro superficial en inventarios viales, resistencias a la compresión por debajo de especificaciones, entre otros. Las brechas entre el estado actual y el estado deseable o normativo están claramente identificadas y dimensionadas. &
& & & \\ \hline
3.3 & Identificación de causas y factores contribuyentes al problema & El diagnóstico identifica, con sustento técnico o bibliográfico, las causas probables y los factores contribuyentes al problema: composición mineralógica del suelo, condiciones climáticas, deficiencias en el proceso constructivo, inadecuación del material de base, ausencia de tratamiento estabilizador, sobrecarga vehicular, drenaje insuficiente, u otros factores técnicamente pertinentes. Las causas identificadas no son especulativas ni genéricas, sino específicas al problema y al contexto descritos. &
& & & \\ \hline
\rowcolor{lightgray}
3.4 & Diagnóstico enmarcado en el dominio del campo de la Ingeniería Civil & El diagnóstico emplea terminología, normativa, clasificaciones y parámetros propios de la Ingeniería Civil y sus especialidades. Las referencias normativas son pertinentes (ASTM, AASHTO, MTC, CE.010, E.050, NTP, Eurocódigos, o normativa local aplicable). Los parámetros técnicos citados corresponden a ensayos y métodos reconocidos por la comunidad especializada. El diagnóstico no se realiza desde un dominio ajeno a la carrera ni emplea terminología imprecisa o ajena a la disciplina. &
& & & \\ \hline
\multicolumn{3}{|r|}{\cellcolor{white}\textbf{Subtotal Dimensión~III (máx.~12 puntos):}} & \multicolumn{4}{c|}{\cellcolor{white}\textbf{/12}} \\ \hline
\end{longtable}
\normalsize

\clearpage
\subsection{Pronóstico y Consecuencias de la No Intervención}
\textit{Se evalúa el pronóstico como segundo componente de la secuencia Diagnóstico--Pronóstico--Control del Pronóstico. El pronóstico debe proyectar, con plausibilidad técnica y trazabilidad lógica desde el diagnóstico, las consecuencias que se derivarían de no abordar el problema identificado, abarcando dimensiones técnicas, estructurales, económicas, ambientales o de seguridad, según corresponda.} Los criterios se aprecian en la tabla~\ref{tab:rubrica-dim4-pronostico}.

\scriptsize
\begin{longtable}{|C{0.55cm}|L{2cm}|L{8cm}|C{0.3cm}|C{0.3cm}|C{0.3cm}|C{0.3cm}|}
\caption{Rúbrica de evaluación de la Dimensión~IV: Pronóstico y consecuencias de la no intervención}\label{tab:rubrica-dim4-pronostico}\\
\hline
\rowcolor{headerblue}
\thd{N.\textsuperscript{o}} & \thd{Criterio} & \thd{Indicador verificable} & \thd{3} & \thd{2} & \thd{1} & \thd{0} \\
\hline\endfirsthead
\hline
\rowcolor{headerblue}
\thd{N.\textsuperscript{o}} & \thd{Criterio} & \thd{Indicador verificable} & \thd{3} & \thd{2} & \thd{1} & \thd{0} \\
\hline\endhead
4.1 & Proyección de consecuencias técnicas y/o estructurales de la no intervención \epref{Pronóstico} & El pronóstico describe las consecuencias técnicas previsibles si el problema no se aborda: progresión del deterioro, pérdida de capacidad portante, incremento de asentamientos diferenciales, aceleración de la fatiga estructural, propagación de fisuras, reducción de la vida útil del pavimento o de la estructura, u otras consecuencias específicas al fenómeno investigado. Las consecuencias están lógicamente derivadas del diagnóstico y no son afirmaciones genéricas desconectadas de la situación descrita. &
& & & \\ \hline
\rowcolor{lightgray}
4.2 & Plausibilidad técnica de los escenarios proyectados & Los escenarios de pronóstico son técnicamente plausibles: se sustentan en principios de mecánica de suelos, mecánica de materiales, ingeniería de pavimentos, análisis estructural u otros fundamentos reconocidos de la Ingeniería Civil. Las proyecciones no son catastrofistas sin fundamento, ni minimizan consecuencias que la evidencia técnica permitiría anticipar. El nivel de severidad proyectado es proporcional a la magnitud del problema diagnosticado. &
& & & \\ \hline
4.3 & Dimensiones múltiples del pronóstico: técnica, económica, ambiental, de seguridad & El pronóstico aborda, según la pertinencia del problema, más de una dimensión de consecuencias: además de las consecuencias técnicas directas, se consideran implicaciones económicas (sobrecostos de mantenimiento, reposición prematura), ambientales (contaminación, consumo de recursos no renovables), de seguridad (riesgo para usuarios o poblaciones cercanas) o de serviciabilidad (interrupción de vías, reducción de niveles de servicio). El pronóstico no es unidimensional cuando la naturaleza del problema demanda una lectura multidimensional. &
& & & \\ \hline
\rowcolor{lightgray}
4.4 & Trazabilidad lógica del pronóstico desde el diagnóstico & Cada consecuencia proyectada en el pronóstico se deriva directamente de al menos una causa, factor o condición identificada en el diagnóstico. No se introducen consecuencias que carezcan de antecedente diagnóstico ni se omiten proyecciones que las causas identificadas harían previsibles. La cadena causal diagnóstico~$\rightarrow$~pronóstico es explícita, verificable y libre de saltos lógicos. &
& & & \\ \hline
\multicolumn{3}{|r|}{\cellcolor{white}\textbf{Subtotal Dimensión~IV (máx.~12 puntos):}} & \multicolumn{4}{c|}{\cellcolor{white}\textbf{/12}} \\ \hline
\end{longtable}
\normalsize

\clearpage
\subsection{Control del Pronóstico y Articulación Pre-Hipótesis}

\textit{Se evalúa el control del pronóstico como tercer componente de la secuencia Diagnóstico--Pronóstico--Control del Pronóstico. Se verifica si el proyecto anticipa la estrategia de investigación que permitirá abordar el problema, si la cadena argumental completa es internamente coherente y si la problematización es suficiente como antecedente lógico de la hipótesis, la justificación y la pregunta de investigación.} Los criterios se aprecian en la tabla~\ref{tab:rubrica-dim5-control-pronostico}.

\scriptsize
\begin{longtable}{|C{0.55cm}|L{2.8cm}|L{8cm}|C{0.3cm}|C{0.3cm}|C{0.3cm}|C{0.3cm}|}
\caption{Rúbrica de evaluación de la Dimensión~V: Control del pronóstico y articulación pre-hipótesis}\label{tab:rubrica-dim5-control-pronostico}\\
\hline
\rowcolor{headerblue}
\thd{N.\textsuperscript{o}} & \thd{Criterio} & \thd{Indicador verificable} & \thd{3} & \thd{2} & \thd{1} & \thd{0} \\
\hline\endfirsthead
\hline
\rowcolor{headerblue}
\thd{N.\textsuperscript{o}} & \thd{Criterio} & \thd{Indicador verificable} & \thd{3} & \thd{2} & \thd{1} & \thd{0} \\
\hline\endhead
5.1 & Anticipación de la estrategia de investigación como respuesta al problema \epref{Control del pronóstico} & El control del pronóstico anticipa, de manera fundamentada, la estrategia general de investigación que permitirá abordar el problema y evitar o mitigar las consecuencias proyectadas: la técnica de estabilización a evaluar, el material alternativo a caracterizar, el procedimiento de refuerzo a comparar, el modelo numérico a validar, o la solución constructiva a optimizar. La estrategia anticipada es coherente con la naturaleza del problema diagnosticado y con las consecuencias pronosticadas, y constituye una respuesta investigativa (no meramente profesional) al problema. &
& & & \\ \hline
\rowcolor{lightgray}
5.2 & Coherencia interna de la cadena diagnóstico--pronóstico--control del pronóstico & La secuencia completa diagnóstico~$\rightarrow$~pronóstico~$\rightarrow$~control del pronóstico presenta coherencia interna verificable: el diagnóstico identifica causas, el pronóstico proyecta consecuencias derivadas de esas causas, y el control del pronóstico propone una estrategia de investigación que actúa sobre las causas identificadas para prevenir las consecuencias proyectadas. No hay contradicciones internas, saltos lógicos, causas huérfanas (sin proyección de consecuencias) ni controles desarticulados (que no se corresponden con diagnóstico y pronóstico). &
& & & \\ \hline
5.3 & Suficiencia de la problematización como antecedente de la hipótesis y la pregunta de investigación & La problematización completa proporciona los elementos lógicos necesarios y suficientes para que la hipótesis, la pregunta de investigación y la justificación se deriven naturalmente de ella. La hipótesis que pueda formularse a partir de esta problematización no resulta arbitraria ni desconectada; la pregunta de investigación que pueda plantearse está contenida lógicamente en la brecha identificada; la justificación que pueda elaborarse encuentra su sustento en el diagnóstico y el pronóstico presentados. La problematización no deja vacíos que obliguen a la hipótesis o a la justificación a introducir elementos argumentales no anticipados. &
& & & \\ \hline
\rowcolor{lightgray}
5.4 & Originalidad y aporte al conocimiento especializado anticipado desde la problematización & La problematización permite identificar, desde su formulación, que la investigación propuesta generará un aporte al conocimiento del campo: un vacío en el estado del arte que se busca llenar, una contradicción empírica que se busca resolver, una técnica que no ha sido evaluada en las condiciones locales, un material que no ha sido caracterizado para el uso propuesto, o una comparación que no ha sido realizada previamente en el contexto de estudio. Este aporte anticipado no es genérico (\enquote{contribuir a la ingeniería}) sino específico y verificable. &
& & & \\ \hline
\multicolumn{3}{|r|}{\cellcolor{white}\textbf{Subtotal Dimensión~V (máx.~12 puntos):}} & \multicolumn{4}{c|}{\cellcolor{white}\textbf{/12}} \\ \hline
\end{longtable}
\normalsize

\newpage
En esta sección se presenta el consolidado de resultados por dimensión y la escala de calificación de la problematización, resumidos en las tablas~\ref{tab:consolidado-problematizacion} y~\ref{tab:escala-calificacion-problematizacion}.

\begin{table}[htbp]
\centering
\scriptsize
\caption{Cuadro consolidado de resultados por dimensión de la problematización}\label{tab:consolidado-problematizacion}
\begin{tabularx}{\textwidth}{|C{0.8cm}|L{7.5cm}|C{1.8cm}|C{1.8cm}|X|}
\hline
\rowcolor{headerblue}
\thd{Dim.} & \thd{Denominación} & \thd{Pts. Máx.} & \thd{Pts. Obt.} & \thd{Obs.} \\ \hline
I & Identificación y Delimitación del Problema & 12 & & \\ \hline
\rowcolor{lightgray}
II & Contextualización y Evidencia del Problema & 12 & & \\ \hline
III & Diagnóstico Técnico de la Situación Problemática & 12 & & \\ \hline
\rowcolor{lightgray}
IV & Pronóstico y Consecuencias de la No Intervención & 12 & & \\ \hline
V & Control del Pronóstico y Articulación Pre-Hipótesis & 12 & & \\ \hline
\rowcolor{white}
& \textbf{TOTAL GENERAL} & \textbf{60} & & \\ \hline
\end{tabularx}
\normalsize
\end{table}

\begin{table}[htbp]
\centering
\scriptsize
\caption{Escala de calificación de la problematización}\label{tab:escala-calificacion-problematizacion}
\begin{tabularx}{\textwidth}{|C{2.5cm}|C{4.2cm}|X|}
\hline
\rowcolor{headerblue}
\thd{Rango} & \thd{Calificación} & \thd{Significado} \\ \hline
\textbf{54--60}\newline(90--100\,\%) &
\textcolor{csuf}{\textbf{PROBLEMATIZACIÓN\newline PLENA}} &
La problematización es integral, técnicamente precisa y lógicamente articulada. El diagnóstico, el pronóstico y el control del pronóstico conforman una cadena argumental coherente, sustentada con evidencia verificable, que proporciona una base sólida y suficiente para la formulación de la hipótesis, la justificación y la pregunta de investigación. \\ \hline
\rowcolor{lightgray}
\textbf{42--53}\newline(70--89\,\%) &
\textcolor{cace}{\textbf{PROBLEMATIZACIÓN\newline SUFICIENTE}} &
La problematización es sustancial y presenta los componentes esenciales de la secuencia diagnóstico--pronóstico--control, con áreas específicas susceptibles de fortalecimiento en cuanto a precisión, sustento o articulación. Las debilidades son corregibles sin reestructuración fundamental. \\ \hline
\textbf{30--41}\newline(50--69\,\%) &
\textcolor{cdef}{\textbf{PROBLEMATIZACIÓN\newline PARCIAL}} &
La problematización es incompleta o predominantemente declarativa. Se evidencian vacíos significativos en el diagnóstico, el pronóstico o el control del pronóstico, afirmaciones sin sustento, o desconexiones lógicas que comprometen la articulación argumental. Requiere fortalecimiento sustancial. \\ \hline
\rowcolor{lightgray}
\textbf{0--29}\newline($<$50\,\%) &
\textcolor{cins}{\textbf{PROBLEMATIZACIÓN\newline INSUFICIENTE}} &
La problematización no demuestra consistencia verificable. El problema carece de identificación precisa, el diagnóstico es superficial o ausente, el pronóstico no se sustenta en el diagnóstico, o el control del pronóstico no anticipa una estrategia coherente. Se requiere una reformulación fundamental del planteamiento del problema. \\ \hline
\end{tabularx}
\normalsize
\end{table}

\clearpage
\section{Evaluación de las Justificaciones}

{\footnotesize
La justificación de un proyecto de tesis constituye el componente argumentativo que demuestra la pertinencia, la necesidad y el valor de la investigación propuesta. Una justificación rigurosa trasciende la simple declaración de intenciones y articula, con evidencia y razonamiento lógico, por qué la investigación merece ser ejecutada, financiada y difundida. La presente rúbrica permite al dictaminante evaluar las cinco dimensiones de la justificación que todo proyecto de investigación debe abordar de manera diferenciada: la \textbf{justificación práctica}, que sustenta la aplicabilidad operacional y el impacto en la resolución de problemas concretos; la \textbf{justificación teórica}, que establece la contribución al avance del conocimiento científico; la \textbf{justificación metodológica}, que defiende la pertinencia, validez y novedad de los métodos empleados; la \textbf{justificación ambiental}, que examina la contribución a la sostenibilidad y la reducción de impactos ecológicos; y la \textbf{justificación tecnológica}, que valora la innovación técnica y el potencial de transferencia a la práctica profesional. Referencias bibliográficas del marco de referencia: Arias Chávez, D. y Cangalaya Sevillano, L.\,M. (2023). \textit{Manual del tesista}. EDUNI, UNI; Bernal Torres, C.\,A. (2006). \textit{Metodología de la investigación} (2.\textsuperscript{a}~ed.). Pearson; Ñaupas Paitán, H. et~al. (2014). \textit{Metodología de la investigación cuantitativa-cualitativa y redacción de la tesis} (4.\textsuperscript{a}~ed.). Ediciones de la U; Vara Horna, A.\,A. (2010). \textit{¿Cómo evaluar la rigurosidad científica de las tesis doctorales?}. USMP. Si bien la normativa interna de la universidad exige la presentación de una justificación de carácter social y, en función de su relevancia, las evaluaciones de justificación analizadas a continuación permiten articular de manera más sistemática y coherente dicho requisito establecido por la normativa institucional.
\par}


Las dimensiones de evaluación y la escala de valoración se resumen en las Tablas~\ref{tab:dimensiones-justificacion} y~\ref{tab:escala-valoracion-justificacion}, respectivamente.

\begin{table}[htbp]
\centering
\scriptsize
\caption{Dimensiones de evaluación de la justificación del proyecto de investigación}
\label{tab:dimensiones-justificacion}
\begin{tabularx}{\textwidth}{|C{0.8cm}|L{4.2cm}|X|}
\hline
\rowcolor{headerblue}
\thd{Dim.} & \thd{Denominación} & \thd{Alcance} \\ \hline
I & Justificación Práctica & Aplicabilidad operacional, beneficiarios, productos tangibles, significancia práctica, ruta de implementación y mejora sobre prácticas vigentes.
\\ \hline
\rowcolor{lightgray}
II & Justificación Teórica & Contribución al conocimiento, vacíos teóricos, razonamiento causal, originalidad conceptual, marcos predictivos y reflexión epistemológica.
\\ \hline
III & Justificación Metodológica & Pertinencia del diseño, defensa de métodos e instrumentos, validez, confiabilidad, replicabilidad, contribución metodológica y gestión de riesgos.
\\ \hline
\rowcolor{lightgray}
IV & Justificación Ambiental & Reducción de huella de carbono, sostenibilidad de materiales, ciclo de vida, protección de ecosistemas y alineamiento con principios de desarrollo sostenible.
\\ \hline
V & Justificación Tecnológica & Innovación técnica, transferencia tecnológica, viabilidad de escalamiento, potencial de estandarización normativa e integración con sistemas existentes.
\\ \hline
\end{tabularx}
\end{table}

\begin{table}[htbp]
\centering
\scriptsize
\caption{Escala de valoración y descriptores operativos para la evaluación de la justificación}
\label{tab:escala-valoracion-justificacion}
\begin{tabularx}{\textwidth}{|C{2.8cm}|C{1cm}|X|}
\hline
\rowcolor{headerblue}
\thd{Nivel} & \thd{Pts.} & \thd{Descriptor operativo} \\ \hline
\textcolor{csuf}{\textbf{Excelente}} & \textbf{3} & La justificación es explícita, fundamentada con evidencia verificable y coherente con el diseño de investigación.
La argumentación es orgánica, específica y cuantificable, con trazabilidad entre el problema, los objetivos, los métodos y los resultados esperados.
\\ \hline
\rowcolor{lightgray}
\textcolor{cace}{\textbf{Aceptable}} & \textbf{2} & La justificación es identificable pero parcial o insuficientemente fundamentada.
El vínculo con el diseño de investigación existe pero carece de especificidad, cuantificación o evidencia de respaldo. Susceptible de fortalecimiento.
\\ \hline
\textcolor{cdef}{\textbf{Deficiente}} & \textbf{1} & La justificación es tangencial, genérica o meramente declarativa.
No hay sustento argumentativo riguroso ni coherencia con el planteamiento del problema o la metodología del proyecto.
\\ \hline
\rowcolor{lightgray}
\textcolor{cins}{\textbf{Ausente}} & \textbf{0} & El proyecto no presenta justificación en la dimensión evaluada, o la redacción es tan vaga que no permite identificar argumento alguno.
\\ \hline
\end{tabularx}
\end{table}

\clearpage
\subsection{Justificación Práctica}

\textit{Se evalúa si el proyecto demuestra que la investigación generará resultados aplicables a la resolución de problemas concretos, identificando beneficiarios, productos tangibles, significancia práctica medible y una ruta plausible desde los hallazgos hasta la implementación en campo (véase la Tabla~\ref{tab:rubrica-justificacion-practica}).}

{\scriptsize
\begin{longtable}{|C{0.55cm}|L{3.8cm}|L{5.8cm}|C{0.65cm}|C{0.65cm}|C{0.65cm}|C{0.65cm}|}
\caption{Rúbrica de evaluación para la dimensión de justificación práctica}\label{tab:rubrica-justificacion-practica}\\
\hline
\rowcolor{headerblue}
\thd{N.\textsuperscript{o}} & \thd{Criterio} & \thd{Indicador verificable} & \thd{3} & \thd{2} & \thd{1} & \thd{0} \\
\hline\endfirsthead
\hline
\rowcolor{headerblue}
\thd{N.\textsuperscript{o}} & \thd{Criterio} & \thd{Indicador verificable} & \thd{3} & \thd{2} & \thd{1} & \thd{0} \\
\hline\endhead
1.1 & Problema operacional y costos tangibles \dimref{Problema y partes interesadas} & El proyecto formula el problema como una condición tangible con costos medibles o consecuencias adversas 
concretas (fallas de cimentación, inestabilidad de taludes, deterioro prematuro de pavimentos, riesgos ocupacionales).
Se diferencia explícitamente entre el dominio del problema operacional y el dominio del problema conceptual, asegurando que los objetivos de investigación permanezcan claramente delimitados.
&
& & & \\ \hline
\rowcolor{lightgray}
1.2 & Beneficiarios y productos esperados \dimref{Claridad de productos} & El proyecto identifica explícitamente a los beneficiarios primarios y secundarios (ingenieros de diseño, contratistas, entidades gubernamentales, sector agrícola) y los vincula a beneficios concretos.
Se describen productos tangibles esperados (diseños óptimos de mezcla, protocolos estandarizados, manuales de aplicación, plantillas de análisis de costo de ciclo de vida) consistentes con el alcance del diseño experimental.
&
& & & \\ \hline
1.3 & Significancia práctica medible \dimref{Impacto práctico medible} & El proyecto especifica cómo se medirá la relevancia práctica mediante umbrales de efecto alineados a estándares técnicos, diferenciando explícitamente entre significancia estadística y relevancia práctica de ingeniería.
Se cuantifican las mejoras esperadas en propiedades del suelo y se compara el rendimiento con métodos de estabilización vigentes (cemento Portland, cal, geosintéticos).
&
& & & \\ \hline
\rowcolor{lightgray}
1.4 & Ruta de implementación y validez externa \dimref{Ruta plausible hacia el uso} & El proyecto describe una trayectoria factible y detallada desde los hallazgos hasta la aplicación real, incluyendo pasos, actores y recursos.
Las barreras previsibles se reconocen junto con estrategias de mitigación.
Se aborda la validez externa, informando sobre la aplicabilidad en tipos de suelo reales, condiciones de campo y variaciones climáticas.
&
& & & \\ \hline
\multicolumn{3}{|r|}{\cellcolor{white}\textbf{Subtotal Dimensión~I (máx.~12 puntos):}} & \multicolumn{4}{c|}{\cellcolor{white}\textbf{/12}} \\ \hline
\end{longtable}
}

\clearpage
\subsection{Justificación Teórica}

\textit{Se evalúa si el proyecto establece el fundamento conceptual y académico que respalda la investigación, articulando la contribución al avance de constructos teóricos, el llenado de vacíos de conocimiento, el razonamiento causal explícito y la originalidad de los marcos conceptuales propuestos (véase la Tabla~\ref{tab:rubrica-justificacion-teorica}).}

{\scriptsize
\begin{longtable}{|C{0.55cm}|L{3.8cm}|L{5.8cm}|C{0.65cm}|C{0.65cm}|C{0.65cm}|C{0.65cm}|}
\caption{Rúbrica de evaluación para la dimensión de justificación teórica}\label{tab:rubrica-justificacion-teorica}\\
\hline
\rowcolor{headerblue}
\thd{N.\textsuperscript{o}} & \thd{Criterio} & \thd{Indicador verificable} & \thd{3} & \thd{2} & \thd{1} & \thd{0} \\
\hline\endfirsthead
\hline
\rowcolor{headerblue}
\thd{N.\textsuperscript{o}} & \thd{Criterio} & \thd{Indicador verificable} & \thd{3} & \thd{2} & \thd{1} & \thd{0} \\
\hline\endhead
2.1 & Fundamento teórico y originalidad \dimref{Contribución al conocimiento} & El proyecto establece 
conexiones claras con teorías preexistentes en la disciplina e introduce elementos novedosos: descripción de comportamientos no documentados, nuevas explicaciones sobre mejoras en propiedades, refutación de resultados previos, o desarrollo de marcos predictivos.
Se evita la simple repetición de estudios anteriores y se demuestra avance más allá del conocimiento establecido.
&
& & & \\ \hline
\rowcolor{lightgray}
2.2 & Razonamiento causal y lógica de hipótesis \dimref{Razonamiento lógico y causal} & El proyecto presenta razonamiento explícito para relaciones causales, considerando variación asociativa, secuenciación de eventos e hipótesis alternativas.
Se establecen los mecanismos fundamentales antes de las hipótesis de moderación e interacción, con fundamento teórico.
Se evalúa críticamente la lógica subyacente a la inclusión o exclusión de variables, protocolos o agentes.
&
& & & \\ \hline
2.3 & Vacío de conocimiento e identificación de la laguna \dimref{Fundamento teórico} & El proyecto identifica de manera convincente la insuficiencia o laguna de investigación existente, articulando con claridad qué nueva comprensión se producirá y para quién es relevante este conocimiento.
Se responde rigurosamente a las preguntas «¿y qué?» y «¿a quién le importa?» con implicaciones teórico-prácticas.
&
& & & \\ \hline
\rowcolor{lightgray}
2.4 & Consistencia, integralidad y límites de generalización \dimref{Consistencia e integración} & El proyecto se alinea con métodos científicos establecidos, se fundamenta en supuestos explícitos y se desarrolla sistemáticamente.
Se delinean los tipos de generalización teórica facilitados por el diseño y aquellos que no lo son.
Se evitan la sobregeneralización sin límites explícitos y las afirmaciones causales sin sustento teórico.
&
& & & \\ \hline
\multicolumn{3}{|r|}{\cellcolor{white}\textbf{Subtotal Dimensión~II (máx.~12 puntos):}} & \multicolumn{4}{c|}{\cellcolor{white}\textbf{/12}} \\ \hline
\end{longtable}
}

\clearpage
\subsection{Justificación Metodológica}

\textit{Se evalúa si el proyecto presenta una defensa rigurosa de la estrategia metodológica adoptada, incluyendo la alineación problema-método, la selección y validación de instrumentos, la replicabilidad, la contribución metodológica novedosa y la planificación de contingencias ante desafíos anticipados (véase la Tabla~\ref{tab:rubrica-justificacion-metodologica}).}

{\scriptsize
\begin{longtable}{|C{0.55cm}|L{3.8cm}|L{5.8cm}|C{0.65cm}|C{0.65cm}|C{0.65cm}|C{0.65cm}|}
\caption{Rúbrica de evaluación para la dimensión de justificación metodológica}\label{tab:rubrica-justificacion-metodologica}\\
\hline
\rowcolor{headerblue}
\thd{N.\textsuperscript{o}} & \thd{Criterio} & \thd{Indicador verificable} & \thd{3} & \thd{2} & \thd{1} & \thd{0} \\
\hline\endfirsthead
\hline
\rowcolor{headerblue}
\thd{N.\textsuperscript{o}} & \thd{Criterio} & \thd{Indicador verificable} & \thd{3} & \thd{2} & \thd{1} & \thd{0} \\
\hline\endhead
3.1 & Alineación problema-método y defensa del diseño \dimref{Alineación y fundamentación} & La metodología adoptada (experimental, 
modelado numérico, campo o integrada) se alinea explícitamente con el problema específico y las preguntas de investigación.
Se articula por qué los enfoques alternativos fueron descartados con argumentos fundamentados en criterios de sostenibilidad, viabilidad económica, desempeño técnico o limitaciones de escala.
La selección de técnicas se justifica en función de su capacidad para cuantificar la eficacia del tratamiento.
&
& & & \\ \hline
\rowcolor{lightgray}
3.2 & Validez, confiabilidad y rigor de instrumentos \dimref{Selección de métodos e instrumentos} & Se establecen medidas concretas de validez interna y externa del diseño.
Se documentan procedimientos de calibración, protocolos de validación previa, justificación estadística del tamaño de muestra, controles negativos y positivos, y procedimientos de asignación.
Se especifica si el estudio replica protocolos normativos (ASTM, ISO, ISSMGE) o introduce innovaciones, articulando el grado de novedad.
&
& & & \\ \hline
3.3 & Contribución metodológica y replicabilidad \dimref{Contribución y reutilización} & Se proporciona una declaración clara sobre la naturaleza de la contribución metodológica (protocolo estandarizado, modelo constitutivo calibrado, procedimiento innovador de ensayo, base de datos de referencia).
Se detalla cómo los protocolos, modelos y datos podrán ser accedidos y reutilizados por otros investigadores.
Se indica si el método es directamente comparable con estudios previos.
&
& & & \\ \hline
\rowcolor{lightgray}
3.4 & Gestión de riesgos y planificación de contingencias \dimref{Practicidad y gestión de riesgos} & Se enumeran explícitamente los desafíos anticipados (degradación del agente, variabilidad mineralógica, dificultades de compactación, variabilidad en el suministro, transición laboratorio-campo) y se articula para cada riesgo una estrategia de mitigación concreta.
Se proporciona evidencia de viabilidad técnica y logística, incluyendo cronograma, disponibilidad de equipos y acceso a materiales.
&
& & & \\ \hline
\multicolumn{3}{|r|}{\cellcolor{white}\textbf{Subtotal Dimensión~III (máx.~12 puntos):}} & \multicolumn{4}{c|}{\cellcolor{white}\textbf{/12}} \\ \hline
\end{longtable}
}

\clearpage
\subsection{Justificación Ambiental}

\textit{Se evalúa si el proyecto fundamenta su contribución a la sostenibilidad ambiental, incluyendo la reducción de la huella de carbono respecto a sustancias convencionales, el uso de materiales de fuentes renovables, el análisis del ciclo de vida, la protección de ecosistemas y la alineación con principios de desarrollo sostenible en la práctica de la ingeniería (véase la Tabla~\ref{tab:rubrica-justificacion-ambiental}).}

{\scriptsize
\begin{longtable}{|C{0.55cm}|L{3.8cm}|L{5.8cm}|C{0.65cm}|C{0.65cm}|C{0.65cm}|C{0.65cm}|}
\caption{Rúbrica de evaluación para la dimensión de justificación ambiental}\label{tab:rubrica-justificacion-ambiental}\\
\hline
\rowcolor{headerblue}
\thd{N.\textsuperscript{o}} & \thd{Criterio} & \thd{Indicador verificable} & \thd{3} & \thd{2} & \thd{1} & \thd{0} \\
\hline\endfirsthead
\hline
\rowcolor{headerblue}
\thd{N.\textsuperscript{o}} & \thd{Criterio} & \thd{Indicador verificable} & \thd{3} & \thd{2} & \thd{1} & 
\thd{0} \\
\hline\endhead
4.1 & Reducción de la huella de carbono y sostenibilidad de materiales \dimref{Mejora sobre práctica vigente} & El proyecto fundamenta cuantitativamente la reducción de la huella de carbono asociada a estabilizantes convencionales (cemento Portland, cal hidratada).
Se documenta el uso de materiales derivados de fuentes renovables o residuos, y se argumenta que los agentes propuestos no compiten con cultivos alimentarios ni generan externalidades negativas.
&
& & & \\ \hline
\rowcolor{lightgray}
4.2 & Análisis del ciclo de vida y economía circular \dimref{Justificación ética y económica} & El proyecto incorpora un enfoque de ciclo de vida o economía circular para la evaluación comparativa del sistema de suelo tratado frente a alternativas convencionales.
Se consideran costos ambientales desde la producción del agente hasta la disposición final, incluyendo la reversibilidad del tratamiento y la reutilización de materiales de desecho locales.
&
& & & \\ \hline
4.3 & Protección de ecosistemas y mitigación de impactos \dimref{Función ética y pragmática} & El proyecto identifica y evalúa los impactos ambientales potenciales de la solución propuesta sobre suelos, cuerpos de agua, biodiversidad y calidad del aire.
Se abordan aspectos como la sedimentación, la erosión superficial, la generación de polvo, la contaminación por lixiviados y la biocompatibilidad del agente de tratamiento con la vegetación y la fauna del entorno.
&
& & & \\ \hline
\rowcolor{lightgray}
4.4 & Resiliencia climática y durabilidad ambiental \dimref{Validez externa} & El proyecto evalúa el desempeño del tratamiento bajo condiciones ambientales variables y cíclicas (ciclos de humedecimiento-secado, congelamiento-descongelamiento, variaciones estacionales de humedad).
Se abordan la durabilidad a largo plazo, la degradación del agente frente a ciclos climáticos y la resiliencia de la infraestructura tratada ante eventos climáticos extremos.
&
& & & \\ \hline
\multicolumn{3}{|r|}{\cellcolor{white}\textbf{Subtotal Dimensión~IV (máx.~12 puntos):}} & \multicolumn{4}{c|}{\cellcolor{white}\textbf{/12}} \\ \hline
\end{longtable}
}

\clearpage
\subsection{Justificación Tecnológica}

\textit{Se evalúa si el proyecto articula la innovación técnica introducida, el potencial de transferencia a la práctica profesional, la viabilidad de escalamiento desde el laboratorio hasta condiciones de campo, la contribución a la estandarización normativa y la integración con los sistemas tecnológicos y productivos existentes (véase la Tabla~\ref{tab:rubrica-justificacion-tecnologica}).}

{\scriptsize
\begin{longtable}{|C{0.55cm}|L{3.8cm}|L{5.8cm}|C{0.65cm}|C{0.65cm}|C{0.65cm}|C{0.65cm}|}
\caption{Rúbrica de evaluación para la dimensión de justificación tecnológica}\label{tab:rubrica-justificacion-tecnologica}\\
\hline
\rowcolor{headerblue}
\thd{N.\textsuperscript{o}} & \thd{Criterio} & \thd{Indicador verificable} & \thd{3} & \thd{2} & \thd{1} & \thd{0} \\
\hline\endfirsthead
\hline
\rowcolor{headerblue}
\thd{N.\textsuperscript{o}} & \thd{Criterio} & \thd{Indicador verificable} & \thd{3} & \thd{2} & \thd{1} & \thd{0} \\
\hline\endhead
5.1 & Innovación técnica y novedad del enfoque \dimref{Novedad 
y aplicabilidad} & El proyecto describe con claridad la innovación técnica introducida (nuevas formulaciones, protocolos de aplicación, tecnologías de tratamiento, sistemas de monitoreo) y la diferencia respecto a soluciones existentes.
Se destacan deficiencias en métodos vigentes de estabilización y se explica cómo el enfoque propuesto las supera mediante perspectivas novedosas, verificables y documentadas.
&
& & & \\ \hline
\rowcolor{lightgray}
5.2 & Transferencia tecnológica y escalamiento \dimref{Ruta plausible hacia el uso} & El proyecto articula la trayectoria de transferencia desde la validación en laboratorio hasta la adopción en campo, identificando actores más importantes (proveedores, contratistas de movimiento de tierras, empresas de remediación), barreras de escalamiento (producción a gran escala, integración en mezcladoras, logística de suministro) y estrategias de superación.
Se evalúa la viabilidad técnica y económica del escalamiento. &
& & & \\ \hline
5.3 & Potencial de estandarización normativa \dimref{Difusión y transferencia} & El proyecto prevé la incorporación de los resultados en especificaciones técnicas de agencias normativas (ASTM, AASHTO, normas nacionales) o en códigos de diseño y normas de práctica de ingeniería.
Se describen los pasos concretos para la validación en tramos de prueba piloto y la generación de cartas de diseño, diagramas de flujo o bases de datos de referencia que faciliten la adopción normativa.
&
& & & \\ \hline
\rowcolor{lightgray}
5.4 & Integración tecnológica y sistemas inteligentes \dimref{Contribución metodológica} & El proyecto evalúa la compatibilidad del tratamiento propuesto con sistemas tecnológicos y productivos existentes (equipos de mezcla convencionales, condiciones de curado en campo, capacitación de operarios).
Cuando es pertinente, se considera la integración con sistemas de monitoreo, sensores, adquisición automatizada de datos o sistemas geotécnicos inteligentes y adaptativos que informen sobre el estado de desempeño.
&
& & & \\ \hline
\multicolumn{3}{|r|}{\cellcolor{white}\textbf{Subtotal Dimensión~V (máx.~12 puntos):}} & \multicolumn{4}{c|}{\cellcolor{white}\textbf{/12}} \\ \hline
\end{longtable}
}


\clearpage
El cuadro consolidado de resultados por dimensión (puntaje máximo y puntaje obtenido) y la escala de calificación global se presentan en las Tablas~\ref{tab:consolidado-justificaciones} y~\ref{tab:escala-calificacion-justificaciones}, respectivamente.

\begin{table}[htbp]
\centering
\scriptsize
\caption{Cuadro consolidado por dimensión para la evaluación de las justificaciones}
\label{tab:consolidado-justificaciones}
\begin{tabularx}{\textwidth}{|C{0.8cm}|L{7.5cm}|C{1.8cm}|C{1.8cm}|X|}
\hline
\rowcolor{headerblue}
\thd{Dim.} & \thd{Denominación} & \thd{Pts.
Máx.} & \thd{Pts. Obt.} & \thd{Obs.} \\ \hline
I & Justificación Práctica & 12 & & \\ \hline
\rowcolor{lightgray}
II & Justificación Teórica & 12 & & \\ \hline
III & Justificación Metodológica & 12 & & \\ \hline
\rowcolor{lightgray}
IV & Justificación Ambiental & 12 & & \\ \hline
V & Justificación Tecnológica & 12 & & \\ \hline
\rowcolor{white}
& \textbf{TOTAL GENERAL} & \textbf{60} & & \\ \hline
\end{tabularx}
\end{table}

\begin{table}[htbp]
\centering
\scriptsize
\caption{Escala de calificación global para las justificaciones del proyecto}
\label{tab:escala-calificacion-justificaciones}
\begin{tabularx}{\textwidth}{|C{2.5cm}|C{3cm}|X|}
\hline
\rowcolor{headerblue}
\thd{Rango} & \thd{Calificación} & \thd{Significado} \\ \hline
\textbf{54--60}\newline(90--100\,\%) &
\textcolor{csuf}{\textbf{JUSTIFICACIÓN\newline PLENA}} &
El proyecto presenta justificaciones explícitas, fundamentadas, cuantificables y coherentes en todas las dimensiones.
La argumentación es orgánica, la trazabilidad problema-método-resultado es verificable, y las contribuciones práctica, teórica, metodológica, ambiental y tecnológica son claramente articuladas y defendibles ante revisión por pares.
\\ \hline
\rowcolor{lightgray}
\textbf{42--53}\newline(70--89\,\%) &
\textcolor{cace}{\textbf{JUSTIFICACIÓN\newline SUFICIENTE}} &
El proyecto evidencia justificaciones sustanciales en la mayoría de las dimensiones, con áreas específicas susceptibles de fortalecimiento.
Los vínculos principales están demostrados, pero se requiere mayor especificidad, cuantificación o fundamentación en al menos una dimensión.
\\ \hline
\textbf{30--41}\newline(50--69\,\%) &
\textcolor{cdef}{\textbf{JUSTIFICACIÓN\newline PARCIAL}} &
Las justificaciones son incompletas o predominantemente declarativas en varias dimensiones.
El proyecto requiere fortalecimiento significativo de la fundamentación argumentativa, la especificación de beneficiarios y productos, la defensa metodológica o la articulación de la contribución ambiental y tecnológica.
\\ \hline
\rowcolor{lightgray}
\textbf{0--29}\newline($<$50\,\%) &
\textcolor{cins}{\textbf{JUSTIFICACIÓN\newline INSUFICIENTE}} &
El proyecto no demuestra justificaciones verificables en la mayoría de las dimensiones.
Las argumentaciones son vagas, genéricas o ausentes, comprometiendo la pertinencia, la viabilidad y la defensa científica de la investigación propuesta.
\\ \hline
\end{tabularx}
\end{table}

\clearpage
\section{Evaluación de la Hipótesis de Investigación}

{\footnotesize
La hipótesis de investigación es la proposición central que anticipa la relación entre variables y orienta el diseño metodológico del proyecto de tesis. En Ingeniería Civil, toda hipótesis bien formulada articula, en un primer nivel, una relación abstracta entre constructos teóricos \textit{(hipótesis conceptual)} y, en un segundo nivel, la traduce en una predicción específica, medible y verificable mediante procedimientos estandarizados \textit{(hipótesis operacional)}. La calidad de la hipótesis no depende únicamente de su correcta redacción gramatical, sino de su solidez epistemológica: debe ser clara, estar fundamentada en teoría y evidencia previa, especificar las variables con precisión operacional, ser susceptible de contrastación empírica mediante métodos estadísticos adecuados, y derivarse con coherencia lógica del problema, la pregunta y los objetivos de investigación declarados. Esta rúbrica permite evaluar, de manera sistemática e independiente de la especialidad de Ingeniería Civil involucrada (geotécnica, estructuras, pavimentos, hidráulica, materiales, mecánica de rocas, infraestructura de transporte u otras), la calidad integral de la hipótesis formulada en el proyecto de tesis de pregrado. La evaluación se estructura en cinco dimensiones complementarias que cubren la claridad enunciativa, la fundamentación teórica, la operacionalización de variables, la verificabilidad empírica y la coherencia con el planteamiento del problema. La rúbrica es aplicable tanto a hipótesis causales, correlacionales y mecanísticas, como a hipótesis experimentales unifactoriales o multifactoriales.\par
}
Las dimensiones de evaluación de la hipótesis se resumen en la Tabla~\ref{tab:dim-hipotesis}, mientras que la escala de valoración y sus descriptores operativos se presentan en la Tabla~\ref{tab:escala-hipotesis}.

\begin{table}[htbp]
\centering
\scriptsize
\caption{Dimensiones de evaluación de la hipótesis de investigación y alcance de cada dimensión.}
\label{tab:dim-hipotesis}
\begin{tabularx}{\textwidth}{|C{0.8cm}|L{4.5cm}|X|}
\hline
\rowcolor{headerblue}
\thd{Dim.} & \thd{Denominación} & \thd{Alcance} \\ \hline
I & Claridad y precisión del enunciado hipotético & Verifica que la hipótesis esté redactada como una proposición afirmativa, no ambigua y gramaticalmente completa, que anticipe explícitamente una relación entre variables sin confundirse con el objetivo, la pregunta de investigación ni la justificación. \\ \hline
\rowcolor{lightgray}
II & Fundamentación teórica y base bibliográfica & Evalúa que la hipótesis se origine en marcos teóricos consolidados y en la síntesis de estudios previos, que la relación propuesta sea plausible a la luz del conocimiento establecido en Ingeniería Civil, y que no postule relaciones arbitrarias o sin sustento técnico-científico. \\ \hline
III & Identificación y operacionalización de variables & Examina que la hipótesis identifique con precisión las variables independiente(s), dependiente(s) y de control, que cada variable esté definida operacionalmente mediante instrumentos, escalas o procedimientos estandarizados, y que se distingan claramente los niveles de abstracción conceptual y operacional. \\ \hline
\rowcolor{lightgray}
IV & Verificabilidad empírica y diseño metodológico & Evalúa que la hipótesis sea contrastable mediante procedimientos experimentales, cuasiexperimentales o correlacionales factibles, que anticipe los métodos estadísticos apropiados para su evaluación, y que la predicción sea suficientemente específica en dirección y, cuando sea posible, en magnitud. \\ \hline
V & Coherencia con el problema, los objetivos y la pregunta de investigación & Verifica que la hipótesis se derive lógicamente del problema identificado, responda directamente a la pregunta de investigación, sea consistente con los objetivos declarados y no introduzca variables, relaciones o alcances no anticipados en el planteamiento. \\ \hline
\end{tabularx}
\end{table}

\begin{table}[htbp]
\centering
\scriptsize
\caption{Escala de valoración (puntaje) y descriptor operativo para la evaluación de la hipótesis.}
\label{tab:escala-hipotesis}
\begin{tabularx}{\textwidth}{|C{2.8cm}|C{1cm}|X|}
\hline
\rowcolor{headerblue}
\thd{Nivel} & \thd{Pts.} & \thd{Descriptor operativo} \\ \hline
\textcolor{csuf}{\textbf{Excelente}} & \textbf{3} & El proyecto satisface plenamente el criterio evaluado. La hipótesis cumple el aspecto evaluado de manera explícita, técnicamente precisa y sin ambigüedades ni vacíos. \\ \hline
\rowcolor{lightgray}
\textcolor{cace}{\textbf{Aceptable}} & \textbf{2} & El criterio es identificable pero se cumple de manera parcial o incompleta. La hipótesis aborda el aspecto evaluado sin la profundidad, la precisión o el sustento suficiente. Susceptible de fortalecimiento con correcciones específicas. \\ \hline
\textcolor{cdef}{\textbf{Deficiente}} & \textbf{1} & El cumplimiento del criterio es tangencial o meramente declarativo. La hipótesis es superficial en el aspecto evaluado, carece de sustento verificable o presenta inconsistencias que debilitan su estructura lógica y metodológica. \\ \hline
\rowcolor{lightgray}
\textcolor{cins}{\textbf{Ausente}} & \textbf{0} & El proyecto no evidencia consideración alguna del criterio evaluado. Existe omisión total del aspecto evaluado o una desconexión que invalida la hipótesis en ese componente. \\ \hline
\end{tabularx}
\end{table}

\clearpage
\subsection{Claridad y Precisión del Enunciado Hipotético}
\textit{Se evalúa si la hipótesis está formulada como una proposición afirmativa, sintácticamente completa y no ambigua, que anticipa una relación entre variables de manera directamente legible. Esta dimensión verifica los requisitos mínimos de forma antes de evaluar el contenido teórico y metodológico.} Véase la Tabla~\ref{tab:rubrica-hip-dim1}.

{\scriptsize
\begin{longtable}{|C{0.55cm}|L{2cm}|L{8cm}|C{0.4cm}|C{0.4cm}|C{0.4cm}|C{0.4cm}|}
\caption{Rúbrica de evaluación de la Dimensión~I: claridad y precisión del enunciado hipotético.}\label{tab:rubrica-hip-dim1}\\
\hline
\rowcolor{headerblue}
\thd{N.\textsuperscript{o}} & \thd{Criterio} & \thd{Indicador verificable} & \thd{3} & \thd{2} & \thd{1} & \thd{0} \\
\hline\endfirsthead
\hline
\rowcolor{headerblue}
\thd{N.\textsuperscript{o}} & \thd{Criterio} & \thd{Indicador verificable} & \thd{3} & \thd{2} & \thd{1} & \thd{0} \\
\hline\endhead
1.1 & Formulación como proposición afirmativa que anticipa una relación entre variables & La hipótesis está redactada como un enunciado declarativo afirmativo que postula la existencia, dirección o naturaleza de una relación entre al menos dos variables. No se confunde con una pregunta de investigación, un objetivo (\enquote{determinar si\ldots}), una descripción del fenómeno, ni una declaración de intención metodológica. La relación postulada es explícita y no requiere inferencia del lector para ser identificada. &
& & & \\ \hline
\rowcolor{lightgray}
1.2 & Ausencia de ambigüedad terminológica en los constructos empleados & Los términos centrales de la hipótesis tienen un significado unívoco en el contexto especializado de la Ingeniería Civil. No se emplean términos vagos o polisémicos sin definición previa (p.\,ej., \enquote{mejorar el suelo}, \enquote{optimizar el sistema}, \enquote{aumentar la calidad}). Cada constructo es reconocible y delimitado en su campo de especialidad: resistencia a la compresión, módulo de elasticidad, factor de seguridad, coeficiente de permeabilidad, deflexión superficial, u otros parámetros técnicos específicos. &
& & & \\ \hline
1.3 & Identificación explícita de la dirección o naturaleza de la relación postulada & La hipótesis especifica la naturaleza de la relación entre variables: positiva, negativa, causal directa, correlacional, de interacción, o de diferencia entre grupos o tratamientos. No se limita a enunciar que \enquote{existe una relación} sin especificar su dirección o tipo, salvo que la teoría disponible no permita anticipar la dirección, en cuyo caso esta limitación debe reconocerse explícitamente. &
& & & \\ \hline
\rowcolor{lightgray}
1.4 & Delimitación del alcance y las condiciones de aplicabilidad de la hipótesis & La hipótesis declara, implícita o explícitamente, las condiciones bajo las cuales se espera que la relación sea válida: tipo de material, sistema estructural, condiciones de carga, horizonte temporal, escala de ensayo (laboratorio, campo, modelo numérico), o contexto geográfico. La hipótesis no pretende generalidad universal cuando el diseño metodológico solo permite conclusiones acotadas al contexto del estudio. &
& & & \\ \hline
\multicolumn{3}{|r|}{\cellcolor{white}\textbf{Subtotal Dimensión~I (máx.~12 puntos):}} & \multicolumn{4}{c|}{\cellcolor{white}\textbf{/12}} \\ \hline
\end{longtable}
}

\clearpage
\subsection{Fundamentación Teórica y Base Bibliográfica}
\textit{Se evalúa si la relación postulada en la hipótesis tiene sustento en marcos teóricos consolidados y en evidencia bibliográfica previa, verificando que no sea una conjetura arbitraria sino una proposición plausible derivada del estado del arte en la especialidad de Ingeniería Civil correspondiente.} Los criterios se aprecian en la Tabla~\ref{tab:rubrica-hip-dim2}.

\scriptsize
\begin{longtable}{|C{0.55cm}|L{2cm}|L{8cm}|C{0.3cm}|C{0.3cm}|C{0.3cm}|C{0.3cm}|}
\caption{Rúbrica de evaluación de la Dimensión~II: fundamentación teórica y base bibliográfica.}\label{tab:rubrica-hip-dim2}\\
\hline
\rowcolor{headerblue}
\thd{N.\textsuperscript{o}} & \thd{Criterio} & \thd{Indicador verificable} & \thd{3} & \thd{2} & \thd{1} & \thd{0} \\
\hline\endfirsthead
\hline
\rowcolor{headerblue}
\thd{N.\textsuperscript{o}} & \thd{Criterio} & \thd{Indicador verificable} & \thd{3} & \thd{2} & \thd{1} & \thd{0} \\
\hline\endhead
2.1 & Derivación de la hipótesis desde un marco teórico de la especialidad explícitamente identificado & La hipótesis se origina en principios, modelos o teorías reconocidos dentro de la Ingeniería Civil y sus especialidades: mecánica del medio continuo, teoría de estructuras, mecánica de suelos, hidráulica, ingeniería de pavimentos, mecánica de rocas, química de materiales de construcción, u otros marcos pertinentes. El marco teórico de referencia es identificable en el proyecto y constituye la base lógica para la relación postulada, no un adorno retórico. &
& & & \\ \hline
\rowcolor{lightgray}
2.2 & Sustento en estudios empíricos o experimentales previos pertinentes & La plausibilidad de la hipótesis está respaldada por hallazgos de investigaciones previas publicadas en revistas indexadas, informes técnicos institucionales o tesis de posgrado reconocidas, que hayan observado relaciones similares o que hayan establecido los antecedentes necesarios para anticipar la relación propuesta. El sustento bibliográfico no es genérico ni decorativo: los estudios citados son directamente pertinentes a las variables y al contexto técnico de la hipótesis. &
& & & \\ \hline
2.3 & Consistencia de la hipótesis con el conocimiento técnico establecido en la disciplina & La relación propuesta es consistente con los principios fundamentales de la Ingeniería Civil aplicables al fenómeno estudiado. Si la hipótesis propone una relación que desafía el conocimiento establecido, el proyecto justifica explícitamente por qué las condiciones específicas del estudio podrían producir un resultado distinto al convencional. No se postulan relaciones que contradigan principios físicos, mecánicos o químicos sin justificación técnica sustentada. &
& & & \\ \hline
\rowcolor{lightgray}
2.4 & Identificación de la brecha o vacío de conocimiento que la hipótesis busca llenar & La hipótesis permite identificar, de manera específica y verificable, qué aspecto del fenómeno estudiado no ha sido suficientemente investigado, qué contradicción empírica persiste en la literatura, o qué relación no ha sido evaluada bajo las condiciones específicas del estudio. El vacío no es genérico (\enquote{existen pocos estudios sobre\ldots}) sino técnicamente acotado y directamente vinculado a la relación que la hipótesis propone investigar. &
& & & \\ \hline
\multicolumn{3}{|r|}{\cellcolor{white}\textbf{Subtotal Dimensión~II (máx.~12 puntos):}} & \multicolumn{4}{c|}{\cellcolor{white}\textbf{/12}} \\ \hline
\end{longtable}
\normalsize

\clearpage
\subsection{Identificación y Operacionalización de Variables}
\textit{Se evalúa si la hipótesis identifica con precisión las variables que intervienen en la relación propuesta y si cada variable está operacionalmente definida mediante instrumentos, escalas o procedimientos estandarizados que permitan su medición o manipulación reproducible en el contexto específico del estudio.} Los criterios se aprecian en la Tabla~\ref{tab:rubrica-hip-dim3}.

\scriptsize
\begin{longtable}{|C{0.55cm}|L{2cm}|L{8cm}|C{0.3cm}|C{0.3cm}|C{0.3cm}|C{0.3cm}|}
\caption{Rúbrica de evaluación de la Dimensión~III: identificación y operacionalización de variables.}\label{tab:rubrica-hip-dim3}\\
\hline
\rowcolor{headerblue}
\thd{N.\textsuperscript{o}} & \thd{Criterio} & \thd{Indicador verificable} & \thd{3} & \thd{2} & \thd{1} & \thd{0} \\
\hline\endfirsthead
\hline
\rowcolor{headerblue}
\thd{N.\textsuperscript{o}} & \thd{Criterio} & \thd{Indicador verificable} & \thd{3} & \thd{2} & \thd{1} & \thd{0} \\
\hline\endhead
3.1 & Identificación explícita y diferenciada de la variable independiente y la variable dependiente & La hipótesis permite identificar sin ambigüedad cuál es la variable independiente (factor manipulado o clasificador) y cuál es la variable dependiente (respuesta medida). En hipótesis correlacionales se identifican ambas variables sin imposición de causalidad. En hipótesis multifactoriales se enuncian todos los factores principales y, cuando sea pertinente, las interacciones anticipadas. No se presentan hipótesis con variables indistinguibles ni con relaciones simétricas no justificadas. &
& & & \\ \hline
\rowcolor{lightgray}
3.2 & Definición operacional de la variable independiente con especificación de niveles o valores & La variable independiente está definida operacionalmente: se especifica el agente, tratamiento, condición o factor en términos técnicos precisos (tipo de aditivo, dosificación en \%, nivel de carga en kN, temperatura en \textdegree C, tiempo de curado en días, tipo de refuerzo, configuración estructural, u otro parámetro manipulable). Cuando corresponde, se declaran los niveles o rangos de variación que se evaluarán en el estudio. La definición es reproducible por otro investigador sin información adicional. &
& & & \\ \hline
3.3 & Definición operacional de la variable dependiente con especificación del método de medición & La variable dependiente está definida operacionalmente: se indica el parámetro de respuesta, la unidad de medida, el procedimiento o norma técnica aplicable para su determinación (ASTM, AASHTO, ACI, ISO, NTP, MTC u otras normas pertinentes), y el equipo o instrumento requerido. No se emplea el nombre del constructo abstracto como definición operacional (p.\,ej., \enquote{estabilidad} sin especificar si se mide como factor de seguridad, índice de resistencia, deflexión u otro parámetro cuantificable). &
& & & \\ \hline
\rowcolor{lightgray}
3.4 & Identificación de variables de control y reconocimiento de la variabilidad inherente & La hipótesis o el diseño metodológico que la acompaña identifica las variables que deben mantenerse constantes o monitorizarse para garantizar la validez interna del estudio: contenido de humedad, temperatura de curado, granulometría del material, densidad de compactación, condiciones de drenaje, estado de esfuerzos, u otras variables de confusión pertinentes. Se reconoce, cuando corresponda, la variabilidad espacial o material inherente al objeto de estudio y su potencial efecto sobre los resultados. &
& & & \\ \hline
\multicolumn{3}{|r|}{\cellcolor{white}\textbf{Subtotal Dimensión~III (máx.~12 puntos):}} & \multicolumn{4}{c|}{\cellcolor{white}\textbf{/12}} \\ \hline
\end{longtable}
\normalsize

\clearpage
\subsection{Verificabilidad Empírica y Diseño Metodológico}
\textit{Se evalúa si la hipótesis es contrastable mediante procedimientos factibles dentro del alcance del proyecto de tesis de pregrado, si anticipa los métodos estadísticos apropiados para su evaluación, y si la predicción tiene el nivel de especificidad requerido para que la prueba de hipótesis sea inequívoca.} Los criterios se aprecian en la Tabla~\ref{tab:rubrica-hip-dim4}.

\scriptsize
\begin{longtable}{|C{0.55cm}|L{2cm}|L{8cm}|C{0.3cm}|C{0.3cm}|C{0.3cm}|C{0.3cm}|}
\caption{Rúbrica de evaluación de la Dimensión~IV: verificabilidad empírica y diseño metodológico.}\label{tab:rubrica-hip-dim4}\\
\hline
\rowcolor{headerblue}
\thd{N.\textsuperscript{o}} & \thd{Criterio} & \thd{Indicador verificable} & \thd{3} & \thd{2} & \thd{1} & \thd{0} \\
\hline\endfirsthead
\hline
\rowcolor{headerblue}
\thd{N.\textsuperscript{o}} & \thd{Criterio} & \thd{Indicador verificable} & \thd{3} & \thd{2} & \thd{1} & \thd{0} \\
\hline\endhead
4.1 & Contrastabilidad empírica de la hipótesis mediante procedimientos factibles & La hipótesis puede ser evaluada empíricamente mediante los recursos, equipos, materiales y plazos disponibles en el contexto del proyecto de tesis de pregrado: ensayos de laboratorio normalizados, levantamientos topográficos, instrumentación geotécnica, modelos numéricos calibrados, análisis de muestras de campo, u otros procedimientos factibles. La hipótesis no requiere para su verificación ensayos o equipos inaccesibles, periodos de observación incompatibles con el cronograma del proyecto, ni recursos económicos desproporcionados. &
& & & \\ \hline
\rowcolor{lightgray}
4.2 & Coherencia entre la naturaleza de la hipótesis y el método estadístico de contrastación & La hipótesis es coherente con el diseño estadístico que le corresponde: hipótesis de diferencia entre grupos (ANOVA, prueba $t$, Kruskal-Wallis); hipótesis de correlación o regresión (correlación de Pearson o Spearman, regresión lineal o polinomial, regresión múltiple); hipótesis de interacción factorial (ANOVA de dos vías, diseños factoriales); hipótesis de ajuste de modelo (bondad de ajuste, análisis de residuos). El método estadístico anticipado es apropiado para el nivel de medición de las variables y para la distribución esperada de los datos. &
& & & \\ \hline
4.3 & Especificidad de la predicción: dirección esperada y, cuando sea pertinente, magnitud anticipada & La predicción es suficientemente específica para que el resultado del estudio pueda calificarse inequívocamente como confirmatorio o refutatorio: se indica la dirección esperada de la relación (positiva, negativa, mayor que, menor que) y, cuando la teoría o los estudios previos lo justifican, se anticipa un rango o umbral de magnitud del efecto. No se formulan hipótesis bidireccionales ambiguas como único recurso cuando la teoría permite anticipar la dirección del efecto. &
& & & \\ \hline
\rowcolor{lightgray}
4.4 & Reconocimiento implícito o explícito de la hipótesis nula como referente de contrastación & La hipótesis está formulada de manera que su refutación sea conceptualmente posible: existe un estado de línea base, un valor de referencia normativo, un grupo de control o una condición de no-efecto respecto del cual la predicción puede evaluarse. No se formulan hipótesis no falsables (p.\,ej., \enquote{el tratamiento mejorará las propiedades del material} sin definición operacional de \enquote{mejorar} ni referente cuantificable de comparación). &
& & & \\ \hline
\multicolumn{3}{|r|}{\cellcolor{white}\textbf{Subtotal Dimensión~IV (máx.~12 puntos):}} & \multicolumn{4}{c|}{\cellcolor{white}\textbf{/12}} \\ \hline
\end{longtable}
\normalsize

\clearpage
\subsection{Coherencia con el Problema, los Objetivos y la Pregunta de Investigación}

\textit{Se evalúa si la hipótesis constituye una respuesta articulada y lógicamente derivada del problema de investigación planteado, si responde de manera directa a la pregunta de investigación formulada y si es consistente con el alcance y los objetivos declarados en el proyecto de tesis.} Los criterios se aprecian en la Tabla~\ref{tab:rubrica-hip-dim5}.

\scriptsize
\begin{longtable}{|C{0.55cm}|L{2cm}|L{8cm}|C{0.3cm}|C{0.3cm}|C{0.3cm}|C{0.3cm}|}
\caption{Rúbrica de evaluación de la Dimensión~V: coherencia con el problema, los objetivos y la pregunta de investigación.}\label{tab:rubrica-hip-dim5}\\
\hline
\rowcolor{headerblue}
\thd{N.\textsuperscript{o}} & \thd{Criterio} & \thd{Indicador verificable} & \thd{3} & \thd{2} & \thd{1} & \thd{0} \\
\hline\endfirsthead
\hline
\rowcolor{headerblue}
\thd{N.\textsuperscript{o}} & \thd{Criterio} & \thd{Indicador verificable} & \thd{3} & \thd{2} & \thd{1} & \thd{0} \\
\hline\endhead
5.1 & Derivación lógica de la hipótesis a partir del problema de investigación identificado & La hipótesis responde al problema de conocimiento planteado en la problematización: las variables de la hipótesis corresponden a los constructos identificados como problemáticos, la relación postulada aborda la brecha, vacío o contradicción descritos en el diagnóstico, y el alcance de la hipótesis es compatible con la delimitación establecida para el problema. No se introduce en la hipótesis ninguna variable o relación que carezca de antecedente en la problematización. &
& & & \\ \hline
\rowcolor{lightgray}
5.2 & Correspondencia biunívoca entre la hipótesis y la pregunta de investigación & Existe correspondencia directa y verificable entre la hipótesis y la pregunta de investigación: si la pregunta interroga por el efecto de una variable sobre otra, la hipótesis predice ese efecto; si la pregunta interroga por la relación entre dos fenómenos, la hipótesis anticipa la naturaleza de esa relación. No hay variables presentes en la pregunta que estén ausentes en la hipótesis, ni variables en la hipótesis sin correlato en la pregunta. &
& & & \\ \hline
5.3 & Consistencia entre el alcance de la hipótesis y los objetivos específicos del proyecto & Los objetivos específicos declarados en el proyecto constituyen, en conjunto, los pasos operacionales necesarios y suficientes para contrastar la hipótesis: la recolección de datos, el procesamiento estadístico y la interpretación de resultados contemplados en los objetivos son suficientes para confirmar o refutar la hipótesis. No existen objetivos desconectados de la hipótesis, ni aspectos de la hipótesis no abordados por ningún objetivo. &
& & & \\ \hline
\rowcolor{lightgray}
5.4 & Ausencia de variables, relaciones o alcances en la hipótesis no anticipados en el planteamiento & La hipótesis no introduce elementos que no hayan sido anticipados en el planteamiento del problema, la pregunta de investigación ni los objetivos: no agrega variables no definidas, no amplía el alcance más allá de la delimitación establecida, y no postula relaciones adicionales a las que el problema y la pregunta de investigación justifican investigar. La hipótesis no es más amplia ni más restringida que el problema que la origina. &
& & & \\ \hline
\multicolumn{3}{|r|}{\cellcolor{white}\textbf{Subtotal Dimensión~V (máx.~12 puntos):}} & \multicolumn{4}{c|}{\cellcolor{white}\textbf{/12}} \\ \hline
\end{longtable}
\normalsize

\newpage
En esta sección se presenta el consolidado de resultados por dimensión y la escala de calificación de la hipótesis, resumidos en las Tablas~\ref{tab:consolidado-hipotesis} y~\ref{tab:escala-calificacion-hipotesis}.

\begin{table}[htbp]
\centering
\scriptsize
\caption{Cuadro consolidado de resultados por dimensión de la hipótesis de investigación.}\label{tab:consolidado-hipotesis}
\begin{tabularx}{\textwidth}{|C{0.8cm}|L{7.5cm}|C{1.8cm}|C{1.8cm}|X|}
\hline
\rowcolor{headerblue}
\thd{Dim.} & \thd{Denominación} & \thd{Pts. Máx.} & \thd{Pts. Obt.} & \thd{Obs.} \\ \hline
I & Claridad y Precisión del Enunciado Hipotético & 12 & & \\ \hline
\rowcolor{lightgray}
II & Fundamentación Teórica y Base Bibliográfica & 12 & & \\ \hline
III & Identificación y Operacionalización de Variables & 12 & & \\ \hline
\rowcolor{lightgray}
IV & Verificabilidad Empírica y Diseño Metodológico & 12 & & \\ \hline
V & Coherencia con el Problema, los Objetivos y la Pregunta & 12 & & \\ \hline
\rowcolor{white}
& \textbf{TOTAL GENERAL} & \textbf{60} & & \\ \hline
\end{tabularx}
\normalsize
\end{table}

\begin{table}[htbp]
\centering
\scriptsize
\caption{Escala de calificación de la hipótesis de investigación.}\label{tab:escala-calificacion-hipotesis}
\begin{tabularx}{\textwidth}{|C{2.5cm}|C{3.2cm}|X|}
\hline
\rowcolor{headerblue}
\thd{Rango} & \thd{Calificación} & \thd{Significado} \\ \hline
\textbf{54--60}\newline(90--100\,\%) &
\textcolor{csuf}{\textbf{HIPÓTESIS\newline PLENA}} &
La hipótesis es clara, técnicamente precisa, teóricamente fundamentada y operacionalmente completa. Las variables están identificadas y definidas sin ambigüedades, la predicción es contrastable mediante métodos estadísticos apropiados y existe correspondencia lógica total con el problema, la pregunta y los objetivos del proyecto. \\ \hline
\rowcolor{lightgray}
\textbf{42--53}\newline(70--89\,\%) &
\textcolor{cace}{\textbf{HIPÓTESIS\newline SUFICIENTE}} &
La hipótesis presenta los componentes esenciales con áreas específicas susceptibles de fortalecimiento en cuanto a precisión operacional, especificidad estadística o articulación con el planteamiento. Las debilidades son corregibles sin reformulación fundamental de la proposición. \\ \hline
\textbf{30--41}\newline(50--69\,\%) &
\textcolor{cdef}{\textbf{HIPÓTESIS\newline PARCIAL}} &
La hipótesis es incompleta o predominantemente declarativa. Se evidencian vacíos significativos en la operacionalización de variables, en el sustento teórico o en la verificabilidad empírica, o existen desconexiones entre la hipótesis y el problema que comprometen la coherencia del proyecto. Requiere fortalecimiento sustancial. \\ \hline
\rowcolor{lightgray}
\textbf{0--29}\newline($<$50\,\%) &
\textcolor{cins}{\textbf{HIPÓTESIS\newline INSUFICIENTE}} &
La hipótesis no satisface los criterios mínimos de calidad científica. Carece de precisión enunciativa, de fundamentación teórica verificable, de operacionalización de variables o de coherencia con el problema de investigación. Se requiere una reformulación fundamental que aborde las deficiencias identificadas en múltiples dimensiones. \\ \hline
\end{tabularx}
\normalsize
\end{table}

\clearpage
\section{Evaluación de la Formulación de la Pregunta de Investigación}

{\footnotesize
La pregunta de investigación constituye el eje conceptual alrededor del cual se estructura toda la investigación empírica, orientando la selección de marcos teóricos, estrategias metodológicas y procedimientos analíticos. Su formulación representa un acto intelectual sustantivo que delimita los confines epistémicos, el alcance operacional y la viabilidad práctica de un estudio. En Ingeniería Civil, la pregunta debe traducir un problema técnico ---identificado en el planteamiento del problema--- en un enunciado interrogativo que especifique variables mensurables, poblaciones o unidades de análisis claramente definidas, condiciones de contorno y una estructura semántica compatible con el diseño metodológico previsto. Una pregunta bien formulada no es genérica ni ambigua: distingue entre preguntas descriptivas, relacionales, comparativas o explicativas, y anticipa la lógica inferencial que guiará la recolección y el análisis de datos. Esta rúbrica permite evaluar la calidad, precisión y completitud de la pregunta de investigación, verificando que sea técnicamente rigurosa, operacionalmente viable, metodológicamente coherente y científicamente pertinente como componente articulador entre la problematización, los objetivos, la hipótesis y el diseño de investigación. La evaluación se centra en la consistencia interna, la especificidad operacional y la suficiencia de la pregunta como antecedente lógico del diseño experimental o metodológico del proyecto de tesis de pregrado en Ingeniería Civil. Las dimensiones de evaluación de la pregunta de investigación se resumen en la Tabla~\ref{tab:dim-pregunta}, mientras que la escala de valoración y sus descriptores operativos se presentan en la Tabla~\ref{tab:escala-pregunta}.\par
}

\begin{table}[htbp]
\centering
\scriptsize
\caption{Dimensiones de evaluación de la pregunta de investigación y alcance de cada dimensión.}
\label{tab:dim-pregunta}
\begin{tabularx}{\textwidth}{|C{0.8cm}|L{4.5cm}|X|}
\hline
\rowcolor{headerblue}
\thd{Dim.} & \thd{Denominación} & \thd{Alcance} \\ \hline
I & Formulación estructural y precisión semántica de la pregunta & Verifica que la pregunta esté formulada como un enunciado interrogativo claro, libre de ambigüedades, con estructura sintáctica que identifique sujeto, predicado y complemento, y que se derive lógicamente de la problematización. \\ \hline
\rowcolor{lightgray}
II & Especificación de variables y operacionalización & Evalúa la identificación explícita de variables dependientes e independientes, la presencia de definiciones operacionales con protocolos de medición normalizados, y la especificación de unidades, rangos y niveles de tratamiento. \\ \hline
III & Coherencia metodológica y alineación con el diseño de investigación & Examina la correspondencia entre la estructura semántica de la pregunta y el diseño metodológico que anticipa: tipo de pregunta (descriptiva, relacional, comparativa o explicativa), lógica inferencial y técnica analítica compatible. \\ \hline
\rowcolor{lightgray}
IV & Viabilidad técnica y delimitación del alcance & Evalúa si la pregunta es respondible con los recursos disponibles, si el alcance espacial, temporal y poblacional está definido, y si la unidad de análisis es coherente con el nivel de inferencia pretendido. \\ \hline
V & Pertinencia científica y aporte al conocimiento del área & Verifica que la pregunta aborde una brecha, vacío o contradicción identificable en el estado del arte, que su resolución genere un aporte específico al conocimiento en Ingeniería Civil y que posea relevancia tanto teórica como práctica. \\ \hline
\end{tabularx}
\end{table}

\begin{table}[htbp]
\centering
\scriptsize
\caption{Escala de valoración (puntaje) y descriptor operativo para la evaluación de la pregunta de investigación.}
\label{tab:escala-pregunta}
\begin{tabularx}{\textwidth}{|C{2.8cm}|C{1cm}|X|}
\hline
\rowcolor{headerblue}
\thd{Nivel} & \thd{Pts.} & \thd{Descriptor operativo} \\ \hline
\textcolor{csuf}{\textbf{Excelente}} & \textbf{3} & El proyecto satisface plenamente el criterio evaluado. La pregunta de investigación es explícita, técnicamente precisa, operacionalmente definida y lógicamente articulada con la problematización y el diseño metodológico. No se identifican ambigüedades, omisiones ni inconsistencias. \\ \hline
\rowcolor{lightgray}
\textcolor{cace}{\textbf{Aceptable}} & \textbf{2} & El criterio es identificable pero se cumple de manera parcial o incompleta. La pregunta aborda el aspecto evaluado sin la profundidad, la especificidad o la coherencia suficiente. Susceptible de fortalecimiento con correcciones específicas. \\ \hline
\textcolor{cdef}{\textbf{Deficiente}} & \textbf{1} & El cumplimiento del criterio es tangencial o meramente declarativo. La pregunta es superficial, carece de especificación operacional o presenta inconsistencias que debilitan su capacidad para orientar el diseño de investigación. \\ \hline
\rowcolor{lightgray}
\textcolor{cins}{\textbf{Ausente}} & \textbf{0} & El proyecto no evidencia consideración alguna del criterio evaluado. Existe omisión total del aspecto evaluado o una desconexión que invalida la pregunta de investigación en ese componente. \\ \hline
\end{tabularx}
\end{table}

\clearpage
\subsection{Formulación Estructural y Precisión Semántica de la Pregunta}

\textit{Se evalúa si la pregunta de investigación está formulada como un enunciado interrogativo genuino, derivado lógicamente de la problematización, con estructura sintáctica clara que identifique las relaciones entre variables, libre de ambigüedades léxicas y con terminología técnica precisa propia de la Ingeniería Civil. Esta dimensión verifica los requisitos formales de la pregunta antes de evaluar su contenido operacional y metodológico.} Véase la Tabla~\ref{tab:rubrica-dim1-formulacion}.

{\scriptsize
\begin{longtable}{|C{0.55cm}|L{2cm}|L{8cm}|C{0.4cm}|C{0.4cm}|C{0.4cm}|C{0.4cm}|}
\caption{Rúbrica de evaluación de la Dimensión~I: formulación estructural y precisión semántica de la pregunta.}\label{tab:rubrica-dim1-formulacion}\\
\hline
\rowcolor{headerblue}
\thd{N.\textsuperscript{o}} & \thd{Criterio} & \thd{Indicador verificable} & \thd{3} & \thd{2} & \thd{1} & \thd{0} \\
\hline\endfirsthead
\hline
\rowcolor{headerblue}
\thd{N.\textsuperscript{o}} & \thd{Criterio} & \thd{Indicador verificable} & \thd{3} & \thd{2} & \thd{1} & \thd{0} \\
\hline\endhead
1.1 & Estructura interrogativa genuina con identificación de relaciones entre variables & La pregunta está formulada como un enunciado interrogativo completo que identifica, como mínimo, una relación verificable entre variables o un fenómeno cuantificable. No se reduce a una afirmación disfrazada de pregunta ni a una declaración de intenciones (p.\,ej., \enquote{se busca analizar el comportamiento del concreto} no es una pregunta de investigación; \enquote{¿En qué medida la adición de fibras de polipropileno al 0,5\,\% y 1,0\,\% modifica la resistencia a la flexión de vigas de concreto f'c\,=\,210\,kg/cm² a los 28 días de curado?} sí lo es). La estructura sintáctica presenta sujeto, predicado y complemento identificables. &
& & & \\ \hline
\rowcolor{lightgray}
1.2 & Derivación lógica de la pregunta desde la problematización & La pregunta de investigación se deriva directamente de la brecha, vacío o contradicción identificada en la problematización. Existe trazabilidad verificable entre el problema diagnosticado, las consecuencias pronosticadas y la interrogante formulada. La pregunta no introduce elementos temáticos, variables o alcances que no hayan sido anticipados en el planteamiento del problema ni omite aspectos centrales de la problematización. La cadena argumental problema~$\rightarrow$~pregunta es explícita y libre de saltos lógicos. &
& & & \\ \hline
1.3 & Precisión terminológica y ausencia de ambigüedades léxicas & La pregunta emplea terminología técnica precisa, consistente con el dominio técnico de la Ingeniería Civil y sus especialidades (estructuras, geotecnia, hidráulica, vías, materiales, construcción). No se utilizan términos genéricos que admitan interpretaciones múltiples (p.\,ej., \enquote{mejorar el suelo} sin especificar el parámetro de mejora; \enquote{optimizar el diseño} sin definir la función objetivo; \enquote{comportamiento del material} sin indicar la propiedad mecánica o física evaluada). Cada término técnico relevante tiene un referente operacional identificable. &
& & & \\ \hline
\rowcolor{lightgray}
1.4 & Consistencia de la pregunta con el título, los objetivos y la hipótesis del proyecto & Existe correspondencia verificable entre la pregunta de investigación y los demás elementos del proyecto: las variables de la pregunta coinciden con las del título y los objetivos; el alcance de la pregunta no excede ni reduce el declarado en la hipótesis; y el tipo de relación implícita en la pregunta (descriptiva, comparativa, causal) es coherente con la formulación de la hipótesis. No hay variables presentes en la pregunta que estén ausentes en los objetivos ni viceversa. &
& & & \\ \hline
\multicolumn{3}{|r|}{\cellcolor{white}\textbf{Subtotal Dimensión~I (máx.~12 puntos):}} & \multicolumn{4}{c|}{\cellcolor{white}\textbf{/12}} \\ \hline
\end{longtable}
}

\clearpage
\subsection{Especificación de Variables y Operacionalización}
\textit{Se evalúa si la pregunta identifica explícitamente las variables dependientes, independientes y de control, si cada variable posee una definición operacional con protocolo de medición normalizado y si se especifican unidades, rangos, niveles de tratamiento y condiciones de ensayo. Esta dimensión verifica la capacidad de la pregunta para orientar un diseño experimental o metodológico replicable.} Los criterios se aprecian en la Tabla~\ref{tab:rubrica-dim2-variables}.

\scriptsize
\begin{longtable}{|C{0.55cm}|L{2.8cm}|L{8cm}|C{0.3cm}|C{0.3cm}|C{0.3cm}|C{0.3cm}|}
\caption{Rúbrica de evaluación de la Dimensión~II: especificación de variables y operacionalización.}\label{tab:rubrica-dim2-variables}\\
\hline
\rowcolor{headerblue}
\thd{N.\textsuperscript{o}} & \thd{Criterio} & \thd{Indicador verificable} & \thd{3} & \thd{2} & \thd{1} & \thd{0} \\
\hline\endfirsthead
\hline
\rowcolor{headerblue}
\thd{N.\textsuperscript{o}} & \thd{Criterio} & \thd{Indicador verificable} & \thd{3} & \thd{2} & \thd{1} & \thd{0} \\
\hline\endhead
2.1 & Identificación explícita de variables dependientes con unidades y protocolos de medición & La pregunta identifica con claridad la(s) variable(s) dependiente(s): la propiedad, el parámetro o el indicador que se medirá como resultado de la investigación. La variable está asociada a unidades técnicas apropiadas (kPa, MPa, \%, cm/s, mm, kN, entre otras) y a un protocolo de ensayo normalizado o reconocido (ASTM, AASHTO, NTP, ACI, EN, ISO u otra normativa pertinente). No se emplean constructos abstractos sin indicador empírico (p.\,ej., \enquote{desempeño} sin especificar resistencia, deformación, durabilidad u otro parámetro cuantificable). &
& & & \\ \hline
\rowcolor{lightgray}
2.2 & Identificación explícita de variables independientes con niveles o rangos de variación & La pregunta especifica la(s) variable(s) independiente(s) ---el factor, tratamiento, intervención o condición que se manipula u observa--- con sus niveles, rangos o categorías de variación. La especificación es operativa: no se limita a nombrar el factor (p.\,ej., \enquote{dosificación de cemento}) sino que indica los valores concretos (p.\,ej., \enquote{dosificaciones del 3\,\%, 5\,\% y 7\,\% respecto al peso seco}), los tipos a comparar o las condiciones a evaluar. La variable independiente es técnicamente manipulable u observable dentro del diseño de investigación anticipado. &
& & & \\ \hline
2.3 & Definición operacional de los constructos técnicos presentes en la pregunta & Cada constructo técnico presente en la pregunta posee una definición operacional verificable que incluye el procedimiento de medición, el instrumento o equipo asociado y el estándar de calibración cuando corresponda. No se incluyen constructos que carezcan de indicadores empíricos (p.\,ej., \enquote{calidad del concreto} sin especificar si se refiere a resistencia a la compresión, trabajabilidad, porosidad o durabilidad). La operacionalización garantiza que dos investigadores independientes, leyendo la pregunta, identificarían los mismos procedimientos de medición. &
& & & \\ \hline
\rowcolor{lightgray}
2.4 & Especificación de condiciones de contorno, variables de control o parámetros fijos & La pregunta explicita, directa o indirectamente, las condiciones bajo las cuales se evaluarán las variables: condiciones de curado (tiempo, temperatura, humedad), energía de compactación, condiciones de carga, escala de ensayo (laboratorio, campo, piloto), o cualquier otro parámetro que permanecerá constante o controlado. La ausencia de condiciones de contorno no se justifica por brevedad; su omisión impide la replicabilidad y genera indeterminación en el alcance del estudio. &
& & & \\ \hline
\multicolumn{3}{|r|}{\cellcolor{white}\textbf{Subtotal Dimensión~II (máx.~12 puntos):}} & \multicolumn{4}{c|}{\cellcolor{white}\textbf{/12}} \\ \hline
\end{longtable}
\normalsize

\clearpage
\subsection{Coherencia Metodológica y Alineación con el Diseño de Investigación}

\textit{Se evalúa la correspondencia entre la estructura semántica de la pregunta y las capacidades inferenciales del diseño metodológico que anticipa. La pregunta debe ser consistente con el tipo de indagación (descriptiva, relacional, comparativa o explicativa), con la lógica inferencial (cuantitativa, cualitativa o mixta) y con la técnica analítica propuesta. Esta dimensión verifica que la pregunta no anticipe conclusiones causales cuando el diseño es observacional, ni se limite a descripciones cuando el diseño es experimental.} Los criterios se aprecian en la Tabla~\ref{tab:rubrica-dim3-coherencia}.

\scriptsize
\begin{longtable}{|C{0.55cm}|L{2.8cm}|L{8cm}|C{0.3cm}|C{0.3cm}|C{0.3cm}|C{0.3cm}|}
\caption{Rúbrica de evaluación de la Dimensión~III: coherencia metodológica y alineación con el diseño de investigación.}\label{tab:rubrica-dim3-coherencia}\\
\hline
\rowcolor{headerblue}
\thd{N.\textsuperscript{o}} & \thd{Criterio} & \thd{Indicador verificable} & \thd{3} & \thd{2} & \thd{1} & \thd{0} \\
\hline\endfirsthead
\hline
\rowcolor{headerblue}
\thd{N.\textsuperscript{o}} & \thd{Criterio} & \thd{Indicador verificable} & \thd{3} & \thd{2} & \thd{1} & \thd{0} \\
\hline\endhead
3.1 & Correspondencia entre el tipo de pregunta y el diseño metodológico anticipado \epref{Tipo de pregunta vs.\ diseño} & La estructura semántica de la pregunta es coherente con el tipo de diseño que el proyecto propone. Las preguntas comparativas (\enquote{¿qué diferencias existen entre\ldots?}) anticipan diseños experimentales o cuasiexperimentales con grupos de comparación; las preguntas descriptivas (\enquote{¿cuál es la distribución de\ldots?}) anticipan diseños observacionales o de caracterización; las preguntas relacionales (\enquote{¿cuál es la relación entre\ldots?}) anticipan diseños correlacionales; y las preguntas explicativas (\enquote{¿por qué ocurre\ldots?}) anticipan diseños que permitan inferencia causal o exploración de mecanismos. No existe incompatibilidad entre lo que la pregunta solicita y lo que el diseño puede responder. &
& & & \\ \hline
\rowcolor{lightgray}
3.2 & Coherencia entre la escala de medición implícita y la técnica analítica propuesta & La pregunta anticipa, a través de su formulación, una escala de medición (nominal, ordinal, de intervalo o de razón) compatible con la técnica estadística o analítica que el proyecto propone emplear. Si la pregunta indaga sobre \enquote{magnitud de incremento} o \enquote{proporción de variación}, la escala de razón y el análisis paramétrico son pertinentes. Si la pregunta indaga sobre \enquote{clasificación} o \enquote{categoría de desempeño}, la escala ordinal o nominal y las técnicas no paramétricas son apropiadas. No se anticipan análisis estadísticos incompatibles con el tipo de dato que la pregunta generará. &
& & & \\ \hline
3.3 & Alineación entre la unidad de análisis de la pregunta y el nivel de inferencia del diseño & La unidad de análisis implícita en la pregunta (probeta, espécimen, tramo, estructura, proyecto, sistema) es coherente con el nivel en el cual se derivarán las conclusiones. La pregunta no formula interrogantes a nivel de proyecto cuando los datos se recolectan a nivel de probeta sin agregación justificada, ni formula interrogantes a nivel de espécimen cuando la inferencia pretende generalizarse a nivel de sistema. La correspondencia entre la unidad conceptual de la pregunta y la unidad empírica del diseño previene falacias ecológicas o atomísticas. &
& & & \\ \hline
\rowcolor{lightgray}
3.4 & Articulación de la pregunta con subpreguntas o jerarquía de indagación cuando corresponde & Si la complejidad del problema lo requiere, la pregunta principal se descompone en subpreguntas que operacionalizan sus componentes: cada subpregunta se asocia con un ensayo, una medición o una fase del diseño experimental. La jerarquía es lógica (pregunta principal~$\rightarrow$~subpreguntas específicas), las subpreguntas cubren las dimensiones relevantes de la pregunta principal sin redundancias ni omisiones, y su resolución conjunta permite responder la pregunta central. Si el alcance de la investigación no requiere subpreguntas, la pregunta principal es autocontenida y suficiente para orientar el diseño completo. &
& & & \\ \hline
\multicolumn{3}{|r|}{\cellcolor{white}\textbf{Subtotal Dimensión~III (máx.~12 puntos):}} & \multicolumn{4}{c|}{\cellcolor{white}\textbf{/12}} \\ \hline
\end{longtable}
\normalsize

\clearpage
\subsection{Viabilidad Técnica y Delimitación del Alcance}
\textit{Se evalúa si la pregunta es respondible dentro de las restricciones de recursos, tiempo, equipamiento y competencias técnicas del proyecto, y si la delimitación espacial, temporal, poblacional y temática es suficiente para acotar el estudio sin trivializarlo. Esta dimensión verifica que la ambición epistémica de la pregunta sea proporcional a las capacidades reales del proyecto de tesis.} Los criterios se aprecian en la Tabla~\ref{tab:rubrica-dim4-viabilidad}.

\scriptsize
\begin{longtable}{|C{0.55cm}|L{2.8cm}|L{8cm}|C{0.3cm}|C{0.3cm}|C{0.3cm}|C{0.3cm}|}
\caption{Rúbrica de evaluación de la Dimensión~IV: viabilidad técnica y delimitación del alcance.}\label{tab:rubrica-dim4-viabilidad}\\
\hline
\rowcolor{headerblue}
\thd{N.\textsuperscript{o}} & \thd{Criterio} & \thd{Indicador verificable} & \thd{3} & \thd{2} & \thd{1} & \thd{0} \\
\hline\endfirsthead
\hline
\rowcolor{headerblue}
\thd{N.\textsuperscript{o}} & \thd{Criterio} & \thd{Indicador verificable} & \thd{3} & \thd{2} & \thd{1} & \thd{0} \\
\hline\endhead
4.1 & Respondibilidad de la pregunta con los recursos, equipos y competencias disponibles \epref{Viabilidad} & La pregunta es respondible con los recursos técnicos, humanos y presupuestarios declarados en el proyecto. Los ensayos, mediciones o procedimientos que la pregunta exige corresponden a equipos accesibles (prensa hidráulica, CBR, consolidómetro, permeámetro, triaxial, equipo topográfico, software de análisis, entre otros); las competencias técnicas del equipo de investigación son suficientes para ejecutar los protocolos requeridos; y el presupuesto estimado cubre materiales, insumos y servicios. La pregunta no demanda técnicas, equipos o infraestructura que excedan las capacidades institucionales sin que se justifique su acceso mediante convenios o servicios externos documentados. &
& & & \\ \hline
\rowcolor{lightgray}
4.2 & Delimitación espacial, temporal y poblacional de la pregunta & La pregunta delimita el alcance del estudio en sus dimensiones pertinentes: \textbf{espacial} (sitio de estudio, tramo vial, zona de proyecto, localización geográfica); \textbf{temporal} (periodos de curado, horizontes de evaluación, vigencia del estudio); y \textbf{poblacional o de unidad muestral} (tipo de suelo, tipo de concreto, clasificación del material, tipo de estructura, especificación de probetas). La delimitación es operativa y no meramente enunciativa: permite determinar el marco de muestreo y los criterios de inclusión y exclusión de las unidades de análisis. &
& & & \\ \hline
4.3 & Calibración del alcance: ni excesivamente amplia ni trivialmente estrecha (prueba Goldilocks) & La pregunta se sitúa en un punto de equilibrio entre la amplitud excesiva, que impide respuestas concluyentes dentro de un proyecto de tesis de pregrado, y la especificidad trivial, que genera observaciones sin valor teórico ni práctico. La pregunta no abarca múltiples vías causales simultáneas, poblaciones heterogéneas irreconciliables o dominios de resultado dispares (p.\,ej., \enquote{¿cómo se comportan todos los suelos ante todos los estabilizadores?}); ni se restringe a un caso tan particular que sus resultados carezcan de transferibilidad (p.\,ej., \enquote{¿cuál es la resistencia de una única probeta?}). Los hallazgos potenciales son publicables, relevantes y generalizables dentro de condiciones razonables. &
& & & \\ \hline
\rowcolor{lightgray}
4.4 & Adecuación del cronograma implícito en la pregunta al marco temporal del proyecto de tesis & Los tiempos de curado, los periodos de monitoreo, las campañas de campo, los ciclos de ensayo u otros requerimientos temporales implícitos en la pregunta son compatibles con el cronograma del proyecto de tesis de pregrado. La pregunta no exige periodos de evaluación de años cuando el proyecto dispone de meses, ni demanda múltiples campañas estacionales cuando solo una es factible. La escala temporal de la pregunta es realista y ha sido verificada contra el calendario académico y las restricciones logísticas del proyecto. &
& & & \\ \hline
\multicolumn{3}{|r|}{\cellcolor{white}\textbf{Subtotal Dimensión~IV (máx.~12 puntos):}} & \multicolumn{4}{c|}{\cellcolor{white}\textbf{/12}} \\ \hline
\end{longtable}
\normalsize

\clearpage
\subsection{Pertinencia Científica y Aporte al Conocimiento de la Disciplina}

\textit{Se evalúa si la pregunta aborda una brecha, vacío, contradicción o insuficiencia identificable en el estado del arte, si su resolución generará un aporte específico y verificable al conocimiento en Ingeniería Civil, y si posee relevancia tanto para el avance teórico como para la aplicación práctica. Esta dimensión verifica que la pregunta no sea una repetición de investigaciones previas sin valor agregado ni una interrogante desconectada del discurso académico vigente.} Los criterios se aprecian en la Tabla~\ref{tab:rubrica-dim5-pertinencia}.

\scriptsize
\begin{longtable}{|C{0.55cm}|L{2.8cm}|L{8cm}|C{0.3cm}|C{0.3cm}|C{0.3cm}|C{0.3cm}|}
\caption{Rúbrica de evaluación de la Dimensión~V: pertinencia científica y aporte al conocimiento científico-técnico.}\label{tab:rubrica-dim5-pertinencia}\\
\hline
\rowcolor{headerblue}
\thd{N.\textsuperscript{o}} & \thd{Criterio} & \thd{Indicador verificable} & \thd{3} & \thd{2} & \thd{1} & \thd{0} \\
\hline\endfirsthead
\hline
\rowcolor{headerblue}
\thd{N.\textsuperscript{o}} & \thd{Criterio} & \thd{Indicador verificable} & \thd{3} & \thd{2} & \thd{1} & \thd{0} \\
\hline\endhead
5.1 & Identificación de la brecha, vacío o contradicción en el estado del arte que la pregunta busca resolver & La pregunta se sustenta en una brecha específica del conocimiento: un fenómeno no caracterizado en las condiciones locales, un material no evaluado para el uso propuesto, una contradicción empírica entre estudios previos, una técnica no comparada con alternativas convencionales, o un parámetro no cuantificado en el contexto de estudio. La brecha no es genérica (\enquote{falta investigación sobre el tema}) sino que se identifica con precisión la insuficiencia que la pregunta pretende subsanar, trazable a la revisión de literatura o a la problematización. &
& & & \\ \hline
\rowcolor{lightgray}
5.2 & Novedad y diferenciación respecto de investigaciones previas & La pregunta se distingue de estudios previamente realizados en al menos una dimensión sustantiva: el tipo de material o sistema investigado, las condiciones de ensayo o exposición, la escala de estudio, el contexto geográfico o ambiental, la combinación de variables evaluadas, o el enfoque metodológico. La diferenciación no es trivial (p.\,ej., cambiar únicamente la marca comercial de un insumo) sino que aporta información nueva que contribuye al avance del conocimiento de la especialidad. El proyecto no duplica investigaciones existentes sin justificación técnica. &
& & & \\ \hline
5.3 & Relevancia práctica de la respuesta esperada para la Ingeniería Civil & La resolución de la pregunta generará resultados con aplicabilidad verificable en el ejercicio de la Ingeniería Civil: optimización de dosificaciones, validación de métodos constructivos, actualización de especificaciones técnicas, comparación de alternativas de diseño, caracterización de materiales para uso normativo, reducción de incertidumbre en parámetros de ingeniería, o fundamentación de decisiones técnicas en proyectos de infraestructura. La relevancia no es declarativa (\enquote{contribuir al desarrollo de la ingeniería}) sino específica, con destinatarios identificables (diseñadores, constructores, entidades normativas, administradores de infraestructura). &
& & & \\ \hline
\rowcolor{lightgray}
5.4 & Adaptabilidad interdisciplinaria y neutralidad terminológica de la pregunta & La pregunta está formulada con un lenguaje que, sin sacrificar el rigor técnico, permite su comprensión por profesionales de disciplinas afines (geología, ciencias ambientales, ciencia de materiales, arquitectura) y facilita la integración de sus resultados en síntesis más amplias o metaanálisis. Los operadores relacionales empleados (\enquote{modificar}, \enquote{variar}, \enquote{comparar}, \enquote{evaluar}) no presuponen mecanismos causales específicos de una sola tradición teórica. Los supuestos teóricos son explícitos y evaluables por diversas comunidades académicas, lo que favorece el diálogo académico acumulativo. &
& & & \\ \hline
\multicolumn{3}{|r|}{\cellcolor{white}\textbf{Subtotal Dimensión~V (máx.~12 puntos):}} & \multicolumn{4}{c|}{\cellcolor{white}\textbf{/12}} \\ \hline
\end{longtable}
\normalsize

\newpage
En esta sección se presenta el consolidado de resultados por dimensión y la escala de calificación de la pregunta de investigación, resumidos en las tablas~\ref{tab:consolidado-pregunta} y~\ref{tab:escala-calificacion-pregunta}.

\begin{table}[htbp]
\centering
\scriptsize
\caption{Cuadro consolidado de resultados por dimensión de la pregunta de investigación}\label{tab:consolidado-pregunta}
\begin{tabularx}{\textwidth}{|C{0.8cm}|L{7.5cm}|C{1.8cm}|C{1.8cm}|X|}
\hline
\rowcolor{headerblue}
\thd{Dim.} & \thd{Denominación} & \thd{Pts. Máx.} & \thd{Pts. Obt.} & \thd{Obs.} \\ \hline
I & Formulación Estructural y Precisión Semántica de la Pregunta & 12 & & \\ \hline
\rowcolor{lightgray}
II & Especificación de Variables y Operacionalización & 12 & & \\ \hline
III & Coherencia Metodológica y Alineación con el Diseño de Investigación & 12 & & \\ \hline
\rowcolor{lightgray}
IV & Viabilidad Técnica y Delimitación del Alcance & 12 & & \\ \hline
V & Pertinencia Científica y Aporte al Conocimiento & 12 & & \\ \hline
\rowcolor{white}
& \textbf{TOTAL GENERAL} & \textbf{60} & & \\ \hline
\end{tabularx}
\normalsize
\end{table}

\begin{table}[htbp]
\centering
\scriptsize
\caption{Escala de calificación de la pregunta de investigación}\label{tab:escala-calificacion-pregunta}
\begin{tabularx}{\textwidth}{|C{2.5cm}|C{3.2cm}|X|}
\hline
\rowcolor{headerblue}
\thd{Rango} & \thd{Calificación} & \thd{Significado} \\ \hline
\textbf{54--60}\newline(90--100\,\%) &
\textcolor{csuf}{\textbf{PREGUNTA\newline PLENA}} &
La pregunta de investigación es precisa, operacionalmente definida, metodológicamente coherente y científicamente pertinente. Identifica variables con especificaciones completas, anticipa un diseño viable y aborda una brecha verificable en el estado del arte. Constituye un eje articulador sólido entre la problematización, los objetivos, la hipótesis y el diseño de investigación. \\ \hline
\rowcolor{lightgray}
\textbf{42--53}\newline(70--89\,\%) &
\textcolor{cace}{\textbf{PREGUNTA\newline SUFICIENTE}} &
La pregunta presenta los componentes esenciales de una interrogante de investigación con áreas específicas susceptibles de fortalecimiento en cuanto a especificidad operacional, coherencia metodológica o delimitación del alcance. Las debilidades son corregibles sin reformulación fundamental de la pregunta. \\ \hline
\textbf{30--41}\newline(50--69\,\%) &
\textcolor{cdef}{\textbf{PREGUNTA\newline PARCIAL}} &
La pregunta es incompleta, ambigua o predominantemente declarativa. Se evidencian vacíos significativos en la especificación de variables, en la coherencia con el diseño metodológico o en la delimitación del alcance. Requiere fortalecimiento sustancial en su formulación antes de orientar el diseño de investigación. \\ \hline
\rowcolor{lightgray}
\textbf{0--29}\newline($<$50\,\%) &
\textcolor{cins}{\textbf{PREGUNTA\newline INSUFICIENTE}} &
La pregunta no demuestra consistencia verificable como eje de investigación. Carece de variables identificables, no se distingue de un problema de aplicación profesional, presenta incompatibilidad con el diseño propuesto o no aborda una brecha en el conocimiento. Se requiere una reformulación fundamental de la pregunta de investigación. \\ \hline
\end{tabularx}
\normalsize
\end{table}

\clearpage
\section{Evaluación de las Variables, Dimensiones e Indicadores}

{\footnotesize\setlength{\baselineskip}{0.85\baselineskip}
La definición, clasificación y operacionalización de variables es un requisito metodológico en las tesis de pregrado en Ingeniería Civil, pues vincula los conceptos teóricos con las mediciones empíricas que permiten contrastar las hipótesis. Una variable es cualquier propiedad, atributo o característica que puede asumir distintos valores o categorías en un dominio de observación; su delimitación rigurosa organiza el diseño experimental, la recolección de datos y la comparación de resultados. Las dimensiones agrupan las variables en categorías conceptualmente coherentes —técnicas, económicas, ambientales, sociales, microestructurales, entre otras— que representan distintos aspectos del fenómeno estudiado, mientras que los indicadores concretan cada variable en resultados verificables mediante procedimientos normalizados. Esta rúbrica evalúa la calidad, consistencia y completitud del sistema \textbf{Variable}~$\rightarrow$~\textbf{Dimensión}~$\rightarrow$~\textbf{Indicador} del proyecto de tesis, verificando la precisión en la identificación de variables, la correcta clasificación según función y escala de medición, la coherencia y organización jerárquica de la estructura dimensional, el respaldo de los indicadores en procedimientos estandarizados de medición y la consistencia interna y suficiencia del sistema para la verificación empírica de las hipótesis. La evaluación se aplica a cualquier subárea de Ingeniería Civil—geotecnia, estructuras, pavimentos, hidráulica, materiales, gestión de la construcción, saneamiento, entre otras—y se centra en la consistencia interna, la precisión técnica y la suficiencia operacional del sistema de variables como base del diseño metodológico.Las dimensiones de evaluación del sistema se resumen en la Tabla~\ref{tab:dim-variables}, y la escala de valoración y sus descriptores operativos en la Tabla~\ref{tab:escala-variables}\par}
\vspace{-12pt}
\begin{table}[htbp]
\centering
\scriptsize
\caption{Dimensiones de evaluación del sistema de variables, dimensiones e indicadores}
\label{tab:dim-variables}
\begin{tabularx}{\textwidth}{|C{0.6cm}|L{3cm}|X|}
\hline
\rowcolor{headerblue}
\thd{Dim.} & \thd{Denominación} & \thd{Alcance} \\ \hline
I & Identificación y definición conceptual de las variables & Verifica que las variables del proyecto estén identificadas con claridad, que posean definiciones conceptuales y operacionales diferenciadas, que su naturaleza corresponda a propiedades observables y medibles, y que exista correspondencia con los objetivos y las hipótesis de la investigación. \\ \hline
\rowcolor{lightgray}
II & Clasificación y función de las variables en el diseño de investigación & Evalúa la correcta clasificación de las variables según su función (independiente, dependiente, interviniente), su naturaleza (cuantitativa continua, cuantitativa discreta, cualitativa) y su escala de medición (nominal, ordinal, de intervalo, de razón), verificando la coherencia entre la clasificación declarada y el tratamiento estadístico propuesto. \\ \hline
III & Estructura dimensional y jerarquía analítica de las variables & Examina la descomposición de las variables en dimensiones analíticas pertinentes, la coherencia conceptual de las agrupaciones propuestas, la identificación de dimensiones simples, compuestas o multifactoriales según corresponda, y la organización jerárquica que vincula el nivel teórico con el nivel empírico. \\ \hline
\rowcolor{lightgray}
IV & Operacionalización mediante indicadores verificables & Evalúa la traducción de cada variable y dimensión en indicadores medibles, la especificación de procedimientos normalizados de medición, la pertinencia de las normas técnicas citadas, la definición de unidades, rangos y criterios de aceptación, y la factibilidad técnica de las mediciones propuestas. \\ \hline
V & Coherencia y articulación del sistema variable--dimensión--indicador & Verifica la consistencia interna del sistema completo, la trazabilidad desde las hipótesis hasta los indicadores, la ausencia de variables huérfanas o indicadores desarticulados, la correspondencia con la matriz de consistencia y la suficiencia del sistema para la verificación empírica de las hipótesis. \\ \hline
\end{tabularx}
\end{table}

\vspace{-12pt}
\begin{table}[htbp]
\centering
\scriptsize
\caption{Escala de valoración para evaluar el sistema de variables, dimensiones e indicadores.}
\label{tab:escala-variables}
\begin{tabularx}{\textwidth}{|C{3cm}|C{1.6cm}|X|}
\hline
\rowcolor{headerblue}
\thd{Nivel} & \thd{Pts.} & \thd{Descriptor operativo} \\ \hline
\textcolor{csuf}{\textbf{Excelente}} & \textbf{3} & El proyecto satisface plenamente el criterio evaluado. Las variables están definidas con precisión, clasificadas correctamente, descompuestas en dimensiones coherentes y operacionalizadas mediante indicadores verificables respaldados por procedimientos normalizados. No se identifican ambigüedades, omisiones ni inconsistencias. \\ \hline
\rowcolor{lightgray}
\textcolor{cace}{\textbf{Aceptable}} & \textbf{2} & El criterio es identificable pero se cumple de manera parcial o incompleta. El sistema de variables aborda el aspecto evaluado sin la profundidad, la precisión clasificatoria o el respaldo normativo suficiente. Susceptible de fortalecimiento con correcciones específicas. \\ \hline
\textcolor{cdef}{\textbf{Deficiente}} & \textbf{1} & El cumplimiento del criterio es tangencial o meramente declarativo. La definición, clasificación u operacionalización es superficial, carece de sustento normativo verificable o presenta inconsistencias que debilitan la estructura del sistema de variables. \\ \hline
\rowcolor{lightgray}
\textcolor{cins}{\textbf{Ausente}} & \textbf{0} & El proyecto no evidencia consideración alguna del criterio evaluado. Existe omisión total del aspecto evaluado o una desconexión que invalida el sistema de variables en ese componente. \\ \hline
\end{tabularx}
\end{table}

\clearpage
\subsection{Identificación y Definición Conceptual de las Variables}

\textit{Se evalúa si el proyecto identifica las variables de investigación con claridad, si las definiciones conceptuales y operacionales están diferenciadas, si la naturaleza de las variables corresponde a propiedades observables y medibles dentro del dominio de la disciplina de la Ingeniería Civil, y si existe correspondencia verificable entre las variables declaradas y los objetivos e hipótesis del proyecto.} Véase la Tabla~\ref{tab:rubrica-ident-def-variables}.
{\scriptsize

\begin{longtable}{|C{0.55cm}|L{2cm}|L{8cm}|C{0.4cm}|C{0.4cm}|C{0.4cm}|C{0.4cm}|}
\caption{Rúbrica de evaluación de la Dimensión~I: identificación y definición conceptual de las variables.}\label{tab:rubrica-ident-def-variables}\\
\hline
\rowcolor{headerblue}
\thd{N.\textsuperscript{o}} & \thd{Criterio} & \thd{Indicador verificable} & \thd{3} & \thd{2} & \thd{1} & \thd{0} \\
\hline\endfirsthead
\hline
\rowcolor{headerblue}
\thd{N.\textsuperscript{o}} & \thd{Criterio} & \thd{Indicador verificable} & \thd{3} & \thd{2} & \thd{1} & \thd{0} \\
\hline\endhead
1.1 & Identificación explícita de las variables de investigación vinculadas a los objetivos e hipótesis & El proyecto identifica de manera explícita todas las variables involucradas en la investigación, estableciendo una correspondencia verificable con los objetivos específicos y las hipótesis formuladas. Cada variable mencionada en las hipótesis aparece declarada en la sección de variables y viceversa; no existen variables implícitas en los objetivos que carezcan de identificación formal, ni variables declaradas que no se vinculen a ningún objetivo o hipótesis. La identificación distingue con claridad entre la variable central del estudio y las variables asociadas del sistema. &
& & & \\ \hline
\rowcolor{lightgray}
1.2 & Definición conceptual precisa y diferenciada de la definición operacional & Cada variable posee una definición conceptual que describe su significado teórico dentro del dominio del área de la Ingeniería Civil (p.\,ej., \enquote{la resistencia al corte es la tensión máxima que un material soporta en un plano de falla antes de la rotura, descrita por el criterio de Mohr-Coulomb}) y una definición operacional que especifica cómo se medirá en la investigación (p.\,ej., \enquote{determinada mediante ensayo triaxial CU según ASTM D4767}). Ambas definiciones son explícitas, diferenciadas y no se confunden entre sí. No se presentan definiciones genéricas que podrían aplicarse a cualquier disciplina sin especificidad técnica. &
& & & \\ \hline
1.3 & Naturaleza observable y medible de las variables dentro del dominio de la Ingeniería Civil & Las variables identificadas corresponden a propiedades, atributos o características susceptibles de observación directa o medición mediante instrumentación y procedimientos reconocidos por la comunidad profesional. Las variables técnicas se expresan mediante magnitudes con unidades dimensionales definidas (resistencia en MPa, permeabilidad en cm/s, deformación en mm, caudal en m$^3$/s, carga en kN, entre otras). Cuando se incluyen constructos abstractos como durabilidad, sostenibilidad o trabajabilidad, estos se acompañan de la descomposición en variables directamente medibles que los operacionalizan. No se presentan como variables conceptos vagos o inmedibles. &
& & & \\ \hline
\rowcolor{lightgray}
1.4 & Consistencia entre las variables declaradas y el título, el problema y el alcance del proyecto & Las variables identificadas guardan correspondencia directa con el título del proyecto, el problema de investigación formulado y el alcance declarado. No existen variables en el título ausentes en la sección de variables, ni variables declaradas que excedan el alcance del problema. La unidad de análisis (tipo de suelo, material, estructura, mezcla, sistema constructivo, tramo vial, cuenca, red de saneamiento, entre otros) está reflejada coherentemente en la definición de las variables. &
& & & \\ \hline
\multicolumn{3}{|r|}{\cellcolor{white}\textbf{Subtotal Dimensión~I (máx.~12 puntos):}} & \multicolumn{4}{c|}{\cellcolor{white}\textbf{/12}} \\ \hline
\end{longtable}
}

\clearpage
\subsection{Clasificación y Función de las Variables en el Diseño de Investigación}

\textit{Se evalúa si las variables están correctamente clasificadas según su función metodológica (independiente, dependiente, interviniente), su naturaleza (cuantitativa continua, cuantitativa discreta, cualitativa) y su escala de medición (nominal, ordinal, de intervalo, de razón), verificando que la clasificación declarada sea coherente con el rol que cada variable desempeña en el diseño experimental y con el tratamiento estadístico propuesto.} Los criterios se aprecian en la Tabla~\ref{tab:rubrica-clasif-funcion-variables}.
\scriptsize
\begin{longtable}{|C{0.55cm}|L{2.8cm}|L{8cm}|C{0.3cm}|C{0.3cm}|C{0.3cm}|C{0.3cm}|}
\caption{Rúbrica de evaluación de la Dimensión~II: clasificación y función de las variables en el diseño de investigación.}\label{tab:rubrica-clasif-funcion-variables}\\
\hline
\rowcolor{headerblue}
\thd{N.\textsuperscript{o}} & \thd{Criterio} & \thd{Indicador verificable} & \thd{3} & \thd{2} & \thd{1} & \thd{0} \\
\hline\endfirsthead
\hline
\rowcolor{headerblue}
\thd{N.\textsuperscript{o}} & \thd{Criterio} & \thd{Indicador verificable} & \thd{3} & \thd{2} & \thd{1} & \thd{0} \\
\hline\endhead
2.1 & Clasificación correcta según la función metodológica: independiente, dependiente e interviniente & Las variables independientes corresponden a los factores o condiciones manipuladas o seleccionadas activamente por el investigador (p.\,ej., dosificación de un aditivo, tipo de refuerzo, método constructivo, espesor de capa). Las variables dependientes representan las respuestas medidas del sistema (p.\,ej., resistencia a compresión, asentamiento, caudal de infiltración, deflexión). Las variables intervinientes identifican condiciones que pueden influir sistemáticamente en la relación entre las independientes y las dependientes (p.\,ej., temperatura de curado, mineralogía del material, condiciones ambientales). La clasificación es coherente con el diseño experimental o metodológico propuesto y no presenta inversiones de rol ni omisiones de variables intervinientes relevantes. &
& & & \\ \hline
\rowcolor{lightgray}
2.2 & Clasificación según la naturaleza de la variable: cuantitativa continua, cuantitativa discreta y cualitativa & Las variables cuantitativas continuas se identifican correctamente como aquellas representables en la recta real (resistencia en MPa, deformación en mm, permeabilidad en cm/s). Las variables cuantitativas discretas se reconocen como restringidas a valores enteros o contables (número de ciclos de carga, número de fisuras, cantidad de capas). Las variables cualitativas se identifican como aquellas que expresan categorías o atributos (tipo de suelo según SUCS, método de compactación, clasificación visual del deterioro). La clasificación según naturaleza es consistente en todo el documento y no se presentan variables cualitativas tratadas como cuantitativas ni viceversa. &
& & & \\ \hline
2.3 & Identificación correcta de la escala de medición de cada variable & Cada variable tiene declarada su escala de medición de forma coherente con su naturaleza: \textbf{nominal} para categorías sin orden inherente (tipo de cemento, clasificación SUCS); \textbf{ordinal} para categorías con jerarquía sin distancias uniformes (grado de deterioro: leve, moderado, severo); \textbf{de intervalo} para magnitudes con distancias constantes pero sin cero absoluto (temperatura en $^\circ$C); \textbf{de razón} para magnitudes con cero absoluto que admiten el cálculo de razones (resistencia en MPa, longitud en m, masa en kg). No se asignan escalas de razón a variables que carecen de cero absoluto ni escalas nominales a variables que admiten ordenamiento. &
& & & \\ \hline
\rowcolor{lightgray}
2.4 & Coherencia entre la clasificación de las variables y el tratamiento estadístico o analítico propuesto & Los métodos de análisis estadístico propuestos son compatibles con la escala de medición y la naturaleza de las variables: pruebas paramétricas (ANOVA, regresión lineal, prueba $t$) se aplican a variables cuantitativas de intervalo o razón con distribución verificable; pruebas no paramétricas (Kruskal-Wallis, Mann-Whitney, chi-cuadrado) se reservan para variables ordinales o cualitativas. Los diseños factoriales declarados son consistentes con el número de niveles de las variables independientes. No se proponen análisis de varianza sobre variables nominales ni correlaciones de Pearson sobre escalas ordinales. &
& & & \\ \hline
\multicolumn{3}{|r|}{\cellcolor{white}\textbf{Subtotal Dimensión~II (máx.~12 puntos):}} & \multicolumn{4}{c|}{\cellcolor{white}\textbf{/12}} \\ \hline
\end{longtable}
\normalsize

\clearpage
\subsection{Estructura Dimensional y Jerarquía Analítica de las Variables}

\textit{Se evalúa si las variables se descomponen en dimensiones analíticas pertinentes que reflejan los diferentes aspectos del fenómeno investigado, si la estructura jerárquica variable--dimensión es conceptualmente coherente, y si la tipología dimensional (simple, compuesta, multifactorial) es apropiada para la complejidad del problema abordado.} Los criterios se aprecian en la Tabla~\ref{tab:rubrica-estructura-dimensional}.
\scriptsize
\begin{longtable}{|C{0.55cm}|L{2.8cm}|L{8cm}|C{0.3cm}|C{0.3cm}|C{0.3cm}|C{0.3cm}|}
\caption{Rúbrica de evaluación de la Dimensión~III: estructura dimensional y jerarquía analítica de las variables.}\label{tab:rubrica-estructura-dimensional}\\
\hline
\rowcolor{headerblue}
\thd{N.\textsuperscript{o}} & \thd{Criterio} & \thd{Indicador verificable} & \thd{3} & \thd{2} & \thd{1} & \thd{0} \\
\hline\endfirsthead
\hline
\rowcolor{headerblue}
\thd{N.\textsuperscript{o}} & \thd{Criterio} & \thd{Indicador verificable} & \thd{3} & \thd{2} & \thd{1} & \thd{0} \\
\hline\endhead
3.1 & Descomposición de las variables en dimensiones analíticas pertinentes al fenómeno investigado & Las variables del proyecto se descomponen en dimensiones que reflejan aspectos distinguibles del fenómeno: dimensiones mecánicas (resistencia, rigidez, tenacidad), hidráulicas (permeabilidad, retención, erosión), constructivas (trabajabilidad, rendimiento, plazo), económicas (costos directos, ciclo de vida), ambientales (emisiones, toxicidad, valorización de residuos), sociales (aceptabilidad, empleo local), u otras pertinentes según la subárea de la Ingeniería Civil. Las dimensiones identificadas son relevantes para el problema y no constituyen categorías arbitrarias o redundantes. Cada dimensión aporta información distinguible y necesaria para la comprensión del fenómeno. &
& & & \\ \hline
\rowcolor{lightgray}
3.2 & Coherencia conceptual de las agrupaciones dimensionales y ausencia de solapamientos & Las variables agrupadas dentro de cada dimensión comparten una base conceptual, física o funcional común que justifica su pertenencia a la misma categoría. No se agrupan en una misma dimensión variables conceptualmente inconexas (p.\,ej., resistencia a compresión y costo de transporte en una dimensión \enquote{técnica}), ni se fragmentan artificialmente variables estrechamente relacionadas en dimensiones separadas sin justificación. Cuando una variable informa sobre múltiples dimensiones (p.\,ej., la durabilidad como indicador técnico y económico simultáneamente), esta transversalidad se reconoce y se justifica explícitamente. &
& & & \\ \hline
3.3 & Identificación de la tipología dimensional: simple, compuesta o multifactorial & El proyecto reconoce la complejidad estructural de las dimensiones empleadas: las dimensiones simples se caracterizan mediante un conjunto reducido de variables directamente relacionadas (p.\,ej., resistencia mecánica descrita por parámetros de corte y compresión); las dimensiones compuestas integran subdimensiones distinguibles pero interdependientes (p.\,ej., durabilidad que combina resistencia a ciclos húmedo-seco, estabilidad química y resistencia a erosión); las dimensiones multifactoriales integran variables heterogéneas que requieren normalización y agregación ponderada (p.\,ej., sostenibilidad que integra criterios técnicos, económicos, ambientales y sociales). La tipología es coherente con la complejidad del problema investigado. &
& & & \\ \hline
\rowcolor{lightgray}
3.4 & Jerarquía analítica que vincula el nivel teórico con el nivel empírico & La estructura dimensional permite transitar desde el nivel conceptual abstracto (variable central del estudio) hacia niveles progresivamente más concretos y medibles (dimensiones, subdimensiones y, finalmente, indicadores cuantificables). La jerarquía es explícita y trazable: cada dimensión se vincula hacia arriba con la variable que descompone y hacia abajo con los indicadores que la operacionalizan. No existen dimensiones desconectadas de la variable a la que pertenecen ni niveles jerárquicos vacíos que interrumpan la cadena de traducción conceptual-empírica. &
& & & \\ \hline
\multicolumn{3}{|r|}{\cellcolor{white}\textbf{Subtotal Dimensión~III (máx.~12 puntos):}} & \multicolumn{4}{c|}{\cellcolor{white}\textbf{/12}} \\ \hline
\end{longtable}
\normalsize

\clearpage
\subsection{Operacionalización Mediante Indicadores Verificables}

\textit{Se evalúa si cada variable y dimensión se traduce en indicadores medibles, si los procedimientos de medición están respaldados por normas técnicas reconocidas, si las unidades, rangos y criterios de aceptación están definidos, y si las mediciones propuestas son técnicamente factibles dentro de las restricciones del proyecto.} Los criterios se aprecian en la Tabla~\ref{tab:rubrica-operacionalizacion-indicadores}.
\scriptsize
\begin{longtable}{|C{0.55cm}|L{2.8cm}|L{8cm}|C{0.3cm}|C{0.3cm}|C{0.3cm}|C{0.3cm}|}
\caption{Rúbrica de evaluación de la Dimensión~IV: operacionalización mediante indicadores verificables.}\label{tab:rubrica-operacionalizacion-indicadores}\\
\hline
\rowcolor{headerblue}
\thd{N.\textsuperscript{o}} & \thd{Criterio} & \thd{Indicador verificable} & \thd{3} & \thd{2} & \thd{1} & \thd{0} \\
\hline\endfirsthead
\hline
\rowcolor{headerblue}
\thd{N.\textsuperscript{o}} & \thd{Criterio} & \thd{Indicador verificable} & \thd{3} & \thd{2} & \thd{1} & \thd{0} \\
\hline\endhead
4.1 & Traducción de cada variable en indicadores medibles con procedimientos normalizados de medición & Cada variable declarada en el proyecto se operacionaliza mediante al menos un indicador cuantificable cuyo procedimiento de obtención está especificado. Los indicadores técnicos se determinan mediante normas reconocidas (ASTM, ISO, AASHTO, NTP, Eurocódigos, o normativa local aplicable): resistencia a compresión según ASTM D2166 o D1633, CBR según ASTM D1883, permeabilidad según ASTM D5084, deflexión según ASTM D4694, resistencia del concreto según ASTM C39, entre otros. No se presentan variables sin indicador asociado ni indicadores sin procedimiento de medición identificable. Los indicadores cualitativos, cuando se emplean, cuentan con escalas de valoración definidas y criterios de clasificación explícitos. &
& & & \\ \hline
\rowcolor{lightgray}
4.2 & Especificación de unidades, rangos, condiciones de ensayo y criterios de aceptación & Los indicadores cuantitativos incluyen las unidades dimensionales de expresión (MPa, cm/s, kN, mm, $\%$, m$^3$/s, entre otras), los rangos esperados o los valores de referencia pertinentes, y las condiciones de ensayo relevantes (velocidad de carga, tiempo de curado, temperatura, gradiente hidráulico, presión de confinamiento, entre otros). Cuando corresponde, se establecen criterios de aceptación basados en normativa o especificaciones técnicas del proyecto (p.\,ej., CBR~$\geq$~12\,\%, UCS~$\geq$~2,0~MPa, asentamiento~$\leq$~25~mm, resistencia del concreto~$\geq$~21~MPa). Las condiciones de ensayo son coherentes con las condiciones de servicio anticipadas. &
& & & \\ \hline
4.3 & Pertinencia y sensibilidad de los indicadores seleccionados para el fenómeno investigado & Los indicadores seleccionados son relevantes para el objetivo de la investigación, sensibles a los cambios que el estudio pretende detectar y específicos para distinguir entre las condiciones experimentales o los tratamientos comparados. No se emplean indicadores de baja sensibilidad como medida principal de un fenómeno que requiere alta resolución (p.\,ej., el límite líquido como indicador principal de la efectividad de un tratamiento cuando la UCS o el módulo resiliente ofrecen mayor sensibilidad). Los indicadores cubren las dimensiones relevantes del fenómeno sin redundancias innecesarias ni omisiones que comprometan la evaluación integral del problema investigado. &
& & & \\ \hline
\rowcolor{lightgray}
4.4 & Factibilidad técnica y económica de las mediciones propuestas & Las mediciones propuestas son realizables con el equipamiento, la infraestructura de laboratorio o campo, el personal técnico y el presupuesto disponibles para el proyecto de tesis. No se proponen indicadores que requieran instrumentación inaccesible (p.\,ej., microscopía electrónica de barrido, tomografía computarizada, columna resonante) sin verificar su disponibilidad, ni ensayos cuyo costo o duración exceda las restricciones del proyecto. Cuando se proponen mediciones de alta complejidad, se justifica su necesidad y se confirma la viabilidad de su ejecución. El programa de ensayos es compatible con el cronograma y los recursos declarados. &
& & & \\ \hline
\multicolumn{3}{|r|}{\cellcolor{white}\textbf{Subtotal Dimensión~IV (máx.~12 puntos):}} & \multicolumn{4}{c|}{\cellcolor{white}\textbf{/12}} \\ \hline
\end{longtable}
\normalsize

\clearpage
\subsection{Coherencia y Articulación del Sistema Variable--Dimensión--Indicador}

\textit{Se evalúa la consistencia interna del sistema completo variable--dimensión--indicador, la trazabilidad desde las hipótesis hasta los indicadores, la correspondencia con la matriz de consistencia, y la suficiencia del sistema para permitir la verificación empírica de las hipótesis y la consecución de los objetivos de la investigación.} Los criterios se aprecian en la Tabla~\ref{tab:rubrica-coherencia-articulacion}.
\scriptsize
\begin{longtable}{|C{0.55cm}|L{2.8cm}|L{8cm}|C{0.3cm}|C{0.3cm}|C{0.3cm}|C{0.3cm}|}
\caption{Rúbrica de evaluación de la Dimensión~V: coherencia y articulación del sistema variable--dimensión--indicador.}\label{tab:rubrica-coherencia-articulacion}\\
\hline
\rowcolor{headerblue}
\thd{N.\textsuperscript{o}} & \thd{Criterio} & \thd{Indicador verificable} & \thd{3} & \thd{2} & \thd{1} & \thd{0} \\
\hline\endfirsthead
\hline
\rowcolor{headerblue}
\thd{N.\textsuperscript{o}} & \thd{Criterio} & \thd{Indicador verificable} & \thd{3} & \thd{2} & \thd{1} & \thd{0} \\
\hline\endhead
5.1 & Trazabilidad completa desde las hipótesis hasta los indicadores de medición & Cada hipótesis del proyecto es contrastable mediante los indicadores declarados: las variables que intervienen en la hipótesis están definidas, descompuestas en dimensiones pertinentes y operacionalizadas mediante indicadores medibles con procedimientos específicos. La cadena hipótesis~$\rightarrow$~variable~$\rightarrow$~dimensión~$\rightarrow$~indicador~$\rightarrow$~procedimiento de medición es explícita y completa para cada hipótesis. No existen hipótesis cuyas variables carezcan de indicadores, ni indicadores que no se vinculen a ninguna hipótesis o pregunta de investigación. &
& & & \\ \hline
\rowcolor{lightgray}
5.2 & Ausencia de variables huérfanas, dimensiones vacías o indicadores desarticulados & El sistema no contiene variables declaradas que carezcan de dimensiones asociadas ni de indicadores que las operacionalicen (variables huérfanas). No existen dimensiones enunciadas que carezcan de indicadores medibles (dimensiones vacías). No se presentan indicadores que no se vinculen a ninguna dimensión ni variable del proyecto (indicadores desarticulados). Todos los elementos del sistema cumplen una función específica y necesaria en la cadena de traducción conceptual-empírica. La revisión del cuadro de operacionalización no revela filas incompletas, celdas vacías ni elementos sin conexión lógica con el resto del sistema. &
& & & \\ \hline
5.3 & Correspondencia verificable con la matriz de consistencia del proyecto & El sistema de variables, dimensiones e indicadores es coherente con la matriz de consistencia: cada objetivo específico tiene variables asociadas claramente definidas, cada variable cuenta con indicadores respaldados por procedimientos normalizados, y los métodos de análisis propuestos son compatibles con la naturaleza y la escala de medición de los datos a recopilar. Las variables independientes y dependientes declaradas en la matriz de consistencia coinciden con las presentadas en la sección de variables. No hay discrepancias entre ambos instrumentos metodológicos. &
& & & \\ \hline
\rowcolor{lightgray}
5.4 & Suficiencia del sistema para la verificación empírica de las hipótesis y el logro de los objetivos & El conjunto de indicadores seleccionados proporciona información necesaria y suficiente para contrastar cada hipótesis del proyecto y alcanzar cada objetivo específico. Los datos que se obtendrán mediante los indicadores declarados permiten la aplicación de los análisis estadísticos propuestos con poder suficiente para detectar efectos de magnitud prácticamente significativa. El sistema no presenta omisiones que obliguen a incorporar mediciones no previstas durante la fase experimental, ni incluye mediciones superfluas que no contribuyan a ningún objetivo. La cobertura del sistema es completa respecto a las preguntas de investigación formuladas. &
& & & \\ \hline
\multicolumn{3}{|r|}{\cellcolor{white}\textbf{Subtotal Dimensión~V (máx.~12 puntos):}} & \multicolumn{4}{c|}{\cellcolor{white}\textbf{/12}} \\ \hline
\end{longtable}
\normalsize
\newpage
En esta sección se presenta el consolidado de resultados por dimensión y la escala de calificación del sistema de variables, dimensiones e indicadores, resumidos en las tablas~\ref{tab:consolidado-variables} y~\ref{tab:escala-calificacion-variables}.
\begin{table}[htbp]
\centering
\scriptsize
\caption{Cuadro consolidado de resultados por dimensión del sistema de variables, dimensiones e indicadores.}\label{tab:consolidado-variables}
\begin{tabularx}{\textwidth}{|C{0.8cm}|L{7.5cm}|C{1.8cm}|C{1.8cm}|X|}
\hline
\rowcolor{headerblue}
\thd{Dim.} & \thd{Denominación} & \thd{Pts. Máx.} & \thd{Pts. Obt.} & \thd{Obs.} \\ \hline
I & Identificación y Definición Conceptual de las Variables & 12 & & \\ \hline
\rowcolor{lightgray}
II & Clasificación y Función de las Variables en el Diseño de Investigación & 12 & & \\ \hline
III & Estructura Dimensional y Jerarquía Analítica de las Variables & 12 & & \\ \hline
\rowcolor{lightgray}
IV & Operacionalización Mediante Indicadores Verificables & 12 & & \\ \hline
V & Coherencia y Articulación del Sistema Variable--Dimensión--Indicador & 12 & & \\ \hline
\rowcolor{white}
& \textbf{TOTAL GENERAL} & \textbf{60} & & \\ \hline
\end{tabularx}
\normalsize
\end{table}
\begin{table}[htbp]
\centering
\scriptsize
\caption{Escala de calificación del sistema de variables, dimensiones e indicadores.}\label{tab:escala-calificacion-variables}
\begin{tabularx}{\textwidth}{|C{2.5cm}|C{4.2cm}|X|}
\hline
\rowcolor{headerblue}
\thd{Rango} & \thd{Calificación} & \thd{Significado} \\ \hline
\textbf{54--60}\newline(90--100\,\%) &
\textcolor{csuf}{\textbf{SISTEMA DE VARIABLES\newline PLENO}} &
El sistema de variables, dimensiones e indicadores es integral, técnicamente preciso y lógicamente articulado. Las variables están definidas con rigor, clasificadas correctamente, descompuestas en dimensiones coherentes y operacionalizadas mediante indicadores verificables respaldados por normas técnicas. La trazabilidad hipótesis--variable--dimensión--indicador es completa y la cobertura del sistema es suficiente para la verificación empírica de las hipótesis. \\ \hline
\rowcolor{lightgray}
\textbf{42--53}\newline(70--89\,\%) &
\textcolor{cace}{\textbf{SISTEMA DE VARIABLES\newline SUFICIENTE}} &
El sistema de variables es sustancial y presenta los componentes esenciales de la cadena de operacionalización, con áreas específicas susceptibles de fortalecimiento en cuanto a precisión clasificatoria, respaldo normativo de indicadores o articulación con la matriz de consistencia. Las debilidades son corregibles sin reestructuración fundamental del sistema. \\ \hline
\textbf{30--41}\newline(50--69\,\%) &
\textcolor{cdef}{\textbf{SISTEMA DE VARIABLES\newline PARCIAL}} &
El sistema de variables es incompleto o predominantemente declarativo. Se evidencian vacíos significativos en la definición, clasificación u operacionalización de variables, indicadores sin respaldo normativo, dimensiones enunciadas pero no desarrolladas, o desconexiones entre las variables declaradas y las hipótesis. Requiere fortalecimiento sustancial. \\ \hline
\rowcolor{lightgray}
\textbf{0--29}\newline($<$50\,\%) &
\textcolor{cins}{\textbf{SISTEMA DE VARIABLES\newline INSUFICIENTE}} &
El sistema de variables no demuestra consistencia verificable. Las variables carecen de definiciones operacionales, la clasificación es ausente o incorrecta, las dimensiones no se identifican, los indicadores no se vinculan a procedimientos de medición, o la articulación con las hipótesis y la matriz de consistencia es inexistente. Se requiere una reformulación fundamental del sistema de variables. \\ \hline
\end{tabularx}
\normalsize
\end{table}

\clearpage
\section{Revisión del Marco Teórico y Conceptual}

{\footnotesize
El marco teórico constituye la estructura conceptual que fundamenta, organiza y delimita el conocimiento existente sobre el cual se sustenta una investigación. En un proyecto de tesis de pregrado en Ingeniería Civil, el marco teórico no se puede concebir como un resumen pasivo de la literatura ni como una recopilación enciclopédica de definiciones, sino como un modelo estructurado que identifica las teorías, los principios, las variables y las relaciones que guían la formulación de hipótesis, la selección metodológica y la interpretación de resultados. Su función epistemológica es proporcionar el \textit{plano conceptual} de la investigación: así como un plano estructural dicta las cargas, los apoyos y las secciones de un elemento, el marco teórico dicta los constructos, los supuestos y los vínculos lógicos que sostienen el diseño investigativo. Esta rúbrica evalúa la calidad, coherencia y suficiencia del marco teórico de un \textbf{proyecto de tesis} de pregrado --- documento que, de ser aprobado, recién se convertirá en una tesis --- verificando que la fundamentación teórica sea pertinente al problema formulado, que las fuentes sean actuales y confiables, que las variables estén conceptualmente definidas, que exista alineación verificable entre el marco teórico y los demás componentes del proyecto, y que el nivel de profundidad sea el apropiado para una investigación de pregrado en Ingeniería Civil. La evaluación se sustenta en criterios derivados de marcos de calidad reconocidos internacionalmente (Lovitts, 2007; Boote \& Beile, 2005; Grant \& Osanloo, 2014), estándares de acreditación (ABET, EUR-ACE) y estudios específicos de evaluación de tesis en programas de Ingeniería Civil latinoamericanos (García-Ramírez, 2023; Del Savio et~al., 2024; Oyewobi et~al., 2024). Las dimensiones de evaluación del marco teórico se resumen en la Tabla~\ref{tab:dim-marco-teorico}, mientras que la escala de valoración y sus descriptores operativos se presentan en la Tabla~\ref{tab:escala-marco-teorico}.\par
}
\vspace{-8pt}
\begin{table}[htbp]
\centering
\scriptsize
\caption{Dimensiones de evaluación del marco teórico y alcance de cada dimensión.}
\label{tab:dim-marco-teorico}
\begin{tabularx}{\textwidth}{|C{0.8cm}|L{4.5cm}|X|}
\hline
\rowcolor{headerblue}
\thd{Dim.} & \thd{Denominación} & \thd{Alcance} \\ \hline
I & Fundamentación teórica y selección de teorías o modelos & Verifica que el proyecto identifique, seleccione y justifique las teorías, los principios o los modelos de la Ingeniería Civil que fundamentan la investigación, y que dicha selección sea pertinente al problema formulado y al nivel de pregrado. \\ \hline
\rowcolor{lightgray}
II & Coherencia interna del marco teórico con el diseño de investigación & Evalúa la alineación verificable entre el marco teórico y los demás componentes del proyecto: problema, objetivos, hipótesis y metodología anticipada, verificando que no existan desconexiones, variables huérfanas ni contradicciones internas. \\ \hline
III & Revisión de literatura y síntesis crítica & Examina la calidad, actualidad, diversidad y tratamiento crítico de las fuentes bibliográficas que sustentan el marco teórico, verificando que la revisión trascienda la descripción secuencial y alcance un nivel de síntesis e identificación de vacíos coherente con el pregrado. \\ \hline
\rowcolor{lightgray}
IV & Claridad conceptual y definición de variables & Verifica que los constructos, las variables y los términos técnicos fundamentales estén definidos conceptualmente con precisión, que las relaciones entre variables estén explicitadas y que las definiciones sean operativas para la investigación propuesta. \\ \hline
V & Rigor académico y pertinencia del campo & Evalúa la profundidad analítica, la pertinencia temática del marco teórico respecto de la Ingeniería Civil, la incorporación de normativa técnica aplicable, y la consideración de factores de sostenibilidad, seguridad o impacto sociotécnico cuando el problema lo requiera. \\ \hline
\end{tabularx}
\end{table}
\begin{table}[htbp]
\centering
\scriptsize
\caption{Escala de valoración para la evaluación del marco teórico.}
\label{tab:escala-marco-teorico}
\begin{tabularx}{\textwidth}{|C{2.8cm}|C{1cm}|X|}
\hline
\rowcolor{headerblue}
\thd{Nivel} & \thd{Pts.} & \thd{Descriptor operativo} \\ \hline
\textcolor{csuf}{\textbf{Excelente}} & \textbf{3} & El proyecto satisface plenamente el criterio evaluado. El marco teórico es explícito, técnicamente preciso, sustentado en fuentes confiables y lógicamente articulado con los demás componentes del proyecto. No se identifican ambigüedades, omisiones críticas ni desconexiones internas. \\ \hline
\rowcolor{lightgray}
\textcolor{cace}{\textbf{Aceptable}} & \textbf{2} & El criterio es identificable pero se cumple de manera parcial o incompleta. El marco teórico aborda el aspecto evaluado sin la profundidad, la precisión o la articulación suficiente. Susceptible de fortalecimiento con correcciones específicas. \\ \hline
\textcolor{cdef}{\textbf{Deficiente}} & \textbf{1} & El cumplimiento del criterio es tangencial o meramente declarativo. La fundamentación es superficial, las fuentes son insuficientes o desactualizadas, o se presentan inconsistencias que debilitan la estructura del marco teórico. \\ \hline
\rowcolor{lightgray}
\textcolor{cins}{\textbf{Ausente}} & \textbf{0} & El proyecto no evidencia consideración alguna del criterio evaluado. Existe omisión total del aspecto evaluado o una desconexión que invalida el marco teórico en ese componente. \\ \hline
\end{tabularx}
\end{table}

\clearpage
\subsection{Fundamentación Teórica y Selección de Teorías o Modelos}

\textit{Se evalúa si el proyecto identifica, selecciona y justifica las teorías, los principios científicos o los modelos de la Ingeniería Civil que fundamentan la investigación propuesta. Esta dimensión verifica que la selección teórica no sea arbitraria ni meramente nominal, sino que responda a la naturaleza del problema formulado y esté sustentada en la literatura revisada. El nivel de exigencia se calibra para un proyecto de tesis de pregrado.} Véase la Tabla~\ref{tab:rubrica-dim1-fundamentacion}.
{\scriptsize
\begin{longtable}{|C{0.55cm}|L{2cm}|L{8cm}|C{0.4cm}|C{0.4cm}|C{0.4cm}|C{0.4cm}|}
\caption{Rúbrica de evaluación de la Dimensión~I: fundamentación teórica y selección de teorías o modelos.}\label{tab:rubrica-dim1-fundamentacion}\\
\hline
\rowcolor{headerblue}
\thd{N.\textsuperscript{o}} & \thd{Criterio} & \thd{Indicador verificable} & \thd{3} & \thd{2} & \thd{1} & \thd{0} \\
\hline\endfirsthead
\hline
\rowcolor{headerblue}
\thd{N.\textsuperscript{o}} & \thd{Criterio} & \thd{Indicador verificable} & \thd{3} & \thd{2} & \thd{1} & \thd{0} \\
\hline\endhead
1.1 & Identificación explícita de las teorías, principios o modelos que fundamentan la investigación & El proyecto identifica de manera explícita las teorías, los principios científicos o los modelos de la Ingeniería Civil sobre los cuales se fundamenta la investigación (p.\,ej., teoría de Mohr-Coulomb para resistencia al corte, modelo constitutivo de consolidación de Terzaghi, teoría multicapa de Burmister para pavimentos, principios de diseño basado en confiabilidad). La identificación no es genérica (\enquote{se basa en la mecánica de suelos}) sino específica: se nombran las teorías, se enuncian sus principios fundamentales y se precisa su relación con el fenómeno investigado. &
& & & \\ \hline
\rowcolor{lightgray}
1.2 & Justificación de la selección teórica en función del problema de investigación & La selección de las teorías o modelos está justificada: el proyecto explica por qué dichas teorías son las más apropiadas para abordar el problema formulado y no otras alternativas disponibles. La justificación no es circular (\enquote{se usa la teoría X porque es la teoría del tema}) sino argumentada: se vincula la naturaleza del fenómeno investigado con los supuestos, el alcance y las condiciones de aplicabilidad de la teoría seleccionada. Se evidencia que la selección responde a un criterio razonado y no a una decisión arbitraria o por conveniencia. &
& & & \\ \hline
1.3 & Presentación de los fundamentos teóricos con precisión técnica y sin distorsiones & Las teorías o modelos presentados se exponen con precisión técnica: las ecuaciones constitutivas, los parámetros, las hipótesis de partida y las condiciones de contorno se describen correctamente. No se presentan versiones simplificadas que distorsionen el modelo original, ni se omiten supuestos fundamentales que delimitan su validez (p.\,ej., presentar la ecuación de Terzaghi sin mencionar las hipótesis de carga centrada, suelo homogéneo y cimentación superficial). Los fundamentos teóricos son trazables a fuentes primarias o textos de referencia reconocidos. &
& & & \\ \hline
\rowcolor{lightgray}
1.4 & Equilibrio entre fuentes seminales y fuentes actuales en la fundamentación & La fundamentación teórica integra un equilibrio verificable entre fuentes seminales o clásicas que establecen los principios fundamentales (textos fundacionales, artículos originales de las teorías empleadas) y fuentes actuales que documentan avances recientes, validaciones experimentales o extensiones de dichas teorías. El marco teórico no se sustenta exclusivamente en textos clásicos desactualizados ni exclusivamente en publicaciones recientes sin anclaje en los fundamentos originales. &
& & & \\ \hline
\multicolumn{3}{|r|}{\cellcolor{white}\textbf{Subtotal Dimensión~I (máx.~12 puntos):}} & \multicolumn{4}{c|}{\cellcolor{white}\textbf{/12}} \\ \hline
\end{longtable}
}

\clearpage
\subsection{Coherencia Interna del Marco Teórico con el Diseño de Investigación}

\textit{Se evalúa la alineación verificable entre el marco teórico y los demás componentes del proyecto de tesis: problema de investigación, objetivos, hipótesis y metodología anticipada. Un marco teórico de calidad no es un componente aislado, sino el eje articulador que conecta el problema con el método y la interpretación de resultados. Esta dimensión verifica que tal articulación sea explícita y libre de contradicciones.} Los criterios se aprecian en la Tabla~\ref{tab:rubrica-dim2-coherencia}.
\scriptsize
\begin{longtable}{|C{0.55cm}|L{2.8cm}|L{8cm}|C{0.3cm}|C{0.3cm}|C{0.3cm}|C{0.3cm}|}
\caption{Rúbrica de evaluación de la Dimensión~II: coherencia interna del marco teórico con el diseño de investigación.}\label{tab:rubrica-dim2-coherencia}\\
\hline
\rowcolor{headerblue}
\thd{N.\textsuperscript{o}} & \thd{Criterio} & \thd{Indicador verificable} & \thd{3} & \thd{2} & \thd{1} & \thd{0} \\
\hline\endfirsthead
\hline
\rowcolor{headerblue}
\thd{N.\textsuperscript{o}} & \thd{Criterio} & \thd{Indicador verificable} & \thd{3} & \thd{2} & \thd{1} & \thd{0} \\
\hline\endhead
2.1 & Alineación entre el marco teórico y el problema de investigación \epref{Alineación problema--teoría} & Las teorías, los principios y los modelos presentados en el marco teórico se vinculan directamente con el problema de investigación formulado. No se incluyen desarrollos teóricos extensos que no guardan relación con el problema ni se omiten fundamentos esenciales para comprenderlo. Cada bloque teórico presentado responde a la pregunta: \enquote{¿qué conocimiento previo es necesario para investigar este problema específico?}. No hay teorías huérfanas (incluidas pero no conectadas con el problema) ni vacíos teóricos (aspectos del problema sin cobertura en el marco). &
& & & \\ \hline
\rowcolor{lightgray}
2.2 & Alineación entre el marco teórico y los objetivos e hipótesis del proyecto & Las variables, los constructos y las relaciones presentadas en el marco teórico se corresponden biunívocamente con los objetivos y la hipótesis del proyecto. Cada variable mencionada en la hipótesis posee sustento teórico en el marco; cada objetivo específico encuentra su fundamentación conceptual en algún componente del marco teórico. No existen variables en la hipótesis ausentes del marco teórico, ni constructos teóricos extensamente desarrollados que no se reflejen en ningún objetivo ni en la hipótesis. &
& & & \\ \hline
2.3 & Alineación entre el marco teórico y la metodología anticipada & El marco teórico provee los fundamentos necesarios para justificar la metodología anticipada en el proyecto: los ensayos seleccionados, los modelos de cálculo, las técnicas de análisis o los procedimientos experimentales encuentran su respaldo teórico en el marco. Si el proyecto anticipa ensayos de resistencia a la compresión no confinada, el marco teórico incluye los principios de la resistencia al corte y los criterios de falla pertinentes. La conexión teoría--método es explícita, no inferida por el lector. &
& & & \\ \hline
\rowcolor{lightgray}
2.4 & Ausencia de contradicciones internas y consistencia lógica del marco teórico & El marco teórico no presenta contradicciones internas: no se postulan simultáneamente principios incompatibles, no se adoptan modelos con supuestos mutuamente excluyentes sin justificación, y no se afirman relaciones entre variables que se contradigan en distintas secciones del documento. La progresión argumental es lógica: los conceptos se presentan en un orden que permite la comprensión acumulativa, desde los fundamentos generales hacia los específicos del problema investigado. &
& & & \\ \hline
\multicolumn{3}{|r|}{\cellcolor{white}\textbf{Subtotal Dimensión~II (máx.~12 puntos):}} & \multicolumn{4}{c|}{\cellcolor{white}\textbf{/12}} \\ \hline
\end{longtable}
\normalsize

\clearpage
\subsection{Revisión de Literatura y Síntesis Crítica}

\textit{Se evalúa la calidad, la actualidad, la diversidad y el tratamiento crítico de las fuentes bibliográficas que sustentan el marco teórico. Una revisión de literatura de calidad no es una enumeración cronológica de autores y hallazgos, sino una síntesis organizada temáticamente que identifica convergencias, divergencias y vacíos en el estado del conocimiento. El nivel de exigencia se calibra para pregrado, donde se espera síntesis temática y posicionamiento frente a la literatura, no necesariamente una revisión sistemática formal.} Los criterios se aprecian en la Tabla~\ref{tab:rubrica-dim3-revision-literatura}.
\scriptsize
\begin{longtable}{|C{0.55cm}|L{2.8cm}|L{8cm}|C{0.3cm}|C{0.3cm}|C{0.3cm}|C{0.3cm}|}
\caption{Rúbrica de evaluación de la Dimensión~III: revisión de literatura y síntesis crítica.}\label{tab:rubrica-dim3-revision-literatura}\\
\hline
\rowcolor{headerblue}
\thd{N.\textsuperscript{o}} & \thd{Criterio} & \thd{Indicador verificable} & \thd{3} & \thd{2} & \thd{1} & \thd{0} \\
\hline\endfirsthead
\hline
\rowcolor{headerblue}
\thd{N.\textsuperscript{o}} & \thd{Criterio} & \thd{Indicador verificable} & \thd{3} & \thd{2} & \thd{1} & \thd{0} \\
\hline\endhead
3.1 & Organización temática y progresión lógica de la revisión de literatura \epref{Organización temática} & La revisión de literatura está organizada temáticamente ---no cronológicamente ni por autor--- siguiendo una progresión lógica que transita desde los conceptos generales hacia los específicos del problema investigado. Los temas se agrupan por afinidad conceptual (p.\,ej., propiedades del suelo, mecanismos de estabilización, agentes estabilizadores, ensayos de caracterización) y cada bloque temático cumple una función argumental identificable dentro del marco teórico. La estructura no es una yuxtaposición de fichas bibliográficas sino un discurso articulado. &
& & & \\ \hline
\rowcolor{lightgray}
3.2 & Actualidad, diversidad y confiabilidad de las fuentes bibliográficas & Las fuentes empleadas son diversas en tipo (artículos de revistas indexadas, normas técnicas, libros de referencia, memorias de congresos, informes técnicos institucionales) y en procedencia (fuentes nacionales e internacionales). La mayoría de las fuentes académicas ---excluyendo textos clásicos de fundamentación--- corresponden al periodo de los últimos diez años. No se sustentan afirmaciones técnicas en fuentes no arbitradas (blogs, páginas web sin autoría identificable, apuntes de clase no publicados) ni se emplean exclusivamente fuentes en un solo idioma cuando la literatura del tema es internacional. &
& & & \\ \hline
3.3 & Síntesis crítica: convergencias, divergencias y vacíos identificados en la literatura & La revisión de literatura no se limita a describir los hallazgos de cada fuente de manera aislada, sino que sintetiza convergencias entre estudios (resultados concordantes, tendencias confirmadas), identifica divergencias (resultados contradictorios, condiciones diferentes que producen desenlaces distintos) y señala vacíos en el estado del conocimiento (aspectos no investigados, condiciones locales no evaluadas, materiales no caracterizados). La identificación de vacíos es específica y se vincula directamente con la justificación de la investigación propuesta. &
& & & \\ \hline
\rowcolor{lightgray}
3.4 & Posicionamiento del proyecto frente al estado del arte revisado & La revisión de literatura concluye con un posicionamiento explícito del proyecto frente al conocimiento existente: se declara qué aspecto específico del estado del arte el proyecto busca complementar, validar, contrastar o extender. Este posicionamiento no es genérico (\enquote{se busca aportar al conocimiento}) sino preciso: identifica la brecha concreta que la investigación abordará y cómo esta se deriva de los vacíos detectados en la revisión. La transición entre la revisión de literatura y la investigación propuesta es lógica y no abrupta. &
& & & \\ \hline
\multicolumn{3}{|r|}{\cellcolor{white}\textbf{Subtotal Dimensión~III (máx.~12 puntos):}} & \multicolumn{4}{c|}{\cellcolor{white}\textbf{/12}} \\ \hline
\end{longtable}
\normalsize

\clearpage
\subsection{Claridad Conceptual y Definición de Variables}

\textit{Se evalúa si los constructos, las variables y los términos técnicos principales están definidos conceptualmente con la precisión necesaria para que el marco teórico funcione como herramienta operativa de la investigación. Un marco teórico con claridad conceptual permite que cualquier lector con formación en Ingeniería Civil comprenda sin ambigüedad qué se mide, qué se controla, qué se manipula y cómo se relacionan las variables del estudio.} Los criterios se aprecian en la Tabla~\ref{tab:rubrica-dim4-claridad-conceptual}.
\scriptsize
\begin{longtable}{|C{0.55cm}|L{2.8cm}|L{8cm}|C{0.3cm}|C{0.3cm}|C{0.3cm}|C{0.3cm}|}
\caption{Rúbrica de evaluación de la Dimensión~IV: claridad conceptual y definición de variables.}\label{tab:rubrica-dim4-claridad-conceptual}\\
\hline
\rowcolor{headerblue}
\thd{N.\textsuperscript{o}} & \thd{Criterio} & \thd{Indicador verificable} & \thd{3} & \thd{2} & \thd{1} & \thd{0} \\
\hline\endfirsthead
\hline
\rowcolor{headerblue}
\thd{N.\textsuperscript{o}} & \thd{Criterio} & \thd{Indicador verificable} & \thd{3} & \thd{2} & \thd{1} & \thd{0} \\
\hline\endhead
4.1 & Definición conceptual precisa de las variables de la investigación \epref{Definición de variables} & Las variables de la investigación ---dependiente(s), independiente(s) e intervinientes, según corresponda--- están definidas conceptualmente dentro del marco teórico con precisión técnica. La definición conceptual establece qué es cada variable en términos teóricos, no solo cómo se mide. Por ejemplo, si la variable dependiente es \enquote{resistencia a la compresión no confinada}, el marco teórico define este constructo desde la teoría de resistencia al corte, no solo como \enquote{el resultado del ensayo UCS}. Las definiciones son consistentes con la literatura citada. &
& & & \\ \hline
\rowcolor{lightgray}
4.2 & Explicitación de las relaciones teóricas entre variables & El marco teórico explicita las relaciones teóricas esperadas entre las variables de la investigación: relaciones causales, correlacionales, de dependencia o de moderación, según la naturaleza del estudio. Las relaciones no se asumen implícitamente ni se postergan al capítulo de metodología; están fundamentadas en los principios teóricos previamente expuestos (p.\,ej., \enquote{la adición de cemite Portland incrementa la resistencia a la compresión de suelos arcillosos por la formación de compuestos puzolánicos, tal como postulan [autores]}). Las relaciones anticipadas son la base teórica de la hipótesis. &
& & & \\ \hline
4.3 & Definición precisa de términos técnicos esenciales y ausencia de ambigüedades terminológicas & Los términos técnicos empleados en el marco teórico están definidos con precisión y se utilizan de manera consistente a lo largo del documento. No se emplean sinónimos intercambiables sin aclaración (p.\,ej., usar indistintamente \enquote{estabilización}, \enquote{mejoramiento} y \enquote{tratamiento} sin definir cada término), ni se introducen acrónimos sin definición previa. La terminología es coherente con la normativa técnica y la literatura especializada del campo. &
& & & \\ \hline
\rowcolor{lightgray}
4.4 & Delimitación de los supuestos y las condiciones de contorno del marco teórico & El marco teórico explicita sus supuestos y sus condiciones de contorno: las limitaciones de las teorías adoptadas, los rangos de validez de los modelos empleados, las condiciones bajo las cuales las relaciones teóricas postuladas se mantienen y aquellas bajo las cuales podrían no cumplirse. Esta delimitación no es genérica (\enquote{toda teoría tiene limitaciones}) sino específica: se identifican los supuestos del modelo adoptado y se evalúa su pertinencia para el contexto de la investigación propuesta. &
& & & \\ \hline
\multicolumn{3}{|r|}{\cellcolor{white}\textbf{Subtotal Dimensión~IV (máx.~12 puntos):}} & \multicolumn{4}{c|}{\cellcolor{white}\textbf{/12}} \\ \hline
\end{longtable}
\normalsize

\clearpage
\subsection{Rigor Académico y Pertinencia para el Campo de Conocimiento}

\textit{Se evalúa la profundidad analítica del marco teórico, su pertinencia respecto del dominio técnico de la Ingeniería Civil, la incorporación de normativa técnica aplicable y la consideración de factores de sostenibilidad, seguridad o impacto sociotécnico que los estándares de acreditación internacional exigen en la formación de ingenieros civiles. Esta dimensión verifica que el marco teórico trascienda la mera recopilación informativa y evidencie pensamiento analítico.} Los criterios se aprecian en la Tabla~\ref{tab:rubrica-dim5-rigor-academico}.
\scriptsize
\begin{longtable}{|C{0.55cm}|L{2.8cm}|L{8cm}|C{0.3cm}|C{0.3cm}|C{0.3cm}|C{0.3cm}|}
\caption{Rúbrica de evaluación de la Dimensión~V: rigor académico y pertinencia temática.}\label{tab:rubrica-dim5-rigor-academico}\\
\hline
\rowcolor{headerblue}
\thd{N.\textsuperscript{o}} & \thd{Criterio} & \thd{Indicador verificable} & \thd{3} & \thd{2} & \thd{1} & \thd{0} \\
\hline\endfirsthead
\hline
\rowcolor{headerblue}
\thd{N.\textsuperscript{o}} & \thd{Criterio} & \thd{Indicador verificable} & \thd{3} & \thd{2} & \thd{1} & \thd{0} \\
\hline\endhead
5.1 & Profundidad analítica: el marco teórico trasciende la descripción y evidencia análisis \epref{Profundidad analítica} & El marco teórico no se limita a describir secuencialmente lo que cada fuente reporta, sino que analiza, compara y evalúa críticamente los aportes de la literatura. Se evidencia pensamiento analítico propio del tesista: comparación de modelos, evaluación de condiciones de aplicabilidad, discusión de ventajas y limitaciones de enfoques alternativos. El texto no es una transcripción parafraseada de fuentes individuales sino una construcción argumental donde el tesista articula el conocimiento existente con voz propia y criterio técnico. &
& & & \\ \hline
\rowcolor{lightgray}
5.2 & Incorporación de normativa técnica, estándares y clasificaciones pertinentes & El marco teórico incorpora la normativa técnica, los estándares de ensayo y las clasificaciones pertinentes al problema de investigación (ASTM, AASHTO, NTP, MTC, ACI, AISC, Eurocódigos, normas CE, E.050, u otras aplicables según la especialidad). La normativa no se cita como referencia decorativa sino que se integra funcionalmente: se especifican los métodos de ensayo normalizados que se emplearán, los criterios de aceptación o rechazo establecidos por la norma, y los parámetros que la normativa define como umbrales de desempeño. &
& & & \\ \hline
5.3 & Pertinencia de la especialidad: el marco teórico se enmarca inequívocamente en la Ingeniería Civil & El contenido del marco teórico se inscribe dentro del dominio de la especialidad de la Ingeniería Civil. Los principios expuestos corresponden a mecánica de suelos, mecánica de materiales, análisis estructural, ingeniería de pavimentos, hidráulica, hidrología, gestión de la construcción, u otra rama reconocida de la especialidad. El marco teórico no está constituido predominantemente por contenido de disciplinas auxiliares (química pura, biología, administración general) sin anclaje en los principios de la Ingeniería Civil. &
& & & \\ \hline
\rowcolor{lightgray}
5.4 & Consideración de factores de sostenibilidad, seguridad o impacto sociotécnico cuando el problema lo requiere & Cuando la naturaleza del problema de investigación lo demanda, el marco teórico incorpora consideraciones de sostenibilidad ambiental (uso de recursos no renovables, huella de carbono, reciclaje de materiales), seguridad pública (estabilidad de taludes, capacidad portante de cimentaciones, durabilidad de estructuras) o impacto sociotécnico (accesibilidad de la solución, disponibilidad de materiales en el contexto local, costos asociados). Estas consideraciones están integradas como parte del fundamento teórico de la investigación, no como adiciones genéricas desconectadas del problema. &
& & & \\ \hline
\multicolumn{3}{|r|}{\cellcolor{white}\textbf{Subtotal Dimensión~V (máx.~12 puntos):}} & \multicolumn{4}{c|}{\cellcolor{white}\textbf{/12}} \\ \hline
\end{longtable}
\normalsize
\newpage
En esta sección se presenta el consolidado de resultados por dimensión y la escala de calificación del marco teórico, resumidos en las tablas~\ref{tab:consolidado-marco-teorico} y~\ref{tab:escala-calificacion-marco-teorico}.
\begin{table}[htbp]
\centering
\scriptsize
\caption{Cuadro consolidado de resultados por dimensión del marco teórico.}\label{tab:consolidado-marco-teorico}
\begin{tabularx}{\textwidth}{|C{0.8cm}|L{7.5cm}|C{1.8cm}|C{1.8cm}|X|}
\hline
\rowcolor{headerblue}
\thd{Dim.} & \thd{Denominación} & \thd{Pts. Máx.} & \thd{Pts. Obt.} & \thd{Obs.} \\ \hline
I & Fundamentación Teórica y Selección de Teorías o Modelos & 12 & & \\ \hline
\rowcolor{lightgray}
II & Coherencia Interna del Marco Teórico con el Diseño de Investigación & 12 & & \\ \hline
III & Revisión de Literatura y Síntesis Crítica & 12 & & \\ \hline
\rowcolor{lightgray}
IV & Claridad Conceptual y Definición de Variables & 12 & & \\ \hline
V & Rigor Académico y Pertinencia Dentro del Campo de Conocimiento & 12 & & \\ \hline
\rowcolor{white}
& \textbf{TOTAL GENERAL} & \textbf{60} & & \\ \hline
\end{tabularx}
\normalsize
\end{table}
\begin{table}[htbp]
\centering
\scriptsize
\caption{Escala de calificación del marco teórico.}\label{tab:escala-calificacion-marco-teorico}
\begin{tabularx}{\textwidth}{|C{2.5cm}|C{3cm}|X|}
\hline
\rowcolor{headerblue}
\thd{Rango} & \thd{Calificación} & \thd{Significado} \\ \hline
\textbf{54--60}\newline(90--100\,\%) &
\textcolor{csuf}{\textbf{MARCO TEÓRICO\newline PLENO}} &
El marco teórico es integral, técnicamente preciso y lógicamente articulado con todos los componentes del proyecto. La fundamentación es pertinente, las fuentes son actuales y diversas, las variables están definidas con claridad, la revisión de literatura evidencia síntesis crítica y la coherencia interna es verificable. El marco teórico constituye una base sólida y suficiente para la ejecución de la investigación propuesta. \\ \hline
\rowcolor{lightgray}
\textbf{42--53}\newline(70--89\,\%) &
\textcolor{cace}{\textbf{MARCO TEÓRICO\newline SUFICIENTE}} &
El marco teórico presenta los componentes esenciales con fundamentación identificable, pero exhibe áreas específicas susceptibles de fortalecimiento en cuanto a profundidad, articulación con otros componentes del proyecto, actualidad de fuentes o precisión en la definición de variables. Las debilidades son corregibles sin reestructuración fundamental del marco. \\ \hline
\textbf{30--41}\newline(50--69\,\%) &
\textcolor{cdef}{\textbf{MARCO TEÓRICO\newline PARCIAL}} &
El marco teórico es incompleto o predominantemente descriptivo. Se evidencian vacíos significativos en la fundamentación, desconexiones entre el marco y otros componentes del proyecto, fuentes insuficientes o desactualizadas, definiciones ambiguas de variables, o ausencia de síntesis crítica. Requiere fortalecimiento sustancial antes de que el proyecto pueda ejecutarse como investigación. \\ \hline
\rowcolor{lightgray}
\textbf{0--29}\newline($<$50\,\%) &
\textcolor{cins}{\textbf{MARCO TEÓRICO\newline INSUFICIENTE}} &
El marco teórico no demuestra consistencia verificable. La fundamentación es ausente o arbitraria, no existe articulación con el problema ni con la metodología, las fuentes son inadecuadas o insuficientes, y las variables no están definidas. Se requiere una reformulación fundamental del marco teórico del proyecto. \\ \hline
\end{tabularx}
\normalsize
\end{table}

\clearpage
\section{Evaluación del Diseño Metodológico}

{\footnotesize
La metodología constituye el componente del proyecto de tesis que operacionaliza la estrategia de investigación: traduce el problema, los objetivos y las hipótesis en un plan de acción técnicamente riguroso, replicable y verificable. En Ingeniería Civil, la metodología se distingue de la tradición de las ciencias sociales en aspectos fundamentales: las variables son magnitudes físicas medibles con instrumentos calibrados; los especímenes son materiales, probetas, muestras de suelo, elementos estructurales o tramos de infraestructura; los protocolos de ensayo están normalizados por organismos como ASTM, AASHTO, NTP, ISO o ACI; y la inferencia se sustenta en principios de mecánica de suelos, mecánica de materiales, hidráulica, análisis estructural u otras ramas reconocidas de la disciplina. Esta rúbrica evalúa la calidad, consistencia, completitud y rigor técnico del diseño metodológico propuesto, verificando que el tesista haya planificado un procedimiento experimental o analítico capaz de responder a la pregunta de investigación con validez interna, trazabilidad de datos y control de las fuentes de variabilidad propias de los materiales y procesos de la Ingeniería Civil.\par
}
Las dimensiones de evaluación del diseño metodológico se resumen en la Tabla~\ref{tab:dim-metodologia}, mientras que la escala de valoración y sus descriptores operativos se presentan en la Tabla~\ref{tab:escala-metodologia}.

\begin{table}[htbp]
\centering
\scriptsize
\caption{Dimensiones de evaluación del diseño metodológico y alcance de cada dimensión.}
\label{tab:dim-metodologia}
\begin{tabularx}{\textwidth}{|C{0.8cm}|L{4.5cm}|X|}
\hline
\rowcolor{headerblue}
\thd{Dim.} & \thd{Denominación} & \thd{Alcance} \\ \hline
I & Tipo, nivel y diseño de investigación & Verifica la coherencia entre el tipo de investigación declarado y el diseño efectivamente propuesto, la adecuación del nivel investigativo a los objetivos, la capacidad del diseño para controlar variables y la posibilidad de inferencia causal cuando el estudio lo requiere. \\ \hline
\rowcolor{lightgray}
II & Operacionalización de variables y definición de tratamientos experimentales & Evalúa la definición cuantitativa de los niveles de las variables independientes, la fundamentación técnica de los tratamientos o condiciones experimentales, la ausencia de indeterminaciones y la presentación de la matriz factorial o esquema de diseño. \\ \hline
III & Muestreo, tamaño muestral y especímenes de ensayo & Examina la justificación del tamaño muestral, la potencia estadística del diseño, los criterios de inclusión y exclusión de especímenes, y la representatividad de las muestras respecto de la población o universo de estudio. \\ \hline
\rowcolor{lightgray}
IV & Técnicas, instrumentos y protocolos de recolección de datos & Evalúa la referencia a protocolos de ensayo normalizados, la calibración y validez de los instrumentos de medición, los procedimientos de control de calidad y la coherencia entre instrumentos y variables declaradas. \\ \hline
V & Plan de análisis de datos y tratamiento de la incertidumbre & Verifica la especificación de los métodos estadísticos a emplear, la cuantificación y propagación de incertidumbres, la pertinencia de los marcos analíticos y la distinción entre efectos reales, variabilidad aleatoria y significancia práctica para la Ingeniería Civil. \\ \hline
\end{tabularx}
\end{table}

\begin{table}[htbp]
\centering
\scriptsize
\caption{Escala de valoración y descriptor operativo para evaluar del diseño metodológico.}
\label{tab:escala-metodologia}
\begin{tabularx}{\textwidth}{|C{2.8cm}|C{1cm}|X|}
\hline
\rowcolor{headerblue}
\thd{Nivel} & \thd{Pts.} & \thd{Descriptor operativo} \\ \hline
\textcolor{csuf}{\textbf{Excelente}} & \textbf{3} & El proyecto satisface plenamente el criterio evaluado. La metodología es explícita, técnicamente precisa, coherente con los objetivos e hipótesis, debidamente sustentada y replicable. No se identifican ambigüedades, omisiones ni inconsistencias procedimentales. \\ \hline
\rowcolor{lightgray}
\textcolor{cace}{\textbf{Aceptable}} & \textbf{2} & El criterio es identificable pero se cumple de manera parcial o incompleta. La metodología aborda el aspecto evaluado sin la profundidad, la precisión o el sustento suficiente. Susceptible de fortalecimiento con correcciones específicas. \\ \hline
\textcolor{cdef}{\textbf{Deficiente}} & \textbf{1} & El cumplimiento del criterio es tangencial o meramente declarativo. La descripción metodológica es superficial, carece de sustento técnico verificable o presenta inconsistencias que comprometen la validez interna del diseño. \\ \hline
\rowcolor{lightgray}
\textcolor{cins}{\textbf{Ausente}} & \textbf{0} & El proyecto no evidencia consideración alguna del criterio evaluado. Existe omisión total del aspecto o una desconexión que invalida la metodología en ese componente. \\ \hline
\end{tabularx}
\end{table}

\clearpage
\subsection{Tipo, Nivel y Diseño de Investigación}

\textit{Se evalúa si el proyecto declara con precisión el tipo y nivel de investigación, si el diseño propuesto es coherente con la declaración y si la estructura experimental o analítica posee la capacidad necesaria para responder a la pregunta de investigación y contrastar las hipótesis formuladas. Esta dimensión verifica que la arquitectura metodológica no sea una etiqueta nominal sino un plan operativo verificable.} Véase la Tabla~\ref{tab:rubrica-dim1-tipo-nivel-diseno}.

\scriptsize
\begin{longtable}{|C{0.55cm}|L{2.8cm}|L{8cm}|C{0.3cm}|C{0.3cm}|C{0.3cm}|C{0.3cm}|}
\caption{Rúbrica de evaluación de la Dimensión~I: tipo, nivel y diseño de investigación.}\label{tab:rubrica-dim1-tipo-nivel-diseno}\\
\hline
\rowcolor{headerblue}
\thd{N.\textsuperscript{o}} & \thd{Criterio} & \thd{Indicador verificable} & \thd{3} & \thd{2} & \thd{1} & \thd{0} \\
\hline\endfirsthead
\hline
\rowcolor{headerblue}
\thd{N.\textsuperscript{o}} & \thd{Criterio} & \thd{Indicador verificable} & \thd{3} & \thd{2} & \thd{1} & \thd{0} \\
\hline\endhead
1.1 & Coherencia entre el tipo de investigación declarado y el diseño efectivamente propuesto \epref{Tipo vs.\ diseño} & El proyecto declara un tipo de investigación (experimental, cuasiexperimental, descriptivo, correlacional, explicativo, u otro) y el diseño metodológico descrito es efectivamente consistente con esa declaración. No se declara \enquote{experimental} cuando en realidad se propone una comparación observacional sin manipulación de variables, ni se declara \enquote{descriptivo} cuando el diseño incluye manipulación deliberada de una variable independiente y medición de su efecto sobre una variable dependiente. La coherencia entre la etiqueta declarada y el procedimiento descrito es verificable en el cuerpo del proyecto. &
& & & \\ \hline
\rowcolor{lightgray}
1.2 & Adecuación del nivel investigativo a los objetivos y a la pregunta de investigación & El nivel de investigación declarado (exploratorio, descriptivo, correlacional o explicativo) es congruente con los objetivos formulados y con la pregunta de investigación. Si el objetivo busca \enquote{determinar el efecto de X sobre Y}, el nivel no puede ser meramente descriptivo. Si el objetivo busca \enquote{caracterizar las propiedades de un material}, el nivel no puede ser explicativo sin una relación causal que explicar. El nivel investigativo refleja la profundidad analítica que el diseño es capaz de alcanzar, no una aspiración desconectada de los procedimientos propuestos. &
& & & \\ \hline
1.3 & Capacidad del diseño para aislar y controlar sistemáticamente las variables & El diseño metodológico describe los mecanismos concretos mediante los cuales se controlarán las variables intervinientes y se aislarán los efectos de la variable independiente sobre la dependiente: condiciones de curado estandarizadas, probetas de control sin tratamiento, procedimientos de compactación uniformes, condiciones de carga controladas, temperatura y humedad reguladas, o cualquier otro protocolo pertinente a la naturaleza del estudio. El control de variables no se declara genéricamente (\enquote{se controlarán las variables externas}) sino que se operacionaliza con procedimientos específicos, verificables y coherentes con la práctica experimental de la Ingeniería Civil. &
& & & \\ \hline
\rowcolor{lightgray}
1.4 & Grupos de control, aleatorización y estructura comparativa del diseño & Cuando el diseño lo requiere, el proyecto prevé grupos de control adecuados (especímenes sin tratamiento, mezclas patrón, condiciones de referencia normativas) que permitan atribuir los efectos observados a la variable independiente y no a factores espurios. La asignación de tratamientos a las unidades experimentales sigue un criterio explícito: aleatorización completa, bloques aleatorizados, asignación sistemática con justificación, u otro esquema documentado. Si el diseño es cuasiexperimental y la aleatorización no es posible, se explicita esta limitación y se describen las estrategias de mitigación del sesgo adoptadas. &
& & & \\ \hline
\multicolumn{3}{|r|}{\cellcolor{white}\textbf{Subtotal Dimensión~I (máx.~12 puntos):}} & \multicolumn{4}{c|}{\cellcolor{white}\textbf{/12}} \\ \hline
\end{longtable}
\normalsize

\clearpage
\subsection{Operacionalización de Variables y Definición de Tratamientos Experimentales}

\textit{Se evalúa si el proyecto traduce las variables de investigación en magnitudes medibles con unidades técnicas precisas, si los niveles o tratamientos de la variable independiente están cuantitativamente definidos y técnicamente fundamentados, y si el esquema experimental está completamente especificado antes de la ejecución. Esta dimensión verifica que la metodología no difiera decisiones críticas a fases empíricas posteriores sin justificación.} Los criterios se aprecian en la Tabla~\ref{tab:rubrica-dim2-operacionalizacion}.

\scriptsize
\begin{longtable}{|C{0.55cm}|L{2cm}|L{9cm}|C{0.3cm}|C{0.3cm}|C{0.3cm}|C{0.3cm}|}
\caption{Rúbrica de evaluación de la Dimensión~II: operacionalización de variables y definición de tratamientos.}\label{tab:rubrica-dim2-operacionalizacion}\\
\hline
\rowcolor{headerblue}
\thd{N.\textsuperscript{o}} & \thd{Criterio} & \thd{Indicador verificable} & \thd{3} & \thd{2} & \thd{1} & \thd{0} \\
\hline\endfirsthead
\hline
\rowcolor{headerblue}
\thd{N.\textsuperscript{o}} & \thd{Criterio} & \thd{Indicador verificable} & \thd{3} & \thd{2} & \thd{1} & \thd{0} \\
\hline\endhead
2.1 & Definición cuantitativa de los niveles de la variable independiente \epref{Operacionalización en ingeniería} & Los niveles o tratamientos de la variable independiente están definidos con precisión cuantitativa y unidades técnicas: porcentajes de adición de agente estabilizante en masa, dosificaciones de cemento en kg/m\textsuperscript{3}, relaciones agua/cemento, espesores de capa en centímetros, niveles de carga en kN, temperaturas de curado en \textdegree C, porcentajes de sustitución de agregados, concentraciones de aditivo, u otras magnitudes pertinentes al diseño. No se emplean expresiones vagas como \enquote{diferentes porcentajes}, \enquote{varias dosificaciones} o \enquote{distintas proporciones} sin especificar los valores numéricos concretos. &
& & & \\ \hline
\rowcolor{lightgray}
2.2 & Fundamentación técnica o bibliográfica de los niveles seleccionados & Los niveles de la variable independiente no se eligen arbitrariamente sino que se fundamentan en al menos una de las siguientes fuentes: antecedentes experimentales reportados en la literatura técnica, recomendaciones de normas o especificaciones técnicas (ASTM, ACI, MTC, NTP), resultados de ensayos preliminares propios, modelos teóricos que predicen rangos óptimos, o análisis de la práctica constructiva regional. La fundamentación es explícita: se citan las fuentes que justifican la selección de cada nivel o rango y se argumenta su pertinencia para el contexto de la investigación. &
& & & \\ \hline
2.3 & Ausencia de indeterminación en las condiciones experimentales & El proyecto no pospone a fases empíricas posteriores la definición de las condiciones experimentales, las cuales deben quedar plenamente determinadas en la etapa de diseño del proyecto de tesis: las dosificaciones, las edades de curado, los tipos de ensayo, las condiciones de preparación de especímenes, los parámetros geométricos de probetas y las condiciones ambientales de ensayo están completamente especificados. Si algún parámetro se determinará mediante ensayos preliminares (p.\,ej., el contenido óptimo de humedad para compactación), el proyecto describe el procedimiento normado que se empleará para determinarlo y en qué momento del cronograma se ejecutará dicha determinación previa. &
& & & \\ \hline
\rowcolor{lightgray}
2.4 & Presentación de la matriz factorial o esquema de diseño experimental & El proyecto presenta un esquema de diseño que explicita todas las combinaciones de tratamientos a ejecutar: una tabla, una matriz factorial, un diagrama de flujo experimental o un esquema equivalente que permita identificar sin ambigüedad el número total de combinaciones, el número de réplicas por combinación, los factores y sus niveles, y la secuencia lógica de ejecución. En diseños con múltiples factores, se especifica si el diseño es factorial completo, factorial fraccionado o de otro tipo. El esquema es suficientemente detallado para que un investigador independiente pueda replicar la estructura experimental. &
& & & \\ \hline
\multicolumn{3}{|r|}{\cellcolor{white}\textbf{Subtotal Dimensión~II (máx.~12 puntos):}} & \multicolumn{4}{c|}{\cellcolor{white}\textbf{/12}} \\ \hline
\end{longtable}
\normalsize

\clearpage
\subsection{Muestreo, Tamaño Muestral y Especímenes de Ensayo}

\textit{Se evalúa si el proyecto justifica rigurosamente el número de especímenes, probetas o unidades de análisis que se ensayarán, si el cálculo del tamaño muestral responde a criterios estadísticos explícitos o a exigencias normativas documentadas, si los criterios de inclusión y exclusión de especímenes están definidos, y si las muestras son representativas del universo de estudio declarado. En Ingeniería Civil, la \enquote{muestra} no es un grupo de personas que responde encuestas, sino un conjunto de probetas, testigos, muestras de suelo, tramos viales, elementos estructurales u otras unidades físicas sometidas a ensayo.} Los criterios se aprecian en la Tabla~\ref{tab:rubrica-dim3-muestreo}.

\scriptsize
\begin{longtable}{|C{0.55cm}|L{2cm}|L{9cm}|C{0.3cm}|C{0.3cm}|C{0.3cm}|C{0.3cm}|}
\caption{Rúbrica de evaluación de la Dimensión~III: muestreo, tamaño muestral y especímenes de ensayo.}\label{tab:rubrica-dim3-muestreo}\\
\hline
\rowcolor{headerblue}
\thd{N.\textsuperscript{o}} & \thd{Criterio} & \thd{Indicador verificable} & \thd{3} & \thd{2} & \thd{1} & \thd{0} \\
\hline\endfirsthead
\hline
\rowcolor{headerblue}
\thd{N.\textsuperscript{o}} & \thd{Criterio} & \thd{Indicador verificable} & \thd{3} & \thd{2} & \thd{1} & \thd{0} \\
\hline\endhead
3.1 & Justificación del tamaño muestral con criterio estadístico o normativo \epref{Tamaño muestral en ingeniería} & El número de especímenes, probetas o unidades experimentales por tratamiento está justificado mediante al menos uno de los siguientes criterios: (a)~cálculo estadístico formal basado en la variabilidad esperada del material, el nivel de significancia, la potencia deseada y el tamaño del efecto mínimo detectable; (b)~requisito normativo explícito (p.\,ej., ASTM C39 exige un mínimo de tres cilindros por edad de curado; MTC E~132 estipula el número mínimo de ensayos CBR); o (c)~fundamentación bibliográfica que documente el tamaño muestral empleado en investigaciones análogas con justificación de su suficiencia. No se acepta la mera declaración de un número sin justificación alguna. &
& & & \\ \hline
\rowcolor{lightgray}
3.2 & Cálculo de la potencia estadística del diseño & El proyecto reporta, cuando el diseño lo requiere, un análisis de potencia estadística (\textit{power analysis}) que demuestre que el tamaño muestral propuesto es suficiente para detectar diferencias de magnitud prácticamente relevante con una probabilidad razonable (típicamente $1-\beta \geq 0{,}80$). El cálculo explicita los parámetros de entrada: nivel de significancia~$\alpha$, tamaño del efecto esperado~$d$ o~$\Delta$, desviación estándar estimada del material~$\sigma$, y número de grupos o tratamientos. Si el diseño se basa en un requisito normativo que fija el número mínimo de ensayos, se documenta dicho requisito como alternativa al cálculo de potencia. &
& & & \\ \hline
3.3 & Criterios de inclusión y exclusión de especímenes & El proyecto define explícitamente los criterios por los cuales un espécimen será incluido o excluido del análisis: rangos aceptables de densidad tras compactación, tolerancias dimensionales de probetas, condiciones de rechazo por defectos visibles (fisuras, segregación, contaminación), criterios de descarte de valores atípicos (\textit{outliers}) basados en procedimientos estadísticos documentados (prueba de Grubbs, criterio de Chauvenet, rango intercuartílico), y condiciones de repetición de ensayos cuando un espécimen falla por causa ajena al fenómeno investigado. Los criterios están definidos antes de la ejecución experimental, no establecidos \textit{post hoc}. &
& & & \\ \hline
\rowcolor{lightgray}
3.4 & Representatividad de las muestras respecto de la población o universo de estudio & Las muestras, especímenes o unidades de análisis son representativas del universo declarado en la delimitación del problema: si el estudio se refiere a suelos arcillosos de una zona geotécnica específica, las muestras provienen de esa zona y se documenta el procedimiento de muestreo en campo (calicatas, sondajes, extracción inalterada o alterada) conforme a normativa aplicable; si el estudio se refiere a un tipo de concreto, los materiales constituyentes son los disponibles en el contexto local y las condiciones de preparación replican las condiciones reales de obra o las condiciones normalizadas de laboratorio, según corresponda. La extrapolabilidad de los resultados está acotada por la representatividad de las muestras. &
& & & \\ \hline
\multicolumn{3}{|r|}{\cellcolor{white}\textbf{Subtotal Dimensión~III (máx.~12 puntos):}} & \multicolumn{4}{c|}{\cellcolor{white}\textbf{/12}} \\ \hline
\end{longtable}
\normalsize

\clearpage
\subsection{Técnicas, Instrumentos y Protocolos de Recolección de Datos}

\textit{Se evalúa si el proyecto identifica los ensayos de laboratorio, mediciones de campo o procedimientos analíticos que se emplearán para obtener los datos, si los protocolos están referenciados a normas técnicas reconocidas, si los instrumentos de medición están sujetos a calibración verificable y si existen procedimientos de control de calidad que garanticen la trazabilidad y confiabilidad de los datos primarios.} Los criterios se aprecian en la Tabla~\ref{tab:rubrica-dim4-instrumentos}.

\scriptsize
\begin{longtable}{|C{0.55cm}|L{2.8cm}|L{8cm}|C{0.3cm}|C{0.3cm}|C{0.3cm}|C{0.3cm}|}
\caption{Rúbrica de evaluación de la Dimensión~IV: técnicas, instrumentos y protocolos de recolección de datos.}\label{tab:rubrica-dim4-instrumentos}\\
\hline
\rowcolor{headerblue}
\thd{N.\textsuperscript{o}} & \thd{Criterio} & \thd{Indicador verificable} & \thd{3} & \thd{2} & \thd{1} & \thd{0} \\
\hline\endfirsthead
\hline
\rowcolor{headerblue}
\thd{N.\textsuperscript{o}} & \thd{Criterio} & \thd{Indicador verificable} & \thd{3} & \thd{2} & \thd{1} & \thd{0} \\
\hline\endhead
4.1 & Referencia a protocolos de ensayo normalizados pertinentes \epref{Protocolos normalizados} & Cada variable dependiente que se medirá mediante ensayo de laboratorio o procedimiento de campo tiene asociado un protocolo normalizado identificado con su designación completa: ASTM (p.\,ej., ASTM D1883 para CBR, ASTM C39 para resistencia a la compresión de cilindros de concreto), NTP, AASHTO, MTC, ISO, ACI, BS, EN, u otra norma técnica reconocida y aplicable. La referencia no es genérica (\enquote{se seguirán normas ASTM}) sino específica: se identifica la designación y el año de la versión vigente de cada norma para cada ensayo. Cuando un ensayo no dispone de norma aplicable, se describe el procedimiento no normalizado con suficiente detalle para garantizar su replicabilidad. &
& & & \\ \hline
\rowcolor{lightgray}
4.2 & Calibración, resolución y validez de los instrumentos de medición & El proyecto identifica los instrumentos de medición principales que se emplearán (prensas hidráulicas, balanzas, hornos de secado, tamices, extensómetros, deformímetros, equipos de ultrasonido, estaciones totales, LVDTs, u otros) y describe las condiciones de calibración vigente o prevista: certificados de calibración emitidos por laboratorios acreditados, periodicidad de calibración, resolución y rango de medición del instrumento, e incertidumbre instrumental declarada. No se asume que los instrumentos del laboratorio están calibrados sin evidencia ni se omite esta información. Si los ensayos se realizarán en un laboratorio externo, se identifica el laboratorio y su condición de acreditación. &
& & & \\ \hline
4.3 & Procedimientos de control de calidad de los ensayos & El proyecto describe los mecanismos de control de calidad que se aplicarán durante la fase experimental: preparación estandarizada de especímenes (moldeo, compactación, curado), verificación de condiciones ambientales del laboratorio, registro fotográfico o videográfico del proceso, hojas de registro prediseñadas para la captura de datos, ensayos de verificación o repetición cuando los resultados de un espécimen individual se desvían significativamente de la tendencia del grupo, y protocolos de conservación y transporte de muestras de campo al laboratorio. Los mecanismos de control son operativos y no declarativos. &
& & & \\ \hline
\rowcolor{lightgray}
4.4 & Coherencia entre los instrumentos declarados y las variables del estudio & Cada variable dependiente e independiente operacionalizada en el proyecto tiene asociado al menos un instrumento, ensayo o procedimiento de medición específico. No existen variables declaradas en los objetivos o hipótesis que carezcan de un instrumento o protocolo de medición en la sección metodológica, ni instrumentos o ensayos listados que no se correspondan con ninguna variable del estudio. La correspondencia variable--instrumento--norma--unidad de medida es completa, verificable y presentada de manera sistematizada (preferiblemente en una tabla de operacionalización de variables). &
& & & \\ \hline
\multicolumn{3}{|r|}{\cellcolor{white}\textbf{Subtotal Dimensión~IV (máx.~12 puntos):}} & \multicolumn{4}{c|}{\cellcolor{white}\textbf{/12}} \\ \hline
\end{longtable}
\normalsize

\clearpage
\subsection{Plan de Análisis de Datos y Tratamiento de la Incertidumbre}

\textit{Se evalúa si el proyecto especifica los métodos estadísticos o analíticos que se emplearán para procesar los datos experimentales, si se prevé la cuantificación y propagación de incertidumbres inherentes a los materiales y a los instrumentos, si los marcos analíticos son pertinentes al tipo de diseño propuesto y si el tesista distingue entre significancia estadística y relevancia práctica para la Ingeniería Civil. Esta dimensión verifica que el plan de análisis no sea una lista genérica de técnicas estadísticas sino un procedimiento articulado con los objetivos, las hipótesis y la naturaleza de los datos.} Los criterios se aprecian en la Tabla~\ref{tab:rubrica-dim5-analisis-datos}.

\scriptsize
\begin{longtable}{|C{0.55cm}|L{2cm}|L{10cm}|C{0.3cm}|C{0.3cm}|C{0.3cm}|C{0.3cm}|}
\caption{Rúbrica de evaluación de la Dimensión~V: plan de análisis de datos y tratamiento de la incertidumbre.}\label{tab:rubrica-dim5-analisis-datos}\\
\hline
\rowcolor{headerblue}
\thd{N.\textsuperscript{o}} & \thd{Criterio} & \thd{Indicador verificable} & \thd{3} & \thd{2} & \thd{1} & \thd{0} \\
\hline\endfirsthead
\hline
\rowcolor{headerblue}
\thd{N.\textsuperscript{o}} & \thd{Criterio} & \thd{Indicador verificable} & \thd{3} & \thd{2} & \thd{1} & \thd{0} \\
\hline\endhead
5.1 & Especificación de los métodos estadísticos vinculados a cada objetivo e hipótesis \epref{Plan de análisis} & El proyecto identifica los métodos estadísticos o analíticos que se emplearán para cada objetivo específico y para la contrastación de cada hipótesis: ANOVA de uno o varios factores cuando se comparan medias de tres o más grupos, prueba~\textit{t} cuando se comparan dos grupos, pruebas \textit{post hoc} (Tukey, Scheffé, Bonferroni, u otra) para identificar diferencias pareadas tras un ANOVA significativo, regresión lineal o no lineal cuando se modelan relaciones funcionales, superficie de respuesta o diseño compuesto central cuando se busca optimización formal. La selección del método se justifica en función del tipo de datos, el número de grupos, la distribución esperada y la naturaleza de la hipótesis. No se listan técnicas estadísticas genéricas sin vincularlas a objetivos específicos. &
& & & \\ \hline
\rowcolor{lightgray}
5.2 & Cuantificación y propagación de incertidumbres de medición y variabilidad del material & El proyecto prevé la cuantificación de las fuentes de incertidumbre relevantes: incertidumbre instrumental (resolución del equipo, incertidumbre expandida del certificado de calibración), variabilidad intrínseca del material (heterogeneidad del suelo, dispersión de resistencias del concreto, variabilidad granulométrica), e incertidumbre asociada al procedimiento de ensayo (repetibilidad y reproducibilidad). Se describe cómo se reportarán los resultados (media $\pm$ desviación estándar, intervalos de confianza, coeficiente de variación) y, cuando la complejidad del modelo lo requiera, se prevé la propagación de incertidumbres mediante métodos analíticos (ley de propagación de incertidumbres, GUM) o computacionales (simulación de Monte Carlo). &
& & & \\ \hline
5.3 & Coherencia entre el objetivo declarado y el método analítico propuesto & Si el proyecto declara como objetivo la \textit{optimización} de una propiedad (p.\,ej., \enquote{optimizar la dosificación de cemento para maximizar la resistencia}), el plan de análisis incluye un modelo formal de optimización: superficie de respuesta, diseño compuesto central, algoritmo de optimización, o método equivalente que permita identificar un óptimo en un dominio continuo. Si el objetivo es \textit{comparar} alternativas de un conjunto finito (p.\,ej., \enquote{comparar tres dosificaciones}), el plan de análisis emplea ANOVA y pruebas \textit{post hoc} sin confundir la selección de la mejor alternativa de un conjunto discreto con la identificación de un óptimo en un continuo. Si el objetivo es \textit{correlacionar}, se especifican coeficientes de correlación y modelos de regresión pertinentes. No existe discrepancia entre el verbo rector del objetivo y la técnica analítica propuesta. &
& & & \\ \hline
\rowcolor{lightgray}
5.4 & Distinción entre significancia estadística, variabilidad aleatoria y relevancia práctica para la Ingeniería Civil & El plan de análisis no se limita a reportar valores~$p$ o a declarar significancia estadística, sino que prevé la evaluación de la relevancia práctica de los resultados: tamaño del efecto, diferencias absolutas comparadas con umbrales normativos o de desempeño (p.\,ej., una diferencia estadísticamente significativa de 0{,}5\,\% en CBR puede ser irrelevante si ambos valores superan el umbral normativo), intervalos de confianza interpretados en el contexto de las especificaciones técnicas, y análisis de la magnitud de las diferencias en relación con la variabilidad típica del material. El tesista demuestra comprensión de que un resultado estadísticamente significativo no es necesariamente relevante para la práctica de la Ingeniería Civil, y viceversa. &
& & & \\ \hline
\multicolumn{3}{|r|}{\cellcolor{white}\textbf{Subtotal Dimensión~V (máx.~12 puntos):}} & \multicolumn{4}{c|}{\cellcolor{white}\textbf{/12}} \\ \hline
\end{longtable}
\normalsize

\newpage
En esta sección se presenta el consolidado de resultados por dimensión y la escala de calificación del diseño metodológico, resumidos en las tablas~\ref{tab:consolidado-metodologia} y~\ref{tab:escala-calificacion-metodologia}.

\begin{table}[htbp]
\centering
\scriptsize
\caption{Cuadro consolidado de resultados por dimensión del diseño metodológico.}\label{tab:consolidado-metodologia}
\begin{tabularx}{\textwidth}{|C{0.8cm}|L{7.5cm}|C{1.8cm}|C{1.8cm}|X|}
\hline
\rowcolor{headerblue}
\thd{Dim.} & \thd{Denominación} & \thd{Pts. Máx.} & \thd{Pts. Obt.} & \thd{Obs.} \\ \hline
I & Tipo, Nivel y Diseño de Investigación & 12 & & \\ \hline
\rowcolor{lightgray}
II & Operacionalización de Variables y Definición de Tratamientos Experimentales & 12 & & \\ \hline
III & Muestreo, Tamaño Muestral y Especímenes de Ensayo & 12 & & \\ \hline
\rowcolor{lightgray}
IV & Técnicas, Instrumentos y Protocolos de Recolección de Datos & 12 & & \\ \hline
V & Plan de Análisis de Datos y Tratamiento de la Incertidumbre & 12 & & \\ \hline
\rowcolor{white}
& \textbf{TOTAL GENERAL} & \textbf{60} & & \\ \hline
\end{tabularx}
\normalsize
\end{table}

\begin{table}[htbp]
\centering
\scriptsize
\caption{Escala de calificación del diseño metodológico.}\label{tab:escala-calificacion-metodologia}
\begin{tabularx}{\textwidth}{|C{2.5cm}|C{4.2cm}|X|}
\hline
\rowcolor{headerblue}
\thd{Rango} & \thd{Calificación} & \thd{Significado} \\ \hline
\textbf{54--60}\newline(90--100\,\%) &
\textcolor{csuf}{\textbf{METODOLOGÍA\newline PLENA}} &
El diseño metodológico es integral, técnicamente preciso y lógicamente articulado con los objetivos, las hipótesis y el marco teórico del proyecto. El tipo y nivel de investigación son coherentes con el diseño propuesto, las variables están operacionalizadas con rigor, el tamaño muestral está justificado, los protocolos de ensayo están normalizados, y el plan de análisis es pertinente y suficiente. La metodología constituye una base sólida y replicable para la ejecución de la investigación. \\ \hline
\rowcolor{lightgray}
\textbf{42--53}\newline(70--89\,\%) &
\textcolor{cace}{\textbf{METODOLOGÍA\newline SUFICIENTE}} &
El diseño metodológico presenta los componentes esenciales con fundamentación identificable, pero exhibe áreas específicas susceptibles de fortalecimiento en cuanto a precisión en la operacionalización, justificación del tamaño muestral, especificación de protocolos o articulación del plan de análisis con los objetivos. Las debilidades son corregibles sin reestructuración fundamental de la metodología. \\ \hline
\textbf{30--41}\newline(50--69\,\%) &
\textcolor{cdef}{\textbf{METODOLOGÍA\newline PARCIAL}} &
El diseño metodológico es incompleto o predominantemente declarativo. Se evidencian vacíos significativos: variables sin operacionalizar, tamaño muestral sin justificación, protocolos de ensayo no identificados, plan de análisis genérico o desarticulado de los objetivos, o inconsistencias entre el tipo de investigación declarado y el diseño descrito. Requiere fortalecimiento sustancial antes de la ejecución. \\ \hline
\rowcolor{lightgray}
\textbf{0--29}\newline($<$50\,\%) &
\textcolor{cins}{\textbf{METODOLOGÍA\newline INSUFICIENTE}} &
El diseño metodológico no demuestra consistencia verificable. La descripción es superficial o ausente, no existe articulación entre la metodología y los demás componentes del proyecto, las variables carecen de operacionalización, los especímenes no están definidos, los protocolos no están identificados y el plan de análisis es inexistente o inadecuado. Se requiere una reformulación fundamental de la sección metodológica. \\ \hline
\end{tabularx}
\normalsize
\end{table}

\clearpage
\section{Evaluación de Cientificidad del Proyecto}

{\footnotesize
La cientificidad de un proyecto de tesis de pregrado en Ingeniería Civil es la propiedad que permite distinguir un planteamiento orientado a la \textbf{generación de conocimiento científico} de uno que se limita a la \textbf{aplicación profesional rutinaria} de normas, procedimientos o tecnologías ya establecidas. Un proyecto de tesis científicamente formulado identifica un vacío de conocimiento genuino, propone hipótesis falsables, diseña una estrategia experimental o analítica con control de variables, anticipa el tratamiento cuantitativo de incertidumbres y aspira a producir resultados transferibles más allá del sitio, material o proyecto específico estudiado. La presente rúbrica operacionaliza diez principios epistemológicos de la contribución científica en ingeniería ---comprensión mecanicista, generalizabilidad, hipótesis falsables, control de variables, cuantificación de incertidumbres, transparencia metodológica, validación cruzada, innovación teórica o metodológica, ampliación del conocimiento y difusión académica--- en cinco dimensiones evaluables con veinte criterios verificables. Su propósito es proporcionar a los evaluadores un mecanismo sistemático y reproducible para determinar si un proyecto de tesis posee el nivel de cientificidad suficiente para ser aprobado como investigación de pregrado en Ingeniería Civil, diferenciándolo de un informe técnico, un expediente de diseño o un trabajo de aplicación profesional. La evaluación se aplica al proyecto \textit{antes} de su ejecución; un proyecto que supere el umbral de cientificidad podrá formularse como tesis de investigación. Este instrumento se fundamenta en los diez principios de la contribución científica en ingeniería (Arbulú, 2026), que establecen las características distintivas de la investigación académica frente a la práctica profesional (véase Tabla~\ref{tab:principios}).\par
}

\begin{table}[htbp]
\centering
\scriptsize
\caption{Principios de la contribución científica en Ingeniería Civil.}
\label{tab:principios}
\begin{tabularx}{\textwidth}{|C{0.8cm}|X|}
\hline
\rowcolor{headerblue}
\thd{N.°} & \thd{Principio} \\ \hline
1 & Comprensión mecanicista más allá de la adecuación empírica: el trabajo busca explicar los mecanismos físicos, químicos o mecánicos subyacentes al comportamiento observado, y no se limita a obtener correlaciones predictivas para fines de diseño. \\ \hline
\rowcolor{lightgray}
2 & Desarrollo de marcos generalizables, no soluciones específicas para un sitio u obra: los resultados son transferibles a clases de materiales, condiciones o dominios problemáticos más amplios que el caso particular estudiado. \\ \hline
3 & Hipótesis nulas explícitas y posible refutación estadística: se formulan hipótesis cuyo resultado no es predecible a priori, se anticipan criterios de refutación y se distingue entre efectos reales y variabilidad aleatoria. \\ \hline
\rowcolor{lightgray}
4 & Aislamiento y control sistemático de variables para establecer causalidad: el diseño experimental permite manipular la variable independiente, controlar las intervinientes e inferir relaciones causales, no meramente correlacionales. \\ \hline
5 & Cuantificación y propagación de incertidumbres mediante marcos probabilísticos: se reconoce, mide y propaga la incertidumbre de materiales, mediciones y condiciones de contorno, trascendiendo los factores de seguridad prescriptivos. \\ \hline
\rowcolor{lightgray}
6 & Trazabilidad mediante transparencia metodológica y detalle suficiente: la documentación permite que un investigador independiente reproduzca los procedimientos experimentales, numéricos o analíticos descritos. \\ \hline
7 & Validación cruzada de múltiples métodos independientes: los resultados se sustentan con evidencia convergente de al menos dos fuentes metodológicas independientes (laboratorio, campo, modelación, soluciones analíticas). \\ \hline
\rowcolor{lightgray}
8 & Avances en marcos metodológicos o teóricos: el trabajo aporta innovación verificable ---nuevos procedimientos, modelos modificados, marcos analíticos novedosos o integración original de enfoques existentes---. \\ \hline
9 & Ampliación del conocimiento a materiales y condiciones no estudiados: la investigación se dirige hacia dominios donde el conocimiento existente es insuficiente, abordando materiales emergentes, condiciones singulares o comportamientos no explicados. \\ \hline
\rowcolor{lightgray}
10 & Difusión a través de medios académicos revisados por pares: los resultados tienen potencial de publicación en revistas indexadas, congresos o medios de difusión científica reconocidos por la comunidad disciplinar. \\ \hline
\end{tabularx}
\end{table}

Las dimensiones de evaluación de la cientificidad se resumen en la Tabla~\ref{tab:dim-cientificidad}, mientras que la escala de valoración y sus descriptores operativos se presentan en la Tabla~\ref{tab:escala-cientificidad}.
\begin{table}[htbp]
\centering
\scriptsize
\caption{Dimensiones de evaluación de la cientificidad y alcance de cada dimensión.}
\label{tab:dim-cientificidad}
\begin{tabularx}{\textwidth}{|C{0.8cm}|L{4.5cm}|X|}
\hline
\rowcolor{headerblue}
\thd{Dim.} & \thd{Denominación} & \thd{Alcance} \\ \hline
I & Orientación epistemológica y alcance científico & Verifica que el proyecto plantee un problema de conocimiento genuino, aspire a la comprensión explicativa del fenómeno y no meramente predictiva, proponga resultados generalizables más allá del caso particular y delimite explícitamente el dominio de validez de sus hallazgos. \\ \hline
\rowcolor{lightgray}
II & Rigor en el diseño metodológico e hipótesis & Evalúa la formulación de hipótesis explícitas y falsables, la capacidad del diseño experimental o analítico para aislar y controlar variables, la diferenciación entre correlación y causalidad, y la suficiencia del diseño para contrastar las hipótesis planteadas. \\ \hline
III & Tratamiento de incertidumbres y análisis cuantitativo & Examina si el proyecto anticipa la cuantificación de incertidumbres en materiales, mediciones y condiciones de contorno, el empleo de marcos probabilísticos o estadísticos, la propagación de incertidumbres y un análisis que trascienda la estadística meramente descriptiva. \\ \hline
\rowcolor{lightgray}
IV & Transparencia metodológica y validación & Evalúa la documentación metodológica suficiente para la reproducibilidad, la validación cruzada con métodos independientes, el compromiso con el reporte transparente de limitaciones y anomalías, y la trazabilidad completa de datos, ensayos y procedimientos. \\ \hline
V & Contribución al conocimiento disciplinar & Verifica la identificación de un vacío de conocimiento genuino en el estado del arte, la anticipación de innovación metodológica o teórica, la ampliación del conocimiento hacia dominios poco estudiados y el potencial de difusión académica de los resultados esperados. \\ \hline
\end{tabularx}
\end{table}
\begin{table}[htbp]
\centering
\scriptsize
\caption{Escala de valoración (puntaje) y descriptor operativo para la evaluación de la cientificidad.}
\label{tab:escala-cientificidad}
\begin{tabularx}{\textwidth}{|C{2.8cm}|C{1cm}|X|}
\hline
\rowcolor{headerblue}
\thd{Nivel} & \thd{Pts.} & \thd{Descriptor operativo} \\ \hline
\textcolor{csuf}{\textbf{Excelente}} & \textbf{3} & El proyecto satisface plenamente el criterio evaluado. El planteamiento es explícito, técnicamente preciso, epistemológicamente fundamentado y coherente con los principios de la investigación científica en ingeniería. No se identifican ambigüedades, omisiones ni vicios metodológicos. \\ \hline
\rowcolor{lightgray}
\textcolor{cace}{\textbf{Aceptable}} & \textbf{2} & El criterio es identificable pero se cumple de manera parcial o incompleta. El proyecto aborda el aspecto evaluado sin la profundidad, la rigurosidad o el sustento suficiente. Susceptible de fortalecimiento con correcciones específicas sin reestructuración fundamental. \\ \hline
\textcolor{cdef}{\textbf{Deficiente}} & \textbf{1} & El cumplimiento del criterio es tangencial o meramente declarativo. El planteamiento presenta vicios metodológicos, confusiones epistemológicas o inconsistencias que comprometen la cientificidad del proyecto en el aspecto evaluado. \\ \hline
\rowcolor{lightgray}
\textcolor{cins}{\textbf{Ausente}} & \textbf{0} & El proyecto no evidencia consideración alguna del criterio evaluado. Existe omisión total del aspecto o un planteamiento que corresponde a práctica profesional estándar sin componente científico identificable. \\ \hline
\end{tabularx}
\end{table}

\clearpage
\subsection{Orientación Epistemológica y Alcance Científico}

\textit{Se evalúa si el proyecto se orienta hacia la generación de conocimiento científico transferible, si aspira a la comprensión explicativa de los fenómenos y no únicamente a la predicción empírica para fines de diseño, y si delimita el dominio de validez de sus hallazgos. Esta dimensión diferencia un proyecto de investigación de un expediente técnico o un informe profesional.} Véase la Tabla~\ref{tab:rubrica-dim1-epistemologica}.
{\scriptsize
\begin{longtable}{|C{0.55cm}|L{2cm}|L{8cm}|C{0.65cm}|C{0.65cm}|C{0.65cm}|C{0.65cm}|}
\caption{Rúbrica de evaluación de la Dimensión~I: orientación epistemológica y alcance científico.}\label{tab:rubrica-dim1-epistemologica}\\
\hline
\rowcolor{headerblue}
\thd{N.\textsuperscript{o}} & \thd{Criterio} & \thd{Indicador verificable} & \thd{3} & \thd{2} & \thd{1} & \thd{0} \\
\hline\endfirsthead
\hline
\rowcolor{headerblue}
\thd{N.\textsuperscript{o}} & \thd{Criterio} & \thd{Indicador verificable} & \thd{3} & \thd{2} & \thd{1} & \thd{0} \\
\hline\endhead
1.1 & Existencia de un problema de conocimiento genuino, diferenciado de un problema de aplicación profesional \epref{Principio~1} & El proyecto formula un problema de conocimiento genuino: una brecha, vacío, contradicción o insuficiencia en el estado del arte de la Ingeniería Civil que requiere investigación para ser resuelta. El problema no se reduce a una necesidad de aplicación profesional (p.\,ej., ``diseñar un pavimento rígido para la Av.\ de la Cultura'' no es un problema de conocimiento; ``determinar la influencia del porcentaje de fibras de polipropileno sobre la tenacidad a la flexión de concretos de alta resistencia en condiciones de curado a~3\,800~m\,s.\,n.\,m.'' sí lo es). Se identifica con claridad qué se desconoce y por qué es necesario investigarlo. &
& & & \\ \hline
\rowcolor{lightgray}
1.2 & Orientación hacia la comprensión explicativa del fenómeno, más allá de la predicción empírica \epref{Principio~1} & El proyecto aspira a comprender los mecanismos físicos, químicos, mecánicos o estructurales subyacentes al comportamiento del material, sistema o fenómeno estudiado, y no se limita a obtener datos empíricos para fines de diseño. El planteamiento distingue entre medir un parámetro y explicar por qué adopta determinados valores; entre verificar el cumplimiento de una norma y entender las causas del comportamiento observado (p.\,ej., no solo medir el CBR sino explicar cómo la reacción puzolánica del estabilizante modifica la microestructura del suelo). &
& & & \\ \hline
1.3 & Generalizabilidad y transferibilidad de los resultados esperados \epref{Principio~2} & El proyecto anticipa que sus resultados serán transferibles a contextos distintos al sitio, obra o proyecto específico estudiado, aplicándose a clases de materiales, condiciones de carga, tipologías estructurales o dominios problemáticos más amplios. El planteamiento no se agota en la solución de un caso particular, sino que busca generar conocimiento aprovechable por la comunidad disciplinar. La transferibilidad propuesta es fundamentada, no genérica ni retórica. &
& & & \\ \hline
\rowcolor{lightgray}
1.4 & Delimitación explícita del dominio de aplicabilidad de los hallazgos esperados \epref{Principio~2} & El proyecto establece explícitamente las condiciones de contorno bajo las cuales sus conclusiones serán válidas: rango de dosificaciones, tipo de material, clasificación del suelo, condiciones climáticas, nivel de solicitación, geometría del elemento, o cualquier otra variable de contexto que defina el dominio de aplicabilidad. El planteamiento no promete resultados universales sin acotación ni se restringe a un solo especimen sin justificación de representatividad. &
& & & \\ \hline
\multicolumn{3}{|r|}{\cellcolor{white}\textbf{Subtotal Dimensión~I (máx.~12 puntos):}} & \multicolumn{4}{c|}{\cellcolor{white}\textbf{/12}} \\ \hline
\end{longtable}
}

\clearpage
\subsection{Rigor en el Diseño Metodológico e Hipótesis}

\textit{Se evalúa si el proyecto formula hipótesis genuinamente falsables, si el diseño experimental o analítico permite aislar y controlar las variables bajo estudio, si distingue entre correlación y causalidad, y si la estrategia metodológica es suficiente para contrastar las hipótesis planteadas. Esta dimensión identifica vicios frecuentes como hipótesis triviales, tautológicas o no falsables, y diseños que no permiten inferencia causal.} Véase la Tabla~\ref{tab:rubrica-dim2-rigor}.
{\scriptsize
\begin{longtable}{|C{0.55cm}|L{2cm}|L{8cm}|C{0.65cm}|C{0.65cm}|C{0.65cm}|C{0.65cm}|}
\caption{Rúbrica de evaluación de la Dimensión~II: rigor en el diseño metodológico e hipótesis.}\label{tab:rubrica-dim2-rigor}\\
\hline
\rowcolor{headerblue}
\thd{N.\textsuperscript{o}} & \thd{Criterio} & \thd{Indicador verificable} & \thd{3} & \thd{2} & \thd{1} & \thd{0} \\
\hline\endfirsthead
\hline
\rowcolor{headerblue}
\thd{N.\textsuperscript{o}} & \thd{Criterio} & \thd{Indicador verificable} & \thd{3} & \thd{2} & \thd{1} & \thd{0} \\
\hline\endhead
2.1 & Formulación de hipótesis explícitas, no triviales y falsables \epref{Principio~3} & El proyecto formula hipótesis cuyo resultado no es predecible a priori por cualquier ingeniero competente. La hipótesis no es trivial (p.\,ej., ``a mayor contenido de cemento, mayor resistencia'' es una tautología), ni es un objetivo disfrazado (p.\,ej., ``se determinará la resistencia del concreto'' no es una hipótesis). La hipótesis es refutable: existe un resultado experimental plausible que la invalidaría. Se anticipan, al menos conceptualmente, la hipótesis nula y la alternativa, aunque la formalización estadística se reserve para la ejecución. &
& & & \\ \hline
\rowcolor{lightgray}
2.2 & Diseño experimental o analítico con aislamiento y control sistemático de variables \epref{Principio~4} & El proyecto propone un diseño que permite aislar el efecto de la variable independiente sobre la dependiente, controlando variables intervinientes. Se identifican las variables a controlar (temperatura de curado, granulometría, relación agua/cemento, energía de compactación, etc.) y se describe cómo se mantendrán constantes o se monitorizarán. El diseño no es meramente observacional cuando la ingeniería permite experimentación controlada; no confunde replicación con repetición; y especifica el número de niveles, tratamientos y réplicas con justificación técnica. &
& & & \\ \hline
2.3 & Diferenciación explícita entre correlación y causalidad en el planteamiento \epref{Principio~4} & El proyecto distingue entre asociaciones estadísticas y relaciones causales. Si se propone establecer causalidad, el diseño experimental lo permite (control de variables, manipulación deliberada de la variable independiente). Si el diseño es observacional o correlacional, el proyecto lo reconoce explícitamente y no anticipa conclusiones causales a partir de coeficientes de correlación o determinación. El planteamiento no incurre en el vicio de asumir que $R^2$ elevado implica causalidad sin control experimental ni base mecanicista. &
& & & \\ \hline
\rowcolor{lightgray}
2.4 & Suficiencia del diseño para la contrastación de las hipótesis planteadas \epref{Principios~3--4} & El diseño metodológico propuesto es coherente con las hipótesis formuladas y suficiente para contrastarlas. Los ensayos, mediciones o análisis seleccionados miden efectivamente las variables involucradas en la hipótesis. El número de muestras, probetas o especímenes permite la aplicación de pruebas estadísticas con potencia adecuada. No se proponen metodologías desconectadas de la hipótesis ni se plantean hipótesis que el diseño experimental no puede evaluar. &
& & & \\ \hline
\multicolumn{3}{|r|}{\cellcolor{white}\textbf{Subtotal Dimensión~II (máx.~12 puntos):}} & \multicolumn{4}{c|}{\cellcolor{white}\textbf{/12}} \\ \hline
\end{longtable}
}

\clearpage
\subsection{Tratamiento de Incertidumbres y Análisis Cuantitativo}

\textit{Se evalúa si el proyecto anticipa la cuantificación de incertidumbres en los materiales, mediciones y condiciones de contorno, si propone marcos probabilísticos o estadísticos rigurosos, y si el tratamiento de datos trasciende la estadística meramente descriptiva. Esta dimensión diferencia la investigación científica ---que reconoce, cuantifica y propaga la incertidumbre--- de la práctica profesional que se limita a factores de seguridad prescriptivos.} Véase la Tabla~\ref{tab:rubrica-dim3-incertidumbres}.
{\scriptsize
\begin{longtable}{|C{0.55cm}|L{2cm}|L{8cm}|C{0.4cm}|C{0.4cm}|C{0.4cm}|C{0.4cm}|}
\caption{Rúbrica de evaluación de la Dimensión~III: tratamiento de incertidumbres y análisis cuantitativo.}\label{tab:rubrica-dim3-incertidumbres}\\
\hline
\rowcolor{headerblue}
\thd{N.\textsuperscript{o}} & \thd{Criterio} & \thd{Indicador verificable} & \thd{3} & \thd{2} & \thd{1} & \thd{0} \\
\hline\endfirsthead
\hline
\rowcolor{headerblue}
\thd{N.\textsuperscript{o}} & \thd{Criterio} & \thd{Indicador verificable} & \thd{3} & \thd{2} & \thd{1} & \thd{0} \\
\hline\endhead
3.1 & Anticipación de la cuantificación explícita de incertidumbres \epref{Principio~5} & El proyecto reconoce las fuentes de incertidumbre inherentes a la investigación propuesta ---variabilidad natural de los materiales, errores de medición, condiciones de contorno no controlables, heterogeneidad del suelo o del concreto, tolerancias de fabricación--- y anticipa su cuantificación mediante indicadores apropiados (desviación estándar, coeficiente de variación, intervalos de confianza, incertidumbre expandida). El planteamiento no asume que las mediciones serán deterministas ni que un solo ensayo es representativo. &
& & & \\ \hline
\rowcolor{lightgray}
3.2 & Propuesta de marcos probabilísticos o estadísticos rigurosos \epref{Principio~5} & El proyecto propone el uso de herramientas estadísticas o probabilísticas apropiadas para la naturaleza de los datos: análisis de varianza (ANOVA), pruebas de hipótesis paramétricas o no paramétricas, regresión con análisis de residuos, diseño factorial, simulación de Monte Carlo, métodos de fiabilidad o inferencia bayesiana, según corresponda. El planteamiento no se limita a promedios y desviaciones estándar sin contrastación formal, ni sustituye el análisis estadístico por factores de seguridad prescriptivos. &
& & & \\ \hline
3.3 & Propagación de incertidumbres a través de los análisis propuestos \epref{Principio~5} & El proyecto anticipa cómo las incertidumbres de las variables de entrada se propagarán a los resultados y conclusiones. Si se propone un modelo analítico, numérico o empírico, se prevé evaluar la sensibilidad de los resultados a la variabilidad de los parámetros de entrada. El planteamiento no trata las variables de entrada como valores únicos deterministas cuando la naturaleza del material o del fenómeno implica variabilidad significativa. &
& & & \\ \hline
\rowcolor{lightgray}
3.4 & Análisis estadístico que trasciende lo descriptivo para la contrastación de hipótesis \epref{Principios~3 y~5} & El proyecto propone un análisis cuantitativo que va más allá de la descripción estadística de los datos (tablas de resumen, gráficos de barras, promedios) y anticipa pruebas inferenciales para la contrastación de las hipótesis formuladas. Se identifican las pruebas estadísticas pertinentes, el nivel de significancia previsto y los criterios de decisión. El tratamiento estadístico no es decorativo sino funcional para distinguir efectos reales de variabilidad aleatoria. &
& & & \\ \hline
\multicolumn{3}{|r|}{\cellcolor{white}\textbf{Subtotal Dimensión~III (máx.~12 puntos):}} & \multicolumn{4}{c|}{\cellcolor{white}\textbf{/12}} \\ \hline
\end{longtable}
}
\clearpage

\subsection{Transparencia Metodológica y Validación}

\textit{Se evalúa si el proyecto anticipa la documentación metodológica suficiente para que un investigador independiente pueda reproducir los procedimientos, si propone validación cruzada con métodos independientes, si se compromete con el reporte transparente de limitaciones y anomalías, y si garantiza la trazabilidad completa de datos y procedimientos. Esta dimensión verifica los requisitos de reproducibilidad y rigor documental propios de la investigación científica.} Véase la Tabla~\ref{tab:rubrica-dim4-transparencia}.
{\scriptsize
\begin{longtable}{|C{0.55cm}|L{2cm}|L{8cm}|C{0.65cm}|C{0.65cm}|C{0.65cm}|C{0.65cm}|}
\caption{Rúbrica de evaluación de la Dimensión~IV: transparencia metodológica y validación.}\label{tab:rubrica-dim4-transparencia}\\
\hline
\rowcolor{headerblue}
\thd{N.\textsuperscript{o}} & \thd{Criterio} & \thd{Indicador verificable} & \thd{3} & \thd{2} & \thd{1} & \thd{0} \\
\hline\endfirsthead
\hline
\rowcolor{headerblue}
\thd{N.\textsuperscript{o}} & \thd{Criterio} & \thd{Indicador verificable} & \thd{3} & \thd{2} & \thd{1} & \thd{0} \\
\hline\endhead
4.1 & Documentación metodológica suficiente para la reproducibilidad \epref{Principio~6} & El proyecto describe los procedimientos experimentales, numéricos o analíticos con el detalle suficiente para que un investigador independiente pueda reproducirlos: normas de ensayo aplicables con su designación completa, equipos e instrumentos con sus especificaciones y rangos, preparación de muestras paso a paso, condiciones de curado o acondicionamiento, parámetros de modelación numérica, software y versiones. La descripción no se limita a enunciar el nombre del ensayo (p.\,ej., ``se realizará el ensayo de CBR'') sin especificar la norma, el método de compactación, las condiciones de inmersión ni los criterios de aceptación. &
& & & \\ \hline
\rowcolor{lightgray}
4.2 & Validación cruzada con métodos independientes \epref{Principio~7} & El proyecto anticipa la validación de sus resultados mediante al menos dos fuentes de evidencia independientes: resultados de laboratorio contrastados con modelación numérica, ensayos de campo verificados con soluciones analíticas, métodos directos comparados con métodos indirectos, o correlación con bases de datos publicadas. El planteamiento no depende de un único enfoque metodológico para sustentar sus conclusiones. Si la validación cruzada no es viable, se justifica técnicamente la limitación y se proponen alternativas parciales. &
& & & \\ \hline
4.3 & Compromiso con el reporte transparente de limitaciones, anomalías y resultados negativos \epref{Principios~6--7} & El proyecto se compromete explícitamente con la transparencia científica: anticipa que se reportarán las limitaciones del estudio, las anomalías observadas, los resultados que contradigan las hipótesis y las discrepancias con las predicciones teóricas. El planteamiento no está diseñado para confirmar una conclusión predeterminada ni omite la posibilidad de resultados negativos. Se reconoce que un resultado que refute la hipótesis es científicamente tan valioso como uno que la confirme. &
& & & \\ \hline
\rowcolor{lightgray}
4.4 & Trazabilidad completa de datos, ensayos y procedimientos \epref{Principio~6} & El proyecto anticipa un sistema de trazabilidad que permita rastrear cada resultado hasta los datos primarios: registro de muestras con identificación única, bitácoras de ensayo, calibración de equipos con certificados vigentes, condiciones ambientales durante los ensayos, y almacenamiento de datos en bruto. La trazabilidad propuesta no se limita a presentar resultados finales sin posibilidad de verificación de los datos de origen. Se distingue entre la trazabilidad propia de la investigación y la mera cita de normas de ensayo. &
& & & \\ \hline
\multicolumn{3}{|r|}{\cellcolor{white}\textbf{Subtotal Dimensión~IV (máx.~12 puntos):}} & \multicolumn{4}{c|}{\cellcolor{white}\textbf{/12}} \\ \hline
\end{longtable}
}
\clearpage
\subsection{Contribución al Conocimiento del Campo de la Ingeniería Civil}
\textit{Se evalúa si el proyecto identifica un vacío de conocimiento genuino en el estado del arte ---distinto de una simple laguna de cobertura geográfica---, si anticipa innovación metodológica o teórica, si amplía el conocimiento hacia dominios poco estudiados y si los resultados esperados tienen potencial de difusión académica. Esta dimensión verifica que el proyecto aspire a contribuir al acervo científico de la Ingeniería Civil y no solo a resolver un problema local de infraestructura.} Véase la Tabla~\ref{tab:rubrica-dim5-contribucion}.
{\scriptsize
\begin{longtable}{|C{0.55cm}|L{2cm}|L{8cm}|C{0.65cm}|C{0.65cm}|C{0.65cm}|C{0.65cm}|}
\caption{Rúbrica de evaluación de la Dimensión~V: contribución al conocimiento disciplinar.}\label{tab:rubrica-dim5-contribucion}\\
\hline
\rowcolor{headerblue}
\thd{N.\textsuperscript{o}} & \thd{Criterio} & \thd{Indicador verificable} & \thd{3} & \thd{2} & \thd{1} & \thd{0} \\
\hline\endfirsthead
\hline
\rowcolor{headerblue}
\thd{N.\textsuperscript{o}} & \thd{Criterio} & \thd{Indicador verificable} & \thd{3} & \thd{2} & \thd{1} & \thd{0} \\
\hline\endhead
5.1 & Identificación de un vacío de conocimiento genuino en el estado del arte \epref{Principios~8--9} & El proyecto identifica críticamente un vacío de conocimiento científico específico: un mecanismo no explicado, un comportamiento no modelado, un material no caracterizado bajo las condiciones propuestas, una contradicción entre modelos existentes, o una limitación reconocida en la literatura. El vacío no se reduce a ``no se ha hecho este estudio en esta localidad'' (laguna geográfica) ni a ``no se ha usado este material en esta obra'' (laguna de aplicación). Se fundamenta bibliográficamente qué se sabe, qué no se sabe y por qué lo que no se sabe es científicamente relevante. &
& & & \\ \hline
\rowcolor{lightgray}
5.2 & Innovación metodológica o teórica anticipada \epref{Principio~8} & El proyecto anticipa algún grado de innovación verificable: un procedimiento experimental adaptado a condiciones singulares, la combinación original de técnicas existentes, la modificación de un modelo constitutivo o analítico, la propuesta de un protocolo de ensayo no estandarizado, o la integración de enfoques previamente separados. La innovación es específica y verificable, no genérica (p.\,ej., no basta con declarar que ``se contribuirá a la ingeniería'' sin precisar en qué consiste la innovación). El nivel de innovación es razonable para una tesis de pregrado. &
& & & \\ \hline
5.3 & Ampliación del conocimiento hacia dominios poco estudiados \epref{Principio~9} & El proyecto se dirige hacia un dominio donde el conocimiento existente es insuficiente o inexistente: materiales emergentes o no convencionales, condiciones extremas de altitud, temperatura, sismicidad o agresividad química, suelos singulares (orgánicos, volcánicos, diatomáceos, colapsables), técnicas constructivas no caracterizadas experimentalmente, o combinaciones de materiales no evaluadas previamente. La singularidad del dominio se justifica técnicamente y no es un artificio retórico para disfrazar un estudio convencional. &
& & & \\ \hline
\rowcolor{lightgray}
5.4 & Potencial de difusión académica y aporte a la comunidad científica \epref{Principio~10} & El proyecto genera resultados con potencial de difusión en medios académicos: los hallazgos esperados son de interés para la comunidad científica más allá del contexto local, la metodología es reproducible por otros investigadores, las conclusiones anticipadas aportan evidencia a debates abiertos en la disciplina o llenan vacíos documentados en la literatura. El proyecto no se limita a producir un documento de archivo institucional sin proyección académica. Se valora la identificación explícita de revistas, congresos o bases de datos donde los resultados podrían difundirse. &
& & & \\ \hline
\multicolumn{3}{|r|}{\cellcolor{white}\textbf{Subtotal Dimensión~V (máx.~12 puntos):}} & \multicolumn{4}{c|}{\cellcolor{white}\textbf{/12}} \\ \hline
\end{longtable}
}
\newpage
En esta sección se presenta el consolidado de resultados por dimensión y la escala de calificación de la cientificidad, resumidos en las tablas~\ref{tab:consolidado-cientificidad} y~\ref{tab:escala-calificacion-cientificidad}.
\begin{table}[htbp]
\centering
\scriptsize
\caption{Cuadro consolidado de resultados por dimensión de la cientificidad}\label{tab:consolidado-cientificidad}
\begin{tabularx}{\textwidth}{|C{0.8cm}|L{7.5cm}|C{1.8cm}|C{1.8cm}|X|}
\hline
\rowcolor{headerblue}
\thd{Dim.} & \thd{Denominación} & \thd{Pts. Máx.} & \thd{Pts. Obt.} & \thd{Obs.} \\ \hline
I & Orientación Epistemológica y Alcance Científico & 12 & & \\ \hline
\rowcolor{lightgray}
II & Rigor en el Diseño Metodológico e Hipótesis & 12 & & \\ \hline
III & Tratamiento de Incertidumbres y Análisis Cuantitativo & 12 & & \\ \hline
\rowcolor{lightgray}
IV & Transparencia Metodológica y Validación & 12 & & \\ \hline
V & Contribución al Conocimiento Disciplinar & 12 & & \\ \hline
\rowcolor{white}
& \textbf{TOTAL GENERAL} & \textbf{60} & & \\ \hline
\end{tabularx}
\normalsize
\end{table}
\begin{table}[htbp]
\centering
\scriptsize
\caption{Escala de calificación de la cientificidad del proyecto de tesis}\label{tab:escala-calificacion-cientificidad}
\begin{tabularx}{\textwidth}{|C{2.5cm}|C{4.2cm}|X|}
\hline
\rowcolor{headerblue}
\thd{Rango} & \thd{Calificación} & \thd{Significado} \\ \hline
\textbf{54--60}\newline(90--100\,\%) &
\textcolor{csuf}{\textbf{CIENTIFICIDAD\newline PLENA}} &
El proyecto de tesis presenta un planteamiento plenamente científico. Formula un problema de conocimiento genuino, propone hipótesis falsables con diseño experimental riguroso, anticipa el tratamiento cuantitativo de incertidumbres, garantiza transparencia metodológica y validación cruzada, e identifica una contribución específica al conocimiento disciplinar. El proyecto está en condiciones de ser aprobado como tesis de investigación. \\ \hline
\rowcolor{lightgray}
\textbf{42--53}\newline(70--89\,\%) &
\textcolor{cace}{\textbf{CIENTIFICIDAD\newline SUFICIENTE}} &
El proyecto presenta los componentes esenciales de una investigación científica en Ingeniería Civil, con áreas específicas susceptibles de fortalecimiento en cuanto a rigurosidad metodológica, tratamiento de incertidumbres o fundamentación del vacío de conocimiento. Las debilidades son corregibles sin reestructuración fundamental. El proyecto puede ser aprobado con correcciones específicas. \\ \hline
\textbf{30--41}\newline(50--69\,\%) &
\textcolor{cdef}{\textbf{CIENTIFICIDAD\newline PARCIAL}} &
El proyecto combina elementos de investigación científica con vicios propios de la práctica profesional estándar: hipótesis triviales o tautológicas, ausencia de control de variables, confusión entre correlación y causalidad, vacío geográfico en lugar de vacío de conocimiento, o marco teórico de tipo catálogo. Requiere fortalecimiento sustancial para ser considerado tesis de investigación. \\ \hline
\rowcolor{lightgray}
\textbf{0--29}\newline($<$50\,\%) &
\textcolor{cins}{\textbf{CIENTIFICIDAD\newline INSUFICIENTE}} &
El proyecto no demuestra cientificidad verificable. El planteamiento corresponde esencialmente a un trabajo de aplicación profesional: diseño de infraestructura, verificación normativa, expediente técnico o informe de campo sin componente de generación de conocimiento. Se requiere una reformulación fundamental del proyecto para que pueda constituir una tesis de investigación. \\ \hline
\end{tabularx}
\normalsize
\end{table}

\clearpage
\section{Evaluación de la Reproducibilidad y Trazabilidad}

{\footnotesize
La reproducibilidad ---la capacidad de que un investigador independiente obtenga resultados consistentes siguiendo el mismo protocolo--- y la trazabilidad ---la posibilidad de rastrear cada resultado hasta los datos primarios, materiales, equipos y condiciones que lo originaron--- constituyen dos pilares inseparables de la calidad científica en Ingeniería Civil experimental. En tesis de pregrado que involucran ensayos de laboratorio, campañas de campo o modelación numérica, la reproducibilidad no se logra con la mera declaración de compromiso ni con la cita genérica de normas de ensayo: exige que el proyecto especifique, con precisión operativa, la identidad química y física de los materiales empleados, el protocolo experimental paso a paso, las condiciones ambientales y de preparación de muestras, el sistema de registro de datos primarios, la suficiencia estadística del plan muestral y la ausencia de ambigüedades que impidan la réplica por terceros. Esta rúbrica evalúa el grado en que el proyecto de tesis proporciona la información técnica necesaria y suficiente para que sus procedimientos sean reproducibles y sus resultados trazables, distinguiéndose de los instrumentos previos que evalúan el compromiso declarativo con la gestión de datos (Sección~III de Observaciones Normativas) y la anticipación genérica de un sistema de trazabilidad (Dimensión~IV de Cientificidad). Aquí se verifica el contenido técnico-operativo que hace posibles la reproducibilidad y la trazabilidad, no su declaración de intención. Las dimensiones de evaluación de la reproducibilidad y trazabilidad se resumen en la Tabla~\ref{tab:dim-reproducibilidad}, mientras que la escala de valoración y sus descriptores operativos se presentan en la Tabla~\ref{tab:escala-reproducibilidad}.\par
}

\begin{table}[htbp]
\centering
\scriptsize
\caption{Dimensiones de evaluación de la reproducibilidad y trazabilidad y sus dimensiones.}
\label{tab:dim-reproducibilidad}
\begin{tabularx}{\textwidth}{|C{0.8cm}|L{4.5cm}|X|}
\hline
\rowcolor{headerblue}
\thd{Dim.} & \thd{Denominación} & \thd{Alcance} \\ \hline
I & Caracterización y trazabilidad de materiales e insumos & Verifica que todos los materiales, reactivos e insumos estén identificados con composición química, pureza, procedencia y justificación de selección, de modo que otro investigador pueda adquirir materiales equivalentes y reproducir las condiciones experimentales. \\ \hline
\rowcolor{lightgray}
II & Especificación del protocolo experimental & Evalúa si el proyecto describe el procedimiento experimental con suficiente detalle operativo: secuencia de pasos, equipos e instrumentos con estado de calibración, condiciones ambientales y protocolos de preparación de muestras y probetas. \\ \hline
III & Registro documental de datos primarios & Examina la estructura del sistema de registro de datos: identificación única de especímenes, formato de las hojas de ensayo, separación entre datos en bruto y datos procesados, y documentación fotográfica o audiovisual del procedimiento. \\ \hline
\rowcolor{lightgray}
IV & Suficiencia estadística y control de variables para la replicación & Verifica que el plan muestral esté justificado estadísticamente, que las variables intervinientes se controlen y registren, que existan criterios explícitos para datos atípicos y que las dosificaciones se expresen en unidades reproducibles. \\ \hline
V & Verificabilidad externa y potencial de replicación independiente & Evalúa si la información proporcionada es suficiente para que un investigador externo replique el estudio sin contacto con el autor: completitud de normas citadas, ausencia de ambigüedades críticas, y disponibilidad declarada de datos y protocolos de soporte. \\ \hline
\end{tabularx}
\end{table}

\begin{table}[htbp]
\centering
\scriptsize
\caption{Escala de valoración (puntaje) y descriptor operativo para la evaluación de la reproducibilidad y trazabilidad.}
\label{tab:escala-reproducibilidad}
\begin{tabularx}{\textwidth}{|C{2.8cm}|C{1cm}|X|}
\hline
\rowcolor{headerblue}
\thd{Nivel} & \thd{Pts.} & \thd{Descriptor operativo} \\ \hline
\textcolor{csuf}{\textbf{Excelente}} & \textbf{3} & El proyecto satisface plenamente el criterio evaluado. La información proporcionada es explícita, técnicamente precisa y suficiente para garantizar la reproducibilidad del procedimiento y la trazabilidad de los resultados sin ambigüedades. \\ \hline
\rowcolor{lightgray}
\textcolor{cace}{\textbf{Aceptable}} & \textbf{2} & El criterio es identificable pero se cumple de manera parcial o incompleta. La información existe pero carece de la precisión o el detalle suficiente para garantizar la reproducibilidad sin consultas adicionales al autor. Susceptible de fortalecimiento con correcciones específicas. \\ \hline
\textcolor{cdef}{\textbf{Deficiente}} & \textbf{1} & El cumplimiento del criterio es tangencial o meramente declarativo. La información es superficial, genérica o presenta omisiones que comprometen la capacidad de un tercero para replicar el procedimiento o verificar los resultados. \\ \hline
\rowcolor{lightgray}
\textcolor{cins}{\textbf{Ausente}} & \textbf{0} & El proyecto no proporciona la información requerida por el criterio evaluado. Existe omisión total del aspecto evaluado o una indefinición que imposibilita la reproducibilidad o la trazabilidad en ese componente. \\ \hline
\end{tabularx}
\end{table}

\clearpage
\subsection{Caracterización y Trazabilidad de Materiales e Insumos}

\textit{Se evalúa si el proyecto identifica todos los materiales, reactivos, aditivos e insumos con la precisión necesaria para que un investigador independiente pueda adquirir productos equivalentes y reproducir las condiciones experimentales. Se verifica la composición química, la pureza, la procedencia, la justificación de selección y la pertinencia del material en la aplicación propuesta. Esta dimensión es crítica en investigaciones que emplean productos comerciales, materiales no convencionales o insumos de formulación variable.} Véase la Tabla~\ref{tab:rubrica-dim1-materiales}.
{\scriptsize
\begin{longtable}{|C{0.55cm}|L{2cm}|L{8cm}|C{0.3cm}|C{0.3cm}|C{0.3cm}|C{0.3cm}|}
\caption{Rúbrica de evaluación de la Dimensión~I: caracterización y trazabilidad de materiales e insumos.}\label{tab:rubrica-dim1-materiales}\\
\hline
\rowcolor{headerblue}
\thd{N.\textsuperscript{o}} & \thd{Criterio} & \thd{Indicador verificable} & \thd{3} & \thd{2} & \thd{1} & \thd{0} \\
\hline\endfirsthead
\hline
\rowcolor{headerblue}
\thd{N.\textsuperscript{o}} & \thd{Criterio} & \thd{Indicador verificable} & \thd{3} & \thd{2} & \thd{1} & \thd{0} \\
\hline\endhead
1.1 & Composición química completa y verificable de todos los insumos empleados & El proyecto presenta la composición química de cada material, aditivo, reactivo o insumo que interviene como variable o como componente de la mezcla experimental: fórmula molecular o composición porcentual, especificación técnica del fabricante (\textit{datasheet} o ficha técnica), o resultados de análisis químico propio (fluorescencia de rayos X, espectroscopía, análisis elemental). No se admiten descripciones genéricas del tipo \enquote{se empleó cemento Portland} sin indicar tipo (IP, I, V, etc.), norma de fabricación, ni contenido de componentes activos. La composición es trazable a un documento verificable. &
& & & \\ \hline
\rowcolor{lightgray}
1.2 & Ausencia de productos comerciales de formulación propietaria no divulgada como variables independientes & Ningún producto cuya formulación sea secreto industrial o información propietaria del fabricante se emplea como variable independiente sin que el proyecto declare explícitamente esta limitación y proponga medidas de mitigación (análisis químico independiente, ensayos de caracterización funcional, o redefinición de la variable en términos de la propiedad medible del producto y no de su nombre comercial). Si se emplea un producto de formulación cerrada, el proyecto reconoce que la reproducibilidad queda condicionada a la disponibilidad de ese producto específico y evalúa las implicaciones para la generalización de los resultados. &
& & & \\ \hline
1.3 & Justificación de la selección de marcas, proveedores o fuentes de material frente a la diversidad de opciones disponibles & El proyecto justifica técnicamente por qué se seleccionó un proveedor, marca o fuente de material específica entre las alternativas disponibles en el mercado: composición química, disponibilidad regional, precedentes en la literatura, costo, representatividad del material respecto de la población de estudio, o condiciones de control de calidad del fabricante. La selección no es arbitraria ni se presenta sin argumentación. Cuando la selección responde a conveniencia logística, se declara explícitamente y se discuten las implicaciones para la reproducibilidad en otros contextos. &
& & & \\ \hline
\rowcolor{lightgray}
1.4 & Identificación del uso original previsto del material y empleo de compuestos de pureza conocida & El proyecto identifica el uso para el cual fue diseñado o comercializado cada material, distinguiéndolo del uso propuesto en la investigación (p.\,ej., un polímero de uso agroindustrial empleado como estabilizador de suelos; un subproducto industrial utilizado como adición puzolánica). Cuando se emplean reactivos de laboratorio, se especifica el grado de pureza (técnico, analítico, ACS). Cuando se emplean productos comerciales en lugar de compuestos puros, se justifica la decisión y se documenta la variabilidad composicional esperada. La pertinencia del material en la aplicación propuesta se argumenta técnicamente. &
& & & \\ \hline
\multicolumn{3}{|r|}{\cellcolor{white}\textbf{Subtotal Dimensión~I (máx.~12 puntos):}} & \multicolumn{4}{c|}{\cellcolor{white}\textbf{/12}} \\ \hline
\end{longtable}
}

\clearpage
\subsection{Especificación del Protocolo Experimental}

\textit{Se evalúa si el proyecto describe el procedimiento experimental con el nivel de detalle operativo necesario para que un investigador independiente pueda ejecutar los mismos ensayos sin recurrir a comunicaciones personales con el autor. Se verifica la secuencia de pasos, la especificación de equipos e instrumentos, las condiciones ambientales durante los ensayos y los protocolos de preparación de muestras y probetas. Esta dimensión se distingue de la mera cita de normas de ensayo: exige la descripción de las decisiones operativas que el tesista adoptó dentro del marco normativo.} Véase la Tabla~\ref{tab:rubrica-dim2-protocolo}.
{\scriptsize
\begin{longtable}{|C{0.55cm}|L{2cm}|L{8cm}|C{0.3cm}|C{0.3cm}|C{0.3cm}|C{0.3cm}|}
\caption{Rúbrica de evaluación de la Dimensión~II: especificación del protocolo experimental.}\label{tab:rubrica-dim2-protocolo}\\
\hline
\rowcolor{headerblue}
\thd{N.\textsuperscript{o}} & \thd{Criterio} & \thd{Indicador verificable} & \thd{3} & \thd{2} & \thd{1} & \thd{0} \\
\hline\endfirsthead
\hline
\rowcolor{headerblue}
\thd{N.\textsuperscript{o}} & \thd{Criterio} & \thd{Indicador verificable} & \thd{3} & \thd{2} & \thd{1} & \thd{0} \\
\hline\endhead
2.1 & Descripción secuencial y operativa del procedimiento experimental & El proyecto describe la secuencia completa de operaciones experimentales con suficiente detalle para su réplica: orden de mezclado, tiempos de curado, velocidades de carga, temperaturas de ensayo, intervalos de lectura, ciclos de exposición u otras variables de procedimiento pertinentes. La descripción no se limita a enunciar \enquote{se realizó el ensayo según la norma ASTM~D\ldots} sino que explicita las decisiones operativas adoptadas dentro del marco normativo: tamaño de molde, energía de compactación, humedad de moldeo, método de saturación, velocidad de deformación, entre otras. &
& & & \\ \hline
\rowcolor{lightgray}
2.2 & Especificación de equipos, instrumentos y estado de calibración & El proyecto identifica los equipos e instrumentos principales empleados en cada ensayo: marca, modelo, capacidad, resolución y rango de medición. Se indica el estado de calibración (certificado vigente, fecha de última calibración, laboratorio acreditado responsable) o, en su defecto, se justifica la ausencia de calibración formal y se describen las verificaciones internas realizadas. No se admiten referencias genéricas del tipo \enquote{se usó una prensa hidráulica} sin especificación de capacidad ni certificación. &
& & & \\ \hline
2.3 & Registro de condiciones ambientales durante los ensayos & El proyecto prevé el registro sistemático de las condiciones ambientales que pueden afectar los resultados: temperatura y humedad relativa del laboratorio, condiciones de curado (temperatura, humedad, tiempo), exposición solar o protección contra intemperie en ensayos de campo, altitud del sitio de ensayo cuando sea relevante para la presión atmosférica o la evaporación. Se especifican los rangos admisibles y el procedimiento ante desviaciones. &
& & & \\ \hline
\rowcolor{lightgray}
2.4 & Protocolo de preparación, moldeo y conservación de muestras y probetas & El proyecto detalla el procedimiento de obtención, preparación, moldeo y conservación de las muestras o probetas: método de muestreo (aleatorio, dirigido, compuesto), técnica de cuarteo o reducción, contenido de humedad de moldeo, método de compactación (estática, dinámica, vibratoria), dimensiones del espécimen, condiciones de curado o almacenamiento, y tiempo transcurrido entre preparación y ensayo. Se identifican las normas técnicas que regulan cada etapa y las desviaciones justificadas respecto de la norma cuando existan. &
& & & \\ \hline
\multicolumn{3}{|r|}{\cellcolor{white}\textbf{Subtotal Dimensión~II (máx.~12 puntos):}} & \multicolumn{4}{c|}{\cellcolor{white}\textbf{/12}} \\ \hline
\end{longtable}
}

\clearpage
\subsection{Registro Documental de Datos Primarios}

\textit{Se evalúa la estructura operativa del sistema de registro de datos que el proyecto propone para documentar la ejecución experimental. Se verifica la existencia de un sistema de identificación única de especímenes, la estructura de los formatos de registro, la separación entre datos en bruto y datos procesados, y la documentación fotográfica o audiovisual del procedimiento. Esta dimensión complementa ---sin duplicar--- el Plan de Gestión de Datos evaluado en las Observaciones Normativas, que verifica la declaración de intención; aquí se evalúa el diseño operativo del registro.} Véase la Tabla~\ref{tab:rubrica-dim3-registro}.
{\scriptsize
\begin{longtable}{|C{0.55cm}|L{2cm}|L{8cm}|C{0.3cm}|C{0.3cm}|C{0.3cm}|C{0.3cm}|}
\caption{Rúbrica de evaluación de la Dimensión~III: registro documental de datos primarios.}\label{tab:rubrica-dim3-registro}\\
\hline
\rowcolor{headerblue}
\thd{N.\textsuperscript{o}} & \thd{Criterio} & \thd{Indicador verificable} & \thd{3} & \thd{2} & \thd{1} & \thd{0} \\
\hline\endfirsthead
\hline
\rowcolor{headerblue}
\thd{N.\textsuperscript{o}} & \thd{Criterio} & \thd{Indicador verificable} & \thd{3} & \thd{2} & \thd{1} & \thd{0} \\
\hline\endhead
3.1 & Sistema de identificación única de muestras, probetas y especímenes & El proyecto define un sistema de codificación que permite identificar de manera inequívoca cada muestra, probeta o espécimen: código alfanumérico que incluya el tipo de material, la dosificación, el número de réplica, la fecha de preparación y la condición de ensayo. El sistema de codificación es consistente, no admite duplicados y permite la trazabilidad inversa desde el resultado hasta el espécimen físico. Se presenta un ejemplo del código y su estructura. &
& & & \\ \hline
\rowcolor{lightgray}
3.2 & Formato y estructura de los registros de ensayo (hojas de laboratorio) & El proyecto incluye o describe el formato de las hojas de registro de ensayo que se emplearán durante la fase experimental: campos obligatorios (fecha, operador, equipo, código de muestra, lecturas directas, observaciones), unidades de medida, resolución de las lecturas y espacio para firmas del ejecutor y del supervisor. Los formatos son consistentes con las normas de ensayo aplicables y permiten reconstruir el ensayo completo a partir de la hoja de registro, sin depender de la memoria del operador. &
& & & \\ \hline
3.3 & Separación explícita entre datos en bruto y datos procesados & El proyecto establece una distinción operativa entre los datos en bruto (lecturas directas del equipo, masas, longitudes, fuerzas, desplazamientos, tiempos) y los datos procesados (esfuerzos calculados, deformaciones unitarias, módulos, índices, porcentajes derivados). Se describe el procedimiento de transformación de datos (fórmulas, factores de corrección, software empleado) y se compromete a preservar ambos conjuntos de datos de manera independiente, de modo que un revisor pueda verificar los cálculos desde las lecturas originales. &
& & & \\ \hline
\rowcolor{lightgray}
3.4 & Documentación fotográfica y audiovisual del procedimiento experimental & El proyecto prevé un protocolo de documentación visual: fotografías de cada etapa crítica del procedimiento (preparación de muestras, montaje del ensayo, modo de falla, estado final del espécimen), con escala de referencia, código de muestra visible y fecha. Cuando el ensayo involucre procesos dinámicos o secuenciales (rotura progresiva, flujo, deformación), se considera el registro en video. La documentación visual no es decorativa sino funcional: sirve para verificar condiciones de ensayo, modos de falla y coherencia con los datos numéricos reportados. &
& & & \\ \hline
\multicolumn{3}{|r|}{\cellcolor{white}\textbf{Subtotal Dimensión~III (máx.~12 puntos):}} & \multicolumn{4}{c|}{\cellcolor{white}\textbf{/12}} \\ \hline
\end{longtable}
}

\clearpage
\subsection{Suficiencia Estadística y Control de Variables para la Replicación}

\textit{Se evalúa si el proyecto anticipa un plan muestral estadísticamente justificado, si las variables intervinientes se identifican y controlan, si existen criterios explícitos para el tratamiento de datos atípicos y si las dosificaciones y proporciones se expresan en unidades que permitan la reproducción exacta del experimento. Esta dimensión verifica que el diseño experimental tenga la robustez cuantitativa necesaria para que los resultados sean replicables y estadísticamente defendibles.} Véase la Tabla~\ref{tab:rubrica-dim4-estadistica}.
{\scriptsize
\begin{longtable}{|C{0.55cm}|L{2cm}|L{8cm}|C{0.3cm}|C{0.3cm}|C{0.3cm}|C{0.3cm}|}
\caption{Rúbrica de evaluación de la Dimensión~IV: suficiencia estadística y control de variables para la replicación.}\label{tab:rubrica-dim4-estadistica}\\
\hline
\rowcolor{headerblue}
\thd{N.\textsuperscript{o}} & \thd{Criterio} & \thd{Indicador verificable} & \thd{3} & \thd{2} & \thd{1} & \thd{0} \\
\hline\endfirsthead
\hline
\rowcolor{headerblue}
\thd{N.\textsuperscript{o}} & \thd{Criterio} & \thd{Indicador verificable} & \thd{3} & \thd{2} & \thd{1} & \thd{0} \\
\hline\endhead
4.1 & Número de réplicas y justificación estadística del tamaño muestral & El proyecto especifica el número de réplicas (probetas, muestras, mediciones) por cada combinación de variables y justifica su suficiencia con criterios estadísticos: norma de ensayo que fije un mínimo, análisis de potencia estadística, precedentes en la literatura para ensayos similares, o restricciones de recursos declaradas con evaluación de su impacto en la confiabilidad de los resultados. No se admite un número de réplicas arbitrario sin justificación (p.\,ej., \enquote{se ensayarán 3 probetas} sin explicar por qué tres y no cinco, ni qué nivel de confianza se alcanza con ese tamaño). &
& & & \\ \hline
\rowcolor{lightgray}
4.2 & Identificación, control y registro de variables intervinientes & El proyecto identifica las variables intervinientes (variables no manipuladas que pueden afectar la variable dependiente: temperatura de curado, humedad ambiental, velocidad de carga, granulometría del agregado, contenido de materia orgánica, edad del espécimen, operador) y describe los mecanismos de control: mantenimiento constante, aleatorización, bloqueo, registro sistemático para análisis posterior. Las variables intervinientes no se ignoran ni se asumen constantes sin verificación documentada. &
& & & \\ \hline
4.3 & Criterios explícitos de aceptación y rechazo de datos atípicos & El proyecto establece, antes de la ejecución experimental, los criterios objetivos para identificar y tratar datos atípicos (\textit{outliers}): test de Grubbs, criterio de Chauvenet, distancia de Cook, rango intercuartílico, o el criterio normativo aplicable (p.\,ej., ASTM~C670 para precisión de ensayos de materiales). Se define si los datos atípicos se excluyen, se reportan por separado o se reensayan, y se compromete a documentar cada exclusión con su justificación. No se admite la eliminación discrecional de datos sin criterio preestablecido. &
& & & \\ \hline
\rowcolor{lightgray}
4.4 & Dosificaciones y proporciones expresadas en unidades absolutas y reproducibles & Las dosificaciones, proporciones de mezcla y concentraciones se expresan en unidades absolutas reproducibles: porcentaje en masa respecto del peso seco del suelo, masa por unidad de volumen, relación agua/cemento en masa, molaridad, o unidades equivalentes que no dependan de condiciones locales variables. No se admiten dosificaciones expresadas exclusivamente en unidades comerciales ambiguas (\enquote{1~bolsa por metro cúbico}), en porcentajes sin base de referencia declarada, o en proporciones volumétricas sin corrección por densidad. Cuando se emplean dosificaciones normativas, se cita la norma y se explicita la conversión a unidades de laboratorio. &
& & & \\ \hline
\multicolumn{3}{|r|}{\cellcolor{white}\textbf{Subtotal Dimensión~IV (máx.~12 puntos):}} & \multicolumn{4}{c|}{\cellcolor{white}\textbf{/12}} \\ \hline
\end{longtable}
}

\clearpage
\subsection{Verificabilidad Externa y Potencial de Replicación Independiente}

\textit{Se evalúa si el conjunto de información proporcionado por el proyecto es suficiente para que un investigador externo, sin contacto con el autor, pueda replicar el estudio en su totalidad. Se verifica la completitud de las normas citadas, la ausencia de ambigüedades críticas en el protocolo, la disponibilidad declarada de datos y protocolos de soporte, y la coherencia global del diseño con los principios de ciencia abierta y verificación independiente.} Véase la Tabla~\ref{tab:rubrica-dim5-verificabilidad}.
{\scriptsize
\begin{longtable}{|C{0.55cm}|L{2cm}|L{8cm}|C{0.3cm}|C{0.3cm}|C{0.3cm}|C{0.3cm}|}
\caption{Rúbrica de evaluación de la Dimensión~V: verificabilidad externa y potencial de replicación independiente.}\label{tab:rubrica-dim5-verificabilidad}\\
\hline
\rowcolor{headerblue}
\thd{N.\textsuperscript{o}} & \thd{Criterio} & \thd{Indicador verificable} & \thd{3} & \thd{2} & \thd{1} & \thd{0} \\
\hline\endfirsthead
\hline
\rowcolor{headerblue}
\thd{N.\textsuperscript{o}} & \thd{Criterio} & \thd{Indicador verificable} & \thd{3} & \thd{2} & \thd{1} & \thd{0} \\
\hline\endhead
5.1 & Suficiencia de información para la replicación por un investigador independiente & Se aplica el \textit{test del investigador externo}: un profesional competente en el área, con acceso a un laboratorio equipado y a los materiales del mercado, podría reproducir el experimento completo ---desde la adquisición de materiales hasta la obtención de resultados--- utilizando exclusivamente la información contenida en el proyecto, sin necesidad de comunicación personal con el autor. No quedan pasos implícitos, decisiones no documentadas ni parámetros sin especificar que obliguen a la interpretación subjetiva del replicador. &
& & & \\ \hline
\rowcolor{lightgray}
5.2 & Referencia completa a normas de ensayo vigentes y especificación de la versión aplicable & Cada ensayo citado en el proyecto se referencia con su designación normativa completa (organismo emisor, número, año de la versión empleada): p.\,ej., ASTM~D1557-12(2021), NTP~339.141:2022, BS~EN~13286-47:2021. No se admiten referencias incompletas (\enquote{norma ASTM de Proctor modificado}), normas sin año de edición, ni la cita de normas derogadas sin justificación explícita del uso de una versión anterior. Cuando el proyecto se aparta del procedimiento normativo estándar, se identifican y justifican las modificaciones. &
& & & \\ \hline
5.3 & Disponibilidad declarada de datos primarios, protocolos y material de soporte & El proyecto declara explícitamente el compromiso de poner a disposición, como anexos o como material complementario en el repositorio institucional, los datos primarios en formato abierto, los protocolos detallados de ensayo, las hojas de cálculo con las transformaciones realizadas, los certificados de calibración y las fichas técnicas de los materiales. La disponibilidad no se limita a la tesis impresa sino que abarca los archivos digitales que permitan la verificación computacional de los resultados por un auditor o replicador. &
& & & \\ \hline
\rowcolor{lightgray}
5.4 & Ausencia de ambigüedades críticas que impidan la reproducción del experimento & El proyecto no contiene ambigüedades que, individualmente o en conjunto, impidan la reproducción del experimento: términos indefinidos (\enquote{se añadió una cantidad adecuada}), parámetros sin valor numérico (\enquote{se curó durante un tiempo suficiente}), materiales sin especificación (\enquote{se utilizó un aditivo comercial}), equipos sin identificación, o procedimientos descritos de forma incompleta. Se verifica que cada decisión experimental esté definida con la precisión necesaria para eliminar la discrecionalidad del replicador. &
& & & \\ \hline
\multicolumn{3}{|r|}{\cellcolor{white}\textbf{Subtotal Dimensión~V (máx.~12 puntos):}} & \multicolumn{4}{c|}{\cellcolor{white}\textbf{/12}} \\ \hline
\end{longtable}
}

\newpage
En esta sección se presenta el consolidado de resultados por dimensión y la escala de calificación de la reproducibilidad y trazabilidad, resumidos en las tablas~\ref{tab:consolidado-reproducibilidad} y~\ref{tab:escala-calificacion-reproducibilidad}.

\begin{table}[htbp]
\centering
\scriptsize
\caption{Cuadro consolidado de resultados por dimensión de la reproducibilidad y trazabilidad}\label{tab:consolidado-reproducibilidad}
\begin{tabularx}{\textwidth}{|C{0.8cm}|L{7.5cm}|C{1.8cm}|C{1.8cm}|X|}
\hline
\rowcolor{headerblue}
\thd{Dim.} & \thd{Denominación} & \thd{Pts. Máx.} & \thd{Pts. Obt.} & \thd{Obs.} \\ \hline
I & Caracterización y Trazabilidad de Materiales e Insumos & 12 & & \\ \hline
\rowcolor{lightgray}
II & Especificación del Protocolo Experimental & 12 & & \\ \hline
III & Registro Documental de Datos Primarios & 12 & & \\ \hline
\rowcolor{lightgray}
IV & Suficiencia Estadística y Control de Variables para la Replicación & 12 & & \\ \hline
V & Verificabilidad Externa y Potencial de Replicación Independiente & 12 & & \\ \hline
\rowcolor{white}
& \textbf{TOTAL GENERAL} & \textbf{60} & & \\ \hline
\end{tabularx}
\normalsize
\end{table}

\begin{table}[htbp]
\centering
\scriptsize
\caption{Escala de calificación de la reproducibilidad y trazabilidad del proyecto de tesis}\label{tab:escala-calificacion-reproducibilidad}
\begin{tabularx}{\textwidth}{|C{2.5cm}|C{4.2cm}|X|}
\hline
\rowcolor{headerblue}
\thd{Rango} & \thd{Calificación} & \thd{Significado} \\ \hline
\textbf{54--60}\newline(90--100\,\%) &
\textcolor{csuf}{\textbf{REPRODUCIBILIDAD\newline PLENA}} &
El proyecto proporciona información técnica completa, precisa y suficiente para que un investigador independiente replique el estudio en su totalidad. Los materiales están plenamente caracterizados, el protocolo experimental es detallado y operativo, el sistema de registro garantiza la trazabilidad de cada resultado hasta los datos primarios, el plan muestral está estadísticamente justificado y no existen ambigüedades que comprometan la verificación externa. \\ \hline
\rowcolor{lightgray}
\textbf{42--53}\newline(70--89\,\%) &
\textcolor{cace}{\textbf{REPRODUCIBILIDAD\newline SUFICIENTE}} &
El proyecto presenta los elementos esenciales para la reproducibilidad y la trazabilidad, con áreas específicas que requieren mayor precisión o detalle. Las deficiencias son corregibles sin reestructuración del diseño experimental y no comprometen la viabilidad de la verificación externa en sus aspectos fundamentales. \\ \hline
\textbf{30--41}\newline(50--69\,\%) &
\textcolor{cdef}{\textbf{REPRODUCIBILIDAD\newline PARCIAL}} &
La información proporcionada es insuficiente para garantizar la replicación completa del estudio. Existen omisiones significativas en la caracterización de materiales, el protocolo experimental, el sistema de registro o el control de variables que obligarían a un replicador a tomar decisiones discrecionales no documentadas. Requiere fortalecimiento sustancial. \\ \hline
\rowcolor{lightgray}
\textbf{0--29}\newline($<$50\,\%) &
\textcolor{cins}{\textbf{REPRODUCIBILIDAD\newline INSUFICIENTE}} &
El proyecto no demuestra las condiciones mínimas para la reproducibilidad ni la trazabilidad. El protocolo experimental es genérico o ausente, los materiales carecen de caracterización verificable, no existe sistema de registro de datos primarios y las dosificaciones o procedimientos presentan ambigüedades que imposibilitan la réplica. Se requiere una reformulación fundamental del componente metodológico-experimental. \\ \hline
\end{tabularx}
\normalsize
\end{table}

\clearpage
\section{Evaluación del Presupuesto y Cronograma}

{\footnotesize\setlength{\baselineskip}{0.85\baselineskip}
El presupuesto y el cronograma constituyen instrumentos de gestión que determinan la viabilidad operativa, económica y temporal de un proyecto de tesis de pregrado en Ingeniería Civil. Un presupuesto correctamente formulado no se limita a una lista de gastos; refleja la comprensión del tesista sobre el alcance real de su investigación, la naturaleza de los recursos requeridos y las restricciones del contexto institucional y del mercado. De forma análoga, un cronograma riguroso evidencia la capacidad de planificar la secuencia lógica de actividades, respetar dependencias técnicas, anticipar tiempos de espera y establecer hitos de control verificables. La planificación financiera y temporal deficiente produce, con frecuencia, retrasos que comprometen la calidad de los datos, matrices experimentales incompletas, omisión de ensayos o análisis sumamente importantes, o la imposibilidad de concluir el proyecto en los plazos académicos establecidos. La evaluación conjunta de ambos componentes permite verificar si el tesista ha considerado la totalidad de los gastos relevantes ---ensayos de laboratorio, ensayos de campo, materiales, software especializado, licencias, herramientas digitales, cuentas de inteligencia artificial, logística, transporte de muestras, servicios externos, consumibles y costos indirectos---, si las estimaciones son realistas y sustentables con cotizaciones o referencias de mercado, si el cronograma es coherente con la secuencia metodológica propuesta y si la planificación global demuestra un nivel de rigurosidad y sustento técnico compatible con la formación de un profesional en Ingeniería Civil. Esta rúbrica se aplica a proyectos de tesis en cualquier especialidad de la Ingeniería Civil: geotecnia, estructuras, pavimentos, hidráulica, recursos hídricos, saneamiento, gestión de la construcción, materiales de construcción, transporte, entre otras.Las dimensiones de evaluación del presupuesto y cronograma se resumen en la Tabla~\ref{tab:dim-presupuesto-cronograma}, mientras que la escala de valoración y sus descriptores operativos se presentan en la Tabla~\ref{tab:escala-presupuesto-cronograma}. \par
}

\begin{table}[htbp]
\centering
\scriptsize
\caption{Dimensiones de evaluación del presupuesto y cronograma}
\label{tab:dim-presupuesto-cronograma}
\begin{tabularx}{\textwidth}{|C{0.6cm}|L{2.5cm}|X|}
\hline
\rowcolor{headerblue}
\thd{Dim.} & \thd{Denominación} & \thd{Alcance} \\ \hline
I & Completitud y desagregación del presupuesto & Verifica que el presupuesto identifique y desagregue la totalidad de los recursos económicos necesarios para la ejecución del proyecto: costos directos (ensayos, materiales, servicios, software, licencias, herramientas digitales), costos indirectos, contingencias y cualquier otro gasto relevante, sin omisiones que comprometan la viabilidad. \\ \hline
\rowcolor{lightgray}
II & Realismo y sustento de las estimaciones de costos & Evalúa si los montos presupuestados son realistas, verificables y defendibles: si se sustentan en cotizaciones, tarifarios institucionales, catálogos de proveedores o referencias de mercado actualizadas, y si las cantidades y precios unitarios son coherentes con el alcance experimental declarado. \\ \hline
III & Estructura y coherencia del cronograma & Examina la calidad del cronograma como instrumento de planificación temporal: la identificación de actividades con duraciones estimadas, la secuencia lógica que respeta dependencias técnicas, la definición de hitos de control, la factibilidad de los plazos y la coherencia entre el cronograma y la metodología propuesta. \\ \hline
\rowcolor{lightgray}
IV & Integración presupuesto-cronograma y gestión de riesgos & Evalúa la articulación entre la planificación financiera y la planificación temporal: la distribución del gasto en el tiempo, la identificación de períodos de concentración de recursos, la consideración de reservas de contingencia y la anticipación de riesgos técnicos, logísticos o económicos que podrían afectar la ejecución del proyecto. \\ \hline
V & Rigor técnico y defendibilidad de la planificación & Verifica el nivel de rigurosidad metodológica de la planificación global: la coherencia interna entre presupuesto, cronograma, metodología y objetivos; la trazabilidad de cada partida hacia actividades específicas del proyecto; la transparencia de los supuestos; y la capacidad del tesista de defender la viabilidad técnica y económica de su propuesta ante un comité evaluador. \\ \hline
\end{tabularx}
\end{table}

\begin{table}[htbp]
\centering
\scriptsize
\caption{Escala de valoración para la evaluación del presupuesto y cronograma.}
\label{tab:escala-presupuesto-cronograma}
\begin{tabularx}{\textwidth}{|C{1.8cm}|C{.6cm}|X|}
\hline
\rowcolor{headerblue}
\thd{Nivel} & \thd{Pts.} & \thd{Descriptor operativo} \\ \hline
\textcolor{csuf}{\textbf{Excelente}} & \textbf{3} & El proyecto satisface plenamente el criterio evaluado. El presupuesto o el cronograma es explícito, desagregado, técnicamente sustentado con evidencia verificable y lógicamente articulado con la metodología y los objetivos. No se identifican omisiones, inconsistencias ni montos injustificados. \\ \hline
\rowcolor{lightgray}
\textcolor{cace}{\textbf{Aceptable}} & \textbf{2} & El criterio es identificable pero se cumple de manera parcial o incompleta. El presupuesto o el cronograma aborda el aspecto evaluado sin la desagregación, la precisión o el sustento suficiente. Susceptible de fortalecimiento con correcciones específicas. \\ \hline
\textcolor{cdef}{\textbf{Deficiente}} & \textbf{1} & El cumplimiento del criterio es tangencial o meramente declarativo. La planificación es superficial, carece de sustento verificable, presenta omisiones significativas o contiene inconsistencias que debilitan la viabilidad del proyecto. \\ \hline
\rowcolor{lightgray}
\textcolor{cins}{\textbf{Ausente}} & \textbf{0} & El proyecto no evidencia consideración alguna del criterio evaluado. Existe omisión total del aspecto evaluado o una desconexión que invalida la planificación en ese componente. \\ \hline
\end{tabularx}
\end{table}

\clearpage
\subsection{Completitud y Desagregación del Presupuesto}

\textit{Se evalúa si el presupuesto identifica y desagrega la totalidad de los recursos económicos requeridos para la ejecución del proyecto, incluyendo costos directos, costos indirectos, contingencias y gastos que a menudo se omiten en tesis de pregrado. La completitud se verifica contra el alcance metodológico declarado: toda actividad descrita en la metodología debe tener un correlato presupuestal, y todo gasto presupuestado debe corresponderse con una actividad del proyecto.} Véase la Tabla~\ref{tab:rubrica-dim1-completitud-presupuesto}.

{\scriptsize
\begin{longtable}{|C{0.55cm}|L{2.5cm}|L{8.8cm}|C{0.3cm}|C{0.3cm}|C{0.3cm}|C{0.3cm}|}
\caption{Rúbrica de evaluación de la Dimensión~I: completitud y desagregación del presupuesto.}\label{tab:rubrica-dim1-completitud-presupuesto}\\
\hline
\rowcolor{headerblue}
\thd{N.\textsuperscript{o}} & \thd{Criterio} & \thd{Indicador verificable} & \thd{3} & \thd{2} & \thd{1} & \thd{0} \\
\hline\endfirsthead
\hline
\rowcolor{headerblue}
\thd{N.\textsuperscript{o}} & \thd{Criterio} & \thd{Indicador verificable} & \thd{3} & \thd{2} & \thd{1} & \thd{0} \\
\hline\endhead
1.1 & Identificación exhaustiva de costos directos vinculados al programa experimental o de campo & El presupuesto desagrega los costos directos asociados a cada actividad técnica del proyecto: ensayos de laboratorio (con identificación de tipo, norma de referencia y número de especímenes o muestras), ensayos de campo (con identificación de tipo, ubicación y número de determinaciones), materiales consumibles (con especificación de tipo, cantidad y unidad), muestreo y transporte de muestras, servicios de análisis externos, alquiler o uso de equipos especializados, y cualquier otro gasto directamente atribuible a la obtención de datos. Cada partida de costo directo es trazable a una actividad específica de la metodología. &
& & & \\ \hline
\rowcolor{lightgray}
1.2 & Inclusión de costos de software, licencias, herramientas digitales y servicios tecnológicos & El presupuesto considera explícitamente los costos de las herramientas tecnológicas necesarias para el proyecto: software de análisis o modelación (elementos finitos, análisis estadístico, diseño asistido, modelación hidrológica, BIM, programación de cronogramas), licencias de uso o suscripciones (anuales o temporales), cuentas de servicios de inteligencia artificial utilizadas como herramientas de apoyo, servicios de almacenamiento en la nube, bases de datos especializadas y cualquier otro recurso digital requerido. Si el tesista declara que utilizará software libre o recursos institucionales de acceso gratuito, se indica explícitamente cuáles son y se justifica que su disponibilidad está asegurada. &
& & & \\ \hline
1.3 & Consideración de costos indirectos, gastos administrativos y logísticos & El presupuesto incluye, según corresponda al contexto institucional, los costos indirectos asociados a la ejecución del proyecto: gastos administrativos institucionales, uso de laboratorios (tarifas por hora, por ensayo o por periodo), consumo de servicios (electricidad, agua), impresión y encuadernación de documentos, viáticos y movilidad para trabajo de campo, gastos de comunicación, gestión de residuos de laboratorio, y cualquier otro gasto de soporte necesario para completar el proyecto. No se omiten categorías de gasto que la experiencia demuestra recurrentes en proyectos de tesis de Ingeniería Civil. &
& & & \\ \hline
\rowcolor{lightgray}
1.4 & Inclusión de reserva de contingencia con justificación del porcentaje adoptado & El presupuesto incorpora una reserva de contingencia explícita, expresada como porcentaje del costo directo total o como monto absoluto calculado, destinada a cubrir imprevistos técnicos, operacionales o económicos. El porcentaje adoptado se justifica cualitativamente (naturaleza experimental del proyecto, disponibilidad de equipos, variabilidad de precios) o cuantitativamente (identificación de riesgos con estimación de probabilidad e impacto). No se trata de un porcentaje arbitrario sin fundamento ni de una omisión total de la incertidumbre inherente a la investigación. &
& & & \\ \hline
\multicolumn{3}{|r|}{\cellcolor{white}\textbf{Subtotal Dimensión~I (máx.~12 puntos):}} & \multicolumn{4}{c|}{\cellcolor{white}\textbf{/12}} \\ \hline
\end{longtable}
}

\clearpage
\subsection{Realismo y Sustento de las Estimaciones de Costos}

\textit{Se evalúa si los montos consignados en el presupuesto son realistas, verificables y coherentes con las condiciones del mercado, las tarifas institucionales y el alcance experimental del proyecto. Esta dimensión verifica que las estimaciones no sean cifras genéricas, redondeadas sin justificación o copiadas de otros proyectos sin adaptación al contexto específico de la tesis.} Los criterios se aprecian en la Tabla~\ref{tab:rubrica-dim2-realismo-sustento}.

\scriptsize
\begin{longtable}{|C{0.55cm}|L{2.5cm}|L{9cm}|C{0.3cm}|C{0.3cm}|C{0.3cm}|C{0.3cm}|}
\caption{Rúbrica de evaluación de la Dimensión~II: realismo y sustento de las estimaciones de costos.}\label{tab:rubrica-dim2-realismo-sustento}\\
\hline
\rowcolor{headerblue}
\thd{N.\textsuperscript{o}} & \thd{Criterio} & \thd{Indicador verificable} & \thd{3} & \thd{2} & \thd{1} & \thd{0} \\
\hline\endfirsthead
\hline
\rowcolor{headerblue}
\thd{N.\textsuperscript{o}} & \thd{Criterio} & \thd{Indicador verificable} & \thd{3} & \thd{2} & \thd{1} & \thd{0} \\
\hline\endhead
2.1 & Sustento documental de precios unitarios mediante cotizaciones, tarifarios o referencias de mercado & Los precios unitarios consignados en el presupuesto se sustentan en al menos una fuente verificable: cotizaciones de proveedores o laboratorios (con fecha de vigencia), tarifarios institucionales oficiales, catálogos de fabricantes, precios de mercado documentados o referencias bibliográficas de costos de ensayos en la disciplina. Los montos no son cifras estimadas ``a ojo'' ni provienen de fuentes obsoletas o de contextos geográficos no comparables. Cuando se utilizan precios de referencia de la literatura, se indica la fuente, el año y, si corresponde, el factor de actualización aplicado. &
& & & \\ \hline
\rowcolor{lightgray}
2.2 & Coherencia entre cantidades presupuestadas y alcance experimental declarado en la metodología & Las cantidades de ensayos, muestras, especímenes, materiales y servicios consignados en el presupuesto son consistentes con el diseño experimental o el plan de trabajo descrito en la metodología. Si la metodología declara, por ejemplo, la elaboración de $n$ especímenes con $k$ réplicas para $m$ condiciones experimentales, el presupuesto refleja exactamente $n \times k \times m$ unidades (o un número debidamente justificado si difiere). No hay discrepancias entre lo que la metodología propone ejecutar y lo que el presupuesto costea. &
& & & \\ \hline
2.3 & Desagregación por precio unitario, cantidad y costo total para las partidas principales & Las partidas principales del presupuesto presentan una estructura de cálculo transparente que permite la verificación: precio unitario $\times$ cantidad $=$ costo total de la partida. La unidad de medida es apropiada y consistente (por espécimen, por ensayo, por muestra, por metro cúbico, por kilogramo, por hora, por licencia, según corresponda). Los subtotales por categoría y el total general son aritméticamente correctos y verificables. No se presentan montos globales (``lump sum'') sin desagregación para partidas que admiten cuantificación unitaria. &
& & & \\ \hline
\rowcolor{lightgray}
2.4 & Ausencia de omisiones críticas o subestimaciones evidentes que comprometan la viabilidad del proyecto & El presupuesto no contiene omisiones de categorías de gasto que resulten evidentes a partir de la lectura de la metodología, ni subestimaciones que, por su magnitud, pongan en riesgo la ejecución completa del programa de trabajo. No se omiten, por ejemplo, los costos de ensayos declarados como necesarios en la metodología, los materiales cuya adquisición es imprescindible para la ejecución experimental, los servicios externos de los que depende la obtención de resultados, ni los costos de actividades de campo cuya logística demanda recursos económicos. La evaluación de este criterio considera la proporcionalidad: el monto total del presupuesto es razonable para el alcance del proyecto propuesto. &
& & & \\ \hline
\multicolumn{3}{|r|}{\cellcolor{white}\textbf{Subtotal Dimensión~II (máx.~12 puntos):}} & \multicolumn{4}{c|}{\cellcolor{white}\textbf{/12}} \\ \hline
\end{longtable}
\normalsize

\clearpage
\subsection{Estructura y Coherencia del Cronograma}
\textit{Se evalúa la calidad del cronograma como instrumento de planificación temporal del proyecto. El cronograma debe demostrar que el tesista comprende la secuencia lógica de las actividades, respeta las dependencias técnicas entre ellas, estima duraciones factibles y establece hitos de control que permitan el seguimiento del avance. Un cronograma no es una simple lista de actividades con fechas arbitrarias, sino una herramienta de gestión que refleja la lógica operativa del proyecto.} Los criterios se aprecian en la Tabla~\ref{tab:rubrica-dim3-cronograma}.

\scriptsize
\begin{longtable}{|C{0.5cm}|L{2.5cm}|L{8.8cm}|C{0.3cm}|C{0.3cm}|C{0.3cm}|C{0.3cm}|}
\caption{Rúbrica de evaluación de la Dimensión~III: estructura y coherencia del cronograma.}\label{tab:rubrica-dim3-cronograma}\\
\hline
\rowcolor{headerblue}
\thd{N.\textsuperscript{o}} & \thd{Criterio} & \thd{Indicador verificable} & \thd{3} & \thd{2} & \thd{1} & \thd{0} \\
\hline\endfirsthead
\hline
\rowcolor{headerblue}
\thd{N.\textsuperscript{o}} & \thd{Criterio} & \thd{Indicador verificable} & \thd{3} & \thd{2} & \thd{1} & \thd{0} \\
\hline\endhead
3.1 & Identificación completa de actividades con duraciones estimadas y representación gráfica adecuada & El cronograma identifica todas las actividades principales del proyecto ---revisión de literatura, gestión administrativa, trabajo de campo, ejecución experimental, análisis de resultados, redacción de la tesis, entre otras--- con duraciones estimadas en unidades temporales apropiadas (días, semanas o meses). La representación es clara: diagrama de Gantt, tabla de actividades con fechas de inicio y fin, o formato equivalente que permita visualizar la distribución temporal del trabajo. Las actividades del cronograma corresponden a las fases descritas en la metodología; no hay actividades metodológicas sin representación temporal ni actividades en el cronograma sin correlato metodológico. &
& & & \\ \hline
\rowcolor{lightgray}
3.2 & Secuencia lógica de actividades que respeta dependencias técnicas y operativas & La secuencia de actividades refleja las dependencias técnicas reales del proyecto: las actividades que requieren resultados previos como insumo están programadas después de sus predecesoras (p.\,ej., los ensayos de resistencia no pueden iniciar antes de que concluya el tiempo de curado; el análisis estadístico no puede preceder a la obtención de datos; la calibración de un modelo numérico no puede realizarse sin datos experimentales previos). Las actividades independientes pueden programarse en paralelo. No se identifican secuencias imposibles, solapamientos técnicamente inviables ni dependencias omitidas que invalidarían la programación. &
& & & \\ \hline
3.3 & Factibilidad de las duraciones estimadas y del plazo total del proyecto & Las duraciones asignadas a las actividades son factibles considerando la naturaleza del trabajo, los recursos declarados y las condiciones institucionales: los tiempos de curado o secado respetan los requisitos normativos, los períodos de trabajo de campo son compatibles con las condiciones climáticas o de accesibilidad del sitio, los tiempos de laboratorio consideran la disponibilidad de equipos compartidos y los tiempos de espera por servicios externos son realistas. El plazo total del proyecto es compatible con el calendario académico de la institución y con los plazos reglamentarios para la sustentación de la tesis. No se asignan duraciones excesivamente optimistas que ignoren restricciones operativas previsibles. &
& & & \\ \hline
\rowcolor{lightgray}
3.4 & Definición de hitos de control y puntos de verificación del avance & El cronograma define hitos de control o puntos de verificación que permiten monitorear el avance del proyecto en momentos clave: finalización de la revisión bibliográfica, completitud del muestreo o trabajo de campo, obtención de resultados experimentales de la primera fase, finalización de todos los ensayos, entrega del borrador de tesis, entre otros. Los hitos están asociados a entregables o productos verificables (no a fechas arbitrarias) y su cumplimiento permite al tesista y al asesor identificar desviaciones respecto al plan original y tomar acciones correctivas oportunas. &
& & & \\ \hline
\multicolumn{3}{|r|}{\cellcolor{white}\textbf{Subtotal Dimensión~III (máx.~12 puntos):}} & \multicolumn{4}{c|}{\cellcolor{white}\textbf{/12}} \\ \hline
\end{longtable}
\normalsize

\clearpage
\subsection{Integración Presupuesto-Cronograma y Gestión de Riesgos}

\textit{Se evalúa la articulación entre la planificación financiera y la planificación temporal, así como la capacidad del tesista de anticipar riesgos que podrían afectar la ejecución del proyecto. Un presupuesto y un cronograma bien elaborados individualmente pero desarticulados entre sí no constituyen una planificación efectiva. Esta dimensión verifica que el tesista haya reflexionado sobre cuándo se producen los gastos, cómo se distribuyen en el tiempo y qué factores de incertidumbre podrían alterar el plan.} Los criterios se aprecian en la Tabla~\ref{tab:rubrica-dim4-integracion-riesgos}.

\scriptsize
\begin{longtable}{|C{0.55cm}|L{2.8cm}|L{8cm}|C{0.3cm}|C{0.3cm}|C{0.3cm}|C{0.3cm}|}
\caption{Rúbrica de evaluación de la Dimensión~IV: integración presupuesto-cronograma y gestión de riesgos.}\label{tab:rubrica-dim4-integracion-riesgos}\\
\hline
\rowcolor{headerblue}
\thd{N.\textsuperscript{o}} & \thd{Criterio} & \thd{Indicador verificable} & \thd{3} & \thd{2} & \thd{1} & \thd{0} \\
\hline\endfirsthead
\hline
\rowcolor{headerblue}
\thd{N.\textsuperscript{o}} & \thd{Criterio} & \thd{Indicador verificable} & \thd{3} & \thd{2} & \thd{1} & \thd{0} \\
\hline\endhead
4.1 & Correspondencia verificable entre las partidas presupuestales y las actividades del cronograma & Cada partida principal del presupuesto es trazable a una o más actividades específicas del cronograma, y cada actividad del cronograma que consume recursos tiene un correlato presupuestal identificable. No existen partidas presupuestales ``huérfanas'' (gastos sin actividad asociada en el cronograma) ni actividades ``no costeadas'' (tareas que demandan recursos pero carecen de asignación presupuestal). La trazabilidad puede verificarse mediante referencias cruzadas explícitas, codificación compartida o una tabla de correspondencia entre partidas y actividades. &
& & & \\ \hline
\rowcolor{lightgray}
4.2 & Distribución temporal del gasto coherente con la secuencia de actividades & El tesista identifica, al menos cualitativamente, cómo se distribuye el gasto a lo largo del cronograma: qué períodos concentran los mayores desembolsos (p.\,ej., adquisición de materiales, ejecución de ensayos de laboratorio, contratación de servicios externos de análisis instrumental), qué períodos tienen bajo consumo de recursos económicos (p.\,ej., revisión bibliográfica, redacción) y si la fuente de financiamiento permite atender los picos de gasto en los momentos requeridos. La distribución temporal del gasto no es uniforme por defecto, sino que refleja la lógica operativa del proyecto. &
& & & \\ \hline
4.3 & Identificación de riesgos técnicos, logísticos o económicos que podrían afectar la ejecución & El tesista identifica al menos los riesgos más probables que podrían afectar la ejecución del presupuesto o del cronograma: falla o indisponibilidad de equipos de laboratorio, retrasos en la entrega de materiales o reactivos, variación de precios, dificultades de acceso al sitio de estudio, tiempos de espera por servicios externos superiores a los estimados, resultados experimentales que requieran ensayos adicionales no previstos, o cualquier otro factor de incertidumbre pertinente al proyecto. Los riesgos identificados no son genéricos (\enquote{imprevistos}) sino específicos al contexto de la investigación propuesta. &
& & & \\ \hline
\rowcolor{lightgray}
4.4 & Estrategias de mitigación o respuesta ante desviaciones previsibles & El tesista anticipa, al menos para los riesgos más críticos, estrategias de respuesta o mitigación: alternativas de laboratorio en caso de indisponibilidad de equipos, proveedores sustitutos para materiales críticos, ajustes al diseño experimental en caso de restricciones presupuestales imprevistas, holguras temporales incorporadas en el cronograma para absorber retrasos moderados, o mecanismos de seguimiento que permitan detectar desviaciones tempranamente. Las estrategias son operativas y proporcionadas a la naturaleza del proyecto de pregrado; no se exigen planes de gestión de riesgos de complejidad empresarial, pero sí evidencia de reflexión anticipatoria sobre la incertidumbre. &
& & & \\ \hline
\multicolumn{3}{|r|}{\cellcolor{white}\textbf{Subtotal Dimensión~IV (máx.~12 puntos):}} & \multicolumn{4}{c|}{\cellcolor{white}\textbf{/12}} \\ \hline
\end{longtable}
\normalsize

\clearpage
\subsection{Rigor Técnico y Defendibilidad de la Planificación}

\textit{Se evalúa el nivel de rigurosidad metodológica y la calidad global de la planificación como instrumento defendible ante un comité evaluador. Esta dimensión integra aspectos de coherencia interna, transparencia de supuestos, alineación con los objetivos del proyecto y la capacidad del presupuesto y cronograma de demostrar que el tesista comprende integralmente el alcance operativo de su investigación.} Los criterios se aprecian en la Tabla~\ref{tab:rubrica-dim5-rigor-defendibilidad}.

\scriptsize
\begin{longtable}{|C{0.55cm}|L{2.8cm}|L{8cm}|C{0.3cm}|C{0.3cm}|C{0.3cm}|C{0.3cm}|}
\caption{Rúbrica de evaluación de la Dimensión~V: rigor técnico y defendibilidad de la planificación.}\label{tab:rubrica-dim5-rigor-defendibilidad}\\
\hline
\rowcolor{headerblue}
\thd{N.\textsuperscript{o}} & \thd{Criterio} & \thd{Indicador verificable} & \thd{3} & \thd{2} & \thd{1} & \thd{0} \\
\hline\endfirsthead
\hline
\rowcolor{headerblue}
\thd{N.\textsuperscript{o}} & \thd{Criterio} & \thd{Indicador verificable} & \thd{3} & \thd{2} & \thd{1} & \thd{0} \\
\hline\endhead
5.1 & Coherencia global entre presupuesto, cronograma, metodología y objetivos del proyecto & El presupuesto y el cronograma son coherentes entre sí y con la metodología y los objetivos del proyecto. El alcance costeado en el presupuesto corresponde al alcance descrito en la metodología; las actividades programadas en el cronograma reflejan las fases metodológicas; los objetivos específicos se traducen en actividades costeadas y programadas; y no hay componentes de la investigación que estén declarados en los objetivos pero ausentes en el presupuesto o en el cronograma. La planificación es un reflejo fiel y operativo del diseño de la investigación. &
& & & \\ \hline
\rowcolor{lightgray}
5.2 & Transparencia de los supuestos y condiciones asumidas en la planificación & Los supuestos que sustentan las estimaciones de costos y duraciones están explícitamente declarados: disponibilidad de equipos de laboratorio en determinados períodos, vigencia de las cotizaciones presentadas, tipo de cambio utilizado (si corresponde), rendimientos de trabajo estimados, condiciones de acceso al sitio de campo, disponibilidad de personal técnico, y cualquier otra condición cuya variación podría alterar los montos o los plazos. La explicitación de supuestos permite evaluar la razonabilidad de la planificación y anticipar los puntos de vulnerabilidad del proyecto. &
& & & \\ \hline
5.3 & Presentación formal ordenada, con totales verificables y estructura que facilite la evaluación & El presupuesto se presenta en formato tabular organizado por categorías de gasto, con columnas para descripción, unidad, cantidad, precio unitario y costo total, subtotales por categoría y total general. El cronograma se presenta en formato gráfico o tabular con actividades, duraciones, fechas y dependencias claramente identificadas. Ambos documentos son legibles, ordenados y libres de errores aritméticos. La presentación permite al evaluador verificar montos, identificar categorías y evaluar la completitud sin necesidad de reconstruir información omitida. &
& & & \\ \hline
\rowcolor{lightgray}
5.4 & Proporcionalidad del presupuesto respecto al alcance del proyecto y al nivel de pregrado & El monto total del presupuesto es proporcionado al alcance real de un proyecto de tesis de pregrado en Ingeniería Civil: no es excesivamente ambicioso para los recursos típicamente disponibles en este nivel académico, ni es tan reducido que evidencie una subestimación generalizada del costo real de la investigación. Las proporciones internas entre categorías de gasto (ensayos, materiales, servicios, software, logística, contingencia) son razonables para la naturaleza del proyecto. El tesista demuestra comprensión del orden de magnitud de los costos involucrados en la ejecución de investigaciones en Ingeniería Civil, acorde a la realidad institucional y al mercado local. &
& & & \\ \hline
\multicolumn{3}{|r|}{\cellcolor{white}\textbf{Subtotal Dimensión~V (máx.~12 puntos):}} & \multicolumn{4}{c|}{\cellcolor{white}\textbf{/12}} \\ \hline
\end{longtable}
\normalsize

\newpage
En esta sección se presenta el consolidado de resultados por dimensión y la escala de calificación del presupuesto y cronograma, resumidos en las tablas~\ref{tab:consolidado-presupuesto-cronograma} y~\ref{tab:escala-calificacion-presupuesto-cronograma}.

\begin{table}[htbp]
\centering
\scriptsize
\caption{Cuadro consolidado de resultados por dimensión del presupuesto y cronograma}\label{tab:consolidado-presupuesto-cronograma}
\begin{tabularx}{\textwidth}{|C{0.8cm}|L{7.5cm}|C{1.8cm}|C{1.8cm}|X|}
\hline
\rowcolor{headerblue}
\thd{Dim.} & \thd{Denominación} & \thd{Pts. Máx.} & \thd{Pts. Obt.} & \thd{Obs.} \\ \hline
I & Completitud y Desagregación del Presupuesto & 12 & & \\ \hline
\rowcolor{lightgray}
II & Realismo y Sustento de las Estimaciones de Costos & 12 & & \\ \hline
III & Estructura y Coherencia del Cronograma & 12 & & \\ \hline
\rowcolor{lightgray}
IV & Integración Presupuesto-Cronograma y Gestión de Riesgos & 12 & & \\ \hline
V & Rigor Técnico y Defendibilidad de la Planificación & 12 & & \\ \hline
\rowcolor{white}
& \textbf{TOTAL GENERAL} & \textbf{60} & & \\ \hline
\end{tabularx}
\normalsize
\end{table}

\begin{table}[htbp]
\centering
\scriptsize
\caption{Escala de calificación del presupuesto y cronograma}\label{tab:escala-calificacion-presupuesto-cronograma}
\begin{tabularx}{\textwidth}{|C{2.5cm}|C{4.2cm}|X|}
\hline
\rowcolor{headerblue}
\thd{Rango} & \thd{Calificación} & \thd{Significado} \\ \hline
\textbf{54--60}\newline(90--100\,\%) &
\textcolor{csuf}{\textbf{PLANIFICACIÓN\newline PLENA}} &
El presupuesto y el cronograma son integrales, técnicamente sustentados y lógicamente articulados con la metodología y los objetivos del proyecto. Los costos están desagregados, verificados contra fuentes de mercado y correctamente cuantificados. El cronograma refleja la secuencia operativa real del proyecto con dependencias, duraciones factibles y hitos de control. La integración entre ambos instrumentos es explícita y la gestión de riesgos es proporcionada y pertinente. \\ \hline
\rowcolor{lightgray}
\textbf{42--53}\newline(70--89\,\%) &
\textcolor{cace}{\textbf{PLANIFICACIÓN\newline SUFICIENTE}} &
El presupuesto y el cronograma presentan los componentes esenciales de una planificación viable, con áreas específicas susceptibles de fortalecimiento en cuanto a desagregación, sustento documental, coherencia temporal o anticipación de riesgos. Las debilidades son corregibles sin reestructuración fundamental del presupuesto o del cronograma. \\ \hline
\textbf{30--41}\newline(50--69\,\%) &
\textcolor{cdef}{\textbf{PLANIFICACIÓN\newline PARCIAL}} &
El presupuesto o el cronograma son incompletos, predominantemente declarativos o presentan inconsistencias significativas. Se evidencian omisiones de categorías de gasto relevantes, montos sin sustento, cronogramas con secuencias ilógicas, duraciones inviables o ausencia de integración entre la planificación financiera y la temporal. Requiere fortalecimiento sustancial. \\ \hline
\rowcolor{lightgray}
\textbf{0--29}\newline($<$50\,\%) &
\textcolor{cins}{\textbf{PLANIFICACIÓN\newline INSUFICIENTE}} &
El presupuesto y el cronograma no demuestran una planificación consistente ni verificable. Los costos carecen de desagregación o sustento, el cronograma es una lista de actividades sin secuencia lógica ni duraciones fundamentadas, o ambos instrumentos están desarticulados de la metodología y los objetivos del proyecto. Se requiere una reformulación fundamental de la planificación. \\ \hline
\end{tabularx}
\normalsize
\end{table}


\newpage
\section{Observaciones Técnicas de Fondo}

La revisión técnica de fondo de una propuesta de tesis en Ingeniería Civil exige que el evaluador examine, en primer lugar, la consistencia entre las normas técnicas citadas y el uso efectivo que el proyecto hace de ellas. Es frecuente que los tesistas invoquen normativa internacional o nacional ---ASTM, AASHTO, NTP, ACI, Eurocódigos, entre otras--- de manera nominal, sin verificar que las condiciones de aplicabilidad establecidas por la propia norma se cumplan en el contexto específico de la investigación. El revisor debe contrastar si los rangos de validez del ensayo o del modelo normativo son coherentes con el tipo de material, la escala de estudio y las condiciones de contorno declaradas en el proyecto. Asimismo, conviene verificar que las versiones citadas de las normas correspondan a ediciones vigentes y que las eventuales modificaciones introducidas respecto del procedimiento estándar estén explícitamente justificadas. Una incongruencia típica consiste en aplicar métodos de ensayo diseñados para un tipo de material a otro cuyas propiedades físicas o químicas difieren sustancialmente, sin discutir las implicaciones de dicha extrapolación sobre la validez de los resultados.

En segundo lugar, el evaluador debe prestar atención a la coherencia entre los principios teóricos invocados en el marco conceptual y las decisiones metodológicas adoptadas en el diseño de la investigación. Un proyecto técnicamente sólido no se limita a enunciar teorías o modelos reconocidos de la disciplina, sino que demuestra que las hipótesis de partida de dichos modelos son compatibles con las condiciones del problema investigado. El revisor debe examinar si los supuestos implícitos de los modelos constitutivos, las ecuaciones de diseño o los criterios de falla seleccionados son pertinentes al fenómeno estudiado o si, por el contrario, se emplean fuera de su dominio de validez. Esta verificación resulta especialmente relevante cuando el proyecto combina enfoques provenientes de distintas subespecialidades ---por ejemplo, mecánica de suelos con química de materiales, o hidráulica con análisis estructural---, ya que la integración de marcos teóricos heterogéneos demanda una justificación explícita de su compatibilidad conceptual y operacional.

En tercer lugar, resulta útil evaluar la trazabilidad lógica entre el planteamiento del problema, los objetivos, la hipótesis y el diseño experimental, con especial énfasis en la detección de inconsistencias entre lo que el proyecto declara investigar y lo que su metodología efectivamente permite concluir. El revisor debe identificar situaciones en las que los objetivos anticipen conclusiones causales, mientras que el diseño solo admite inferencias correlacionales, o en las que la hipótesis postule relaciones entre variables que la matriz experimental propuesta no es capaz de contrastar con la potencia estadística suficiente. De igual manera, conviene verificar que las variables independientes estén definidas con la precisión necesaria para garantizar la reproducibilidad del experimento y que los niveles o tratamientos seleccionados no sean arbitrarios, sino técnicamente fundamentados. La ausencia de grupos de control, la omisión de variables intervinientes relevantes o la confusión entre la significancia estadística y la relevancia práctica para la ingeniería constituyen debilidades frecuentes que el evaluador debe señalar con claridad.

En cuarto lugar, el revisor debe examinar si los parámetros, las magnitudes y las unidades empleadas a lo largo del proyecto son internamente consistentes y si las afirmaciones técnicas están respaldadas por evidencia verificable. Es habitual encontrar proyectos que presentan valores de propiedades mecánicas, hidráulicas o térmicas sin indicar las condiciones de ensayo bajo las cuales fueron obtenidos, o que citan resultados de la literatura sin evaluar si las condiciones experimentales de los estudios de referencia son comparables con las del contexto local. El evaluador debe verificar que los órdenes de magnitud reportados sean plausibles para el tipo de material y las condiciones declaradas, que las conversiones entre sistemas de unidades sean correctas y que los criterios de aceptación o umbrales normativos contra los cuales se compararán los resultados estén explícitamente identificados. La detección de cifras incongruentes, unidades omitidas o valores fuera de los rangos típicos reportados en la literatura especializada constituye un indicador temprano de debilidades conceptuales o de falta de dominio técnico por parte del tesista.

En quinto lugar, el evaluador debe valorar si el proyecto reconoce las limitaciones inherentes a su enfoque técnico y si anticipa las fuentes de incertidumbre que podrían afectar la interpretación de los resultados. Un proyecto que no identifica las restricciones de generalización derivadas de la escala de ensayo, de la variabilidad natural de los geomateriales, de las simplificaciones introducidas en los modelos numéricos o de las condiciones ambientales no controladas revela una comprensión superficial del problema investigado. El revisor debe distinguir entre las limitaciones declarativas —aquellas que el tesista enuncia de manera genérica sin consecuencias operativas— y las limitaciones integradas, que se reflejan en decisiones concretas del diseño metodológico, como la inclusión de réplicas suficientes, la adopción de intervalos de confianza, la propagación de incertidumbres o la validación cruzada con métodos independientes. La ausencia de esta reflexión crítica compromete la credibilidad científica de la propuesta y debe ser observada como una debilidad técnica de fondo.

Finalmente, el evaluador debe considerar que las observaciones técnicas de fondo no se agotan en la verificación de errores puntuales, sino que abarcan la evaluación integral de la solidez conceptual del proyecto como contribución al conocimiento de la Ingeniería Civil. Esto implica examinar si el proyecto distingue con claridad entre un problema de conocimiento genuino y un problema de aplicación profesional rutinaria, si la estrategia investigativa propuesta tiene el potencial de generar resultados transferibles más allá del caso particular estudiado y si el diseño metodológico refleja una comprensión mecanicista del fenómeno que trasciende la mera adecuación empírica. Las incongruencias técnicas más profundas no residen en errores de cálculo o en citas normativas incorrectas, sino en la desarticulación entre los fundamentos científicos de la disciplina y la implementación operativa del proyecto. El revisor, al formular sus observaciones, debe orientar al tesista hacia la corrección de estas debilidades estructurales con indicaciones específicas y constructivas que fortalezcan la calidad científica del trabajo propuesto.

\newpage
\section{Estructura Documental del Proyecto de Tesis}

\vspace{-6pt}
{\footnotesize\linespread{1}\selectfont
La estructura documental de un proyecto de tesis de pregrado en Ingeniería Civil no constituye un requisito meramente formal ni una exigencia burocrática, sino un indicador verificable de la madurez metodológica del tesista y del grado de institucionalización de su propuesta dentro del marco normativo que la universidad ha establecido para la producción de conocimiento. Un documento estructuralmente completo, ordenado y conforme a los lineamientos institucionales facilita la evaluación académica, garantiza la trazabilidad de los componentes del proyecto, permite la verificación de la coherencia interna entre secciones y asegura que el trabajo podrá ser registrado, archivado y difundido conforme a las exigencias del sistema nacional de repositorios y grados académicos. La presente rúbrica evalúa la conformidad del proyecto de tesis con el esquema aprobado por la autoridad competente, la completitud y corrección de los elementos preliminares (carátula, metadatos, índices), la presencia, secuencia y articulación de los elementos del cuerpo (planteamiento del problema, marco teórico, hipótesis, variables, metodología, aspectos administrativos), la inclusión y calidad de los elementos finales (referencias, anexos, cronograma, presupuesto, matriz de consistencia) y la coherencia global de la presentación documental como unidad integrada. La evaluación se sustenta en la normativa institucional vigente: el esquema de proyecto de tesis aprobado por Consejo de Facultad conforme al Art.~67 de la R.~N.\textsuperscript{o}~529-CU-2025-UAC, el Reglamento Marco de Grados y Títulos (R\_CU-256-2024), y las disposiciones específicas de la Escuela Profesional de Ingeniería Civil de la Universidad Andina del Cusco. La rúbrica es aplicable a proyectos de tesis de cualquier especialidad dentro de la Ingeniería Civil y se centra en la verificación objetiva de la presencia, completitud, conformidad normativa y articulación de los componentes documentales exigidos. Las dimensiones de evaluación de la estructura documental se resumen en la Tabla~\ref{tab:dim-estructura-documental}, mientras que la escala de valoración y sus descriptores operativos se presentan en la Tabla~\ref{tab:escala-estructura-documental}.\par
}

\vspace{-10pt}
\begin{table}[htbp]
\centering
\scriptsize
\caption{Dimensiones de evaluación de la estructura documental y alcance de cada dimensión.}
\label{tab:dim-estructura-documental}
\begin{tabularx}{\textwidth}{|C{0.8cm}|L{2.5cm}|X|}
\hline
\rowcolor{headerblue}
\thd{Dim.} & \thd{Denominación} & \thd{Alcance} \\ \hline
I & Adecuación al esquema aprobado y marco normativo & Verifica la conformidad del proyecto con el esquema de tesis aprobado por Consejo de Facultad (Art.~67, R.~N.\textsuperscript{o}~529-CU-2025-UAC), la inclusión de todas las secciones obligatorias según la Escuela Profesional, y el cumplimiento de los formatos de metadatos e informe de similitud del Reglamento Marco R\_CU-256-2024. \\ \hline
\rowcolor{lightgray}
II & Elementos preliminares del documento & Evalúa la completitud y conformidad de la carátula institucional, la ficha de metadatos, las declaraciones formales exigidas y los índices de contenido, tablas y figuras, verificando que cada elemento cumpla con el formato y la información requeridos por la normativa vigente. \\ \hline
III & Elementos del cuerpo del proyecto & Examina la presencia, secuencia correcta y numeración jerárquica de los capítulos y secciones sustantivas del proyecto: planteamiento del problema, marco teórico, hipótesis y variables, diseño metodológico, aspectos administrativos, y la existencia de la matriz de consistencia como instrumento articulador. \\ \hline
\rowcolor{lightgray}
IV & Elementos finales y complementarios & Verifica la inclusión, completitud y corrección formal de las referencias bibliográficas, los anexos, el cronograma de ejecución, el presupuesto y financiamiento, y la matriz de consistencia como componente integrador del proyecto. \\ \hline
V & Coherencia global y calidad de la presentación documental & Evalúa la unidad e integración del documento como totalidad: consistencia terminológica entre secciones, ausencia de contradicciones internas, calidad tipográfica y editorial, y conformidad general con los estándares de presentación de documentos académicos en Ingeniería Civil. \\ \hline
\end{tabularx}
\end{table}

\vspace{-10pt}
\begin{table}[htbp]
\centering
\scriptsize
\caption{Escala de valoración para la evaluación de la estructura documental.}
\label{tab:escala-estructura-documental}
\begin{tabularx}{\textwidth}{|C{1.8cm}|C{.6cm}|X|}
\hline
\rowcolor{headerblue}
\thd{Nivel} & \thd{Pts.} & \thd{Descriptor operativo} \\ \hline
\textcolor{csuf}{\textbf{Excelente}} & \textbf{3} & El proyecto satisface plenamente el criterio evaluado. El componente documental está presente, completo, conforme a la normativa institucional vigente y articulado coherentemente con el resto del documento. No se identifican omisiones, errores formales ni inconsistencias. \\ \hline
\rowcolor{lightgray}
\textcolor{cace}{\textbf{Aceptable}} & \textbf{2} & El criterio es identificable pero se cumple de manera parcial o incompleta. El componente documental existe pero presenta deficiencias menores en formato, contenido o conformidad normativa que no comprometen la viabilidad del proyecto. Susceptible de corrección con ajustes específicos. \\ \hline
\textcolor{cdef}{\textbf{Deficiente}} & \textbf{1} & El cumplimiento del criterio es tangencial o presenta errores significativos. El componente documental existe pero con omisiones graves, errores de formato que dificultan la evaluación, o incumplimientos normativos que requieren subsanación obligatoria antes de la aprobación del proyecto. \\ \hline
\rowcolor{lightgray}
\textcolor{cins}{\textbf{Ausente}} & \textbf{0} & El proyecto no incluye el componente documental evaluado, o su inclusión es tan deficiente que equivale a una omisión. Se requiere la incorporación del elemento faltante antes de que el proyecto pueda ser evaluado en sus aspectos sustantivos. \\ \hline
\end{tabularx}
\end{table}

\clearpage
\subsection{Adecuación al Esquema Aprobado y Marco Normativo}
\textit{Se evalúa si el proyecto de tesis se ajusta al esquema oficial aprobado por Consejo de Facultad, si incluye la totalidad de las secciones obligatorias exigidas por la Escuela Profesional de Ingeniería Civil y si cumple con los requisitos documentales establecidos en el Reglamento Marco de Grados y Títulos. Esta dimensión verifica la conformidad normativa del documento antes de evaluar el contenido de cada sección.} Véase la Tabla~\ref{tab:rubrica-dim1-adecuacion-esquema}.

{\scriptsize
\begin{longtable}{|C{0.4cm}|L{2.8cm}|L{8.2cm}|C{0.4cm}|C{0.4cm}|C{0.4cm}|C{0.4cm}|}
\caption{Rúbrica de evaluación de la Dimensión~I: adecuación al esquema aprobado y marco normativo.}\label{tab:rubrica-dim1-adecuacion-esquema}\\
\hline
\rowcolor{headerblue}
\thd{N.\textsuperscript{o}} & \thd{Criterio} & \thd{Indicador verificable} & \thd{3} & \thd{2} & \thd{1} & \thd{0} \\
\hline\endfirsthead
\hline
\rowcolor{headerblue}
\thd{N.\textsuperscript{o}} & \thd{Criterio} & \thd{Indicador verificable} & \thd{3} & \thd{2} & \thd{1} & \thd{0} \\
\hline\endhead
1.1 & Conformidad con el esquema de proyecto de tesis aprobado por Consejo de Facultad & El proyecto se estructura conforme al esquema oficial aprobado por Consejo de Facultad según el Art.~67 de la R.~N.\textsuperscript{o}~529-CU-2025-UAC. La organización de capítulos, secciones y subsecciones sigue el orden y la denominación establecidos en dicho esquema. No se omiten secciones obligatorias, no se alteran las denominaciones oficiales sin justificación reglamentaria y no se introducen secciones no previstas en posiciones que alteren la estructura aprobada. La conformidad es verificable mediante cotejo directo entre el índice del proyecto y el esquema de referencia institucional. &
& & & \\ \hline
\rowcolor{lightgray}
1.2 & Inclusión de todas las secciones obligatorias según la Escuela Profesional de Ingeniería Civil & El proyecto incluye la totalidad de las secciones exigidas por la Escuela Profesional: planteamiento del problema (descripción, formulación, justificación, objetivos, delimitación), marco teórico (antecedentes, bases teóricas, hipótesis, variables, definición de términos), diseño metodológico (tipo, nivel, diseño, población, muestra, técnicas, instrumentos, plan de análisis) y aspectos administrativos (cronograma, presupuesto, matriz de consistencia). No se identifican secciones obligatorias ausentes ni secciones fusionadas que impidan la evaluación independiente de cada componente. &
& & & \\ \hline
1.3 & Conformidad con los formatos de metadatos e informe de similitud del Reglamento Marco R\_CU-256-2024 & El proyecto incluye la ficha de metadatos conforme al formato establecido en el Reglamento Marco de Grados y Títulos (R\_CU-256-2024): campos de identificación del trabajo (título, autor, asesor, programa, tipo de investigación, línea de investigación), resumen y palabras clave en español e inglés, y declaración de acceso. Asimismo, se adjunta o se prevé el informe de similitud conforme a las disposiciones vigentes del reglamento, con indicación del software de detección utilizado y el porcentaje de coincidencia obtenido o previsto. Los formatos son los oficiales de la universidad y no versiones libres o adaptaciones no autorizadas. &
& & & \\ \hline
\rowcolor{lightgray}
1.4 & Cumplimiento de las disposiciones reglamentarias sobre formato, extensión y presentación & El proyecto cumple con las disposiciones reglamentarias sobre formato de presentación: tipo y tamaño de fuente, márgenes, interlineado, formato de tablas y figuras, sistema de numeración de páginas, y cualquier otra especificación establecida por la Escuela Profesional o por el Reglamento Marco. La extensión del documento es proporcionada al nivel de pregrado y al alcance del proyecto. El documento no presenta deficiencias de formato que dificulten su lectura, evaluación o archivo institucional. &
& & & \\ \hline
\multicolumn{3}{|r|}{\cellcolor{white}\textbf{Subtotal Dimensión~I (máx.~12 puntos):}} & \multicolumn{4}{c|}{\cellcolor{white}\textbf{/12}} \\ \hline
\end{longtable}
}

\clearpage
\subsection{Elementos Preliminares del Documento}

\textit{Se evalúa la completitud, conformidad y corrección de los elementos que preceden al cuerpo del proyecto de tesis: carátula institucional, ficha de metadatos, declaraciones formales, índice general, índice de tablas e índice de figuras. Los elementos preliminares constituyen la primera impresión documental del proyecto y condicionan su registro, archivo y difusión institucional; su evaluación verifica que la información consignada sea completa, exacta y conforme al formato oficial.} Los criterios se aprecian en la Tabla~\ref{tab:rubrica-dim2-elementos-preliminares}.

\scriptsize
\begin{longtable}{|C{0.5cm}|L{1.8cm}|L{9cm}|C{0.3cm}|C{0.3cm}|C{0.3cm}|C{0.3cm}|}
\caption{Rúbrica de evaluación de la Dimensión~II: elementos preliminares del documento.}\label{tab:rubrica-dim2-elementos-preliminares}\\
\hline
\rowcolor{headerblue}
\thd{N.\textsuperscript{o}} & \thd{Criterio} & \thd{Indicador verificable} & \thd{3} & \thd{2} & \thd{1} & \thd{0} \\
\hline\endfirsthead
\hline
\rowcolor{headerblue}
\thd{N.\textsuperscript{o}} & \thd{Criterio} & \thd{Indicador verificable} & \thd{3} & \thd{2} & \thd{1} & \thd{0} \\
\hline\endhead
2.1 & Carátula conforme al formato institucional con información completa y verificable & La carátula del proyecto incluye todos los elementos exigidos por el formato institucional de la Universidad Andina del Cusco: logotipo oficial de la universidad, denominación de la Facultad de Ingeniería y Arquitectura, denominación de la Escuela Profesional de Ingeniería Civil, título completo del proyecto de tesis, nombre(s) del autor(es) con su(s) identificador(es) ORCID, nombre del asesor con su identificador ORCID, lugar (Cusco -- Perú) y año de presentación. Cada dato es exacto y verificable: los identificadores ORCID corresponden efectivamente a las personas declaradas, el título coincide con el registrado en la resolución de aprobación del proyecto, y la denominación institucional está actualizada. No se emplean logotipos desactualizados ni denominaciones incorrectas de la Facultad o Escuela. &
& & & \\ \hline
\rowcolor{lightgray}
2.2 & Ficha de metadatos completa conforme al formato del Reglamento Marco & El proyecto incluye la ficha de metadatos con todos los campos exigidos por el Reglamento Marco (R\_CU-256-2024): título en español y en inglés, autor(es), asesor, escuela profesional, línea de investigación, sublínea (si aplica), tipo de investigación, enfoque, nivel, diseño, resumen en español (con extensión dentro del rango permitido), abstract en inglés, palabras clave en español e inglés (número y pertinencia conformes al formato), fecha prevista de sustentación, y declaración de tipo de acceso (abierto, restringido, embargado) con justificación cuando corresponda. No se identifican campos vacíos, incompletos o con información genérica que impida la catalogación del trabajo en el repositorio institucional y en el RENATI. &
& & & \\ \hline
2.3 & Índice general con numeración de páginas correcta y completa & El proyecto incluye un índice general (tabla de contenidos) que refleja fielmente la estructura del documento: todos los capítulos, secciones y subsecciones aparecen listados con su numeración jerárquica y su número de página correspondiente. La numeración de páginas del índice coincide exactamente con la paginación real del documento; no hay desfases, páginas sin número ni secciones omitidas. El índice se genera automáticamente (preferiblemente) o, si se elabora manualmente, es verificablemente exacto. La profundidad del índice (hasta qué nivel de subsección se incluye) es coherente con la complejidad del documento y con las disposiciones de la Escuela Profesional. &
& & & \\ \hline
\rowcolor{lightgray}
2.4 & Índice de tablas e índice de figuras completos y correctamente referenciados & El proyecto incluye un índice de tablas y un índice de figuras (o lista de ilustraciones) independientes del índice general. Cada tabla y cada figura que aparece en el cuerpo del documento está registrada en su índice correspondiente con su número, su título (caption) y su número de página. No hay tablas ni figuras en el documento que estén ausentes de sus respectivos índices, ni entradas en los índices que no correspondan a elementos existentes en el documento. La numeración de tablas y figuras es consecutiva y coherente con el sistema de numeración adoptado (numeración corrida o numeración por capítulo). Los títulos consignados en los índices coinciden literalmente con los títulos que aparecen en el cuerpo del documento. &
& & & \\ \hline
\multicolumn{3}{|r|}{\cellcolor{white}\textbf{Subtotal Dimensión~II (máx.~12 puntos):}} & \multicolumn{4}{c|}{\cellcolor{white}\textbf{/12}} \\ \hline
\end{longtable}
\normalsize

\clearpage
\subsection{Elementos del Cuerpo del Proyecto}

\textit{Se evalúa la presencia, la secuencia correcta y la articulación de los capítulos y secciones sustantivas que conforman el cuerpo del proyecto de tesis. El cuerpo del documento es el componente que alberga el contenido científico-técnico de la propuesta y su organización debe reflejar la lógica investigativa del proyecto: desde la identificación del problema hasta la planificación metodológica. Esta dimensión verifica la arquitectura del documento, no el contenido de cada sección ---que se evalúa en las rúbricas temáticas correspondientes---.} Los criterios se aprecian en la Tabla~\ref{tab:rubrica-dim3-elementos-cuerpo}.

\scriptsize
\begin{longtable}{|C{0.5cm}|L{1.8cm}|L{9cm}|C{0.3cm}|C{0.3cm}|C{0.3cm}|C{0.3cm}|}
\caption{Rúbrica de evaluación de la Dimensión~III: elementos del cuerpo del proyecto.}\label{tab:rubrica-dim3-elementos-cuerpo}\\
\hline
\rowcolor{headerblue}
\thd{N.\textsuperscript{o}} & \thd{Criterio} & \thd{Indicador verificable} & \thd{3} & \thd{2} & \thd{1} & \thd{0} \\
\hline\endfirsthead
\hline
\rowcolor{headerblue}
\thd{N.\textsuperscript{o}} & \thd{Criterio} & \thd{Indicador verificable} & \thd{3} & \thd{2} & \thd{1} & \thd{0} \\
\hline\endhead
3.1 & Presencia y secuencia correcta de las secciones sustantivas del proyecto & El cuerpo del proyecto contiene, en el orden establecido por el esquema aprobado, las siguientes secciones sustantivas: (a)~planteamiento del problema (descripción de la realidad problemática, formulación del problema general y específicos, justificación, objetivos general y específicos, delimitación del estudio); (b)~marco teórico (antecedentes de la investigación, bases teóricas, hipótesis, variables, definición de términos básicos); (c)~diseño metodológico (tipo, nivel, enfoque y diseño de investigación, población y muestra, técnicas e instrumentos de recolección de datos, plan de análisis de datos); y (d)~aspectos administrativos (cronograma, presupuesto, financiamiento). Ninguna de estas secciones está ausente, fusionada con otra de manera que impida su identificación individual, ni ubicada fuera de la secuencia establecida. &
& & & \\ \hline
\rowcolor{lightgray}
3.2 & Numeración jerárquica consistente de capítulos, secciones y subsecciones & El documento emplea un sistema de numeración jerárquica coherente y sostenido a lo largo de todo el proyecto: capítulos numerados con cifras romanas o arábigas según el formato institucional, secciones con numeración decimal o equivalente (1.1, 1.2, 1.3\ldots), y subsecciones con tercer nivel cuando corresponda (1.1.1, 1.1.2\ldots). La numeración no presenta saltos (p.\,ej., de 2.1 a 2.3 sin 2.2), duplicaciones (dos secciones con el mismo número), ni inconsistencias entre el índice y el cuerpo del documento. Los niveles jerárquicos son distinguibles tipográficamente (tamaño de fuente, negritas, sangría) y permiten al lector identificar la posición de cada sección dentro de la estructura del proyecto sin ambigüedad. &
& & & \\ \hline
3.3 & Existencia de la matriz de consistencia como instrumento articulador del proyecto & El proyecto incluye una matriz de consistencia que sintetiza, en formato tabular, la correspondencia entre los problemas (general y específicos), los objetivos (general y específicos), las hipótesis (general y específicas), las variables con sus dimensiones e indicadores, y la metodología (diseño, población, muestra, técnicas, instrumentos). La matriz permite verificar, en una sola vista, que cada problema tiene un objetivo correspondiente, cada objetivo tiene una hipótesis asociada, cada hipótesis identifica variables operacionalizadas y cada variable tiene un método de medición declarado. La matriz no es un elemento decorativo sino un instrumento operativo de verificación de la coherencia interna del proyecto. &
& & & \\ \hline
\rowcolor{lightgray}
3.4 & Correspondencia verificable entre el índice del documento y el contenido efectivo de cada sección & Cada sección y subsección listada en el índice general existe efectivamente en el cuerpo del documento con la denominación, la numeración y la paginación declaradas. No existen secciones enunciadas en el índice que estén vacías o que contengan únicamente un título sin desarrollo; no hay secciones desarrolladas en el cuerpo que no aparezcan en el índice; y no se detectan discrepancias entre los títulos del índice y los títulos que encabezan las secciones correspondientes en el cuerpo del proyecto. La correspondencia índice-contenido es exacta y verificable mediante revisión cruzada. &
& & & \\ \hline
\multicolumn{3}{|r|}{\cellcolor{white}\textbf{Subtotal Dimensión~III (máx.~12 puntos):}} & \multicolumn{4}{c|}{\cellcolor{white}\textbf{/12}} \\ \hline
\end{longtable}
\normalsize

\clearpage
\subsection{Elementos Finales y Complementarios}

\textit{Se evalúa la inclusión, completitud y corrección formal de los elementos que cierran el documento del proyecto de tesis: referencias bibliográficas, anexos, cronograma de ejecución, presupuesto y financiamiento. Estos elementos, aunque frecuentemente tratados como secundarios, condicionan la viabilidad operativa del proyecto (cronograma, presupuesto), la verificabilidad de las fuentes (referencias) y la completitud documental exigida para el registro y la aprobación institucional (anexos).} Los criterios se aprecian en la Tabla~\ref{tab:rubrica-dim4-elementos-finales}.

\scriptsize
\begin{longtable}{|C{0.5cm}|L{1.8cm}|L{9.5cm}|C{0.3cm}|C{0.3cm}|C{0.3cm}|C{0.3cm}|}
\caption{Rúbrica de evaluación de la Dimensión~IV: elementos finales y complementarios.}\label{tab:rubrica-dim4-elementos-finales}\\
\hline
\rowcolor{headerblue}
\thd{N.\textsuperscript{o}} & \thd{Criterio} & \thd{Indicador verificable} & \thd{3} & \thd{2} & \thd{1} & \thd{0} \\
\hline\endfirsthead
\hline
\rowcolor{headerblue}
\thd{N.\textsuperscript{o}} & \thd{Criterio} & \thd{Indicador verificable} & \thd{3} & \thd{2} & \thd{1} & \thd{0} \\
\hline\endhead
4.1 & Referencias bibliográficas completas, verificables y conforme al sistema de citación exigido & {\linespread{0.92}\selectfont El proyecto incluye una sección de referencias bibliográficas que lista todas las fuentes citadas en el cuerpo del documento, sin omisiones ni entradas huérfanas (referencias listadas que no se citan en el texto, o citas en el texto sin entrada correspondiente en la lista de referencias). Cada referencia contiene los datos bibliográficos completos exigidos por el sistema de citación adoptado (APA, Vancouver, ISO~690, u otro según el reglamento institucional): autor(es), año, título, fuente (revista, editorial, institución), volumen, páginas, DOI o URL cuando corresponda. El formato es uniforme a lo largo de toda la lista y no se mezclan estilos de citación. Las fuentes son verificables y trazables; no se incluyen referencias inventadas, incompletas al punto de ser irrastreables, ni fuentes de calidad inadmisible para un trabajo académico de pregrado.} &
& & & \\ \hline
\rowcolor{lightgray}
4.2 & Anexos correctamente numerados, titulados y referenciados & {\linespread{0.92}\selectfont Los anexos están numerados en forma consecutiva (Anexo~A, Anexo~B\ldots{} o Anexo~1, Anexo~2\ldots) y tienen un título descriptivo que permite identificar su contenido sin recurrir al cuerpo del documento. Cada anexo se cita al menos una vez en el texto principal, en el lugar pertinente (p.\,ej., \enquote{véase Anexo~C: Fichas de ensayo de laboratorio}); no existen anexos sin referencia cruzada ni referencias a anexos inexistentes. Su contenido es relevante para el proyecto: instrumentos de recolección de datos, fichas de ensayo, certificados de calibración, planos, mapas, autorizaciones, fichas técnicas de materiales u otros documentos de soporte verificable.} &
& & & \\ \hline
4.3 & Cronograma de ejecución con estructura temporal coherente y representación gráfica adecuada & {\linespread{0.92}\selectfont El proyecto incluye un cronograma de ejecución que identifica las actividades principales del proyecto con sus duraciones estimadas y su distribución en el tiempo, presentado en formato de diagrama de Gantt, tabla de actividades con fechas o representación equivalente. El cronograma es coherente con la secuencia metodológica descrita en el proyecto: las actividades siguen un orden lógico que respeta las dependencias técnicas (p.\,ej., los ensayos no preceden al muestreo; el análisis estadístico no precede a la obtención de datos). Las duraciones estimadas son plausibles para el alcance del proyecto y compatibles con el calendario académico institucional. El cronograma aparece en la sección designada por el esquema aprobado y está referenciado en el cuerpo del documento.} &
& & & \\ \hline
\rowcolor{lightgray}
4.4 & Presupuesto y financiamiento con desagregación, sustento y fuente de recursos identificada & {\linespread{0.92}\selectfont El proyecto incluye un presupuesto que desagrega los costos asociados a la ejecución de la investigación: ensayos de laboratorio, materiales, transporte, servicios externos, software, impresión y otros gastos pertinentes. Las partidas presentan, como mínimo, descripción, unidad, cantidad, precio unitario y costo total. El presupuesto es coherente con el alcance metodológico declarado: los ensayos presupuestados corresponden a los descritos en la metodología y las cantidades son consistentes con el número de especímenes o muestras previsto. Se identifica la fuente de financiamiento (recursos propios, financiamiento institucional, convenios, fondos concursables) y el monto total es proporcionado al nivel de un proyecto de tesis de pregrado. El presupuesto no es una lista genérica de rubros sin cuantificación ni un monto global sin desagregación verificable.} &
& & & \\ \hline
\multicolumn{3}{|r|}{\cellcolor{white}\textbf{Subtotal Dimensión~IV (máx.~12 puntos):}} & \multicolumn{4}{c|}{\cellcolor{white}\textbf{/12}} \\ \hline
\end{longtable}
\normalsize

\clearpage
\subsection{Coherencia Global y Calidad de la Presentación Documental}

\textit{Se evalúa la unidad e integración del documento del proyecto de tesis como totalidad, verificando que los distintos componentes documentales formen un conjunto coherente, libre de contradicciones internas, con calidad tipográfica y editorial apropiada para un documento académico de pregrado en Ingeniería Civil. Esta dimensión trasciende la verificación individual de cada sección y examina la articulación del documento como producto integrado.} Los criterios se aprecian en la Tabla~\ref{tab:rubrica-dim5-coherencia-global}.

\scriptsize
\begin{longtable}{|C{0.5cm}|L{1.8cm}|L{9.5cm}|C{0.3cm}|C{0.3cm}|C{0.3cm}|C{0.3cm}|}
\caption{Rúbrica de evaluación de la Dimensión~V: coherencia global y calidad de la presentación documental.}\label{tab:rubrica-dim5-coherencia-global}\\
\hline
\rowcolor{headerblue}
\thd{N.\textsuperscript{o}} & \thd{Criterio} & \thd{Indicador verificable} & \thd{3} & \thd{2} & \thd{1} & \thd{0} \\
\hline\endfirsthead
\hline
\rowcolor{headerblue}
\thd{N.\textsuperscript{o}} & \thd{Criterio} & \thd{Indicador verificable} & \thd{3} & \thd{2} & \thd{1} & \thd{0} \\
\hline\endhead
5.1 & Consistencia terminológica y conceptual entre secciones del proyecto & Los términos técnicos, las denominaciones de variables, los nombres de ensayos, las siglas y los acrónimos se emplean de manera uniforme a lo largo de todo el documento. No se utilizan sinónimos intercambiables sin definición previa (p.\,ej., \enquote{estabilización}, \enquote{mejoramiento} y \enquote{tratamiento} usados indistintamente para referirse al mismo proceso), ni se cambian las denominaciones de las variables entre la sección de hipótesis, la operacionalización y la metodología. Las siglas se definen en su primera aparición y se usan consistentemente en adelante. El título del proyecto, los objetivos, las hipótesis y las variables emplean la misma terminología sin discrepancias que generen ambigüedad. &
& & & \\ \hline
\rowcolor{lightgray}
5.2 & Ausencia de contradicciones internas entre secciones del documento & El documento no presenta contradicciones entre sus secciones: el número de objetivos específicos coincide en el planteamiento del problema, en la matriz de consistencia y en la metodología; las variables declaradas en la hipótesis coinciden con las operacionalizadas en el capítulo de variables; los ensayos descritos en la metodología coinciden con los presupuestados; el número de muestras o especímenes es consistente entre la metodología, el presupuesto y el cronograma; y la delimitación temporal del problema es coherente con el horizonte del cronograma. No se detectan cifras, porcentajes, dosificaciones, plazos o denominaciones que difieran entre secciones al referirse al mismo elemento del proyecto. &
& & & \\ \hline
5.3 & Calidad tipográfica, editorial y de presentación visual del documento & El documento exhibe calidad editorial apropiada para un trabajo académico de pregrado: las tablas y figuras presentan títulos (captions) descriptivos, numeración consecutiva y fuente cuando corresponde; las ecuaciones están numeradas y referenciadas en el texto; los márgenes, el interlineado, el tipo y tamaño de fuente son uniformes y conformes al formato institucional; la ortografía y la redacción son correctas en español académico; y la maquetación no presenta viudas, huérfanas, tablas cortadas innecesariamente, figuras sin título o de resolución insuficiente. El documento evidencia revisión editorial previa a la presentación y no da la impresión de un borrador sin pulir. &
& & & \\ \hline
\rowcolor{lightgray}
5.4 & Integración del documento como unidad coherente y evaluable & El proyecto de tesis se lee como una unidad integrada y no como una yuxtaposición de secciones independientes. Existe hilo conductor entre el planteamiento del problema, el marco teórico, las hipótesis, la operacionalización de variables y la metodología: cada sección se construye sobre la anterior y anticipa la siguiente. Las transiciones entre capítulos son lógicas, las referencias cruzadas internas son correctas (p.\,ej., \enquote{como se formuló en la Sección~1.2\ldots}, \enquote{véase la Tabla~4.1\ldots}), y el lector puede reconstruir la lógica completa del proyecto sin necesidad de releer secciones previas para identificar conexiones no explicitadas. El documento demuestra que el tesista comprende su proyecto como totalidad articulada y no como una colección de requisitos formales cumplidos por separado. &
& & & \\ \hline
\multicolumn{3}{|r|}{\cellcolor{white}\textbf{Subtotal Dimensión~V (máx.~12 puntos):}} & \multicolumn{4}{c|}{\cellcolor{white}\textbf{/12}} \\ \hline
\end{longtable}
\normalsize

\newpage
En esta sección se presenta el consolidado de resultados por dimensión y la escala de calificación de la estructura documental, resumidos en las tablas~\ref{tab:consolidado-estructura-documental} y~\ref{tab:escala-calificacion-estructura-documental}.

\begin{table}[htbp]
\centering
\scriptsize
\caption{Cuadro consolidado de resultados por dimensión de la estructura documental}\label{tab:consolidado-estructura-documental}
\begin{tabularx}{\textwidth}{|C{0.8cm}|L{7.5cm}|C{1.8cm}|C{1.8cm}|X|}
\hline
\rowcolor{headerblue}
\thd{Dim.} & \thd{Denominación} & \thd{Pts. Máx.} & \thd{Pts. Obt.} & \thd{Obs.} \\ \hline
I & Adecuación al Esquema Aprobado y Marco Normativo & 12 & & \\ \hline
\rowcolor{lightgray}
II & Elementos Preliminares del Documento & 12 & & \\ \hline
III & Elementos del Cuerpo del Proyecto & 12 & & \\ \hline
\rowcolor{lightgray}
IV & Elementos Finales y Complementarios & 12 & & \\ \hline
V & Coherencia Global y Calidad de la Presentación Documental & 12 & & \\ \hline
\rowcolor{white}
& \textbf{TOTAL GENERAL} & \textbf{60} & & \\ \hline
\end{tabularx}
\normalsize
\end{table}

\begin{table}[htbp]
\centering
\scriptsize
\caption{Escala de calificación de la estructura documental del proyecto de tesis}\label{tab:escala-calificacion-estructura-documental}
\begin{tabularx}{\textwidth}{|C{2.5cm}|C{4.2cm}|X|}
\hline
\rowcolor{headerblue}
\thd{Rango} & \thd{Calificación} & \thd{Significado} \\ \hline
\textbf{54--60}\newline(90--100\,\%) &
\textcolor{csuf}{\textbf{ESTRUCTURA\newline DOCUMENTAL PLENA}} &
El proyecto presenta una estructura documental integral, conforme a la normativa institucional vigente y lógicamente articulada. Todos los elementos preliminares, del cuerpo y finales están presentes, completos y correctamente formateados. La numeración es consistente, las referencias cruzadas son exactas, la matriz de consistencia es operativa y el documento se lee como una unidad coherente. La estructura documental constituye una base sólida para la evaluación de los contenidos sustantivos del proyecto. \\ \hline
\rowcolor{lightgray}
\textbf{42--53}\newline(70--89\,\%) &
\textcolor{cace}{\textbf{ESTRUCTURA\newline DOCUMENTAL SUFICIENTE}} &
El proyecto presenta los componentes documentales esenciales con conformidad normativa sustancial, pero exhibe deficiencias menores en cuanto a completitud de metadatos, exactitud de índices, consistencia de numeración o articulación entre secciones. Las debilidades son corregibles con ajustes específicos sin reestructuración fundamental del documento. \\ \hline
\textbf{30--41}\newline(50--69\,\%) &
\textcolor{cdef}{\textbf{ESTRUCTURA\newline DOCUMENTAL PARCIAL}} &
La estructura documental es incompleta o presenta deficiencias significativas: secciones obligatorias ausentes, elementos preliminares incompletos, numeración inconsistente, ausencia de matriz de consistencia, referencias bibliográficas deficientes o contradicciones internas entre secciones. El proyecto requiere una revisión documental sustancial antes de que sus contenidos puedan ser evaluados con rigor. \\ \hline
\rowcolor{lightgray}
\textbf{0--29}\newline($<$50\,\%) &
\textcolor{cins}{\textbf{ESTRUCTURA\newline DOCUMENTAL INSUFICIENTE}} &
El proyecto no cumple con los requisitos mínimos de estructura documental exigidos por la normativa institucional. Faltan secciones obligatorias, la organización del documento no corresponde al esquema aprobado, los elementos preliminares son inexistentes o gravemente deficientes, y la presentación impide una evaluación académica formal del proyecto. Se requiere una reformulación fundamental de la estructura del documento. \\ \hline
\end{tabularx}
\normalsize
\end{table}

\clearpage
\section{Calidad de la Redacción Académica}

{\footnotesize
La calidad de la redacción académica en un proyecto de tesis de Ingeniería Civil es condición necesaria para la comunicación, la verificabilidad y la credibilidad científica de la propuesta. Un proyecto con redacción imprecisa, ambigua o gramaticalmente deficiente dificulta la evaluación de los dictaminantes y compromete la transmisión exacta de las ideas, procedimientos y hallazgos que constituyen el aporte del tesista. La escritura científica en Ingeniería Civil se caracteriza por la precisión terminológica —cada concepto técnico tiene un significado unívoco en la disciplina—, el rigor discursivo —las afirmaciones se sustentan en evidencia y se articulan mediante una lógica argumentativa verificable— y la corrección lingüística —el texto se ajusta a las normas ortográficas, gramaticales y de puntuación del español académico—. A estas propiedades se suman la cohesión textual —articulación fluida y lógica entre oraciones, párrafos y secciones— y la adecuación al registro académico formal —adopción de convenciones de escritura científica que distinguen el discurso técnico del lenguaje coloquial, periodístico o literario—. Esta rúbrica evalúa la calidad de la redacción académica del proyecto de tesis en cinco dimensiones complementarias que abarcan la competencia comunicativa escrita exigible a un tesista de pregrado en Ingeniería Civil. La evaluación no considera el estilo personal del autor —legítimamente variable—, sino el cumplimiento de los estándares mínimos de claridad, precisión, corrección y rigor que la comunidad académica y profesional exige para que un documento técnico sea publicable, evaluable y reproducible.Las dimensiones de evaluación de la calidad de la redacción académica se resumen en la Tabla~\ref{tab:dim-redaccion-academica}, mientras que la escala de valoración y sus descriptores operativos se presentan en la Tabla~\ref{tab:escala-redaccion-academica}.\par
}

\vspace{-10pt}
\begin{table}[htbp]
\centering
\scriptsize
\caption{Dimensiones de evaluación de la calidad de la redacción académica.}
\label{tab:dim-redaccion-academica}
\begin{tabularx}{\textwidth}{|C{0.8cm}|L{3cm}|X|}
\hline
\rowcolor{headerblue}
\thd{Dim.} & \thd{Denominación} & \thd{Alcance} \\ \hline
I & Precisión terminológica y consistencia léxica & Verifica el uso correcto y unívoco de la terminología técnica de la disciplina, la ausencia de sinonimia no controlada, la consistencia en la denominación de materiales, ensayos y procedimientos a lo largo del documento, y la uniformidad ortográfica en nombres propios, marcas comerciales y designaciones normativas. \\ \hline
\rowcolor{lightgray}
II & Rigor discursivo y sustento argumentativo & Evalúa la diferenciación entre afirmaciones con respaldo bibliográfico y opiniones del autor, la ausencia de lenguaje retórico sin sustento técnico, el uso adecuado de conectores lógicos y transiciones argumentativas, y la construcción de argumentos verificables y libres de falacias lógicas. \\ \hline
III & Corrección lingüística y normativa del idioma & Examina el cumplimiento de las normas ortográficas, gramaticales, de puntuación y de sintaxis del español académico, incluyendo la concordancia de género, número y tiempo verbal, la correcta segmentación de oraciones y la ausencia de errores que dificulten la comprensión del texto. \\ \hline
\rowcolor{lightgray}
IV & Cohesión textual y estructura del discurso & Evalúa la articulación lógica entre oraciones, párrafos y secciones, la progresión temática del texto, la ausencia de saltos argumentativos, la organización interna de los párrafos y la capacidad del documento de guiar al lector a través de la argumentación sin rupturas de continuidad. \\ \hline
V & Registro académico y convenciones de la escritura científica & Verifica la adopción del registro formal propio del discurso científico-técnico, la impersonalidad del texto, el uso apropiado de la voz pasiva refleja, la incorporación de convenciones tipográficas de la escritura técnica y la diferenciación del discurso académico respecto del lenguaje coloquial, periodístico o comercial. \\ \hline
\end{tabularx}
\end{table}

\vspace{-10pt}
\begin{table}[htbp]
\centering
\scriptsize
\caption{Escala de valoración para la evaluación de la calidad de la redacción académica.}
\label{tab:escala-redaccion-academica}
\begin{tabularx}{\textwidth}{|C{1.8cm}|C{.6cm}|X|}
\hline
\rowcolor{headerblue}
\thd{Nivel} & \thd{Pts.} & \thd{Descriptor operativo} \\ \hline
\textcolor{csuf}{\textbf{Excelente}} & \textbf{3} & El proyecto satisface plenamente el criterio evaluado. La redacción es precisa, técnicamente correcta, lingüísticamente impecable y plenamente conforme a las convenciones de la escritura científica en Ingeniería Civil. No se identifican ambigüedades terminológicas, errores gramaticales significativos ni deficiencias discursivas. \\ \hline
\rowcolor{lightgray}
\textcolor{cace}{\textbf{Aceptable}} & \textbf{2} & El criterio es identificable pero se cumple de manera parcial o incompleta. La redacción presenta deficiencias menores que no comprometen la comprensión global del texto pero que afectan la precisión, la fluidez o la conformidad con el registro académico. Susceptible de corrección con revisión editorial específica. \\ \hline
\textcolor{cdef}{\textbf{Deficiente}} & \textbf{1} & El cumplimiento del criterio es tangencial o presenta errores significativos. La redacción exhibe imprecisiones terminológicas, errores gramaticales recurrentes, deficiencias argumentativas o incumplimientos del registro académico que dificultan la lectura, generan ambigüedad o comprometen la credibilidad del documento. \\ \hline
\rowcolor{lightgray}
\textcolor{cins}{\textbf{Ausente}} & \textbf{0} & El proyecto no evidencia consideración alguna del criterio evaluado. La redacción es incompatible con los estándares mínimos de la escritura académica y técnica, al punto de impedir la evaluación rigurosa del contenido científico del proyecto. \\ \hline
\end{tabularx}
\end{table}

\clearpage
\subsection{Precisión Terminológica y Consistencia Léxica}

\textit{Se evalúa si el proyecto emplea la terminología técnica de la Ingeniería Civil de manera correcta, unívoca y consistente a lo largo de todo el documento. La precisión terminológica es un requisito fundamental de la comunicación científica: cada término técnico posee un significado específico dentro de la disciplina y su uso incorrecto, ambiguo o inconsistente compromete la interpretabilidad del texto y la reproducibilidad de los procedimientos descritos. Esta dimensión verifica que el vocabulario técnico del proyecto corresponda al léxico normalizado de la especialidad y que su empleo sea uniforme desde la primera hasta la última página.} Véase la Tabla~\ref{tab:rubrica-dim1-precision-terminologica}.

{\scriptsize
\begin{longtable}{|C{0.4cm}|L{2cm}|L{9.5cm}|C{0.4cm}|C{0.4cm}|C{0.4cm}|C{0.4cm}|}
\caption{Rúbrica de evaluación de la Dimensión~I: precisión terminológica y consistencia léxica.}\label{tab:rubrica-dim1-precision-terminologica}\\
\hline
\rowcolor{headerblue}
\thd{N.\textsuperscript{o}} & \thd{Criterio} & \thd{Indicador verificable} & \thd{3} & \thd{2} & \thd{1} & \thd{0} \\
\hline\endfirsthead
\hline
\rowcolor{headerblue}
\thd{N.\textsuperscript{o}} & \thd{Criterio} & \thd{Indicador verificable} & \thd{3} & \thd{2} & \thd{1} & \thd{0} \\
\hline\endhead
1.1 & Uso correcto y unívoco de la terminología técnica de la disciplina & Los términos técnicos empleados en el proyecto corresponden al léxico normalizado de la Ingeniería Civil y sus especialidades: los nombres de ensayos se designan conforme a la nomenclatura de la norma de referencia (p.\,ej., \enquote{ensayo de compresión no confinada} y no \enquote{prueba de resistencia del suelo}; \enquote{límite líquido} y no \enquote{punto de fluidez del material}); los parámetros mecánicos, hidráulicos o geotécnicos se denominan con su designación técnica precisa (cohesión, ángulo de fricción interna, módulo resiliente, coeficiente de permeabilidad, índice CBR); y los conceptos teóricos se nombran según la convención disciplinar (criterio de Mohr-Coulomb, teoría de consolidación de Terzaghi, modelo de Burmister). No se emplean denominaciones coloquiales, imprecisas o ajenas al campo técnico cuando existe un término normalizado. &
& & & \\ \hline
\rowcolor{lightgray}
1.2 & Ausencia de sinonimia no controlada que genere ambigüedad conceptual & El proyecto no emplea sinónimos no definidos para referirse al mismo concepto, variable, material o procedimiento. Si se utiliza más de un término para designar una misma entidad, la equivalencia se declara explícitamente en la primera aparición (p.\,ej., \enquote{resistencia a la compresión simple, en adelante UCS por sus siglas en inglés}). No se alternan indistintamente denominaciones como \enquote{suelo estabilizado}, \enquote{suelo tratado}, \enquote{suelo mejorado} y \enquote{suelo modificado} sin definir cada término y sin justificar su uso diferenciado o su equivalencia declarada. La sinonimia no controlada genera ambigüedad respecto de si el autor se refiere al mismo fenómeno o a fenómenos distintos, y debe estar resuelta explícitamente en el texto. &
& & & \\ \hline
1.3 & Consistencia en la denominación de materiales, ensayos y procedimientos a lo largo del documento & Los materiales, los ensayos y los procedimientos experimentales se denominan de manera idéntica en todas las secciones del proyecto: el nombre del material empleado como variable independiente es el mismo en el título, en la hipótesis, en la operacionalización de variables, en la metodología, en el presupuesto y en la matriz de consistencia. La designación de los ensayos de laboratorio mantiene la misma referencia normativa en todas sus apariciones (p.\,ej., no se cita \enquote{ASTM~D1883} en una sección y \enquote{ensayo CBR según norma MTC~E~132} en otra sin aclarar la equivalencia). Los nombres de procedimientos, dosificaciones y condiciones experimentales no varían entre secciones al referirse a la misma operación o tratamiento. &
& & & \\ \hline
\rowcolor{lightgray}
1.4 & Consistencia ortográfica en nombres propios, marcas comerciales y designaciones normativas & Los nombres propios de autores citados, las marcas comerciales de productos o equipos y las designaciones de normas técnicas se escriben correctamente y de manera uniforme a lo largo del documento. No se detectan variaciones ortográficas del mismo nombre propio (p.\,ej., \enquote{Terzaghi} / \enquote{Terzagi} / \enquote{Tersaghi}), ni de marcas comerciales (p.\,ej., \enquote{GeoStudio} / \enquote{Geo-Studio} / \enquote{Geostudio}), ni de designaciones normativas (p.\,ej., \enquote{ASTM D-2166} / \enquote{ASTM D2166} / \enquote{ASTM~D~2166-16}). Las denominaciones de instituciones (SUNEDU, CONCYTEC, INGEMMET, MTC) se escriben con la sigla oficial y se expanden en su primera aparición. La uniformidad ortográfica es verificable en todo el documento. &
& & & \\ \hline
\multicolumn{3}{|r|}{\cellcolor{white}\textbf{Subtotal Dimensión~I (máx.~12 puntos):}} & \multicolumn{4}{c|}{\cellcolor{white}\textbf{/12}} \\ \hline
\end{longtable}
}

\clearpage
\subsection{Rigor Discursivo y Sustento Argumentativo}

\textit{Se evalúa si el discurso del proyecto distingue con claridad entre afirmaciones sustentadas bibliográficamente y juicios del autor, si la argumentación se construye sobre evidencia verificable y no sobre retórica vacía, si los conectores lógicos articulan adecuadamente las relaciones entre proposiciones y si la cadena argumentativa es libre de falacias lógicas. El rigor discursivo garantiza que el lector pueda evaluar la solidez de cada afirmación y reconstruir la lógica del razonamiento sin ambigüedad.} Los criterios se aprecian en la Tabla~\ref{tab:rubrica-dim2-rigor-discursivo}.

\scriptsize
\begin{longtable}{|C{0.55cm}|L{2cm}|L{10cm}|C{0.3cm}|C{0.3cm}|C{0.3cm}|C{0.3cm}|}
\caption{Rúbrica de evaluación de la Dimensión~II: rigor discursivo y sustento argumentativo.}\label{tab:rubrica-dim2-rigor-discursivo}\\
\hline
\rowcolor{headerblue}
\thd{N.\textsuperscript{o}} & \thd{Criterio} & \thd{Indicador verificable} & \thd{3} & \thd{2} & \thd{1} & \thd{0} \\
\hline\endfirsthead
\hline
\rowcolor{headerblue}
\thd{N.\textsuperscript{o}} & \thd{Criterio} & \thd{Indicador verificable} & \thd{3} & \thd{2} & \thd{1} & \thd{0} \\
\hline\endhead
2.1 & Diferenciación entre afirmaciones con respaldo bibliográfico y juicios o interpretaciones del autor & El texto distingue con claridad las afirmaciones que provienen de fuentes externas ---marcadas con la cita correspondiente--- de las interpretaciones, inferencias o juicios que el tesista formula con voz propia. No se presentan como hechos establecidos proposiciones que constituyen opinión del autor sin evidencia de respaldo (p.\,ej., \enquote{este método es el más eficiente para estabilizar suelos arcillosos} sin cita ni datos comparativos), ni se diluye la autoría intelectual de ideas ajenas mediante paráfrasis sin atribución. Cuando el tesista formula una interpretación propia, esta se presenta explícitamente como tal (\enquote{a juicio del autor\ldots}, \enquote{se infiere que\ldots}, \enquote{los datos sugieren que\ldots}) y se diferencia del conocimiento reportado en la literatura. &
& & & \\ \hline
\rowcolor{lightgray}
2.2 & Ausencia de lenguaje retórico, hiperbólico o valorativo sin sustento técnico & El proyecto no emplea expresiones retóricas que carezcan de sustento técnico verificable: adjetivos valorativos sin base empírica (\enquote{excelentes resultados}, \enquote{mejora significativa} sin prueba estadística, \enquote{solución óptima} sin análisis de optimización), superlativos no justificados (\enquote{el método más avanzado}, \enquote{la técnica más innovadora}), generalizaciones infundadas (\enquote{es ampliamente conocido que\ldots} sin referencia, \enquote{todos los estudios coinciden en que\ldots} sin revisión sistemática), ni apelaciones emocionales (\enquote{la urgente necesidad de\ldots}, \enquote{la grave situación de\ldots}) que sustituyan la argumentación técnica. Cada calificativo que implique magnitud, frecuencia o valoración está respaldado por evidencia cuantificable o por una referencia verificable. &
& & & \\ \hline
2.3 & Uso adecuado de conectores lógicos y transiciones argumentativas & Los conectores discursivos empleados en el texto corresponden a la relación lógica que efectivamente vincula las proposiciones: los conectores causales (\enquote{por lo tanto}, \enquote{en consecuencia}, \enquote{debido a}) se emplean cuando existe una relación de causa-efecto sustentada; los conectores adversativos (\enquote{sin embargo}, \enquote{no obstante}, \enquote{en contraste}) se reservan para proposiciones que efectivamente se oponen o matizan; y los conectores de adición (\enquote{asimismo}, \enquote{además}, \enquote{de manera complementaria}) se utilizan cuando la nueva proposición efectivamente complementa la anterior. No se emplean conectores causales para vincular proposiciones sin relación causal verificable, ni se usan conectores adversativos donde no existe oposición real, ni se yuxtaponen oraciones sin conexión lógica explícita. &
& & & \\ \hline
\rowcolor{lightgray}
2.4 & Construcción de argumentos verificables y ausencia de falacias lógicas & La argumentación del proyecto se construye mediante proposiciones verificables encadenadas lógicamente: premisas sustentadas en evidencia conducen a conclusiones que se derivan de ellas sin saltos lógicos. El texto no incurre en falacias comunes en la escritura de tesis: argumento de autoridad sin evidencia (\enquote{según el prestigioso autor X, el método Y es superior} sin datos de respaldo), generalización apresurada (extrapolar desde un caso o estudio particular a una conclusión universal), petición de principio (asumir como premisa lo que se pretende demostrar), falsa dicotomía (presentar solo dos alternativas cuando existen más), o non sequitur (conclusiones que no se derivan lógicamente de las premisas presentadas). Cada paso del razonamiento es reconstruible y evaluable por el lector. &
& & & \\ \hline
\multicolumn{3}{|r|}{\cellcolor{white}\textbf{Subtotal Dimensión~II (máx.~12 puntos):}} & \multicolumn{4}{c|}{\cellcolor{white}\textbf{/12}} \\ \hline
\end{longtable}
\normalsize

\clearpage

\subsection{Corrección Lingüística y Normativa del Idioma}

\textit{Se evalúa el cumplimiento de las normas ortográficas, gramaticales, de puntuación y de sintaxis del español académico. Los errores lingüísticos en un documento de tesis no son meras faltas estéticas: una coma mal colocada puede alterar el significado de una especificación técnica, una discordancia verbal puede generar ambigüedad temporal en la descripción de un procedimiento, y una oración sintácticamente defectuosa puede tornar incomprensible una hipótesis o un objetivo de investigación. La corrección lingüística es una condición necesaria para que el contenido técnico sea comunicado sin distorsión.} Los criterios se aprecian en la Tabla~\ref{tab:rubrica-dim3-correccion-linguistica}.

\scriptsize
\begin{longtable}{|C{0.55cm}|L{2cm}|L{10cm}|C{0.3cm}|C{0.3cm}|C{0.3cm}|C{0.3cm}|}
\caption{Rúbrica de evaluación de la Dimensión~III: corrección lingüística y normativa del idioma.}\label{tab:rubrica-dim3-correccion-linguistica}\\
\hline
\rowcolor{headerblue}
\thd{N.\textsuperscript{o}} & \thd{Criterio} & \thd{Indicador verificable} & \thd{3} & \thd{2} & \thd{1} & \thd{0} \\
\hline\endfirsthead
\hline
\rowcolor{headerblue}
\thd{N.\textsuperscript{o}} & \thd{Criterio} & \thd{Indicador verificable} & \thd{3} & \thd{2} & \thd{1} & \thd{0} \\
\hline\endhead
3.1 & Ortografía correcta conforme a las normas vigentes de la Real Academia Española & El texto no presenta errores ortográficos: las palabras se escriben conforme a la normativa vigente de la RAE, las tildes se colocan según las reglas de acentuación (incluyendo las palabras técnicas hispanizadas y los extranjerismos adaptados), y las mayúsculas se emplean según las convenciones del español (no se escriben con mayúscula inicial los nombres de ensayos, los nombres de disciplinas ni los adjetivos gentilicios, salvo cuando forman parte de un nombre propio institucional). Se distingue correctamente entre homófonos que difieren en su grafía (\enquote{haya} / \enquote{halla} / \enquote{allá}; \enquote{echo} / \enquote{hecho}; \enquote{a ver} / \enquote{haber}). Los errores ortográficos aislados, si existen, no son recurrentes ni afectan la comprensión del texto. &
& & & \\ \hline
\rowcolor{lightgray}
3.2 & Puntuación correcta que garantice la interpretación unívoca del texto técnico & La puntuación del texto es correcta y contribuye a la claridad del discurso: las comas separan elementos enumerativos, incisos, oraciones subordinadas y complementos circunstanciales desplazados; los puntos y comas articulan proposiciones relacionadas dentro de un mismo enunciado complejo; los dos puntos introducen enumeraciones, explicaciones o citas; y los puntos finales cierran enunciados completos. No se incurre en errores de puntuación que alteren el significado técnico del texto (p.\,ej., la ausencia de coma antes de una oración de relativo explicativa puede convertirla en especificativa y cambiar el alcance de la afirmación). Las oraciones no son excesivamente largas al punto de perder la referencia del sujeto, ni se fragmentan innecesariamente con puntos que interrumpan la unidad sintáctica. &
& & & \\ \hline
3.3 & Concordancia correcta de género, número y tiempo verbal & El texto mantiene concordancia gramatical correcta en todas sus manifestaciones: concordancia de género y número entre sustantivo y adjetivo, entre sujeto y verbo, y entre pronombre y antecedente. Los tiempos verbales son coherentes dentro de cada sección y entre secciones: el planteamiento del problema y el marco teórico emplean predominantemente el presente indicativo para hechos establecidos y el pretérito para hallazgos de estudios previos; la metodología emplea el futuro o el condicional para acciones planificadas, conforme a la convención adoptada. No se producen cambios injustificados de tiempo verbal dentro de un mismo párrafo ni discordancias entre sujeto y verbo que generen ambigüedad sobre quién ejecuta la acción descrita. &
& & & \\ \hline
\rowcolor{lightgray}
3.4 & Sintaxis clara y estructuras oracionales que faciliten la comprensión del contenido técnico & Las oraciones del proyecto presentan una estructura sintáctica clara, con sujeto, verbo y complementos identificables. No se abusa de oraciones subordinadas encadenadas que dificulten la comprensión (oraciones de más de cinco líneas con múltiples incisos y subordinadas relativas anidadas). Las construcciones pasivas se emplean cuando el agente es irrelevante o desconocido, no como recurso sistemático que oscurezca la redacción. Los gerundios se utilizan correctamente (simultaneidad o anterioridad inmediata) y no se incurre en el gerundio de posterioridad, de cualidad ni copulativo. Los participios concuerdan con sus referentes. La estructura sintáctica permite que un lector con formación en Ingeniería Civil comprenda cada oración en una sola lectura, sin necesidad de releer para descifrar la estructura gramatical. &
& & & \\ \hline
\multicolumn{3}{|r|}{\cellcolor{white}\textbf{Subtotal Dimensión~III (máx.~12 puntos):}} & \multicolumn{4}{c|}{\cellcolor{white}\textbf{/12}} \\ \hline
\end{longtable}
\normalsize

\clearpage
\subsection{Cohesión Textual y Estructura del Discurso}

\textit{Se evalúa la articulación interna del texto como discurso continuo y coherente: la progresión temática entre oraciones y entre párrafos, la organización lógica dentro de cada párrafo, la ausencia de saltos argumentativos y la capacidad del documento de guiar al lector a través de la argumentación sin rupturas de continuidad ni información descontextualizada. La cohesión textual es la propiedad que distingue un texto articulado de una colección de oraciones gramaticalmente correctas pero conceptualmente inconexas.} Los criterios se aprecian en la Tabla~\ref{tab:rubrica-dim4-cohesion-textual}.

\scriptsize
\begin{longtable}{|C{0.55cm}|L{2.8cm}|L{9cm}|C{0.3cm}|C{0.3cm}|C{0.3cm}|C{0.3cm}|}
\caption{Rúbrica de evaluación de la Dimensión~IV: cohesión textual y estructura del discurso.}\label{tab:rubrica-dim4-cohesion-textual}\\
\hline
\rowcolor{headerblue}
\thd{N.\textsuperscript{o}} & \thd{Criterio} & \thd{Indicador verificable} & \thd{3} & \thd{2} & \thd{1} & \thd{0} \\
\hline\endfirsthead
\hline
\rowcolor{headerblue}
\thd{N.\textsuperscript{o}} & \thd{Criterio} & \thd{Indicador verificable} & \thd{3} & \thd{2} & \thd{1} & \thd{0} \\
\hline\endhead
4.1 & Progresión temática lógica entre oraciones y entre párrafos & El texto progresa temáticamente de manera lógica: cada oración se construye sobre la información proporcionada por la anterior (progresión lineal) o retoma un tema previamente introducido para desarrollarlo con mayor profundidad (progresión de tema constante o derivado). No se introducen conceptos, datos o afirmaciones sin antecedente textual que los contextualice; no se producen saltos abruptos entre temas sin transición; y cada párrafo retoma, implícita o explícitamente, la información del párrafo precedente antes de avanzar hacia nueva información. El lector puede seguir el hilo argumentativo sin perder la referencia conceptual que conecta cada proposición con las anteriores. &
& & & \\ \hline
\rowcolor{lightgray}
4.2 & Organización interna de los párrafos: unidad temática, desarrollo y cierre & Cada párrafo del proyecto constituye una unidad temática identificable: presenta una idea central (oración tópica), la desarrolla mediante evidencia, explicación o ejemplificación, y concluye con una proposición que cierra el argumento o prepara la transición al párrafo siguiente. No se detectan párrafos que aborden múltiples temas inconexos dentro del mismo bloque textual, ni párrafos de una sola oración (salvo casos justificados como transiciones entre secciones), ni párrafos excesivamente extensos que acumulen información heterogénea sin estructura interna. La extensión de los párrafos es proporcionada a la complejidad del contenido: ni telegráficamente breves ni desproporcionadamente largos. &
& & & \\ \hline
4.3 & Ausencia de información descontextualizada, redundancias y repeticiones innecesarias & El texto no contiene información descontextualizada ---datos, definiciones o afirmaciones que aparecen sin conexión con el argumento en curso y sin justificación de su presencia--- ni redundancias que repitan sustancialmente la misma idea en diferentes secciones sin aportar nueva información o perspectiva. Cuando una información se reitera (p.\,ej., la mención de una variable en el planteamiento del problema y en la operacionalización), la repetición es funcional y no literal: en cada aparición el concepto se aborda desde la perspectiva propia de la sección correspondiente. El texto no presenta oraciones o párrafos que puedan eliminarse sin afectar la coherencia del argumento. &
& & & \\ \hline
\rowcolor{lightgray}
4.4 & Transiciones efectivas entre secciones y entre capítulos del proyecto & Las transiciones entre secciones y entre capítulos no son abruptas: el cierre de una sección anticipa el contenido de la siguiente, y la apertura de cada nueva sección contextualiza su propósito dentro de la estructura global del proyecto. No se inician capítulos o secciones con contenido técnico sin una presentación que oriente al lector sobre qué se abordará y por qué (p.\,ej., un capítulo de marco teórico no comienza directamente con una definición sin explicar previamente la lógica de organización de la revisión). Las transiciones reflejan la arquitectura argumentativa del proyecto y permiten al lector comprender la función de cada sección dentro del conjunto del documento sin necesidad de consultar el índice. &
& & & \\ \hline
\multicolumn{3}{|r|}{\cellcolor{white}\textbf{Subtotal Dimensión~IV (máx.~12 puntos):}} & \multicolumn{4}{c|}{\cellcolor{white}\textbf{/12}} \\ \hline
\end{longtable}
\normalsize

\clearpage
\subsection{Registro Académico y Convenciones de la Escritura Científica}

\textit{Se evalúa si el proyecto adopta el registro formal propio del discurso científico-técnico, si respeta las convenciones de impersonalidad y objetividad, si emplea las formas gramaticales apropiadas para la escritura académica en español y si incorpora las convenciones tipográficas y notacionales que la comunidad de Ingeniería Civil reconoce como estándares de la comunicación técnica. Esta dimensión distingue un documento académico de un texto informativo, periodístico o comercial sobre el mismo tema.} Los criterios se aprecian en la Tabla~\ref{tab:rubrica-dim5-registro-academico}.

\scriptsize
\begin{longtable}{|C{0.55cm}|L{2cm}|L{10cm}|C{0.3cm}|C{0.3cm}|C{0.3cm}|C{0.3cm}|}
\caption{Rúbrica de evaluación de la Dimensión~V: registro académico y convenciones de la escritura científica.}\label{tab:rubrica-dim5-registro-academico}\\
\hline
\rowcolor{headerblue}
\thd{N.\textsuperscript{o}} & \thd{Criterio} & \thd{Indicador verificable} & \thd{3} & \thd{2} & \thd{1} & \thd{0} \\
\hline\endfirsthead
\hline
\rowcolor{headerblue}
\thd{N.\textsuperscript{o}} & \thd{Criterio} & \thd{Indicador verificable} & \thd{3} & \thd{2} & \thd{1} & \thd{0} \\
\hline\endhead
5.1 & Impersonalidad y objetividad del texto conforme al registro científico & El proyecto emplea un registro impersonal conforme a las convenciones del discurso científico en español: se utiliza preferentemente la voz pasiva refleja (\enquote{se determinó la resistencia\ldots}, \enquote{se evaluaron las propiedades\ldots}), la tercera persona impersonal (\enquote{se propone un diseño experimental\ldots}) o la primera persona del plural cuando la convención institucional lo admite (\enquote{proponemos evaluar\ldots}). No se emplea la primera persona del singular de manera recurrente (\enquote{yo considero}, \enquote{en mi opinión}), ni se adopta un tono subjetivo o emocional que comprometa la apariencia de objetividad del texto. Las valoraciones personales, cuando existen, se presentan como interpretaciones fundamentadas y no como afirmaciones categóricas. &
& & & \\ \hline
\rowcolor{lightgray}
5.2 & Ausencia de lenguaje coloquial, periodístico o comercial en el discurso técnico & El texto no emplea expresiones coloquiales (\enquote{hoy en día}, \enquote{a lo largo y ancho de}, \enquote{desde tiempos inmemoriales}), periodísticas (\enquote{la alarmante cifra de}, \enquote{el sorprendente hallazgo de}, \enquote{la revolucionaria técnica}), ni comerciales (\enquote{el producto líder en el mercado}, \enquote{la solución más avanzada disponible}, \enquote{tecnología de punta}). El registro es uniformemente académico: sobrio, preciso, descriptivo y orientado a la transmisión verificable de información técnica. Los términos evaluativos, cuando se emplean, son propios del léxico científico (\enquote{estadísticamente significativo}, \enquote{técnicamente viable}, \enquote{normativamente conforme}) y no del lenguaje publicitario o de divulgación popular. &
& & & \\ \hline
5.3 & Uso correcto de convenciones tipográficas y notacionales de la escritura técnica & El proyecto respeta las convenciones tipográficas de la escritura científica y técnica: los símbolos de unidades del Sistema Internacional se escriben conforme a la norma ISO~80000 (sin punto abreviativo, con espacio entre la cifra y la unidad, sin pluralización); las expresiones matemáticas y las variables se presentan en cursiva; los vectores y matrices se distinguen tipográficamente según la convención adoptada; los números decimales se expresan con el separador correcto según la convención institucional (coma decimal en español académico); los miles se separan con espacio fino y no con punto ni coma; y las abreviaturas de normas técnicas mantienen el formato oficial de la entidad emisora. Las tablas y figuras presentan leyendas con la terminología y las unidades correctas. &
& & & \\ \hline
\rowcolor{lightgray}
5.4 & Tratamiento apropiado de extranjerismos, siglas y acrónimos en el texto académico & Los extranjerismos técnicos que no poseen equivalente normalizado en español se presentan en cursiva en su primera aparición y se acompañan de una definición o explicación entre paréntesis cuando el término no es de uso generalizado en la comunidad de Ingeniería Civil hispanohablante (p.\,ej., \textit{shotcrete}, \textit{slurry}, \textit{grouting}). Los extranjerismos que poseen equivalente castellano establecido se escriben en español (\enquote{hormigón} o \enquote{concreto}, no \enquote{concrete}; \enquote{esfuerzo cortante}, no \enquote{shear stress}, salvo en contextos donde la denominación en inglés sea la convención disciplinar aceptada). Las siglas y acrónimos se expanden en su primera aparición (\enquote{Índice de Soporte de California (CBR, por sus siglas en inglés: California Bearing Ratio)}) y se utilizan de manera consistente en el resto del documento. No se acumulan siglas no expandidas que dificulten la lectura. &
& & & \\ \hline
\multicolumn{3}{|r|}{\cellcolor{white}\textbf{Subtotal Dimensión~V (máx.~12 puntos):}} & \multicolumn{4}{c|}{\cellcolor{white}\textbf{/12}} \\ \hline
\end{longtable}
\normalsize

\newpage
En esta sección se presenta el consolidado de resultados por dimensión y la escala de calificación de la calidad de la redacción académica, resumidos en las tablas~\ref{tab:consolidado-redaccion-academica} y~\ref{tab:escala-calificacion-redaccion-academica}.

\begin{table}[htbp]
\centering
\scriptsize
\caption{Cuadro consolidado de resultados por dimensión de la calidad de la redacción académica}\label{tab:consolidado-redaccion-academica}
\begin{tabularx}{\textwidth}{|C{0.8cm}|L{7.5cm}|C{1.8cm}|C{1.8cm}|X|}
\hline
\rowcolor{headerblue}
\thd{Dim.} & \thd{Denominación} & \thd{Pts. Máx.} & \thd{Pts. Obt.} & \thd{Obs.} \\ \hline
I & Precisión Terminológica y Consistencia Léxica & 12 & & \\ \hline
\rowcolor{lightgray}
II & Rigor Discursivo y Sustento Argumentativo & 12 & & \\ \hline
III & Corrección Lingüística y Normativa del Idioma & 12 & & \\ \hline
\rowcolor{lightgray}
IV & Cohesión Textual y Estructura del Discurso & 12 & & \\ \hline
V & Registro Académico y Convenciones de la Escritura Científica & 12 & & \\ \hline
\rowcolor{white}
& \textbf{TOTAL GENERAL} & \textbf{60} & & \\ \hline
\end{tabularx}
\normalsize
\end{table}

\begin{table}[htbp]
\centering
\scriptsize
\caption{Escala de calificación de la calidad de la redacción académica}\label{tab:escala-calificacion-redaccion-academica}
\begin{tabularx}{\textwidth}{|C{2.5cm}|C{4.2cm}|X|}
\hline
\rowcolor{headerblue}
\thd{Rango} & \thd{Calificación} & \thd{Significado} \\ \hline
\textbf{54--60}\newline(90--100\,\%) &
\textcolor{csuf}{\textbf{REDACCIÓN\newline ACADÉMICA PLENA}} &
El proyecto exhibe una redacción académica integral: la terminología es precisa y consistente, el discurso es riguroso y verificable, la corrección lingüística es impecable, la cohesión textual permite una lectura fluida y el registro es uniformemente científico. El documento cumple con los estándares de comunicación técnica exigibles a un trabajo de pregrado en Ingeniería Civil y está en condiciones de ser evaluado exclusivamente por su contenido científico, sin que la redacción constituya un obstáculo para la comprensión. \\ \hline
\rowcolor{lightgray}
\textbf{42--53}\newline(70--89\,\%) &
\textcolor{cace}{\textbf{REDACCIÓN\newline ACADÉMICA SUFICIENTE}} &
La redacción presenta los atributos esenciales del discurso científico-técnico, con deficiencias menores y localizadas en precisión terminológica, corrección gramatical, cohesión o registro que no comprometen la comprensión global del texto. Las debilidades son corregibles mediante revisión editorial sin necesidad de reescritura fundamental del documento. \\ \hline
\textbf{30--41}\newline(50--69\,\%) &
\textcolor{cdef}{\textbf{REDACCIÓN\newline ACADÉMICA PARCIAL}} &
La redacción presenta deficiencias significativas en múltiples dimensiones: imprecisiones terminológicas recurrentes, errores gramaticales que afectan la comprensión, discurso retórico sin sustento, falta de cohesión entre secciones o desviaciones frecuentes del registro académico. El documento requiere una revisión editorial profunda y, en algunas secciones, la reescritura de párrafos o pasajes completos. \\ \hline
\rowcolor{lightgray}
\textbf{0--29}\newline($<$50\,\%) &
\textcolor{cins}{\textbf{REDACCIÓN\newline ACADÉMICA INSUFICIENTE}} &
La redacción no alcanza los estándares mínimos de la escritura académica y técnica. Los errores lingüísticos, las imprecisiones terminológicas, la ausencia de rigor discursivo y las deficiencias de cohesión son tan generalizados que impiden la evaluación confiable del contenido científico del proyecto. Se requiere una reescritura fundamental del documento con acompañamiento en competencias de escritura académica. \\ \hline
\end{tabularx}
\normalsize
\end{table}


\clearpage
\section{Formato y Presentación}

{\footnotesize
El formato y la presentación de un proyecto de tesis no constituyen un aspecto meramente estético, sino un requisito de comunicación técnica que garantiza la legibilidad, la trazabilidad y la verificabilidad del documento como producto académico. En el marco del Reglamento Específico para Obtener el Grado Académico de Bachiller y el Título Profesional en la Facultad de Ingeniería y Arquitectura de la Universidad Andina del Cusco (R.\,N.\textsuperscript{o}~529-CU-2025-UAC), el proyecto de tesis debe ajustarse a los esquemas aprobados por Consejo de Facultad para cada Escuela Profesional (Arts.~67 y~89), emplear la carátula institucional oficial, incluir los metadatos y el informe de similitud conforme al Reglamento Marco R\_CU-256-2024 (Anexos~2 y~3), y cumplir con las Consideraciones Generales de la Propuesta para el Contenido del Proyecto de Tesis de la Escuela Profesional de Ingeniería Civil: hoja A4, fuente Times New Roman tamaño~12, interlineado de 1,5~líneas, texto justificado en ambos márgenes, márgenes de 2,5\,cm, paginación en arábigos en el extremo superior derecho (excepto la carátula), jerarquía de hasta cinco niveles de encabezados con formato diferenciado, y referencias en norma APA. La presente rúbrica evalúa la calidad formal del documento en cinco dimensiones: estructura documental y elementos obligatorios, normas tipográficas y maquetación, presentación de tablas y figuras, tratamiento de ecuaciones y expresiones matemáticas, y uso de unidades de medida y convenciones numéricas. Cada dimensión comprende cuatro criterios evaluados en una escala de cuatro niveles (0--3), para un total de 20~criterios y 60~puntos máximos. Las dimensiones de evaluación del formato y la presentación se resumen en la Tabla~\ref{tab:dim-formato}, mientras que la escala de valoración y sus descriptores operativos se presentan en la Tabla~\ref{tab:escala-formato}.\par
}

\vspace{-10pt}
\begin{table}[htbp]
\centering
\scriptsize
\caption{Dimensiones de evaluación del formato y la presentación, y alcance de cada dimensión.}
\label{tab:dim-formato}
\begin{tabularx}{\textwidth}{|C{0.8cm}|L{2.5cm}|X|}
\hline
\rowcolor{headerblue}
\thd{Dim.} & \thd{Denominación} & \thd{Alcance} \\ \hline
I & Estructura documental y elementos obligatorios & Verifica la presencia, completitud y conformidad de los elementos obligatorios del documento: carátula institucional con los campos prescritos (incluyendo línea de investigación y códigos ORCID), metadatos, declaración de autenticidad, informe de similitud, índices generales, índice de tablas, índice de figuras y organización del cuerpo conforme al esquema de cuatro capítulos aprobado por Consejo de Facultad (Planteamiento del Problema, Marco Teórico, Método, Aspectos Operativos) con Referencias en norma APA. \\ \hline
\rowcolor{lightgray}
II & Normas tipográficas y maquetación & Evalúa el cumplimiento de las especificaciones tipográficas institucionales: fuente Times New Roman tamaño~12, interlineado de 1,5~líneas, texto justificado en ambos márgenes, hoja A4, márgenes de 2,5\,cm en los cuatro lados, paginación en números arábigos en el extremo superior derecho (excluyendo la carátula) y jerarquía de cinco niveles de encabezados con el formato prescrito. \\ \hline
III & Presentación de tablas y figuras & Examina la calidad formal de los elementos gráficos y tabulares: numeración secuencial, títulos descriptivos, referencia cruzada en el texto, citación de fuentes, legibilidad, resolución de imágenes y cumplimiento de las convenciones de la escritura técnica en Ingeniería Civil. \\ \hline
\rowcolor{lightgray}
IV & Ecuaciones y expresiones matemáticas & Evalúa la presentación formal de las expresiones matemáticas: numeración secuencial, definición explícita de variables y parámetros, consistencia notacional, referencia cruzada y conformidad con las convenciones de la comunidad disciplinar. \\ \hline
V & Unidades de medida y convenciones numéricas & Verifica el empleo correcto y consistente del Sistema Internacional de Unidades, la notación numérica conforme a las convenciones técnicas (coma decimal, separador de miles con espacio fino), la coherencia dimensional a lo largo de todo el documento y el tratamiento adecuado de cifras significativas y precisión numérica. \\ \hline
\end{tabularx}
\end{table}

\vspace{-10pt}
\begin{table}[htbp]
\centering
\scriptsize
\caption{Escala de valoración (puntaje) y descriptor operativo para la evaluación del formato y la presentación.}
\label{tab:escala-formato}
\begin{tabularx}{\textwidth}{|C{2.8cm}|C{1cm}|X|}
\hline
\rowcolor{headerblue}
\thd{Nivel} & \thd{Pts.} & \thd{Descriptor operativo} \\ \hline
\textcolor{csuf}{\textbf{Excelente}} & \textbf{3} & El proyecto satisface plenamente el criterio evaluado. El elemento formal es completo, correcto, consistente y conforme a la normativa institucional y a las convenciones de la escritura técnica en Ingeniería Civil. No se identifican omisiones ni inconsistencias. \\ \hline
\rowcolor{lightgray}
\textcolor{cace}{\textbf{Aceptable}} & \textbf{2} & El criterio es identificable pero se cumple de manera parcial o con inconsistencias menores. El elemento formal está presente y es funcional, aunque requiere correcciones específicas para alcanzar la conformidad plena. \\ \hline
\textcolor{cdef}{\textbf{Deficiente}} & \textbf{1} & El cumplimiento del criterio es incompleto o presenta errores recurrentes. El elemento formal existe pero con deficiencias que comprometen la legibilidad, la trazabilidad o la conformidad normativa del documento. \\ \hline
\rowcolor{lightgray}
\textcolor{cins}{\textbf{Ausente}} & \textbf{0} & El proyecto no evidencia consideración alguna del criterio evaluado. El elemento formal está ausente, es ilegible o presenta una desconexión tal con la normativa institucional que equivale a una omisión. \\ \hline
\end{tabularx}
\end{table}

\clearpage
\subsection{Estructura Documental y Elementos Obligatorios}

\textit{Se evalúa si el proyecto de tesis contiene todos los elementos documentales obligatorios establecidos por la R.\,N.\textsuperscript{o}~529-CU-2025-UAC y el Reglamento Marco R\_CU-256-2024, si dichos elementos se presentan en el orden y con el formato institucional exigido, y si la organización del cuerpo del documento responde al esquema aprobado por Consejo de Facultad para la Escuela Profesional de Ingeniería Civil.} Véase la Tabla~\ref{tab:rubrica-dim1-estructura-documental}.

{\scriptsize
\begin{longtable}{|C{0.55cm}|L{2cm}|L{10cm}|C{0.3cm}|C{0.3cm}|C{0.3cm}|C{0.3cm}|}
\caption{Rúbrica de evaluación de la Dimensión~I: estructura documental y elementos obligatorios.}\label{tab:rubrica-dim1-estructura-documental}\\
\hline
\rowcolor{headerblue}
\thd{N.\textsuperscript{o}} & \thd{Criterio} & \thd{Indicador verificable} & \thd{3} & \thd{2} & \thd{1} & \thd{0} \\
\hline\endfirsthead
\hline
\rowcolor{headerblue}
\thd{N.\textsuperscript{o}} & \thd{Criterio} & \thd{Indicador verificable} & \thd{3} & \thd{2} & \thd{1} & \thd{0} \\
\hline\endhead
1.1 & Carátula institucional conforme al formato oficial aprobado & El proyecto presenta la carátula conforme al modelo establecido en la R.\,N.\textsuperscript{o}~529-CU-2025-UAC y en la Propuesta para el Contenido del Proyecto de Tesis de la Escuela Profesional de Ingeniería Civil, conteniendo los siguientes elementos en el orden prescrito: (1)~denominación ``Universidad Andina del Cusco''; (2)~``Facultad de Ingeniería y Arquitectura''; (3)~``Escuela Profesional de Ingeniería Civil''; (4)~logo institucional de la Universidad; (5)~la indicación ``Proyecto de Tesis''; (6)~título de la tesis enmarcado entre una línea superior y otra inferior, sin comillas, con solo la primera letra en mayúscula y el resto en minúsculas (salvo nombres propios); (7)~línea de investigación aprobada para Ingeniería Civil en la que se inscribe la investigación; (8)~la mención ``Presentado por:'' seguida del nombre del o los tesistas con sus códigos ORCID; (9)~la frase ``Para optar al Título Profesional de Ingeniero Civil''; (10)~nombre del asesor (y del co-asesor si corresponde) con su código ORCID; y (11)~``Cusco -- Perú'' con el año de presentación. Todos los campos están completos, correctamente ubicados y sin errores tipográficos. La carátula no se numera. &
& & & \\ \hline
\rowcolor{lightgray}
1.2 & Metadatos, declaración de autenticidad e informe de similitud & El documento incluye, inmediatamente después de la carátula y en el orden institucional exigido: (a)~la hoja de metadatos conforme al Anexo~2 del Reglamento Marco R\_CU-256-2024; (b)~la Declaración Personal de Autenticidad y de No Plagio (Formato~T1, R.\,N.\textsuperscript{o}~529-CU-2025-UAC) firmada por el o los tesistas; y (c)~el informe de similitud emitido mediante el software oficial de la Universidad (Turnitin), con el porcentaje de coincidencia claramente visible, conforme a los Arts.~64.f y~91 del reglamento específico. La ausencia de cualquiera de estos tres elementos constituye una omisión grave. &
& & & \\ \hline
1.3 & Índices generales, de tablas y de figuras & El documento contiene un índice general (tabla de contenidos) que refleja fielmente la estructura jerárquica del proyecto con numeración de página correcta, un índice de tablas y un índice de figuras, cada uno con la numeración, el título abreviado y la página correspondiente. Los tres índices son consistentes con el contenido real del documento: no existen entradas huérfanas (listadas pero inexistentes en el cuerpo), entradas omitidas (presentes en el cuerpo pero ausentes del índice), ni discrepancias en la numeración de página. &
& & & \\ \hline
\rowcolor{lightgray}
1.4 & Organización del cuerpo conforme al esquema aprobado por Consejo de Facultad & La estructura del cuerpo del proyecto de tesis sigue el esquema aprobado mediante resolución de Consejo de Facultad para la Escuela Profesional de Ingeniería Civil (Art.~67, R.\,N.\textsuperscript{o}~529-CU-2025-UAC), cuya secuencia obligatoria es: Carátula, Cuerpo del proyecto de tesis (Índices, Capítulo~I: Planteamiento del Problema, Capítulo~II: Marco Teórico, Capítulo~III: Método, Capítulo~IV: Aspectos Operativos) y Referencias. Todos los capítulos y secciones obligatorios están presentes, se identifican con la denominación y la numeración establecidas y aparecen en el orden prescrito. Las referencias bibliográficas se elaboran conforme a la norma APA vigente, como lo exige el esquema institucional. No se incluyen secciones ajenas al esquema sin justificación ni se omiten las obligatorias. La jerarquía de encabezados (capítulo, sección, subsección) es consistente y respeta la estructura aprobada. &
& & & \\ \hline
\multicolumn{3}{|r|}{\cellcolor{white}\textbf{Subtotal Dimensión~I (máx.~12 puntos):}} & \multicolumn{4}{c|}{\cellcolor{white}\textbf{/12}} \\ \hline
\end{longtable}
}

\clearpage
\subsection{Normas Tipográficas y Maquetación}

\textit{Se evalúa si el documento cumple con las especificaciones tipográficas y de maquetación establecidas en las Consideraciones Generales del esquema aprobado para la Escuela Profesional de Ingeniería Civil: fuente Times New Roman tamaño~12, interlineado de 1,5~líneas, texto justificado en ambos márgenes, hoja tamaño~A4, márgenes de 2,5\,cm, paginación en arábigos en el extremo superior derecho y jerarquía de cinco niveles de encabezados con formato diferenciado.} Véase la Tabla~\ref{tab:rubrica-dim2-normas-tipograficas}.

{\scriptsize
\begin{longtable}{|C{0.55cm}|L{2cm}|L{10.5cm}|C{0.3cm}|C{0.3cm}|C{0.3cm}|C{0.3cm}|}
\caption{Rúbrica de evaluación de la Dimensión~II: normas tipográficas y maquetación.}\label{tab:rubrica-dim2-normas-tipograficas}\\
\hline
\rowcolor{headerblue}
\thd{N.\textsuperscript{o}} & \thd{Criterio} & \thd{Indicador verificable} & \thd{3} & \thd{2} & \thd{1} & \thd{0} \\
\hline\endfirsthead
\hline
\rowcolor{headerblue}
\thd{N.\textsuperscript{o}} & \thd{Criterio} & \thd{Indicador verificable} & \thd{3} & \thd{2} & \thd{1} & \thd{0} \\
\hline\endhead
2.1 & Tipo de fuente, tamaño, espaciado y justificación conforme a la normativa institucional & El documento se presenta en hoja tamaño A4 y emplea la fuente \textbf{Times New Roman en tamaño~12} para el cuerpo del texto, con \textbf{interlineado de 1,5~líneas} y \textbf{texto justificado en ambos márgenes}, conforme a las Consideraciones Generales del esquema aprobado para la Escuela Profesional de Ingeniería Civil. El tipo de fuente es uniforme en todo el cuerpo del texto; las excepciones admitidas se limitan a tablas, figuras, notas al pie o código fuente, donde puede emplearse un tamaño menor manteniendo la misma familia tipográfica. La redacción se realiza en tiempo presente o futuro, según lo exige el esquema. Todas las palabras que requieren acentuación llevan tilde sin excepción, independientemente de que estén en mayúsculas o minúsculas. No se observan variaciones injustificadas de fuente, tamaño o espaciado entre secciones o capítulos del documento. &
& & & \\ \hline
\rowcolor{lightgray}
2.2 & Márgenes de 2,5\,cm en los cuatro lados conforme al esquema aprobado & Los márgenes del documento son de \textbf{2,5\,cm} en los cuatro lados (\textbf{superior, inferior, derecho e izquierdo}), conforme a las Consideraciones Generales del esquema aprobado para la Escuela Profesional de Ingeniería Civil. Los márgenes son uniformes a lo largo de todo el documento y no se alteran para forzar la distribución del contenido en un número determinado de páginas. Los elementos gráficos (tablas, figuras, ecuaciones) no exceden los límites de los márgenes establecidos. En modo de impresión física, el documento se presenta impreso en ambas caras de cada hoja (a excepción de la carátula), por lo que los márgenes simétricos de 2,5\,cm se mantienen sin ajuste diferencial para encuadernación, salvo instrucción institucional específica posterior. &
& & & \\ \hline
2.3 & Paginación en números arábigos en el extremo superior derecho, excluyendo la carátula & La numeración de páginas cumple con las Consideraciones Generales del esquema aprobado: \textbf{todas las hojas se numeran con números arábigos, sin excepción, a excepción de la carátula}. La \textbf{numeración se ubica en el extremo superior derecho} de cada página. La secuencia numérica es continua y correcta a lo largo de todo el documento, desde la primera página posterior a la carátula hasta la última página de las referencias o anexos. No se observan saltos, duplicaciones ni omisiones en la secuencia numérica. No se emplea numeración romana en ninguna sección del documento, conforme a la disposición institucional que establece el uso exclusivo de números arábigos para toda la paginación. &
& & & \\ \hline
\rowcolor{lightgray}
2.4 & Jerarquía de cinco niveles de encabezados conforme al formato institucional & Los encabezados del documento respetan la jerarquía de cinco niveles establecida en las Consideraciones Generales del esquema aprobado: \textbf{Nivel~1} (sin numeración): centrado, negrita, cada palabra inicia en mayúscula, el resto en minúsculas; \textbf{Nivel~2} (1.1): alineado a la izquierda, negrita, cada palabra inicia en mayúscula, el texto comienza en nuevo párrafo; \textbf{Nivel~3} (1.1.1): alineado a la izquierda, negrita y cursiva, cada palabra inicia en mayúscula, el texto en nuevo párrafo; \textbf{Nivel~4} (1.1.1.1): con sangría de \textonehalf\textquotedblright{} (1,27\,cm), negrita, cada palabra inicia en mayúscula, punto final, el texto inicia en la misma línea; \textbf{Nivel~5} (1.1.1.1.1): con sangría de \textonehalf\textquotedblright{} (1,27\,cm), negrita y cursiva, cada palabra inicia en mayúscula, punto final, el texto inicia en la misma línea. En ningún caso se presenta un título o subtítulo escrito enteramente en mayúsculas. La jerarquía se mantiene uniforme en todos los capítulos del documento: un mismo nivel jerárquico recibe siempre el mismo tratamiento tipográfico. Se recomienda no emplear más de cinco niveles de profundidad. &
& & & \\ \hline
\multicolumn{3}{|r|}{\cellcolor{white}\textbf{Subtotal Dimensión~II (máx.~12 puntos):}} & \multicolumn{4}{c|}{\cellcolor{white}\textbf{/12}} \\ \hline
\end{longtable}
}

\clearpage
\subsection{Presentación de Tablas y Figuras}

\textit{Se evalúa la calidad formal de los elementos tabulares y gráficos del proyecto: numeración secuencial, títulos descriptivos ubicados según la convención (tablas con título superior, figuras con título inferior), referencia cruzada obligatoria en el texto, citación de fuente, legibilidad y resolución. Una tabla o figura que no se menciona en el texto no cumple función comunicativa; una figura ilegible o una tabla sin título descriptivo comprometen la verificabilidad del documento.} Véase la Tabla~\ref{tab:rubrica-dim3-tablas-figuras}.

{\scriptsize
\begin{longtable}{|C{0.55cm}|L{2cm}|L{9cm}|C{0.3cm}|C{0.3cm}|C{0.3cm}|C{0.3cm}|}
\caption{Rúbrica de evaluación de la Dimensión~III: presentación de tablas y figuras.}\label{tab:rubrica-dim3-tablas-figuras}\\
\hline
\rowcolor{headerblue}
\thd{N.\textsuperscript{o}} & \thd{Criterio} & \thd{Indicador verificable} & \thd{3} & \thd{2} & \thd{1} & \thd{0} \\
\hline\endfirsthead
\hline
\rowcolor{headerblue}
\thd{N.\textsuperscript{o}} & \thd{Criterio} & \thd{Indicador verificable} & \thd{3} & \thd{2} & \thd{1} & \thd{0} \\
\hline\endhead
3.1 & Numeración secuencial y títulos descriptivos de tablas y figuras & Todas las tablas y figuras del documento están numeradas de manera secuencial y continua (numeración simple correlativa o compuesta capítulo-número, según la convención adoptada). Cada tabla lleva un título descriptivo ubicado en la parte superior que indica de forma concisa y precisa su contenido, las variables presentadas y, cuando corresponda, las condiciones del ensayo o la fuente de los datos. Cada figura lleva un título descriptivo ubicado en la parte inferior con la misma precisión. Los títulos no son genéricos (p.\,ej., ``Tabla~1: Resultados'' es insuficiente; ``Tabla~1: Resultados del ensayo de compresión simple a 7, 14 y 28~días para probetas con 4\,\%, 6\,\% y 8\,\% de cemento Portland tipo~IP'' es adecuado). &
& & & \\ \hline
\rowcolor{lightgray}
3.2 & Referencia cruzada en el texto para cada tabla y figura presentada & Toda tabla y toda figura incluida en el documento se mencionan explícitamente en el cuerpo del texto mediante referencia cruzada (p.\,ej., ``como se observa en la Tabla~3'', ``la Figura~7 muestra\ldots''). No existen tablas ni figuras huérfanas (presentes en el documento pero nunca referenciadas en el texto). La referencia cruzada precede o acompaña al elemento gráfico y cumple una función interpretativa: no se limita a señalar la existencia de la tabla o figura, sino que la integra en el hilo argumentativo del texto. La numeración citada en el texto coincide exactamente con la numeración del elemento referido. &
& & & \\ \hline
3.3 & Citación de fuente, nota al pie y créditos en tablas y figuras & Toda tabla y toda figura que reproduzca, adapte o sintetice información proveniente de una fuente externa indica su procedencia de manera explícita conforme a la norma APA vigente (sistema de citación obligatorio según el esquema institucional): ``Fuente: [referencia bibliográfica en formato APA]'' para reproducciones textuales; ``Adaptado de [referencia]'' para modificaciones; ``Elaboración propia'' para tablas y figuras originales del tesista, indicando cuando corresponda ``a partir de datos de [fuente]''. Las notas al pie de tabla aclaran abreviaturas, símbolos o condiciones particulares que no son evidentes desde el título o los encabezados de columna. No existen tablas ni figuras de fuente externa sin atribución ni elementos originales sin la mención de elaboración propia. &
& & & \\ \hline
\rowcolor{lightgray}
3.4 & Legibilidad, resolución y calidad gráfica de tablas y figuras & Las tablas son legibles en el formato de impresión previsto: el tamaño de fuente no es inferior al mínimo aceptable (generalmente 8~pt), las columnas tienen ancho suficiente para su contenido, los encabezados de columna son claros e incluyen las unidades de medida correspondientes, y las líneas de separación facilitan la lectura. Las figuras (fotografías, gráficos, planos, diagramas, mapas) tienen resolución suficiente para su reproducción sin pérdida de legibilidad (mínimo 300~dpi para impresión), los ejes de gráficos están rotulados con la magnitud y la unidad, las escalas son apropiadas, las leyendas son completas y los elementos gráficos son distinguibles tanto en color como en escala de grises. &
& & & \\ \hline
\multicolumn{3}{|r|}{\cellcolor{white}\textbf{Subtotal Dimensión~III (máx.~12 puntos):}} & \multicolumn{4}{c|}{\cellcolor{white}\textbf{/12}} \\ \hline
\end{longtable}
}

\clearpage
\subsection{Ecuaciones y Expresiones Matemáticas}

\textit{Se evalúa la presentación formal de las ecuaciones y expresiones matemáticas del proyecto de tesis: numeración secuencial, definición explícita de todas las variables y parámetros, consistencia notacional y referencia cruzada en el texto. En Ingeniería Civil, las ecuaciones constituyen la expresión formal de los modelos, criterios de diseño y relaciones constitutivas que sustentan la investigación; su presentación deficiente compromete la reproducibilidad y la verificabilidad del trabajo.} Véase la Tabla~\ref{tab:rubrica-dim4-ecuaciones}.

{\scriptsize
\begin{longtable}{|C{0.55cm}|L{2.8cm}|L{8cm}|C{0.3cm}|C{0.3cm}|C{0.3cm}|C{0.3cm}|}
\caption{Rúbrica de evaluación de la Dimensión~IV: ecuaciones y expresiones matemáticas.}\label{tab:rubrica-dim4-ecuaciones}\\
\hline
\rowcolor{headerblue}
\thd{N.\textsuperscript{o}} & \thd{Criterio} & \thd{Indicador verificable} & \thd{3} & \thd{2} & \thd{1} & \thd{0} \\
\hline\endfirsthead
\hline
\rowcolor{headerblue}
\thd{N.\textsuperscript{o}} & \thd{Criterio} & \thd{Indicador verificable} & \thd{3} & \thd{2} & \thd{1} & \thd{0} \\
\hline\endhead
4.1 & Numeración secuencial y referencia cruzada de ecuaciones en el texto & Todas las ecuaciones relevantes del documento están numeradas de manera secuencial (numeración simple correlativa o compuesta capítulo-número, según la convención adoptada) y alineadas conforme a un criterio uniforme (centradas con número a la derecha es la convención más extendida). Cada ecuación numerada se menciona en el texto mediante referencia cruzada (p.\,ej., ``conforme a la Ec.~(3.2)\ldots'', ``sustituyendo en la expresión~(7)\ldots''). No existen ecuaciones numeradas que nunca se citen ni ecuaciones citadas cuya numeración no corresponda. Las expresiones matemáticas menores integradas en el texto (en línea) no requieren numeración, pero las relaciones constitutivas, los modelos y las fórmulas de diseño sí la exigen. &
& & & \\ \hline
\rowcolor{lightgray}
4.2 & Definición explícita de todas las variables y parámetros empleados & Cada ecuación va acompañada, inmediatamente después de su presentación, de una lista de definición en la que se identifica cada variable y cada parámetro con su nombre, su significado físico, su unidad de medida y, cuando corresponda, su valor numérico o su rango admisible. No se emplean variables no definidas, parámetros sin unidades, ni símbolos introducidos sin explicación previa. La lista de definición emplea un formato consistente a lo largo del documento (p.\,ej., ``donde: $\sigma$ = esfuerzo axial [kPa]; $A$ = área de la sección transversal [m$^{2}$]; \ldots'') y no omite ningún símbolo presente en la ecuación. &
& & & \\ \hline
4.3 & Consistencia notacional a lo largo del documento & La notación matemática es consistente en todo el documento: una misma variable se representa siempre con el mismo símbolo; un mismo símbolo no se emplea para designar magnitudes diferentes en distintas secciones; los subíndices y superíndices siguen un criterio uniforme; las variables se presentan en cursiva y las unidades en redonda, conforme a la convención ISO~80000; los vectores y matrices se distinguen tipográficamente del escalar mediante negrita, flecha u otra convención declarada y sostenida. No se observan conflictos de notación entre ecuaciones del marco teórico y ecuaciones de la metodología o del análisis de resultados. &
& & & \\ \hline
\rowcolor{lightgray}
4.4 & Procedencia, trazabilidad y correcta transcripción de las ecuaciones & Toda ecuación proveniente de una fuente externa (norma técnica, artículo científico, manual de diseño) indica su procedencia mediante referencia bibliográfica explícita. Las ecuaciones desarrolladas por el tesista como parte de la investigación se identifican claramente como tales. Las ecuaciones transcritas de fuentes externas están correctamente reproducidas: no se detectan errores de transcripción en signos, exponentes, subíndices, coeficientes o variables. Cuando una ecuación se presenta en una forma simplificada o adaptada respecto de la fuente original, la adaptación se declara y se justifica. &
& & & \\ \hline
\multicolumn{3}{|r|}{\cellcolor{white}\textbf{Subtotal Dimensión~IV (máx.~12 puntos):}} & \multicolumn{4}{c|}{\cellcolor{white}\textbf{/12}} \\ \hline
\end{longtable}
}

\clearpage
\subsection{Unidades de Medida y Convenciones Numéricas}

\textit{Se evalúa el empleo correcto y consistente del Sistema Internacional de Unidades (SI), las convenciones numéricas propias de la escritura técnica en español (coma decimal, espacio fino como separador de miles), la coherencia dimensional a lo largo de todo el documento y el tratamiento de cifras significativas. El incumplimiento de estas convenciones introduce ambigüedad en la interpretación de los resultados y compromete la reproducibilidad de la investigación.} Véase la Tabla~\ref{tab:rubrica-dim5-unidades}.

{\scriptsize
\begin{longtable}{|C{0.55cm}|L{2.8cm}|L{8cm}|C{0.3cm}|C{0.3cm}|C{0.3cm}|C{0.3cm}|}
\caption{Rúbrica de evaluación de la Dimensión~V: unidades de medida y convenciones numéricas.}\label{tab:rubrica-dim5-unidades}\\
\hline
\rowcolor{headerblue}
\thd{N.\textsuperscript{o}} & \thd{Criterio} & \thd{Indicador verificable} & \thd{3} & \thd{2} & \thd{1} & \thd{0} \\
\hline\endfirsthead
\hline
\rowcolor{headerblue}
\thd{N.\textsuperscript{o}} & \thd{Criterio} & \thd{Indicador verificable} & \thd{3} & \thd{2} & \thd{1} & \thd{0} \\
\hline\endhead
5.1 & Empleo del Sistema Internacional de Unidades (SI) conforme a la norma ISO~80000 & El documento emplea las unidades del Sistema Internacional de Unidades conforme a la norma ISO~80000: los símbolos de unidades se escriben en redonda, sin punto abreviativo y sin pluralización (``kg'', no ``kgs.'' ni ``Kg''); se deja un espacio entre la cifra y el símbolo de la unidad (``25~MPa'', no ``25MPa''); el símbolo de grado Celsius se escribe ``\textdegree C'' con espacio entre la cifra y el símbolo; los prefijos del SI se emplean correctamente (``kN'', ``MPa'', ``mm'', ``$\mu$m''). Cuando se emplean unidades ajenas al SI de uso extendido en la práctica profesional local (p.\,ej., psi, pulgadas, libras), se acompaña su equivalencia en unidades SI entre paréntesis o en nota. &
& & & \\ \hline
\rowcolor{lightgray}
5.2 & Consistencia de las unidades a lo largo de todo el documento & Una misma magnitud se expresa en la misma unidad a lo largo de todo el documento, salvo justificación explícita del cambio. No se emplean simultáneamente kilogramos-fuerza y newtons para expresar fuerzas, ni megapascales y kilogramos por centímetro cuadrado para presiones, sin declarar y justificar la equivalencia. Las tablas comparativas presentan todos los valores en las mismas unidades. Los resultados de ensayos se reportan en las unidades que la norma técnica de referencia establece. No se detectan conversiones implícitas, inconsistencias dimensionales ni discrepancias de unidades entre el texto, las tablas, las figuras y las ecuaciones. &
& & & \\ \hline
5.3 & Notación numérica conforme a las convenciones técnicas en español & El documento emplea la coma como separador decimal (``2,54'' y no ``2.54''), conforme a la convención establecida por la norma ISO~80000 para la escritura en español y a la práctica académica latinoamericana. Los millares se separan con espacio fino (``12\,500'' y no ``12.500'' ni ``12,500'') para evitar ambigüedad con el separador decimal. Los números de cuatro cifras pueden escribirse sin separador (``2500'' o ``2\,500''). Estas convenciones se mantienen de manera consistente en el texto, las tablas, las figuras y las ecuaciones. No se observa alternancia entre punto decimal y coma decimal en distintas secciones del documento. &
& & & \\ \hline
\rowcolor{lightgray}
5.4 & Cifras significativas, precisión numérica y coherencia dimensional en resultados & Los valores numéricos presentados en el documento reflejan una precisión coherente con la resolución del instrumento de medición, el método de ensayo empleado y la incertidumbre inherente a la magnitud reportada. No se reportan valores con una precisión ficticia (p.\,ej., ``CBR = 4,3271\,\%'' cuando el ensayo tiene una resolución de $\pm$0,5\,\%), ni se presentan resultados con menos cifras significativas que las que el método permite obtener. Los cálculos intermedios y finales mantienen coherencia dimensional verificable: las unidades resultantes de las operaciones algebraicas son las esperadas y no se producen inconsistencias dimensionales. Los redondeos siguen un criterio uniforme declarado o implícito en las normas técnicas de referencia. &
& & & \\ \hline
\multicolumn{3}{|r|}{\cellcolor{white}\textbf{Subtotal Dimensión~V (máx.~12 puntos):}} & \multicolumn{4}{c|}{\cellcolor{white}\textbf{/12}} \\ \hline
\end{longtable}
}
\normalsize

\newpage
En esta sección se presenta el consolidado de resultados por dimensión y la escala de calificación del formato y la presentación, resumidos en las tablas~\ref{tab:consolidado-formato} y~\ref{tab:escala-calificacion-formato}.

\begin{table}[htbp]
\centering
\scriptsize
\caption{Cuadro consolidado de resultados por dimensión del formato y la presentación}\label{tab:consolidado-formato}
\begin{tabularx}{\textwidth}{|C{0.8cm}|L{7.5cm}|C{1.8cm}|C{1.8cm}|X|}
\hline
\rowcolor{headerblue}
\thd{Dim.} & \thd{Denominación} & \thd{Pts. Máx.} & \thd{Pts. Obt.} & \thd{Obs.} \\ \hline
I & Estructura Documental y Elementos Obligatorios & 12 & & \\ \hline
\rowcolor{lightgray}
II & Normas Tipográficas y Maquetación & 12 & & \\ \hline
III & Presentación de Tablas y Figuras & 12 & & \\ \hline
\rowcolor{lightgray}
IV & Ecuaciones y Expresiones Matemáticas & 12 & & \\ \hline
V & Unidades de Medida y Convenciones Numéricas & 12 & & \\ \hline
\rowcolor{white}
& \textbf{TOTAL GENERAL} & \textbf{60} & & \\ \hline
\end{tabularx}
\normalsize
\end{table}

\begin{table}[htbp]
\centering
\scriptsize
\caption{Escala de calificación del formato y la presentación}\label{tab:escala-calificacion-formato}
\begin{tabularx}{\textwidth}{|C{2.5cm}|C{4.2cm}|X|}
\hline
\rowcolor{headerblue}
\thd{Rango} & \thd{Calificación} & \thd{Significado} \\ \hline
\textbf{54--60}\newline(90--100\,\%) &
\textcolor{csuf}{\textbf{FORMATO Y\newline PRESENTACIÓN PLENOS}} &
El proyecto cumple íntegramente con los requisitos formales institucionales y con las convenciones de la escritura técnica en Ingeniería Civil. La estructura documental es completa, la tipografía y la maquetación son conformes, las tablas, figuras y ecuaciones se presentan con rigor, y las unidades y convenciones numéricas son consistentes. El documento está en condiciones de ser evaluado exclusivamente por su contenido científico, sin que el formato constituya un obstáculo para la lectura, la interpretación o la verificabilidad. \\ \hline
\rowcolor{lightgray}
\textbf{42--53}\newline(70--89\,\%) &
\textcolor{cace}{\textbf{FORMATO Y\newline PRESENTACIÓN\newline SUFICIENTES}} &
El proyecto presenta los elementos formales esenciales con deficiencias menores y localizadas: inconsistencias tipográficas aisladas, tablas o figuras con títulos incompletos, ecuaciones sin definición parcial de variables o convenciones numéricas con alternancia ocasional. Las debilidades no comprometen la comprensión global del documento y son corregibles mediante revisión editorial sin reestructuración. \\ \hline
\textbf{30--41}\newline(50--69\,\%) &
\textcolor{cdef}{\textbf{FORMATO Y\newline PRESENTACIÓN\newline PARCIALES}} &
El proyecto presenta deficiencias formales significativas en múltiples dimensiones: elementos obligatorios ausentes o incompletos, inconsistencias tipográficas recurrentes, tablas o figuras sin numeración o sin referencia cruzada, ecuaciones con variables no definidas o unidades empleadas de manera inconsistente. El documento requiere una revisión formal profunda que afecta a la presentación de varias secciones. \\ \hline
\rowcolor{lightgray}
\textbf{0--29}\newline($<$50\,\%) &
\textcolor{cins}{\textbf{FORMATO Y\newline PRESENTACIÓN\newline INSUFICIENTES}} &
El formato y la presentación no alcanzan los estándares mínimos exigibles. La carátula no se ajusta al modelo institucional, los elementos obligatorios están mayoritariamente ausentes, la tipografía y la maquetación son inconsistentes, las tablas y figuras carecen de títulos o numeración, las ecuaciones no se definen y las unidades de medida se emplean sin criterio uniforme. Las deficiencias formales impiden la evaluación confiable del contenido científico del proyecto. \\ \hline
\end{tabularx}
\normalsize
\end{table}

\clearpage
\section{Referencias Bibliográficas}

{\footnotesize
Las referencias bibliográficas constituyen el aparato erudito que sustenta la investigación y permiten al evaluador verificar la trazabilidad, la originalidad y la fundamentación del proyecto de tesis. En el marco del Reglamento Específico para Obtener el Grado Académico de Bachiller y el Título Profesional en la Facultad de Ingeniería y Arquitectura de la Universidad Andina del Cusco (R.\,N.\textsuperscript{o}~529-CU-2025-UAC), el proyecto de tesis debe emplear un sistema de citación normalizado; las Consideraciones Generales de la Propuesta para el Contenido del Proyecto de Tesis de la Escuela Profesional de Ingeniería Civil establecen de manera explícita: \textbf{``REFERENCIAS: USAR NORMA APA''}. Por su parte, el Art.~91 del reglamento específico exige el informe de similitud emitido por el asesor (Turnitin), lo que presupone un manejo riguroso y verificable de las fuentes. La Rúbrica de Integridad Científica y Ética de la Investigación (Sección~II del instrumento de evaluación normativa) ya evalúa la ausencia de plagio, la corresponsabilidad del asesor y el sistema de citación normalizado; la presente rúbrica complementa aquella evaluación centrándose en la \textit{calidad técnica}, la \textit{pertinencia disciplinar}, la \textit{actualidad}, la \textit{cobertura normativa} y la \textit{corrección formal} de las referencias bibliográficas como componente esencial de un proyecto de investigación en Ingeniería Civil. La rúbrica se estructura en cinco dimensiones con cuatro criterios cada una, para un total de 20~criterios y 60~puntos máximos. Las dimensiones de evaluación de las referencias bibliográficas se resumen en la Tabla~\ref{tab:dim-referencias}, mientras que la escala de valoración y sus descriptores operativos se presentan en la Tabla~\ref{tab:escala-referencias}.\par
}

\vspace{-10pt}
\begin{table}[htbp]
\centering
\scriptsize
\caption{Dimensiones de evaluación de las referencias bibliográficas y alcance de cada dimensión.}
\label{tab:dim-referencias}
\begin{tabularx}{\textwidth}{|C{0.8cm}|L{2.5cm}|X|}
\hline
\rowcolor{headerblue}
\thd{Dim.} & \thd{Denominación} & \thd{Alcance} \\ \hline
I & Sistema de citación y correspondencia texto-referencias & Verifica el empleo consistente de la norma APA (sistema obligatorio según el esquema aprobado para la Escuela Profesional de Ingeniería Civil), la correspondencia biunívoca entre las citas en el texto y las entradas en la lista de referencias, la ausencia de referencias fantasma y huérfanas, y la corrección formal de las citas dentro del texto. \\ \hline
\rowcolor{lightgray}
II & Calidad y jerarquía de las fuentes & Evalúa el predominio de fuentes primarias (artículos en revistas indexadas, normas técnicas, reportes institucionales oficiales) sobre fuentes secundarias o grises, la pertinencia y relevancia técnica de cada fuente citada en relación con el problema de investigación, y la diversidad de tipos documentales que sustenta el proyecto. \\ \hline
III & Actualidad, vigencia y cobertura temporal & Examina la actualidad de las referencias (proporción de fuentes de los últimos cinco años), la inclusión justificada de fuentes clásicas o seminales, la vigencia de las normas técnicas citadas (versiones actualizadas) y la coherencia de la cobertura temporal con las exigencias del estado del arte en la especialidad. \\ \hline
\rowcolor{lightgray}
IV & Normativa técnica y fuentes especializadas en Ingeniería Civil & Verifica la inclusión de las normas técnicas internacionales y nacionales pertinentes al problema de investigación (ASTM, ISO, NTP, MTC, CE.010, E.050, Eurocódigos, entre otras), la referencia a fuentes institucionales y bases de datos oficiales, y la presencia de artículos de la especialidad en revistas reconocidas por la comunidad disciplinar. \\ \hline
V & Formato, completitud y uniformidad de la lista de referencias & Evalúa el ordenamiento de la lista (alfabético conforme a la norma APA), la completitud de los datos bibliográficos de cada entrada (autores, año, título, revista/editorial, volumen, páginas, DOI), la uniformidad tipográfica de todas las entradas y la inclusión de identificadores digitales que garanticen la accesibilidad y verificabilidad de las fuentes. \\ \hline
\end{tabularx}
\end{table}

\vspace{-10pt}
\begin{table}[htbp]
\centering
\scriptsize
\caption{Escala de valoración (puntaje) y descriptor operativo para la evaluación de las referencias bibliográficas.}
\label{tab:escala-referencias}
\begin{tabularx}{\textwidth}{|C{2.8cm}|C{1cm}|X|}
\hline
\rowcolor{headerblue}
\thd{Nivel} & \thd{Pts.} & \thd{Descriptor operativo} \\ \hline
\textcolor{csuf}{\textbf{Excelente}} & \textbf{3} & El proyecto satisface plenamente el criterio evaluado. Las referencias son completas, pertinentes, actualizadas, correctamente formateadas en norma APA y verificables. No se identifican omisiones, inconsistencias ni errores de citación. \\ \hline
\rowcolor{lightgray}
\textcolor{cace}{\textbf{Aceptable}} & \textbf{2} & El criterio es identificable pero se cumple de manera parcial o con inconsistencias menores. Las referencias son funcionales y mayoritariamente correctas, aunque presentan deficiencias localizadas susceptibles de corrección editorial. \\ \hline
\textcolor{cdef}{\textbf{Deficiente}} & \textbf{1} & El cumplimiento del criterio es incompleto o presenta errores recurrentes. Las referencias existen pero con deficiencias que comprometen la trazabilidad, la verificabilidad o la conformidad con la norma APA. \\ \hline
\rowcolor{lightgray}
\textcolor{cins}{\textbf{Ausente}} & \textbf{0} & El proyecto no evidencia consideración alguna del criterio evaluado. Las referencias están ausentes, son inverificables o presentan una desconexión tal con la norma APA y con las exigencias de rigor académico que equivale a una omisión. \\ \hline
\end{tabularx}
\end{table}

\clearpage
\subsection{Sistema de Citación y Correspondencia Texto-Referencias}

\textit{Se evalúa si el proyecto emplea de manera consistente la norma APA ---sistema de citación obligatorio según las Consideraciones Generales del esquema aprobado para la Escuela Profesional de Ingeniería Civil---, si existe correspondencia biunívoca entre las citas en el texto y las entradas en la lista de referencias, y si las citas dentro del texto cumplen con el formato prescrito por dicha norma.} Véase la Tabla~\ref{tab:rubrica-dim1-sistema-citacion}.

{\scriptsize
\begin{longtable}{|C{0.55cm}|L{2cm}|L{9cm}|C{0.3cm}|C{0.3cm}|C{0.3cm}|C{0.3cm}|}
\caption{Rúbrica de evaluación de la Dimensión~I: sistema de citación y correspondencia texto-referencias.}\label{tab:rubrica-dim1-sistema-citacion}\\
\hline
\rowcolor{headerblue}
\thd{N.\textsuperscript{o}} & \thd{Criterio} & \thd{Indicador verificable} & \thd{3} & \thd{2} & \thd{1} & \thd{0} \\
\hline\endfirsthead
\hline
\rowcolor{headerblue}
\thd{N.\textsuperscript{o}} & \thd{Criterio} & \thd{Indicador verificable} & \thd{3} & \thd{2} & \thd{1} & \thd{0} \\
\hline\endhead
1.1 & Empleo consistente y exclusivo de la norma APA vigente como sistema de citación & El proyecto emplea de manera consistente y exclusiva la \textbf{norma APA vigente} (7.\textsuperscript{a}~edición o la versión que la Universidad adopte oficialmente), conforme a la disposición explícita de las Consideraciones Generales del esquema aprobado para la Escuela Profesional de Ingeniería Civil: ``REFERENCIAS: USAR NORMA APA''. No se mezclan formatos de diferentes sistemas de citación (APA con IEEE, Vancouver, ISO~690 u otros) en ninguna sección del documento. Las citas en el texto emplean el sistema autor-año propio de la norma APA (p.\,ej., ``(Terzaghi et~al., 1996)'', ``según Das (2019)\ldots'') y no el sistema numérico ni el sistema nota al pie. La consistencia se mantiene en todo el documento: planteamiento del problema, marco teórico, método, aspectos operativos y anexos. &
& & & \\ \hline
\rowcolor{lightgray}
1.2 & Correspondencia biunívoca entre citas en el texto y entradas en la lista de referencias & Existe correspondencia biunívoca verificable entre las citas incluidas en el cuerpo del texto y las entradas registradas en la lista de referencias: toda fuente citada en el texto aparece en la lista de referencias con los datos completos, y toda entrada de la lista de referencias ha sido citada al menos una vez en el cuerpo del texto. No se detectan discrepancias de apellido, año de publicación o número de autores entre la cita en el texto y la entrada correspondiente en la lista. La verificación es exhaustiva: se abarca el planteamiento del problema, el marco teórico, el método y los aspectos operativos, así como los pies de tabla y figura. &
& & & \\ \hline
1.3 & Ausencia de referencias fantasma y de referencias huérfanas & El proyecto no contiene \textit{referencias fantasma} (fuentes citadas en el texto que no aparecen en la lista de referencias) ni \textit{referencias huérfanas} (entradas en la lista de referencias que no son citadas en ningún lugar del texto). La verificación cruzada entre el cuerpo del documento y la lista de referencias es completa y no arroja discrepancias. La ausencia de ambos tipos de inconsistencia constituye un requisito de integridad documental exigido por la norma APA y por el principio de trazabilidad del Código Nacional de Integridad Científica (Res.~035-2024-CONCYTEC-P). Cada fuente listada cumple una función argumentativa, metodológica o normativa identificable en el texto. &
& & & \\ \hline
\rowcolor{lightgray}
1.4 & Corrección formal de las citas dentro del texto conforme a la norma APA & Las citas dentro del texto cumplen con las reglas de formato de la norma APA vigente: citas textuales cortas (menos de 40~palabras) se integran en el párrafo entre comillas con indicación de autor, año y página; citas textuales largas (40~palabras o más) se presentan en bloque con sangría, sin comillas, con indicación de autor, año y página; las citas parafraseadas incluyen autor y año; las citas con tres o más autores emplean ``et~al.'' desde la primera mención; las citas de fuentes institucionales sin autor personal emplean el nombre de la institución; las comunicaciones personales se citan en el texto pero no se incluyen en la lista de referencias; y las citas secundarias se identifican como tales (``como se citó en\ldots''). No se detectan errores sistemáticos en la aplicación de estas reglas. &
& & & \\ \hline
\multicolumn{3}{|r|}{\cellcolor{white}\textbf{Subtotal Dimensión~I (máx.~12 puntos):}} & \multicolumn{4}{c|}{\cellcolor{white}\textbf{/12}} \\ \hline
\end{longtable}
}

\clearpage
\subsection{Calidad y Jerarquía de las Fuentes}

\textit{Se evalúa la calidad intrínseca de las fuentes que sustentan el proyecto: el predominio de fuentes primarias sobre fuentes secundarias o grises, la pertinencia y relevancia técnica de cada fuente en relación con el problema de investigación en Ingeniería Civil, la diversidad de tipos documentales y la exclusión de fuentes no verificables o no académicas como sustento de afirmaciones técnicas.} Véase la Tabla~\ref{tab:rubrica-dim2-calidad-fuentes}.

{\scriptsize
\begin{longtable}{|C{0.55cm}|L{2cm}|L{9cm}|C{0.3cm}|C{0.3cm}|C{0.3cm}|C{0.3cm}|}
\caption{Rúbrica de evaluación de la Dimensión~II: calidad y jerarquía de las fuentes.}\label{tab:rubrica-dim2-calidad-fuentes}\\
\hline
\rowcolor{headerblue}
\thd{N.\textsuperscript{o}} & \thd{Criterio} & \thd{Indicador verificable} & \thd{3} & \thd{2} & \thd{1} & \thd{0} \\
\hline\endfirsthead
\hline
\rowcolor{headerblue}
\thd{N.\textsuperscript{o}} & \thd{Criterio} & \thd{Indicador verificable} & \thd{3} & \thd{2} & \thd{1} & \thd{0} \\
\hline\endhead
2.1 & Predominio de fuentes primarias: artículos en revistas indexadas, normas técnicas y reportes institucionales oficiales & La mayoría de las fuentes citadas corresponde a fuentes primarias verificables: artículos publicados en revistas indexadas en bases de datos reconocidas (Scopus, Web of Science, SciELO, Latindex Catálogo~2.0), normas técnicas emitidas por organismos de normalización (ASTM, ISO, AASHTO, INACAL/NTP, BSI, Eurocódigos), reportes técnicos institucionales oficiales (MTC, INGEMMET, CENEPRED, ACI, FHWA) y libros de texto de autores reconocidos en la disciplina. Las fuentes primarias representan al menos el 70\,\% del total de referencias. No se sustenta el marco teórico ni la metodología exclusivamente con tesis de pregrado, páginas web de divulgación o fuentes no arbitradas. &
& & & \\ \hline
\rowcolor{lightgray}
2.2 & Pertinencia y relevancia técnica de cada fuente citada respecto del problema de investigación & Cada fuente citada tiene una conexión verificable con el problema de investigación, las variables del estudio, la metodología propuesta o el marco normativo aplicable. No se incluyen referencias de relleno cuya relación con el proyecto sea tangencial, genérica o inexistente. Las fuentes del marco teórico abordan directamente los conceptos, modelos, ensayos o fenómenos que el proyecto investiga. Las fuentes metodológicas corresponden a los procedimientos, normas de ensayo o técnicas analíticas que el proyecto declara emplear. No se detectan referencias que, por su temática, alcance o disciplina, resulten ajenas al campo de la Ingeniería Civil o a la especialidad del estudio. &
& & & \\ \hline
2.3 & Diversidad de tipos documentales que sustenta el proyecto & El conjunto de referencias presenta diversidad de tipos documentales apropiada para un proyecto de investigación en Ingeniería Civil: artículos de investigación original en revistas indexadas (para el estado del arte), normas técnicas (para el marco normativo y los procedimientos de ensayo), libros de texto (para los fundamentos teóricos), tesis de posgrado (como antecedentes directos, empleadas con discernimiento), reportes técnicos institucionales (para datos contextuales) y, cuando corresponda, manuales de diseño y guías profesionales. No se observa una dependencia excesiva de un solo tipo de fuente (p.\,ej., exclusivamente tesis de pregrado o exclusivamente libros de texto) ni se emplean tipos documentales inadecuados como sustento principal (blogs, páginas comerciales, enciclopedias abiertas). &
& & & \\ \hline
\rowcolor{lightgray}
2.4 & Exclusión de fuentes no verificables, no académicas o de dudosa rigurosidad como sustento técnico & El proyecto no emplea como sustento de afirmaciones técnicas, metodológicas o normativas fuentes no verificables (páginas web sin autoría, sin fecha o sin respaldo institucional), no académicas (blogs personales, foros, artículos de Wikipedia, publicaciones de divulgación popular, folletos comerciales de fabricantes sin respaldo técnico documentado) ni de dudosa rigurosidad (revistas depredadoras, publicaciones sin revisión por pares presentadas como fuentes arbitradas). Cuando excepcionalmente se cita una fuente no convencional (catálogo de fabricante, ficha técnica comercial), se identifica su naturaleza y se emplea exclusivamente para información factual específica (dosificaciones del fabricante, especificaciones de producto), nunca como sustento teórico o metodológico. &
& & & \\ \hline
\multicolumn{3}{|r|}{\cellcolor{white}\textbf{Subtotal Dimensión~II (máx.~12 puntos):}} & \multicolumn{4}{c|}{\cellcolor{white}\textbf{/12}} \\ \hline
\end{longtable}
}

\clearpage
\subsection{Actualidad, Vigencia y Cobertura Temporal}

\textit{Se evalúa la distribución temporal de las referencias: la proporción de fuentes recientes que refleja el estado actual del conocimiento, la inclusión justificada de fuentes clásicas o seminales cuando el fundamento teórico lo exige, la vigencia de las normas técnicas citadas (versiones actualizadas, no derogadas) y la coherencia de la cobertura temporal con la naturaleza del problema de investigación.} Véase la Tabla~\ref{tab:rubrica-dim3-actualidad-vigencia}.

{\scriptsize
\begin{longtable}{|C{0.55cm}|L{2cm}|L{9cm}|C{0.3cm}|C{0.3cm}|C{0.3cm}|C{0.3cm}|}
\caption{Rúbrica de evaluación de la Dimensión~III: actualidad, vigencia y cobertura temporal.}\label{tab:rubrica-dim3-actualidad-vigencia}\\
\hline
\rowcolor{headerblue}
\thd{N.\textsuperscript{o}} & \thd{Criterio} & \thd{Indicador verificable} & \thd{3} & \thd{2} & \thd{1} & \thd{0} \\
\hline\endfirsthead
\hline
\rowcolor{headerblue}
\thd{N.\textsuperscript{o}} & \thd{Criterio} & \thd{Indicador verificable} & \thd{3} & \thd{2} & \thd{1} & \thd{0} \\
\hline\endhead
3.1 & Actualidad de las referencias: al menos el 50\,\% de los últimos cinco años para el estado del arte & Al menos el \textbf{50\,\%} de las referencias que sustentan el estado del arte, los antecedentes y la revisión de literatura corresponde a publicaciones de los \textbf{últimos cinco años} contados desde la fecha de presentación del proyecto. Esta proporción se verifica sobre las fuentes empleadas en la construcción del conocimiento actual del problema (marco teórico, antecedentes, discusión de técnicas), excluyendo del cómputo las normas técnicas, los libros clásicos de fundamentos y las fuentes históricas cuya antigüedad está justificada por su carácter seminal. La actualidad de las referencias demuestra que el tesista ha realizado una revisión de literatura vigente y no se limita a reproducir bibliografía obsoleta o desactualizada. &
& & & \\ \hline
\rowcolor{lightgray}
3.2 & Inclusión justificada de fuentes clásicas o seminales cuando el fundamento teórico lo exige & Cuando el proyecto cita fuentes anteriores al periodo de cinco años (obras clásicas, publicaciones seminales, contribuciones fundacionales de la disciplina), su inclusión está justificada por la relevancia intrínseca de la fuente para el marco teórico: teorías fundamentales (Terzaghi, Meyerhof, Mohr-Coulomb), métodos de diseño consolidados (AASHTO~93, método PCA), modelos constitutivos de referencia u otros aportes cuya vigencia trasciende su antigüedad. Estas fuentes no se emplean como sustituto de la revisión actualizada sino como complemento que proporciona los cimientos teóricos del estudio. No se detectan fuentes antiguas cuya inclusión carezca de justificación disciplinar ni fuentes desactualizadas empleadas para describir el estado actual del conocimiento. &
& & & \\ \hline
3.3 & Vigencia de las normas técnicas citadas: versiones actualizadas y no derogadas & Todas las normas técnicas citadas en el proyecto corresponden a la \textbf{versión vigente} al momento de la presentación del proyecto. No se citan normas derogadas, sustituidas o canceladas como si estuvieran en vigor (p.\,ej., no se cita la ASTM~D~1557-12 si ya existe la ASTM~D~1557-12(2021); no se cita una NTP derogada por INACAL sin indicar la norma sustituta). Cuando la investigación requiere referirse a una versión anterior de una norma (por ejemplo, para comparar evoluciones normativas o porque los datos de antecedentes se obtuvieron bajo esa versión), se indica explícitamente que se trata de una versión anterior y se cita también la versión vigente. &
& & & \\ \hline
\rowcolor{lightgray}
3.4 & Cobertura temporal coherente con la naturaleza del problema y la especialidad del estudio & La distribución temporal de las referencias es coherente con la naturaleza del problema de investigación y con el ritmo de producción científica de la especialidad. Temas de rápida evolución tecnológica (aditivos químicos, fibras sintéticas, geosintéticos, BIM, modelación numérica) exigen una concentración mayor de fuentes recientes; temas de fundamentos consolidados (mecánica de suelos clásica, resistencia de materiales) admiten una proporción mayor de fuentes establecidas. La cobertura temporal no presenta vacíos injustificados (décadas sin referencias en un campo con producción continua) ni concentraciones artificiales en un solo año o bienio que sugieran una revisión de literatura superficial o circunstancial. &
& & & \\ \hline
\multicolumn{3}{|r|}{\cellcolor{white}\textbf{Subtotal Dimensión~III (máx.~12 puntos):}} & \multicolumn{4}{c|}{\cellcolor{white}\textbf{/12}} \\ \hline
\end{longtable}
}

\clearpage
\subsection{Normativa Técnica y Fuentes Especializadas en Ingeniería Civil}

\textit{Se evalúa si el proyecto incluye las normas técnicas internacionales y nacionales pertinentes al problema de investigación, si se referencian las fuentes institucionales y bases de datos oficiales relevantes, y si el aparato bibliográfico demuestra inserción en la comunidad disciplinar de la Ingeniería Civil mediante la citación de artículos publicados en revistas reconocidas de la especialidad.} Véase la Tabla~\ref{tab:rubrica-dim4-normativa-tecnica}.

{\scriptsize
\begin{longtable}{|C{0.55cm}|L{2cm}|L{9.5cm}|C{0.3cm}|C{0.3cm}|C{0.3cm}|C{0.3cm}|}
\caption{Rúbrica de evaluación de la Dimensión~IV: normativa técnica y fuentes especializadas en Ingeniería Civil.}\label{tab:rubrica-dim4-normativa-tecnica}\\
\hline
\rowcolor{headerblue}
\thd{N.\textsuperscript{o}} & \thd{Criterio} & \thd{Indicador verificable} & \thd{3} & \thd{2} & \thd{1} & \thd{0} \\
\hline\endfirsthead
\hline
\rowcolor{headerblue}
\thd{N.\textsuperscript{o}} & \thd{Criterio} & \thd{Indicador verificable} & \thd{3} & \thd{2} & \thd{1} & \thd{0} \\
\hline\endhead
4.1 & Inclusión de normas técnicas internacionales pertinentes al problema de investigación & El proyecto cita las normas técnicas internacionales directamente pertinentes a los ensayos, clasificaciones, diseños o procedimientos que la investigación emplea o evalúa: normas ASTM (para ensayos de materiales, suelos, concreto, asfalto), ISO (para sistemas de gestión, terminología, procedimientos normalizados), AASHTO (para diseño de pavimentos y clasificación de suelos), Eurocódigos (para diseño estructural, cuando la investigación lo requiera), NZS o IS (cuando el marco comparativo lo exija), u otras normas internacionales reconocidas por la comunidad de Ingeniería Civil. No se omiten normas de ensayo que el proyecto declara utilizar en su metodología ni normas de clasificación cuyos resultados se reportan en el diagnóstico o en los resultados esperados. &
& & & \\ \hline
\rowcolor{lightgray}
4.2 & Inclusión de normativa nacional peruana aplicable al contexto del estudio & El proyecto cita la normativa técnica nacional vigente pertinente al problema de investigación: Normas Técnicas Peruanas (NTP) emitidas por INACAL, normas del Reglamento Nacional de Edificaciones (CE.010 para pavimentos urbanos, E.050 para suelos y cimentaciones, E.060 para concreto armado, E.030 para diseño sismorresistente, u otras según la especialidad), manuales y directivas del Ministerio de Transportes y Comunicaciones (Manual de Carreteras: Suelos, Geología, Geotecnia y Pavimentos; Manual de Ensayo de Materiales; Especificaciones Técnicas Generales para Construcción), y normativa sectorial adicional cuando el problema lo requiera (normativa ambiental, normativa de gestión de riesgos). No se omiten normas nacionales que rigen los procedimientos o criterios de diseño que el proyecto emplea. &
& & & \\ \hline
4.3 & Referencia a fuentes institucionales y bases de datos oficiales pertinentes & El proyecto cita fuentes institucionales y bases de datos oficiales pertinentes para la contextualización y el sustento empírico del problema: registros del SINPAD/CENEPRED (para datos de riesgo y desastres), INGEMMET (para información geológica y geotécnica), SENAMHI (para datos climáticos e hidrológicos), INEI (para datos demográficos y socioeconómicos), Provías Nacional o Descentralizado (para inventarios viales y datos de tráfico), repositorios institucionales de universidades (para antecedentes directos), y bases de datos técnicas de organismos internacionales (Banco Mundial, BID, CEPAL) cuando el contexto del estudio lo requiera. Las fuentes institucionales se emplean para sustentar datos factuales verificables y no se confunden con opiniones editoriales o documentos de difusión sin rigor técnico. &
& & & \\ \hline
\rowcolor{lightgray}
4.4 & Presencia de artículos de investigación en revistas reconocidas de la especialidad & El aparato bibliográfico incluye artículos de investigación original publicados en revistas reconocidas por la comunidad de Ingeniería Civil, tales como: \textit{Géotechnique}, \textit{Canadian Geotechnical Journal}, \textit{Journal of Geotechnical and Geoenvironmental Engineering} (ASCE), \textit{Transportation Research Record}, \textit{Construction and Building Materials}, \textit{Cement and Concrete Research}, \textit{Soils and Foundations}, \textit{Engineering Geology}, \textit{Revista Ingeniería de Construcción}, u otras revistas indexadas en Scopus, Web of Science o SciELO pertinentes a la temática del proyecto. La presencia de artículos en estas revistas demuestra que el tesista ha consultado la producción científica actual y arbitrada de la disciplina, y no ha limitado su revisión a libros de texto y tesis previas. &
& & & \\ \hline
\multicolumn{3}{|r|}{\cellcolor{white}\textbf{Subtotal Dimensión~IV (máx.~12 puntos):}} & \multicolumn{4}{c|}{\cellcolor{white}\textbf{/12}} \\ \hline
\end{longtable}
}

\clearpage
\subsection{Formato, Completitud y Uniformidad de la Lista de Referencias}

\textit{Se evalúa la presentación formal de la lista de referencias: el ordenamiento conforme a la norma APA (alfabético por apellido del primer autor), la completitud de los datos bibliográficos de cada entrada, la uniformidad tipográfica a lo largo de toda la lista y la inclusión de identificadores digitales que garanticen la accesibilidad y verificabilidad de cada fuente.} Véase la Tabla~\ref{tab:rubrica-dim5-formato-lista}.

{\scriptsize
\begin{longtable}{|C{0.55cm}|L{2cm}|L{9.5cm}|C{0.3cm}|C{0.3cm}|C{0.3cm}|C{0.3cm}|}
\caption{Rúbrica de evaluación de la Dimensión~V: formato, completitud y uniformidad de la lista de referencias.}\label{tab:rubrica-dim5-formato-lista}\\
\hline
\rowcolor{headerblue}
\thd{N.\textsuperscript{o}} & \thd{Criterio} & \thd{Indicador verificable} & \thd{3} & \thd{2} & \thd{1} & \thd{0} \\
\hline\endfirsthead
\hline
\rowcolor{headerblue}
\thd{N.\textsuperscript{o}} & \thd{Criterio} & \thd{Indicador verificable} & \thd{3} & \thd{2} & \thd{1} & \thd{0} \\
\hline\endhead
5.1 & Orden alfabético de la lista de referencias conforme a la norma APA & La lista de referencias se presenta bajo el encabezado ``Referencias'' y está ordenada \textbf{alfabéticamente por el apellido del primer autor} de cada entrada, conforme a la norma APA. Las entradas de un mismo autor se ordenan cronológicamente (de la más antigua a la más reciente); las entradas de un mismo autor y un mismo año se distinguen con sufijos alfabéticos (2019a, 2019b). Las obras institucionales sin autor personal se alfabetizan por el nombre de la institución. El ordenamiento es riguroso y no presenta alteraciones (entradas fuera de secuencia, agrupamientos temáticos no conformes con APA, o mezcla de criterios de ordenación). &
& & & \\ \hline
\rowcolor{lightgray}
5.2 & Completitud de los datos bibliográficos de cada entrada & Cada entrada de la lista de referencias contiene \textbf{todos los datos bibliográficos} exigidos por la norma APA según el tipo de fuente: para \textit{artículos de revista}: autor(es), año, título del artículo, nombre de la revista en cursiva, volumen en cursiva, número entre paréntesis, páginas y DOI; para \textit{libros}: autor(es), año, título en cursiva, editorial y DOI si existe; para \textit{normas técnicas}: organismo emisor, año, designación y título de la norma; para \textit{tesis}: autor, año, título en cursiva, tipo de tesis y universidad, y URL del repositorio; para \textit{reportes institucionales}: institución, año, título y URL. No se detectan entradas incompletas (sin año, sin editorial, sin número de volumen, sin páginas, sin URL de acceso para documentos digitales) que impidan la localización de la fuente. &
& & & \\ \hline
5.3 & Uniformidad tipográfica de todas las entradas de la lista de referencias & Todas las entradas de la lista de referencias presentan \textbf{uniformidad tipográfica} conforme a la norma APA: sangría francesa (primera línea al margen, líneas siguientes con sangría de 1,27\,cm); títulos de revistas y de libros en cursiva; títulos de artículos y de capítulos en redonda; uso consistente de mayúsculas y minúsculas (solo la primera palabra del título y los nombres propios en mayúscula, conforme a APA en español); puntuación uniforme (punto después del año entre paréntesis, punto final de cada entrada); interlineado consistente con el del cuerpo del texto (1,5~líneas según el esquema institucional) o a espacio simple si la normativa institucional lo permite para la lista de referencias. No se observan variaciones de formato entre entradas de un mismo tipo documental. &
& & & \\ \hline
\rowcolor{lightgray}
5.4 & Inclusión de DOI, URL y elementos de accesibilidad digital para fuentes verificables & Las entradas de la lista de referencias incluyen el \textbf{DOI} (\textit{Digital Object Identifier}) cuando la fuente lo posee, presentado como enlace completo (https://doi.org/10.xxxx/xxxxx) conforme a la norma APA vigente. Para fuentes digitales sin DOI, se incluye la URL de acceso directo al documento. Para normas técnicas adquiridas mediante suscripción institucional, se indica la designación completa que permita su localización en los catálogos de los organismos emisores (ASTM, ISO, INACAL). Para tesis y reportes institucionales disponibles en repositorios digitales, se incluye la URL del registro en el repositorio institucional o en ALICIA/RENATI conforme a la Ley~N.\textsuperscript{o}~30035. La inclusión de estos identificadores permite la verificabilidad de las fuentes sin depender exclusivamente de la consulta física de los documentos. &
& & & \\ \hline
\multicolumn{3}{|r|}{\cellcolor{white}\textbf{Subtotal Dimensión~V (máx.~12 puntos):}} & \multicolumn{4}{c|}{\cellcolor{white}\textbf{/12}} \\ \hline
\end{longtable}
}
\normalsize

\newpage
En esta sección se presenta el consolidado de resultados por dimensión y la escala de calificación de las referencias bibliográficas, resumidos en las tablas~\ref{tab:consolidado-referencias} y~\ref{tab:escala-calificacion-referencias}.

\begin{table}[htbp]
\centering
\scriptsize
\caption{Cuadro consolidado de resultados por dimensión de las referencias bibliográficas}\label{tab:consolidado-referencias}
\begin{tabularx}{\textwidth}{|C{0.8cm}|L{7.5cm}|C{1.8cm}|C{1.8cm}|X|}
\hline
\rowcolor{headerblue}
\thd{Dim.} & \thd{Denominación} & \thd{Pts. Máx.} & \thd{Pts. Obt.} & \thd{Obs.} \\ \hline
I & Sistema de Citación y Correspondencia Texto-Referencias & 12 & & \\ \hline
\rowcolor{lightgray}
II & Calidad y Jerarquía de las Fuentes & 12 & & \\ \hline
III & Actualidad, Vigencia y Cobertura Temporal & 12 & & \\ \hline
\rowcolor{lightgray}
IV & Normativa Técnica y Fuentes Especializadas en Ingeniería Civil & 12 & & \\ \hline
V & Formato, Completitud y Uniformidad de la Lista de Referencias & 12 & & \\ \hline
\rowcolor{white}
& \textbf{TOTAL GENERAL} & \textbf{60} & & \\ \hline
\end{tabularx}
\normalsize
\end{table}

\begin{table}[htbp]
\centering
\scriptsize
\caption{Escala de calificación de las referencias bibliográficas}\label{tab:escala-calificacion-referencias}
\begin{tabularx}{\textwidth}{|C{2.5cm}|C{4.2cm}|X|}
\hline
\rowcolor{headerblue}
\thd{Rango} & \thd{Calificación} & \thd{Significado} \\ \hline
\textbf{54--60}\newline(90--100\,\%) &
\textcolor{csuf}{\textbf{REFERENCIAS\newline BIBLIOGRÁFICAS\newline PLENAS}} &
Las referencias son integrales: el sistema de citación APA se aplica de manera rigurosa y consistente, la correspondencia texto-referencias es biunívoca, las fuentes son primarias, pertinentes y actualizadas, la normativa técnica nacional e internacional está completa, y la lista de referencias cumple con todos los requisitos formales. El aparato bibliográfico demuestra una revisión de literatura exhaustiva, actualizada y disciplinarmente pertinente que sustenta de manera verificable todas las afirmaciones del proyecto. \\ \hline
\rowcolor{lightgray}
\textbf{42--53}\newline(70--89\,\%) &
\textcolor{cace}{\textbf{REFERENCIAS\newline BIBLIOGRÁFICAS\newline SUFICIENTES}} &
Las referencias presentan los atributos esenciales de un aparato bibliográfico riguroso, con deficiencias menores y localizadas: inconsistencias aisladas en el formato APA, alguna referencia fantasma u huérfana, proporción de fuentes recientes ligeramente inferior al umbral, omisión de una norma técnica secundaria, o datos bibliográficos incompletos en entradas puntuales. Las debilidades son corregibles mediante revisión editorial sin necesidad de reformular la revisión de literatura. \\ \hline
\textbf{30--41}\newline(50--69\,\%) &
\textcolor{cdef}{\textbf{REFERENCIAS\newline BIBLIOGRÁFICAS\newline PARCIALES}} &
Las referencias presentan deficiencias significativas en múltiples dimensiones: mezcla de sistemas de citación, proporción elevada de referencias fantasma o huérfanas, predominio de fuentes secundarias o no arbitradas, antigüedad generalizada de las fuentes, omisión de normas técnicas fundamentales, o formato de la lista de referencias con errores recurrentes. El aparato bibliográfico requiere una revisión profunda que puede implicar la ampliación de la revisión de literatura. \\ \hline
\rowcolor{lightgray}
\textbf{0--29}\newline($<$50\,\%) &
\textcolor{cins}{\textbf{REFERENCIAS\newline BIBLIOGRÁFICAS\newline INSUFICIENTES}} &
Las referencias no alcanzan los estándares mínimos exigibles. El sistema de citación no se identifica o se aplica de manera errática, no existe correspondencia entre citas y lista de referencias, las fuentes son predominantemente no académicas o inverificables, las normas técnicas están ausentes y la lista de referencias carece de los datos mínimos para localizar las fuentes. Se requiere una reformulación fundamental de la revisión de literatura y del aparato bibliográfico del proyecto. \\ \hline
\end{tabularx}
\normalsize
\end{table}

\end{document}